\begin{bt}%[1H8V2-2]
% nguon: Archive/SP_SauLuoi_2024-2025/Sp_GHKII/data/GHKII/1-TN-DS-TLN-TL-THPT-ThiXaQuangTri-QuangTri-GHKII-NH24-25.tex #146
 %[Dự án đề GHKII THPT Thị Xã Quảng Trị - Quảng Trị - 24-25][Lê Phúc]
	Cho hình chóp $S.ABCD$ có đáy $ABCD$ là một hình vuông và $SA\perp(ABCD)$.
	\begin{enumerate}
		\item Chứng minh rằng $BC\perp(SAB)$.
	\item Gọi $H$ là hình chiếu vuông góc của $A$ lên mặt phẳng $(SBD)$. Chứng minh rằng $H$ là trực tâm của tam giác $SBD$.	
	\end{enumerate}
	\loigiai{
		\begin{center}
			\begin{tikzpicture}[scale=1, font=\footnotesize, line join=round, line cap=round, >=stealth]
				\coordinate (A) at (0,0);
				\coordinate (B) at (-2,-3);
				\coordinate (D) at (7,0);
				\coordinate (C) at ($(B)+(D)-(A)$);
				\coordinate (S) at ($(A)+(0,5)$);
				\coordinate (O) at (intersection of A--C and B--D);
				\coordinate (H) at ($(S)!2/3!(O)$);
				\draw(S)--(B) (S)--(C) (S)--(D) (B)--(C)--(D);
				\draw[dashed,thin](S)--(A)--(C) (B)--(A)--(H) (A)--(D)--(B) (S)--(O) (D)--(H);
				\pic[draw,thin,angle radius=3mm] {right angle = S--A--D} pic[draw,thin,angle radius=3mm] {right angle = S--A--B} pic[draw,thin,angle radius=3mm] {right angle = A--H--O};
				\foreach \i/\g in {S/90,A/-90,B/-90,C/-90,D/0,O/-90,H/45}{\draw[fill=black](\i) circle (1pt) ($(\i)+(\g:3mm)$) node[scale=1]{$\i$};}
			\end{tikzpicture}
		\end{center}
	\begin{enumerate}
		\item Ta có $\heva{&BC\perp AB \ (ABCD \ \text{là hình vuông})\\&BC\perp SA \ (SA\perp (ABCD))}\Rightarrow BC \perp (SAB)$.
	\item 	Gọi $O$ là tâm hình vuông $ABCD$.\\
	Kẻ $AH\perp SO$.\\
	\begin{itemize}
		\item $\heva{&BD\perp AC\\&BD\perp SA}\Rightarrow BD\perp (SAC)\Rightarrow BD\perp AH$.\\
	\item  $\heva{&AH\perp SO\\&AH\perp BD}$\\
	$\Rightarrow AH\perp (SBD)$ hay $H$ là hình chiếu của $A$ lên mặt phẳng $(SBD)$.
	\item $\heva{&AH\perp(SBD)\\&SB\subset(SBD)}\Rightarrow AH\perp SB$.\\
	\item $\heva{&BC\perp (SAB)\\&BC\parallel AD}\Rightarrow AD\perp (SAB) \Rightarrow AD\perp SB$.
	\item $\heva{&SB\perp AH\\&SB\perp AD}\Rightarrow SB\perp (AHD)\Rightarrow SB\perp DH$.
	\end{itemize}
	Gọi $E$ là giao điểm của $DH$ và $SB$.\\
	Suy ra $DE\perp SB$.\\
	Xét $\triangle SBD$ có $SO$ và $DE$ là hai đường cao cắt nhau tại $H$.\\
	Vậy $H$ là trực tâm của $\triangle SBD$.
	\end{enumerate}

	}
\end{bt}

\begin{bt}%[1H8H1-3]
% nguon: Archive/SP_SauLuoi_2024-2025/Sp_GHKII/data/GHKII/1-TN-DS-TLN-TL-THPT-ThiXaQuangTri-QuangTri-GHKII-NH24-25.tex #147
 %[Dự án đề GHKII THPT Thị Xã Quảng Trị - Quảng Trị - 24-25][Lê Phúc]
\immini{Khung của một mái nhà được ghép các thanh gỗ như hình bên. Cho biết tam giác $ABC$ vuông cân tại $A$. Tính góc giữa hai thanh gỗ $a$ và $b$ $(b\equiv BC)$.}
{
\begin{tikzpicture}[scale=0.6, font=\footnotesize, line join=round, line cap=round, >=stealth]
	\coordinate (A1) at (-3,1);
	\coordinate (A2) at (1,-1);
	\coordinate (A3) at (6,1);
	\coordinate (A4) at (2,2);
	\coordinate (A) at (4,5);
	\coordinate (B) at (3,3);
	\coordinate (D) at (6,3.3);
	\coordinate (A7) at (-1,5);
	\coordinate (C) at (intersection of A--A3 and B--D);
	\coordinate (A5) at (3,0.5);
	\coordinate (A6) at (1,-2);
	\coordinate (E) at ($(A7)!0.6!(A)$);
	\coordinate (F) at ($(A1)!0.6!(A2)$);
	\coordinate (G) at ($(A7)!0.3!(A)$);
	\coordinate (H) at ($(A1)!0.3!(A2)$);
	\coordinate (K) at ($(A7)!0.75!(A)$);
	\coordinate (L) at ($(A1)!0.75!(A2)$);
	\coordinate (A9) at ($(A7)!0.15!(A)$);
	\coordinate (A10) at ($(A1)!0.15!(A2)$);
	\coordinate (A11) at ($(A7)!0.45!(A)$);
	\coordinate (A12) at ($(A1)!0.45!(A2)$);
	\coordinate (A13) at ($(A7)!0.85!(A)$);
	\coordinate (A14) at ($(A1)!0.85!(A2)$);
	\coordinate (A15) at ($(A4)!0.15!(A3)$);
	\coordinate (A16) at ($(A4)!0.3!(A3)$);
	\coordinate (A17) at ($(A4)!0.6!(A3)$);
	\coordinate (A18) at ($(A4)!0.75!(A3)$);
	\coordinate (A19) at ($(A4)!0.85!(A3)$);
	\draw[brown,fill=lightgray](A1)--(A2)--(A3)--(A4)--cycle;
	\draw[fill=magenta](0,1)--(-0.5,1)--(-0.5,5.5)--(0,5.5)--cycle;
	\draw[brown](A)--(A7) (A7)--(A1)  (B)--(D) (A4)--(A7) (A)--(A2) (A3)--(A) (G)--(H) (K)--(L) (A9)--(A10) (A11)--(A12) (A13)--(A14) (A17)--(E) (A16)--(G) (A18)--(K) (A15)--(A9) (A19)--(A13);
	\draw[blue,very thick](F)--(E) node [above right]{$a$};
	\draw[blue,very thick](A3)--(A) (A2)--(A);
	\draw[blue,very thick] (B) -- (C) node[midway, above] {$b$};
	\foreach \i/\g in {B/180,C/45,A/90}{\draw[fill=blue](\i) circle (1pt) ($(\i)+(\g:3mm)$) node[scale=1]{$\i$};}
\end{tikzpicture}
}
	\loigiai{
	Do tam giác $ABC$ vuông cân tại $A$ nên $\widehat{ABC}=45^\circ$.\\
	Suy ra $(AB,BC)=45^\circ$.\\
	Do đó $(a,b)=(AB,BC)=45^\circ$.
	}
\end{bt}

\begin{bt}%[1H8V4-3]
% nguon: Archive/SP_SauLuoi_2024-2025/Sp_GHKII/data/GHKII/1-TN-DS-TLN-TL-THPT-NguyenHuuHuan-HCM-GHKII-NH24-25.tex #153
%[Đề thi GHKII THPT Nguyễn Hữu Huân - NH24-25][Nguyễn Ngọc Huy Trường]
	Cho hình chóp $S.ABCD$ có đáy $ABCD$ là hình thang vuông tại $A$ và $B$ có $AD=2a$ và $AB=BC=a$, $SA=a$, $SA\perp(ABCD)$. Gọi $I$ là trung điểm của $AD$.
	\begin{enumerate}
		\item Chứng minh rằng $BC\perp(SAB)$.
		\item Chứng minh rằng $(SAC)\perp(SBI)$.
		\item Tính góc tạo bởi mặt phẳng $(SBC)$ và $(SAD)$.
	\end{enumerate}
	\loigiai{
		\begin{center}
			\begin{tikzpicture}[scale=0.8, font=\footnotesize,line join=round, line cap=round, >=stealth]
				\path 
				(0,0) coordinate (A)
				++(-150:3) coordinate (B)
				++(0:4) coordinate (C)
				($(A)+2*(C)-2*(B)$) coordinate (D)
				($(A)+(0,4)$) coordinate (S)
				($(A)!0.5!(D)$) coordinate (I)
				($(S)+(4,0)$) coordinate (x)
				;
				\foreach \i in{B,C,D}{\draw (S)--(\i);};
				\draw (B)--(C)--(D) (S)--(x);
				\draw[dashed] (C)--(A)--(B)--(I)--(C) (S)--(A)--(D);
				\pic[draw,angle eccentricity=1.8,angle radius=2mm]{right angle=B--A--D};
				\pic[draw,angle eccentricity=1.8,angle radius=2mm]{right angle=C--B--A};
				\pic[draw,angle eccentricity=1.8,angle radius=2mm]{right angle=D--A--S};
				\foreach \i/\g in {B/-120,C/-90,D/-90,A/150,S/90,I/90}
				\fill[black] (\i) circle(1pt)+(\g:4mm)node[scale=1]{$\i$};
				\node[above] at (x) {$x$};
			\end{tikzpicture}
		\end{center}
		\begin{enumerate}
			\item Ta có $BC\perp AB$, $BC\perp SA$ ($SA\perp (ABCD)$) suy ra $BC\perp (SAB)$.
			\item Ta có $AI = \dfrac{1}{2}AD = \dfrac{1}{2}\cdot 2a = a = BC$. Mà $AI\parallel BC$ nên $AIBC$ là hình bình hành. \\
			Lại có $\widehat{ABC} =90^\circ$ và $AB = BC$ nên $AIBC$ là hình vuông. \\
			Ta có $BI\perp AC$, $BI\perp SA$ suy ra $BI\perp (SAC)$. Do đó $(SAC)\perp (SBI)$.
			\item Ta có $S\in (SBC)\cap (SAD)$, $AD\parallel BC$ suy ra $(SBC)\cap (SAD) = Sx$ với $Sx\parallel AD\parallel BC$. \\
			Ta có $BC\perp (SAB)$ suy ra $BC\perp SB$ mà $BC\parallel Sx$ nên $SB\perp Sx$. \\
			Ta cũng có $SA\perp AD$, $AD\parallel Sx$ suy ra $SA\perp Sx$. \\
			Ta có $\heva{&(SBC)\cap (SAD)=Sx \\ &SA\perp Sx,\, SA\subset (SAD) \\ &SB\perp Sx,\, SB\subset (SBC)}\Rightarrow ((SBC),(SAD)) = (SB,SA) = \widehat{BSA}$. \\
			Xét $\triangle SAB$ vuông tại $A$, ta có 
			$\tan\widehat{BSA} = \dfrac{AB}{AS} = \dfrac{a}{a} = 1$. \\
			Suy ra $\widehat{BSA} = 45^\circ$. \\
			Vậy $((SBC),(SAD)) = (SB,SA) = \widehat{BSA}=45^\circ$.
		\end{enumerate}
	}
\end{bt}

\begin{bt}%[1H8H2-2]
% nguon: Archive/SP_SauLuoi_2024-2025/Sp_GHKII/data/GHKII/1-TN-DS-TLN-TL-THPT-HoaHoi-BRVT-GHKII-NH24-25.tex #163
%[Dự án tex đề GHKII THPT Hoà Hội - Bà Rịa Vũng Tàu- 24-25][Đỗ Minh Vũ]
	 			Cho hình chóp $S.ABCD$ có đáy $ABCD$ là hình chữ nhật, $SA$ vuông góc với mặt phẳng $(ABCD)$, $SA = AB$, gọi $H$ là hình chiếu vuông góc của $A$ lên đường thẳng $SB$, $M$ là trung điểm $SC$.
	 			\begin{enumerate}
	 				\item Chứng minh rằng đường thẳng $BC$ vuông góc với mặt phẳng $(SAB)$.
	 				\item Chứng minh rằng đường thẳng $MH$ vuông góc với đường thẳng $SA$.
	 			\end{enumerate}
	 		\loigiai{
	 			\begin{center}
	 				\begin{tikzpicture}[scale=1, font=\footnotesize, line join=round, line cap=round]
	 					\foreach \x\y\t in {0/3/S,0/0/A,-1.3/-1.6/B,2.5/-1.6/C}
	 					\coordinate (\t) at (\x,\y);
	 					\coordinate (D) at ($(A)+(C)-(B)$);
	 					\coordinate (H) at ($(S)!1/2!(B)$);
	 					\coordinate (M) at ($(S)!1/2!(C)$);
	 					\draw (M)--(H) (S)--(C)--(B)--(S)--(D)--(C);
	 					\draw[dashed](S)--(A)  (B)--(A)--(D) (H)--(A);
	 					\path pic[draw,angle radius=6]{right angle=S--A--C};
	 					\foreach \t/\g in {S/90,A/-90,B/-130,C/-60,D/0,H/160,M/40}
	 					\draw[fill=black] (\t) circle(1.5pt)
	 					node[shift={(\g:8pt)}]{$\t$};
	 				\end{tikzpicture}
	 			\end{center}
	 			\begin{enumerate}
	 				\item Chứng minh $BC\perp (SAB)$\\
	 				Ta có \\
	 				$BC \perp AB$ (do $ABCD$ là hình chữ nhật)\\
	 				$BC \perp SA$ (do $SA \perp (ABCD)$)\\
	 				$\Rightarrow BC \perp (SAB)$.
	 				\item Chứng minh $MH \perp SA$\\
	 				Tam giác $SAB$ vuông cân tại $A$ (do $SA=AB$) nên $AH$ vừa là đường cao vừa là trung tuyến $\Rightarrow H$ là trung điểm $SB$.\\
	 				Và $M$ là trung điểm $SC$ (giả thiết)\\
	 				Nên $HM \parallel BC$ (tính chất đường trung bình).\\
	 				$\Rightarrow HM \perp (SAB) \Rightarrow HM \perp SA$ (do $SA \subset (SAB)$).
	 			\end{enumerate}
	 		}
	 	\end{bt}

\begin{bt}%[1H8H6-1]
% nguon: Archive/SP_SauLuoi_2024-2025/Sp_GHKII/data/GHKII/1-TN-DS-TLN-TL-THPT-AnLuongDong-Hue-GHKII-NH24-25.tex #171
%[Dự án đề kiểm tra NH24-25 -đợt 9 - Nhật Thiện]%[THPT An Lương Đông - Huế]
	\immini{Một tấm ván hình chữ nhật $ABCD$ được dùng làm mặt phẳng nghiêng để kéo một vật lên khỏi hố sâu $4$ m (như hình vẽ). Cho biết $AB=2$ m, $AD=5$ m. Tính góc giữa đường thẳng $AC$ và đáy hố.}{
	\begin{tikzpicture}[scale=1, font=\footnotesize, line join=round, line cap=round, >=stealth]
		\path 
		(0,0) coordinate (A)
		(4,0) coordinate (B)
		(2,0.5) coordinate (K)
		($(B)+(K)-(A)$) coordinate (H)
		(K)+(0,2) coordinate (D)
		(H)+(0,2) coordinate (C)
		;
		\draw (A)--node[midway,sloped,above]{$5\text{ m}$}(D)--(C)--(B)--node[midway,sloped,below]{$2\text{ m}$}cycle (A)--(C) (C)--node[midway,sloped,above]{$4\text{ m}$}(H)--(B);
		\draw[dashed] (A)--(K)--(H) (K)--(D);
		\foreach \p/\r in {A/-90,B/-90,H/-60,C/60,D/90,K/-90}
		\fill (\p) circle (1pt) node[shift={(\r:3mm)}]{$\p$};
		\foreach \x/\dinh/\y in {C/H/B} \draw[fill = gray!50] ($(\dinh)!3pt!(\x)$)--($(\dinh)!3pt!(\x)+(\dinh)!3pt!(\y)-(\dinh)$)--($(\dinh)!3pt!(\y)$)--(\dinh)--cycle; 
	\end{tikzpicture}
}
	\loigiai{
	\immini{Do $CH\perp (ABKH)$ suy ra $AH$ là hình chiếu của $AC$ lên $(ABKH)$.\\
	Nên góc giữa $AC$ và đáy hố là góc $\widehat{CAH}$.\\
	Ta có $AC=\sqrt{AD^2+DC^2}=\sqrt{29}$.\\
	Xét tam giác $ACH$ vuông tại $H$.\\
	Ta có $\sin \widehat{CAH}=\dfrac{CH}{AC}=\dfrac{4}{\sqrt{29}}$.\\
	Suy ra $\widehat{CAH}\approx 47,97^\circ$.}{
\begin{tikzpicture}[scale=1, font=\footnotesize, line join=round, line cap=round, >=stealth]
	\path 
	(0,0) coordinate (A)
	(4,0) coordinate (B)
	(2,0.5) coordinate (K)
	($(B)+(K)-(A)$) coordinate (H)
	(K)+(0,2) coordinate (D)
	(H)+(0,2) coordinate (C)
	;
	\draw (A)--node[midway,sloped,above]{$5\text{ m}$}(D)--(C)--(B)--node[midway,sloped,below]{$2\text{ m}$}cycle (A)--(C) (C)--node[midway,sloped,above]{$4\text{ m}$}(H)--(B);
	\draw[dashed] (A)--(K)--(H)--cycle (K)--(D);
	\foreach \p/\r in {A/-90,B/-90,H/-60,C/60,D/90,K/45}
	\fill (\p) circle (1pt) node[shift={(\r:3mm)}]{$\p$};
	\foreach \x/\dinh/\y in {C/H/B} \draw[fill = gray!50] ($(\dinh)!3pt!(\x)$)--($(\dinh)!3pt!(\x)+(\dinh)!3pt!(\y)-(\dinh)$)--($(\dinh)!3pt!(\y)$)--(\dinh)--cycle; 
\end{tikzpicture}
}
}
\end{bt}

\begin{bt}%[1H8H6-2]
% nguon: Archive/SP_SauLuoi_2024-2025/Sp_GHKII/data/GHKII/1-TN-DS-TLN-THPT-NgocLac-ThanhHoa-GHKII-NH24-25.tex #179
%[KNTT - Lớp 11 - Ôn tập giữa học kì 2 -THPT Ngọc Lạc - Thanh Hoá]%[TheHung Nguyễn]
	Cho hình chóp $S.ABCD$ có đáy $ABCD$ là hình vuông cạnh $a$, $SA \perp (ABCD)$ và $SA = \dfrac{a\sqrt{2}}{2}$. Tính số đo của góc nhị diện $[S, BD, C]$ (đơn vị: độ).
	\loigiai{
		\immini{Gọi $O$ là giao điểm của $AC$ và $BD$, khi đó $CO \perp BD$, $SO \perp B D$.\\
			Do đó, góc phẳng nhị diện $[S, B D, C]=\widehat {SOC}$.\\
			Xét $\triangle SAO$, có $A O=\dfrac{a \sqrt{2}}{2}=S A$ và góc $SAO$ là góc vuông nên tam giác $SAO$ là tam giác vuông cân tại $A$, suy ra $\widehat{SOA}=45^{\circ}$; $\widehat{SOC} =135^{\circ}$.\\			
			Vậy số đo của góc nhị diện $\left[S, B D, C\right]$ bằng $135^{\circ}$.}{\begin{tikzpicture}[scale=.75, font=\footnotesize, line join=round, line cap=round, >=stealth]
				\def\a{3}
				\def\b{2.5}
				\def\h{3}
				\def\g{120}
				\path
				(0,0) coordinate (A)
				(0:\a) coordinate (D)
				--++(-\g:\b) coordinate (C)
				--++(-180:\a)coordinate (B)
				;
				\coordinate (O) at ($(A)!1/2!(C)$); 
				\coordinate (S) at ($(A)+(0,\h)$); 
				\coordinate (H) at ($(S)!(A)!(D)$);
				\draw (B)--(C)--(D)--(S)--(B) (S)--(C) ;
				\draw[dashed] (O)--(S)--(A)--(D)--(B)--(A)--(C) ;
				\foreach \x/\g in {A/45,B/-90,C/-90,D/0,S/90,O/-90}  \fill (\x) circle (1pt)+(\g:.45)node {$\x$};
		\end{tikzpicture}}
			
		
	}
\end{bt}

\begin{bt}%[1H8V2-4]
% nguon: Archive/SP_SauLuoi_2024-2025/Sp_GHKII/data/GHKII/1-TN-DS-TLN-TL-THPT-HauNghia-LongAn-GHKII-NH24-25.tex #214
%[Dự án đề kiểm tra Toán 11 GHKII NH24-25 - Đợt 9 - Huỳnh Đức Vũ]%[THPT Hậu Nghĩa - Long An]
Cho hình chóp $S.ABCD$ có đáy $ABCD$ là hình vuông và $SA\perp(ABCD)$. Gọi $I$, $J$, $K$ lần lượt là trung điểm của $A B$, $BC$ và $SB$. Chứng minh
\begin{enumerate}
	\item $ CD\perp\left(SAD\right).$
\item  $ BD\perp\left(\text{IJ}K\right).$
\end{enumerate}
\loigiai{
	\immini{
\begin{enumerate}
	\item Ta có $$\heva{&
CD\perp AD\\&CD\subset (ABCD)\\&CD\perp SA}\Rightarrow SA\perp ABCD.$$
\item Ta có $IK$, $IJ$ lần lượt là đường trung bình của các tam giác $SAB$, $ABC$ nên $IK \parallel SA$, $I J \parallel AC$.\\
Suy ra $(I J K) \parallel  (SAC)$.\\
Khi đó $\heva{&(IJK) \parallel (SAC)\\& B D\perp(S A C)}\Rightarrow B D\perp(IJK).$
\end{enumerate}}{\begin{tikzpicture}[scale=1, font=\footnotesize, line join=round, line cap=round, >=stealth]
\def\bc{4} % cạnh BC
\def\ba{2} % cạnh BA
\def\h{4} % đường cao
\def\gocB{30} % góc B của đáy
\coordinate[label=below left:$B$] (B) at (0,0);
\coordinate[label=above right:$A$] (A) at (\gocB:\ba);
\coordinate[label=below:$C$] (C) at (\bc,0);
\coordinate[label=right:$D$] (D) at ($(C)-(B)+(A)$);
\coordinate[label=above:$S$] (S) at ($(A)+(90:\h)$);
\coordinate[label=above left:$I$] (I) at ($(A)!0.5!(B)$);
\coordinate[label=below:$J$] (J) at ($(C)!0.5!(B)$);
\coordinate[label=above left:$K$] (K) at ($(S)!0.5!(B)$);
\draw (B)--(C)--(D)--(S)--cycle (S)--(C) (K)--(J);
\draw[dashed] (A)--(D) (S)--(A)--(B) (K)--(I)--(J) (A)--(C) (B)--(D)
;
\foreach \diem in {A,B,C,D,S,I,J,K}	\fill (\diem)circle(1.0pt);
\newcommand{\gocv}[4][black]{\draw[#1] ($(#3)!5pt!(#2)$)--($(#3)!2!($($(#3)!5pt!(#2)$)!.5!($(#3)!5pt!(#4)$)$)$)--($(#3)!5pt!(#4)$);}
\gocv{S}{A}{D}
\end{tikzpicture}
}
}
\end{bt}

\begin{bt}%[1H8H2-2]
% nguon: Archive/SP_SauLuoi_2024-2025/Sp_GHKII/data/GHKII/1-TN-DS-TLN-TL-THPT-NinhBinh-BacLieu-GHKII-NH24-25.tex #230
%[Dự án đề kiểm tra Toán 11 GHKII NH24-25-Lê Minh Thiện Anh]%[THPT Ninh Bình-Bạc Liêu]
	Cho hình chóp $S.ABCD$ có đáy $ABCD$ là hình vuông, cạnh bên $SA$ vuông góc với đáy. Chứng minh $BD \perp (SAC)$.	
	\loigiai{
		\immini{Ta có $BD \perp AC$ ($ABCD$ là hình vuông).\\
		$BD \perp SA$ ($SA \perp (ABCD)$) nên $BD \perp (SAC)$.}
		{\begin{tikzpicture}[font=\footnotesize, line join=round, line cap=round, >=stealth,scale=0.7]
		\path
		(0,0) coordinate (A) 
		(-2,-2) coordinate (B)
		(3,-2) coordinate (C)
		(5,0) coordinate (D)
		(0,4) coordinate (S)
		;
		\draw 
		(S)--(B)--(C)--(D)--cycle
		(S)--(C);
		\draw [dashed] (S)--(A)--(D)--(B)--(A)--(C); 
		\pic[draw,angle radius=1.5mm,angle eccentricity=1.5] {right angle = S--A--B};
		\pic[draw,angle radius=1.5mm,angle eccentricity=1.5] {right angle = S--A--D};		
		\foreach \x/\g in
		{A/150,B/180,C/0,D/0,S/90}
		\fill[black](\x) circle (1pt)
		($(\x)+(\g:3mm)$) node{$\x$};
		\end{tikzpicture}}
	}
\end{bt}

\begin{bt}%[1H8V4-7]
% nguon: Archive/SP_SauLuoi_2024-2025/Sp_GHKII/data/GHKII/1-TN-DS-TLN-TL-THPT-PhanBoiChau-BinhThuan-GHKII-NH24-25.tex #239
%[GK2-K11-THPT Phan Bội Châu-Bình Thuận NH2425]%[Hieu Phan]
Kim tự tháp Kheops tại Ai Cập được xây dựng vào khoảng thời gian từ năm $2580-2560$ trước công nguyên có dạng hình chóp tứ giác đều với chiều cao khoảng $149$ m và cạnh đáy dài khoảng $230$ m. Tính tan góc giữa mặt bên và mặt đáy của kim tự tháp.
\loigiai
{ \begin{center}
		\begin{tikzpicture}[scale=0.8, font=\footnotesize,line join=round, line cap=round, >=stealth]
			\path 
			(0,0) coordinate (A)
			++(-140:2) coordinate (B)
			++(0:3.5) coordinate (C)
			($(A)+(C)-(B)$) coordinate (D)
			($(A)!1/2!(C)$) coordinate (O)
			($(O)+(0,3.5)$) coordinate (S)
			($(C)!1/2!(D)$) coordinate (M)
			;
			\foreach \i in{B,C,D}{\draw (S)--(\i);};
			\draw (B)--(C)--(D);
			\draw[dashed] (S)--(A)--(B) (C)--(A)--(D)--(B) (S)--(O)--(M)--cycle;
			\pic[draw,angle eccentricity=1.8,angle radius=2mm]{right angle=S--O--D};
			\foreach \i/\g in {A/-90,B/-90,C/-90,D/-90,O/-90,S/90,M/-90}
			\fill[black] (\i) circle(1pt)+(\g:3mm)node[scale=1]{$\i$};
		\end{tikzpicture}
	\end{center}
	Gọi $O$ là tâm của hình vuông $ABCD$; $ M$ là trung điểm của $ CD$.\\
Ta có $\left\{\begin{aligned}
&\left(SCD\right)\cap\left(ABCD\right)=CD\\ 
& OM\perp CD,\text{OM}\subset\left(ABCD\right)\\ 
& SM\perp CD,\text{SM}\subset\left(SCD\right).\\ 
\end{aligned}\right.$\\
Góc giữa $\left(SCD\right)$ và $\left(ABCD\right)$ là góc $\widehat{SMO}$.\\
Ta có $ OM=\dfrac{1}{2}AD=\dfrac{230}{2}=115$ m.\\
Xét tam giác $ SOM$ vuông tại $ O$, ta có $\tan\widehat{SMO}=\dfrac{SO}{OM}$$=\dfrac{149}{115}$.
}
\end{bt}

\begin{bt}%[1H8H2-2]
% nguon: Archive/SP_SauLuoi_2024-2025/Sp_GHKII/data/GHKII/1-TN-DS-TLN-TL-THPT-LeLoi-KonTum-GHKII-NH24-25.tex #251
%[Dự án A đợt 9, THPT Lê Lợi Kon Tum GHK2 2024-2025]%[BCTuan]
Cho hình chóp $S.ABCD$ có đáy $ABCD$ là hình chữ nhật, $SA\perp(ABCD)$.
Chứng minh $BC\perp (SAB)$.
\loigiai{
\immini
{
Ta có $\heva{&BC\perp AB \\& BC \perp SA\ (\text{do}\ SA\perp (ABCD))} \Rightarrow BC\perp(SAB)$.
}
{
	\begin{tikzpicture}[scale=0.8, font=\footnotesize, line join=round, line cap=round, >=stealth]
		\def\bc{4} % cạnh BC
		\def\ba{2} % cạnh BA
		\def\h{3} % đường cao
		\def\gocB{30} % góc B của đáy
		\coordinate[label=below left:$B$] (B) at (0,0);
		\coordinate[label=above left:$A$] (A) at (\gocB:\ba);
		\coordinate[label=below:$C$] (C) at (\bc,0);
		\coordinate[label=right:$D$] (D) at ($(C)-(B)+(A)$);
		\coordinate[label=above:$S$] (S) at ($(A)+(90:\h)$);
		% \coordinate[label=above right:$H$] (H) at ($(S)!0.5!(D)$);
		\draw (B)--(C)--(D)--(S)--cycle (S)--(C);
		\draw[dashed] (A)--(D) (S)--(A)--(B);
		\foreach \diem in {A,B,C,D,S}	\fill (\diem)circle(1pt);
	\end{tikzpicture}
}
}
\end{bt}

\begin{bt}%[1H8V1-3]
% nguon: Archive/SP_SauLuoi_2024-2025/Sp_GHKII/data/GHKII/1-TN-DS-TLN-TL-THPT-LeLoi-KonTum-GHKII-NH24-25.tex #252
%[Dự án A đợt 9, THPT Lê Lợi Kon Tum GHK2 2024-2025]%[BCTuan]
Cho hình tứ diện $S.ABC$ có tất cả các cạnh bằng nhau và bằng $a$. Điểm $K$ thuộc cạnh $SB$ sao cho $SK=\dfrac{2}{3}SB$. Tính côsin của góc giữa hai đường thẳng $AB$ và $CK$.
\loigiai{
\begin{center}
	\begin{tikzpicture}[scale=1, font=\footnotesize, line join=round, line cap=round, >=stealth]
		\def\ac{4} % cạnh AC
		\def\ab{2} % cạnh AB
		\def\h{4} % chiều cao
		\def\gocA{50} % góc A của đáy
		\coordinate[label=left:$A$] (A) at (0,0);
		\coordinate[label=right:$C$] (C) at (\ac,0);
		\coordinate[label=below left:$B$] (B) at (-\gocA:\ab);
		\coordinate[label=above:$S$] (S) at ($(A)+(70:\h)$);
		\coordinate[label=above right:$K$] (K) at ($(S)!2/3!(B)$);
		\coordinate[label=above left:$I$] (I) at ($(S)!2/3!(A)$);
		% \coordinate[label=above right:$H$] (H) at ($(S)!0.5!(B)$);
		\draw (A)--(B)--(C)--(S)--(A) (S)--(B) (I)--(K)--(C);
		\draw[dashed] (A)--(C) (I)--(C);
		\foreach \diem in {A,B,C,S,I,K}	\fill (\diem)circle(1pt);
	\end{tikzpicture}	
\end{center}
Trong $(SAB)$ gọi $I\in SA$ sao cho $KI\parallel AB \Rightarrow \dfrac{SI}{SA}=\dfrac{2}{3}$. Do đó $(AB,KC)=(IK,KC)$.\\
Xét tam giác $BKC$, áp dụng định lí côsin có:
\begin{eqnarray*}
KC^2&=&KB^2+BC^2-2\cdot KB\cdot BC\cdot\cos\widehat{KBC}\\
&=&\left(\dfrac{a}{3}\right)^2+a^2-2\cdot \dfrac{a}{3}\cdot a\cdot\cos 60^{\circ}\\
&=&\dfrac{7a^2}{9} \Rightarrow KC=\dfrac{a\sqrt{7}}{3}.	
\end{eqnarray*}
Tương tự tính được $IC=\dfrac{a\sqrt{7}}{3}$; $IK=\dfrac{2a}{3}$.\\
Áp dụng định lí côsin cho tam giác $IKC$ có 
$$\cos\widehat{IKC}=\dfrac{IK^2+KC^2-IC^2}{2\cdot IK\cdot KC}=\dfrac{\left(\dfrac{a\sqrt{7}}{3}\right)^2+\left(\dfrac{2a}{3}\right)^2-\left(\dfrac{a\sqrt{7}}{3}\right)^2}{2\cdot\dfrac{a\sqrt{7}}{3} \cdot\dfrac{2a}{3}}=\dfrac{\sqrt{7}}{7}.$$
Vậy $\cos(AB,KC)=\dfrac{\sqrt{7}}{7}$.
}
\end{bt}

\begin{bt}%[1H8H2-2]
% nguon: Archive/SP_SauLuoi_2024-2025/Sp_GHKII/data/GHKII/1-TN-DS-TLN-TL-THPT-Edison-HaiPhong-GKII-NH24-25.tex #260
%[Dự án đề kiểm tra Toán 10 GHKII NH24-25- Quan Ón]%[THPT Edison - Hải Phòng]
	Cho tứ diện $OABC$ có $OA$ vuông góc với mặt phẳng $(OBC)$ và có $A'$, $B'$, $C'$ lần lượt là trung điểm của $OA$, $AB$, $AC$. Vẽ $OH$ là đường cao của tam giác $OBC$. Chứng minh rằng $B'C' \perp (OAH)$.
	\loigiai{
		\immini{
			Ta có $\heva{&OA \perp (OBC)\\&BC \subset (OBC)} \Rightarrow OA \perp BC$.\\
			Vì $\heva{&BC\perp OA \textrm{ (cmt)}\\&BC\perp OH \textrm{ ($OH$ là đường cao của $\triangle OBC$)}\\&OA \cap OH = \{O\}}$\\
			$\Rightarrow BC \perp (OAH)$.\\
			Mà tam giác $ABC$ có $B'C'$ là đường trung bình nên $B'C' \parallel BC$.\\
			Ta có $\heva{&BC \perp (OAH) \textrm{ (cmt)}\\&B'C' \parallel BC \textrm{ (cmt)}}$\\
			$\Rightarrow B'C' \perp (OAH)$.
		}{
			\begin{tikzpicture}[font=\footnotesize,line join=round, line cap=round, >=stealth,scale=1] 
				\path  
				(0,0) coordinate (O)
				(1,-1.6) coordinate (B)
				(0,2.8) coordinate (A)
				(3.5,0) coordinate (C)
				($(O)!0.5!(A)$) coordinate (A')
				($(A)!0.5!(B)$) coordinate (B')
				($(A)!0.5!(C)$) coordinate (C')
				($(B)!0.35!(C)$) coordinate (H)
				; 
				\draw (A)--(O)--(B)--(C)--(A)--(B) (A')--(B')--(C') (A)--(H);
				\draw[dashed] (A')--(C') (C)--(O)--(H);
				\foreach \x/\g in {A/90, B/-90, C/-90, O/-135,H/-45,A'/180,B'/-20,C'/45} \fill (\x) circle(1pt) node[{shift=(\g:0.3)}]{$\x$}; 
				\foreach \a/\b/\c in {O/H/B} {\draw pic [draw, angle radius=2mm] {right angle=\a--\b--\c};}
			\end{tikzpicture}
		}
	}
\end{bt}

\begin{bt}%[1H8N2-2]
% nguon: Archive/SP_SauLuoi_2024-2025/Sp_GHKII/data/GHKII/1-TN-DS-TLN-TL-THPT-NguyenHue-Hue-GHKII-NH24-25.tex #268
%[Dự án đề kiểm tra Toán 11 Giữa HKII NH24-25- Xuan Vy Pham%[THPT Nguyễn Huệ - Huế]
Cho lăng trụ đứng tam giác $ABC.A'B'C'$, biết $CA=CB$. Gọi $I$ là trung điểm của $AB$. Chứng minh rằng $CI \perp (ABB'A')$.	
\loigiai{\begin{center}
		\begin{tikzpicture}[line join=round, line cap=round,>=stealth,font=\footnotesize,scale=1]
			\def\a{3.4} 
			\def\b{2}
			\def\h{3.1}
			\path 	(0:0) coordinate (A)
			++(0:\a) coordinate (C)
			(A)	++(-50:\b) coordinate (B)
			($(A)+(90:\h)$) coordinate (A')
			($(B)+(90:\h)$) coordinate (B')
			($(C)+(90:\h)$) coordinate (C')
			($(A)!0.5!(B)$) coordinate (I);
			\draw[dashed,thick] 	(A)--(C)--(I);
			\draw[thick] (C)--(C')--(B')--(A') (B)--(B') 	(B)--(C) (C')--(A')--(A)--(B) ;
			\foreach \x/\g in {A/180,B/-60,C/0,A'/180,B'/-20,C'/0,I/-90}
			\fill[black] 	(\x) circle (1pt)
			($(\g:4mm)+(\x)$) node {$\x$};	
		\end{tikzpicture}
\end{center}
Xét tam giác $ACB$ có $CA=CB$ nên tam giác $CAB$ cân tại $C$.\\
Mà $CI$ là đường trung tuyến nên cũng là đường cao. Do đó $CI \perp AB$.\\
Ta có $\heva{&CI \perp AA' \; \text{(vì $AA' \perp (ACB)$)}\\&CI \perp AB} \Rightarrow CI \perp (ABB'A')$.

}
\end{bt}

\begin{bt}%[1H8V6-2]
% nguon: Archive/SP_SauLuoi_2024-2025/Sp_GHKII/data/GHKII/1-TN-DS-TLN-TL-THPT-VinhLinh-QuangTri-GHKII-NH24-25.tex #292
%[Lớp 11 - Ôn tập giữa học kì 2 - THPT Vĩnh Linh-Quảng Trị]%[Trần Văn Hùng]
	Cho hình chóp $S.ABCD$ có đáy $ABCD$ là hình vuông cạnh $a$, $SA$ vuông góc với mặt phẳng $(ABCD)$, $SA = a\sqrt{2}$.
	\begin{enumEX}{1}
	\item Chứng minh $BD \perp (SAC)$.
	\item Gọi $M$, $N$ lần lượt là hình chiếu vuông góc của điểm $A$ lên các cạnh $SB$, $SD$. Chứng minh $MN \perp SC$.
	\item Tính số đo góc nhị diện $[S, BD, C]$.
	\end{enumEX}
	\loigiai{
	\begin{center}
			\begin{tikzpicture}[scale=1, font=\footnotesize, line join=round, line cap=round, >=stealth]
			\coordinate (A) at (0,0);
			\coordinate (B) at (-2,-2);
			\coordinate (D) at (4,0);
			\coordinate (C) at ($(B)+(D)-(A)$);
			\coordinate (O) at ($(A)!0.5!(C)$);
			\coordinate (S) at ($(A)+(0,5)$);
			\coordinate (M) at ($(S)!2/3!(B)$);
			\coordinate (N) at ($(S)!2/3!(D)$);
			\draw(S)--(B) (S)--(C)--(D)--(S) (B)--(C);
			\draw[dashed,thin](A)--(B) (A)--(C)--(D) (S)--(A)--(D)--(B) (A)--(M)--(N)--(A) (S)--(O);
			\pic[draw,thin,angle radius=2mm] {right angle = S--A--D} 
			pic[draw,thin,angle radius=2mm] {right angle = S--A--B}
			pic[draw,thin,angle radius=2mm] {right angle = A--O--B}
			pic[draw,thin,angle radius=2mm] {right angle = A--M--B}
			pic[draw,thin,angle radius=2mm] {right angle = A--N--D}
			pic[draw,thin,angle radius=2mm] {right angle = A--B--C}
			pic[draw,thin,angle radius=2mm] {right angle = B--A--D};
			\foreach \i/\g in {S/90,A/180,B/-90,C/-90,D/0,M/150,N/30,O/-90}{\draw[fill=black](\i) circle (1pt) ($(\i)+(\g:3mm)$) node[scale=1]{$\i$};}
		\end{tikzpicture}
		\end{center}
		\begin{enumEX}{1}
		\item Vì $\heva{&BD\perp AC\\&BD\perp SA\, \left(\text{do } SA\perp (ABCD)\right)}$\\
		 nên $BD\perp (SAC)$.
		\item Tam giác $SAB$ vuông tại $A$, $AM$ là đường cao nên 
		$$SM\cdot SB=SA^2\Rightarrow \dfrac{SM}{SB}=\dfrac{SA^2}{SB^2}=\dfrac{2a^2}{2a^2+a^2}=\dfrac{2}{3}\quad (1).$$
		Tam giác $SAD$ vuông tại $A$, $AN$ là đường cao nên 
		$$SN\cdot SD=SA^2\Rightarrow \dfrac{SN}{SD}=\dfrac{SA^2}{SD^2}=\dfrac{2a^2}{2a^2+a^2}=\dfrac{2}{3}\quad (2).$$
		Từ $(1)$ và $(2)$ suy ra $\dfrac{SM}{SB}=\dfrac{SN}{SD}=\dfrac{2}{3}\Rightarrow MN\parallel BD$.\\
		Mặt khác $BD\perp (SAC)$ nên $MN\perp (SAC)\Rightarrow MN\perp SC$.
		\item Gọi $O$ là tâm hình vuông $ABCD$.\\
		Vì $\heva{&SO\perp BD\\&OC\perp BD}$ nên $[S,BD,C]=\widehat{SOC}=180^\circ-\widehat{SOA}$.\\
		Do $ABCD$ là hình vuông cạnh $a$ nên $AC=a\sqrt{2}\Rightarrow AO=\dfrac{a\sqrt{2}}{2}$.\\
		Tam giác $SAO$ vuông tại $A$, suy ra 
		$$ \tan\widehat{SAO}=\dfrac{SA}{AO}=\dfrac{a\sqrt{2}}{\dfrac{a\sqrt{2}}{2}}=2\Rightarrow \widehat{SAO}\approx 63{,}4^\circ\Rightarrow \widehat{SOC}\approx 116{,}6^\circ.$$
		Vậy số đo góc nhị diện $[S,BD,C]$ là $116{,}6^\circ$.
	\end{enumEX}
	}
\end{bt}

\begin{bt}%[1H8V6-2]
% nguon: Archive/SP_SauLuoi_2024-2025/Sp_GHKII/data/GHKII/1-TN-DS-TLN-TL-THPT-VinhLinh-QuangTri-GHKII-NH24-25.tex #293
%[Lớp 11 - Ôn tập giữa học kì 2 - THPT Vĩnh Linh-Quảng Trị]%[Trần Văn Hùng]
	Để chuẩn bị cho đợt Hội trại 26/3 sắp tới, bạn Nam lớp 11A đã thiết kế trại theo hình vẽ sau
	\begin{center}
	\begin{tikzpicture}[scale=1.2,font=\footnotesize, line join=round, line cap=round, >=stealth]
		\coordinate (A) at (0,0);\coordinate (B) at (1,1);\coordinate (C) at (4,1);	\coordinate (D) at ($(A)+(C)-(B)$);
		\coordinate (A') at ($(A)-(0,2)$);
		\coordinate (B') at ($(A')+(B)-(A)$);
		\coordinate (C') at ($(B')+(C)-(B)$);
		\coordinate (D') at ($(A')+(D)-(A)$);
		\coordinate (N) at ($(A)!0.5!(D)$);\coordinate (E) at ($(A')!0.5!(D')$);
		\coordinate (O) at ($(C)!0.5!(A)$);\coordinate (M) at ($(C)!0.5!(D)$);
		\coordinate (S) at ($(O)+(0,1)$);
		\draw (D)--(A)--(B) (C)--(D)--(D')--(C')--(C) (A)--(A')--(D') (A)--(S)--(D) (C)--(S)--(B) (E)--(N)--(S)--(M);
		\draw[dashed] (A')--(B')--(C') (B')--(B)--(C)--(A) (B)--(D) (S)--(O)--(M) (O)--(N);
		\foreach \x/\g in {A/180,B/90,C/0,D/0,A'/-90,B'/120,C'/0,D'/-90,S/90,O/-90,M/0,N/-60,E/-90} \fill[black](\x) circle (1pt) ($(\x)+(\g:3mm)$) node{$\x$};
	\end{tikzpicture}
	\end{center}
	Cho biết $SA=SB=SC=SD$, $ABCD.A'B'C'D'$ là hình hộp chữ nhật có kích thước $AB=2{,}6$, $AD = 4{,}8$ và $AA'=2{,}3$, đồng thời góc giữa hai mặt phẳng $(SCD)$ và $(ABCD)$ bằng $30^\circ$. Tính số đo góc nhị diện $[S,AD,A']$.
	\loigiai{
	Gọi $O=AC\cap BD$ và $M$, $N$ lần lượt là trung điểm của $CD$, $AD \Rightarrow SO\perp (ABCD)$, $OM\perp CD$, $ON\perp AD$.\\
	Vì $\heva{&CD\perp OM\\&CD\perp SO}$ nên $CD\perp (SOM)\Rightarrow $ góc giữa hai mặt phẳng $(SCD)$ và $(ABCD)$ bằng $\widehat{SOM}$.\\
	Theo giả thiết suy ra $\widehat{SMO}=30^\circ$.\\
	Tam giác $SOM$ vuông tại $O$, ta có $SO=OM\cdot \tan 30^\circ=2{,}4\cdot \dfrac{1}{\sqrt{3}}=\dfrac{4\sqrt{3}}{5}$.\\
	Gọi $E$ là trung điểm $A'D'$, suy ra $AD\perp NE$.\\
	Vì $\heva{&AD\perp SN\\&AD\perp NE}$ nên $[S,AD,A']=\widehat{SNE}=90^\circ+\widehat{SNO}$.\\
	Tam giác $SNO$ vuông tại $O$, ta có $\tan\widehat{SNO}=\dfrac{SO}{ON}=\dfrac{\dfrac{4\sqrt{3}}{5}}{1{,3}}=\dfrac{8\sqrt{3}}{13}\Rightarrow \widehat{SNO}\approx 46{,}8^\circ$.\\
	Vậy số đo góc nhị diện $[S,AD,A']$ là $136{,}8^\circ$.
	}
\end{bt}

\begin{bt}%[1H8H6-1]
% nguon: Archive/SP_SauLuoi_2024-2025/Sp_GHKII/data/GHKII/1-TL-TN-TLN-THPT-TranCaoVan-KhanhHoa-GHKII-NH24-25.tex #299
%[Dự án đề kiểm tra Toán 11 GHKII NH24-25- TranXuanHoa]%[THPT Trần Cao Vân - Khánh Hòa]
	\immini{Cho hình chóp $S.ABCD$ có đáy $ABCD$ là hình chữ nhật, với $SA=AB=a$, $AD=a\sqrt{2}$, $SA\perp (ABCD)$. Gọi $H$ là hình chiếu vuông góc của $A$ trên $SD$.
	\begin{enumerate}
		\item Chứng minh $AH\perp (SCD)$.
		\item Tính góc giữa $SC$ và mặt phẳng $(ABCD)$.
	\end{enumerate}}
	{\begin{tikzpicture}[scale=0.7, font=\footnotesize,line join=round, line cap=round, >=stealth]
			\path 
			(0,0) coordinate (A)
			++(-140:2) coordinate (B)
			++(0:3.5) coordinate (C)
			($(A)+(C)-(B)$) coordinate (D)
			($(A)+(0,3)$) coordinate (S)
			($(S)!1/4!(D)$) coordinate (H)
			;
			\foreach \i in{B,C,D}{\draw (S)--(\i);};
			\draw (B)--(C)--(D);
			\draw[dashed] (S)--(A)--(B) (A)--(D) (A)--(H);
			\pic[draw,angle eccentricity=1.8,angle radius=2mm]{right angle=S--A--D};
			\pic[draw,angle eccentricity=1.8,angle radius=2mm]{right angle=S--H--A};
			\foreach \i/\g in {A/-90,B/-90,C/-90,D/-90,S/90,H/60}
			\fill[black] (\i) circle(1pt)+(\g:3mm)node[scale=1]{$\i$};
	\end{tikzpicture}}
	\loigiai{
	\begin{enumerate}[a)]
		\item Ta có $\heva{&CD\perp AD\subset (SAD)\\ &CD\perp SA\subset (SAD)}$ nên $CD\perp (SAD)$. Do đó $CD\perp AH$.\\
		Ta có $\heva{&AH\perp CD\subset (SCD)\\& AH\perp SD\subset (SCD)}$ nên $AH\perp (SCD)$.
		\item Ta có $SA\perp (ABCD)$ nên góc giữa $SC$ và mặt phẳng đáy là góc $SCA$.\\
		Xét tam giác $SAC$ vuông tại $A$. Ta có $AC=a\sqrt{3}$ và $\tan \widehat{SCA}=\dfrac{SA}{CA}=\dfrac{1}{\sqrt{3}}$.\\
		Do đó $\widehat{SCA}=30^\circ$.
	\end{enumerate}
	}
\end{bt}

\begin{bt}%[1H8H1-3]
% nguon: Archive/2024/TOAN_11_XE/TOAN_11_XE/data4/11C7/KNTT-7.8.tex #1369
%[Tex hóa SGK11, Kiều Ngân]
	Cho hình chóp $S.ABC$ có $SA\perp (ABC)$, $AC\perp BC$, $SA=BC=a\sqrt{3}$, $AC=a$.
	\immini{
		\begin{enumerate}
			\item Tính góc giữa hai đường thẳng $SA$ và $BC$.
			\item Tính góc giữa đường thẳng $SC$ và mặt phẳng $(ABC)$.
			\item Tính số đo của góc nhị diện $[B,SA,C]$.
			\item Tính khoảng cách từ $B$ đến mặt phẳng $(SAC)$.
			\item Tính khoảng cách giữa hai đường thẳng $SA$ và $BC$.
			\item Tính thể tích của khối chóp $S.ABC$.
		\end{enumerate}
	}{
		\begin{tikzpicture}[scale=1, font=\footnotesize, line join=round, line cap=round, >=stealth]
			\path
			(0,0) coordinate (A)
			(1,-1) coordinate (C)
			(3,0) coordinate (B)
			($(O)+(0,2.5)$) coordinate (S)
			;
			\draw (S)--(A)--(C)--(B)--(S)--(C);
			\draw[dashed] (A)--(B);
			\foreach \x/\g in {S/90,A/180,C/-90,B/0} \draw[fill=black] (\x) circle(1pt)++(\g:0.3)node{$\x$};
			\path ($(S)!0.5!(A)$) node[left]{$a\sqrt{3}$};
			\path ($(A)!0.5!(C)$) node[below left]{$a$};
			\path ($(B)!0.5!(C)$) node[below right]{$a\sqrt{3}$};
			\draw pic[draw=black, angle eccentricity=2, angle radius=0.25cm]{right angle=B--A--S};
			\draw pic[draw=black, angle eccentricity=2, angle radius=0.25cm]{right angle=C--A--S};
			\draw pic[draw=black, angle eccentricity=2, angle radius=0.25cm]{right angle=B--C--A};
		\end{tikzpicture}
	}
	\loigiai{
		\begin{center}
			\begin{tikzpicture}[scale=1, font=\footnotesize, line join=round, line cap=round, >=stealth]
				\path
				(0,0) coordinate (A)
				(1,-1) coordinate (C)
				(3,0) coordinate (B)
				($(O)+(0,2.5)$) coordinate (S)
				;
				\draw (S)--(A)--(C)--(B)--(S)--(C);
				\draw[dashed] (A)--(B);
				\foreach \x/\g in {S/90,A/180,C/-90,B/0} \draw[fill=black] (\x) circle(1pt)++(\g:0.3)node{$\x$};
				\path ($(S)!0.5!(A)$) node[left]{$a\sqrt{3}$};
				\path ($(A)!0.5!(C)$) node[below left]{$a$};
				\path ($(B)!0.5!(C)$) node[below right]{$a\sqrt{3}$};
				\draw pic[draw=black, angle eccentricity=2, angle radius=0.25cm]{right angle=B--A--S};
				\draw pic[draw=black, angle eccentricity=2, angle radius=0.25cm]{right angle=C--A--S};
				\draw pic[draw=black, angle eccentricity=2, angle radius=0.25cm]{right angle=B--C--A};
			\end{tikzpicture}
		\end{center}
		\begin{enumerate}
			\item Tính góc giữa hai đường thẳng $SA$ và $BC$.\\
			Ta có $SA\perp (ABC)\Rightarrow SA\perp BC$.\\
			Do đó góc giữa hai đường thẳng $SA$ và $BC$ là $90^\circ$.
			\item Tính góc giữa đường thẳng $SC$ và mặt phẳng $(ABC)$.\\
			Ta có $SA\perp (ABC)$, suy ra $AC$ là hình chiếu vuông góc của $SC$ lên mặt phẳng $(ABC)$.\\
			Do đó góc giữa đường thẳng $SC$ và mặt phẳng $(ABC)$ là $(SC,AC)=\widehat{SCA}$.\\
			Ta có $\tan \widehat{SCA}=\dfrac{a\sqrt{3}}{a}=\sqrt{3}$.\\
			Vậy $\widehat{SCA}=60^\circ$.
			\item Tính số đo của góc nhị diện $[B,SA,C]$.\\
			Ta có $\heva{&CA\perp SA\\&BA\perp SA}$, suy ra góc nhị diện $[B,SA,C]$ là góc $\widehat{BAC}$.\\
			Ta có $\tan \widehat{BAC}=\dfrac{BC}{AC}=\dfrac{a\sqrt{3}}{a}=\sqrt{3}\Rightarrow \widehat{BAC}=60^\circ$.
			\item Tính khoảng cách từ $B$ đến mặt phẳng $(SAC)$.\\
			Ta có $\heva{&BC\perp SA\\&BC\perp AC}\Rightarrow BC\perp (SAC)$, suy ra $\mathrm{d}(B,(SAC))=BC=a\sqrt{3}$.
			\item Tính khoảng cách giữa hai đường thẳng $SA$ và $BC$.\\
			Ta có $\heva{&SA\perp AC\\&BC\perp AC}$, suy ra $AC$ là đoạn vuông góc chung của $SA$ và $BC$.\\
			Vậy $\mathrm{d}(SA,BC)=AC=a$.
			\item Tính thể tích của khối chóp $S.ABC$.\\
			Diện tích tam giác $ABC$ là $S_{\triangle {ABC}}=\dfrac{1}{2}\cdot AC\cdot BC=\dfrac{a^2\sqrt{3}}{2}$.\\
			Thể tích khối chóp $S.ABC$ là $V_{S.ABC}=\dfrac{1}{3}\cdot S_{\triangle {ABC}}\cdot SA=\dfrac{1}{3}\cdot \dfrac{a^2\sqrt{3}}{2}\cdot a\sqrt{3}=\dfrac{a^3}{2}$.
		\end{enumerate}
	}
\end{bt}

\begin{bt}%[1H8H1-3]
% nguon: Archive/2024/TOAN_11_XE/TOAN_11_XE/data4/11C7/KNTT-7.8.tex #1370
%[Tex hóa SGK11, Kiều Ngân]
	Cho hình lăng trụ tam giác đều $ABC.A'B'C'$ có tất cả các cạnh bằng $a$. Gọi $M$ là trung điểm của $AB$.
	\immini{
		\begin{enumerate}
			\item Tính góc giữa hai đường thẳng $AB$ và $B'C'$.
			\item Tính góc giữa đường thẳng $A'B$ và mặt phẳng $(ABC)$.
			\item Tính số đo của góc nhị diện $[B,CC',M]$.
			\item Chứng minh rằng $CC'\parallel (ABB'A')$. Tính khoảng cách giữa đường thẳng $CC'$ và mặt phẳng $(ABB'A')$.
			\item Chứng minh rằng $CM\perp (ABB'A')$. Tính khoảng cách giữa hai đường thẳng $CC'$ và $A'M$.
			\item Tính thể tích của khối lăng trụ tam giác đều $ABC.A'B'C'$ và thể tích khối chóp $A'.MBC$.
		\end{enumerate}
	}{
		\begin{tikzpicture}[scale=1, font=\footnotesize, line join=round, line cap=round, >=stealth]
			\path
			(0,0) coordinate (A)
			(1.5,-1) coordinate (B)
			(2.5,0) coordinate (C)
			($(A)+(0,3)$) coordinate (A')
			($(A')-(A)+(B)$) coordinate (B')
			($(A')-(A)+(C)$) coordinate (C')
			($(A)!0.5!(B)$) coordinate (M)
			;
			\draw (A')--(B')--(C')--(A') (A')--(A) (B)--(B') (C)--(C') (A)--(B)--(C) (A')--(M) (A')--(B);
			\draw[dashed] (A)--(C)--(M) (A')--(C);
			\foreach \x/\g in {A/180,B/-90,C/0,A'/180,B'/90,C'/0,M/-150} \draw[fill=black] (\x) circle(1pt)++(\g:0.3)node{$\x$};
			\path ($(A)!0.5!(A')$) node[left]{$a$};
			\path ($(A')!0.5!(B')$) node[above right]{$a$};
		\end{tikzpicture}
	}
	\loigiai{
		\begin{center}
			\begin{tikzpicture}[scale=1, font=\footnotesize, line join=round, line cap=round, >=stealth]
				\path
				(0,0) coordinate (A)
				(1.5,-1) coordinate (B)
				(2.5,0) coordinate (C)
				($(A)+(0,3)$) coordinate (A')
				($(A')-(A)+(B)$) coordinate (B')
				($(A')-(A)+(C)$) coordinate (C')
				($(A)!0.5!(B)$) coordinate (M)
				;
				\draw (A')--(B')--(C')--(A') (A')--(A) (B)--(B') (C)--(C') (A)--(B)--(C) (A')--(M) (A')--(B);
				\draw[dashed] (A)--(C)--(M) (A')--(C);
				\foreach \x/\g in {A/180,B/-90,C/0,A'/180,B'/90,C'/0,M/-150} \draw[fill=black] (\x) circle(1pt)++(\g:0.3)node{$\x$};
				\path ($(A)!0.5!(A')$) node[left]{$a$};
				\path ($(A')!0.5!(B')$) node[above right]{$a$};
				\draw pic[draw=black, angle eccentricity=2, angle radius=0.25cm]{right angle=B--M--C};
				\draw pic[draw=black, angle eccentricity=2, angle radius=0.25cm]{right angle=C'--C--B};
				\draw pic[draw=black, angle eccentricity=2, angle radius=0.25cm]{right angle=C'--C--M};
				\draw pic[draw=black, angle eccentricity=2, angle radius=0.5cm]{angle=M--C--B};
			\end{tikzpicture}
		\end{center}
		\begin{enumerate}
			\item Tính góc giữa hai đường thẳng $AB$ và $B'C'$.\\
			Ta có $B'C'\parallel BC$, suy ra $(AB,B'C')=(AB,BC)=\widehat{ABC}=60^\circ$ (vì tam giác $ABC$ đều).
			\item Tính góc giữa đường thẳng $A'B$ và mặt phẳng $(ABC)$.\\
			Ta có $AA'\perp (ABC)$, suy ra $AB$ là hình chiếu vuông góc của $A'B$ lên mặt phẳng $(ABC)$.\\
			Do đó góc giữa đường thẳng $A'B$ và mặt phẳng $(ABC)$ là $(A'B,AB)=\widehat{A'BA}=45^\circ$ (vì tam giác $A'AB$ vuông cân tại $A$).
			\item Tính số đo của góc nhị diện $[B,CC',M]$.\\
			Ta có $CC'\perp (ABC)\Rightarrow \heva{&CC'\perp CM\\&CC'\perp CB}$, suy ra góc nhị diện $[B,CC',M]$ là $\widehat{BCM}=30^\circ$ (vì $CM$ là phân giác trong tam giác $ABC$ đều).
			\item Chứng minh rằng $CC'\parallel (ABB'A')$. Tính khoảng cách giữa đường thẳng $CC'$ và mặt phẳng $(ABB'A')$.\\
			Ta có $\heva{&CC'\parallel BB'\\&BB'\subset (ABB'A')\\&CC'\not\subset (ABB'A')}\Rightarrow CC'\parallel (ABB'A')$.\\ Do đó $\mathrm{d}(CC',(ABB'A'))=\mathrm{d}(C,(ABB'A'))$.\\
			Lại có $\heva{&AA'\perp (ABC)\\&CM\subset (ABC)}\Rightarrow AA'\perp CM$.\\
			Khi đó $\heva{&CM\perp AB\\&CM\perp AA'}\Rightarrow CM\perp (ABB'A')$.\\
			Vậy $\mathrm{d}(CC',(ABB'A'))=\mathrm{d}(C,(ABB'A'))=CM=\dfrac{a\sqrt{3}}{2}$.
			\item Chứng minh rằng $CM\perp (ABB'A')$. Tính khoảng cách giữa hai đường thẳng $CC'$ và $A'M$.\\
			Ta có $\heva{&AA'\perp (ABC)\\&CM\subset (ABC)}\Rightarrow AA'\perp CM$.\\
			Khi đó $\heva{&CM\perp AB\\&CM\perp AA'}\Rightarrow CM\perp (ABB'A')$.\\
			Ta có $\heva{&CC'\parallel (ABB'A')\\&A'M\subset (ABB'A')}$\\
			$\Rightarrow \mathrm{d}(CC',A'M)=\mathrm{d}(CC',(ABB'A'))=\mathrm{d}(C,(ABB'A'))=CM=\dfrac{a\sqrt{3}}{2}$.
			\item Tính thể tích của khối lăng trụ tam giác đều $ABC.A'B'C'$ và thể tích khối chóp $A'.MBC$.\\
			Diện tích tam giác $ABC$ là $S_{\triangle {ABC}}=\dfrac{a^2\sqrt{3}}{4}$.\\
			Thể tích của khối lăng trụ tam giác đều $ABC.A'B'C'$ là $V_{ABC.A'B'C'}=S_{\triangle {ABC}}\cdot AA'=\dfrac{a^3\sqrt{3}}{4}$.\\
			Diện tích tam giác $MBC$ là $S_{\triangle {MBC}}=\dfrac{1}{2}S_{\triangle{ABC}}=\dfrac{a^2\sqrt{3}}{8}$.\\
			Thể tích khối chóp $A'.MBC$ là $V_{A'.MBC}=\dfrac{1}{3}\cdot S_{\triangle{MBC}}\cdot AA'=\dfrac{1}{3}\cdot \dfrac{a^2\sqrt{3}}{8}\cdot a=\dfrac{a^3\sqrt{3}}{24}$.
		\end{enumerate}
	}
\end{bt}

\begin{bt}%[1H8H2-2]
% nguon: Archive/2024/TOAN_11_XE/TOAN_11_XE/data4/11C7/KNTT-7.2.tex #1381

	Cho hình chóp $S.ABCD$ có $SA \perp(ABCD)$. Cho biết $ABCD$ là hình thang vuông tại $A$ và $D$, $AB=2AD$.
	\begin{enumerate}
		\item Chứng minh $CD \perp(SAD)$;
		\item Gọi $M$ là trung điểm của $AB$. Chứng minh $CM \perp (SAB)$.
	\end{enumerate}
	\loigiai{
		\immini{
			\begin{enumerate}
				\item Ta có $CD \perp AD \quad (1)$.\\
				Mặt khác $CD \perp SA \quad (2)$ do $SA \perp (ABCD) $.\\
				Từ $(1)$ và $(2)$ suy ra $CD \perp (SAD)$.
				\item  Ta có $SA \perp CM \quad (3)$ do $SA \perp (ABCD) $.\\
				Mặt khác, $CM \parallel AD \Rightarrow CM \perp AB \quad (4)$.\\
				Từ $(3)$ và $(4)$ suy ra $(SAB)$.
			\end{enumerate}
		}
		{
			\begin{tikzpicture}[scale=0.5, font=\footnotesize, line join=round, line cap=round, >=stealth]
				\coordinate (A) at (0,0);
				\coordinate (D) at (-2,-3);
				\coordinate (B) at (8,0);
				\coordinate (C) at ($(D)+(4.5,0)-(A)$);
				\coordinate (S) at ($(A)+(0,5)$);
				\coordinate (M) at ($(A)!0.5!(B)$);
				\draw(S)--(D) (S)--(C) (S)--(B) (D)--(C)--(B);
				\draw[dashed,thin](S)--(A) (A)--(D) (A)--(B) (C)--(M);
				\pic[draw,thin,angle radius=3mm] {right angle = S--A--B};
				\foreach \i/\g in {S/90,A/-90,B/-90,C/-90,D/-90,M/90}{\draw[fill=black](\i) circle (1.5pt) ($(\i)+(\g:3mm)$) node[scale=1]{$\i$};}
			\end{tikzpicture}
		}
	}
\end{bt}

\begin{bt}%[1H8H2-2]
% nguon: Archive/2024/TOAN_11_XE/TOAN_11_XE/data4/11C7/KNTT-7.2.tex #1382

	Cho hình vuông $ABCD$. Gọi $H$, $K$ lần lượt là trung điểm của $AB$, $AD$. Trên đường thẳng vuông góc với $(ABCD)$ tại $H$, lấy điểm $S$. Chứng minh rằng:
	\begin{enumEX}{2}
		\item $AC \perp (SHK)$;
		\item $CK \perp (SDH)$.
	\end{enumEX}
	\loigiai{
		\begin{center}
			\begin{tikzpicture}[scale=.65, font=\footnotesize, line join=round, line cap=round, >=stealth]
				\coordinate (A) at (0,0);
				\coordinate (B) at (-2,-3);
				\coordinate (D) at (7,0);
				\coordinate (C) at ($(B)+(D)-(A)$);
				\coordinate (S) at ($(-1,0)+(0,5)$);
				\coordinate (H) at ($(A)!0.5!(B)$);
				\coordinate (K) at ($(A)!0.5!(D)$);
				\coordinate (I) at (intersection of H--D and C--K);
				\draw(S)--(B) (S)--(C) (S)--(D) (B)--(C)--(D);
				\draw[dashed,thin](S)--(A) (A)--(B) (A)--(D) (S)--(H) (S)--(K)--(H) (A)--(C) (H)--(D) (C)--(K);
				\foreach \i/\g in {S/90,A/-90,B/-90,C/-90,D/0,H/180,K/45,I/-30}{\draw[fill=black](\i) circle (1.5pt) ($(\i)+(\g:3mm)$) node[scale=1]{$\i$};}
			\end{tikzpicture}
		\end{center}
		\begin{enumerate}
			\item Theo đề bài ta có $SH \perp (ABCD)$ mà $AC \subset (ABCD)$ nên $SH \perp AC \quad (1)$.\\
			Vì $HK$ là đường trung bình của $\triangle ABD \Rightarrow HK \parallel BD$.\\
			Mà $BD \perp AC$ (Vì $BD$ và $AC$ là hai đường chéo của hình vuông $ABCD$).\\
			Suy ra $HK \perp AC \quad (2)$.\\
			Từ $(1)$ và $(2)$ ta được $\heva{&AC \perp SH \subset (SHK)\\ & AC \perp HK \subset (SHK)\\ & SH \cup HK =H} \Rightarrow AC \perp (SHK)$.
			\item Gọi $I=CK \cap DH$.\\
			Suy ra, $\triangle IDC$ có $\widehat{IDC}+\widehat{ICD}=\widehat{IDC}+\widehat{ADH}=90^\circ \Rightarrow CK \perp DH \quad (3)$.\\
			Mà $AH \perp (ABCD) \Rightarrow AH \perp CK \quad (4)$.\\
			Từ $(3)$ và $(4)$ suy ra $CK \perp (SDH)$.
		\end{enumerate}
	}
\end{bt}

\begin{bt}%[1H8H2-2]
% nguon: Archive/2024/TOAN_11_XE/TOAN_11_XE/data4/11C7/KNTT-7.2.tex #1383

	Cho tứ diện $ABCD$ có $AB\perp (BCD)$, các tam giác $BCD$ và $ACD$ là những tam giác nhọn. Gọi $H$ và $K$ lần lượt là trực tâm của các tam giác $BCD$, $ACD$. Chứng minh rằng:
	\begin{enumerate}
		\item $CD\perp (ABH)$.
		\item $CD\perp (ABK)$.
		\item Ba đường thẳng $AK$, $BH$, $CD$ cùng đi qua một điểm.
	\end{enumerate}
	\loigiai{
		\immini{\begin{enumerate}
				\item Ta có $CD\perp BH$ (vì $H$ là trực tâm $\triangle BCD$)\\
				$CD\perp AB$ (vì $AB\perp (BCD)$)\\
				$BH,AB\subset (ABH)$.\\
				$\Rightarrow CD\perp (ABH)$.
				\item Ta có $CD\perp AK$ (vì $K$ là trực tâm $\triangle ACD$)\\
				$CD\perp AB$ (vì $AB\perp (BCD)$)\\
				$AK,AB\subset (ABK)$.\\
				$\Rightarrow CD\perp(ABK)$.
				\item Ta có hai mặt phẳng $(ABH)$ và $(ABK)$ cùng vuông góc với $CD$ nên chúng song song hoặc trùng nhau. Hơn nữa chúng có chung 2 điểm $A$, $B$ nên trùng nhau.\\
				Hay $4$ điểm $A,B,H,K$ đồng phẳng.\\
				Gọi $I$ là giao điểm của $BH$ và $CD$.\\
				Ta có $I$ vừa thuộc $CD$, vừa thuộc $BH$ nên nằm trên giao tuyến $AK$ của hai mặt phẳng $(ABK)$ và $(ACD)$.\\
				Vậy ba đường thẳng $BH$, $AK$, $CD$ đồng quy.
		\end{enumerate}}{\begin{tikzpicture}[scale=1,line join=round,line cap=round]
				\tkzDefPoints{0/0/B,1.2/-1.5/C,4/0/D}
				\coordinate (A) at ($(B)+(0,3)$);
				\coordinate (I) at ($(C)!0.4!(D)$);
				\coordinate (H) at ($(B)!2/3!(I)$);
				\coordinate (K) at ($(A)!2/3!(I)$);		
				\tkzDrawPolygon(A,B,C,D)
				\tkzDrawSegments(A,B A,C A,I)
				\tkzDrawSegments[dashed](B,D B,I)
				\tkzDrawPoints[fill=black,size=4](A,B,C,D,H,K,I)
				\tkzLabelPoints[above](A,K)
				\tkzLabelPoints[below](B,H,I)
				\tkzLabelPoints[below](C)
				\tkzLabelPoints[right](D)
		\end{tikzpicture}}	 
	}
\end{bt}

\begin{bt}%[1H8V2-3]
% nguon: Archive/2024/TOAN_11_XE/TOAN_11_XE/data4/11C7/KNTT-7.2.tex #1384

	Cho tứ diện $ABCD$ có $AB\perp (BCD)$, $BC\perp CD$. Gọi $M$, $N$ lần lượt là hình chiếu vuông góc của $B$ trên $AC$ và $AD$. Chứng minh rằng:
	\begin{enumerate}
		\item $CD\perp BM$;
		\item $BM\perp MN$.
	\end{enumerate}
	\loigiai{	
		\immini{\begin{enumerate}
				\item Ta có $CD\perp AB$ (vì $AB\perp (BCD)$)\\
				$CD\perp BC$ (giả thiết)\\
				$AB, BC\subset (ABC)$\\
				Suy ra $CD\perp (ABC)$.\\
				Suy ra $CD \perp BM$.
				\item Ta có $BM\perp AC$ ($M$ là hình chiếu của $B$ trên $AC$).\\
				$BM\perp CD$ (chứng minh trên)\\
				$\Rightarrow BM\perp (ACD)$\\
				$\Rightarrow BM\perp MN$.
		\end{enumerate}}{\begin{tikzpicture}[scale=1,line join=round,line cap=round]
				\tkzDefPoints{0/0/B,1.2/-1.5/C,4/0/D}
				\coordinate (A) at ($(B)+(0,3)$);
				\coordinate (M) at ($(A)!0.5!(C)$);
				\coordinate (N) at ($(A)!0.6!(D)$);	
				\tkzDrawPolygon(A,B,C,D)
				\tkzDrawSegments(A,B A,C B,M M,N)
				\tkzDrawSegments[dashed](B,D B,N)
				\tkzDrawPoints[fill=black,size=4](A,B,C,D,M,N)
				\tkzLabelPoints[above](A)
				\tkzLabelPoints[below](B,C)
				\tkzLabelPoints[above right](M,N)
				\tkzLabelPoints[right](D)
		\end{tikzpicture}}
	}
\end{bt}

\begin{bt}%[1H8H2-3]
% nguon: Archive/2024/TOAN_11_XE/TOAN_11_XE/data4/11C7/KNTT-7.2.tex #1385

	Cho hình chóp $O.ABC$ có $\widehat{AOB}=\widehat{BOC}=\widehat{COA}=90^{\circ}$. Chứng minh rằng:
	\begin{enumerate}
		\item $BC\perp OA$;
		\item $CA\perp OB$;
		\item $AB\perp OC$.
	\end{enumerate}
	\loigiai{
		\immini{\begin{enumerate}
				\item Ta có $\heva{&OA\perp OC\\&OA\perp OB\\&OB,OC\subset (OBC)}\Rightarrow OA\perp (OBC)\Rightarrow OA\perp BC$.
				\item Ta có $\heva{&OB\perp OA\\&OB\perp OC\\&OA,OC\subset (OAC)}\Rightarrow OB\perp (OAC)\Rightarrow OB\perp AC$.
				\item Tương tự, ta chứng minh được $OC\perp (OAB)\Rightarrow OC\perp AB$.
		\end{enumerate}}{\begin{tikzpicture}[scale=1, line join=round, line cap=round]
				\tkzDefPoints{0/0/A,1.3/-1.6/B,4.5/0/C,1/3.5/O}
				\tkzDrawPolygon(A,B,C,O)
				\tkzDrawSegments(O,B)
				\tkzDrawSegments[dashed](A,C)
				\tkzDrawPoints[fill=black,size=4](A,B,C,O)
				\tkzLabelPoints[above](O)
				\tkzLabelPoints[below](B)
				\tkzLabelPoints[left](A)
				\tkzLabelPoints[right](C)
				\tkzMarkRightAngles[size=0.3](A,O,B B,O,C A,O,C)
		\end{tikzpicture}}
	}
\end{bt}

\begin{bt}%[1H8H2-3]
% nguon: Archive/2024/TOAN_11_XE/TOAN_11_XE/data4/11C7/KNTT-7.1.tex #1386

	Cho hình hộp $ABCD.A'B'C'D'$ có $6$ mặt đều là hình vuông. Chứng minh rằng $AB\perp CC'$, $AC \perp B'D'$.
	\loigiai{
		Ta có $CC' \parallel BB'$, suy ra $\left(AB, CC'\right)=\left(AB, BB'\right)=\widehat{ABB'}=90^{\circ}$. Vậy $A B \perp C C'$.\\
		Ta có $B'D' \parallel BD$, suy ra $\left(AC, B'D'\right)=(AC, BD)=90^{\circ}$ (hai đường chéo của hình vuông luôn vuông góc với nhau). Vậy $AC \perp B'D'$.
	}
\end{bt}

\begin{bt}%[1H8H1-3]
% nguon: Archive/2024/TOAN_11_XE/TOAN_11_XE/data4/11C7/KNTT-7.1.tex #1387
%[Tex hóa SGK 11 Cánh Diều T2, Lương Như Quỳnh]
	\immini{Cho hình chóp $S.ABCD$ có đáy $A B C D$ là hình bình hành và $\widehat{SAB}=100^{\circ}$ (Hình bên). Tính góc giữa hai đường thẳng:
		\begin{listEX}
			\item $SA$ và $AB$,
			\item $SA$ và $CD$.
	\end{listEX}}
	{
		\begin{tikzpicture}[scale=1.2, font=\footnotesize, line join=round, line cap=round, >=stealth]
			\def\bc{4} % cạnh BC
			\def\ba{2} % cạnh BA
			\def\h{4} % đường cao
			\def\gocB{35} % góc B của đáy
			\coordinate[label=below left:$B$] (B) at (0,0);
			\coordinate[label=above right:$A$] (A) at (\gocB:\ba);
			\coordinate[label=below:$C$] (C) at (\bc,0);
			\coordinate[label=above right:$D$] (D) at ($(C)-(B)+(A)$);
			\coordinate [label=above:$S$](S) at (60:5);
			\draw (S)--(D)--(C)--(B)--cycle (S)--(C);
			\draw[dashed] (A)--(D)--(B)--cycle (A)--(C) (S)--(A);
			
			\foreach \i in {A,B,C,D,S} \fill[black] (\i) circle (1.5pt);
			\draw pic["$ 100^\circ $",draw=black, angle eccentricity=1.7, angle radius=.4cm]
			{angle=S--A--B}; 
		\end{tikzpicture}
	}
	\loigiai{
		\begin{listEX}
			\item Góc giữa hai đường thẳng $SA$ và $AB$ là $\widehat{SAB}=100^{\circ}$.
			\item Vì $CD\parallel AB$ nên góc giữa hai đường thẳng $SA$ và $CD$ bằng góc giữa hai đường thẳng $SA$ và $AB$. Suy ra $ (SA,CD)=100^{\circ} $.
		\end{listEX}
	}
\end{bt}

\begin{bt}%[1H8H2-1]
% nguon: Archive/2024/TOAN_11_XE/TOAN_11_XE/data4/11C7/KNTT-7.1.tex #1388
%[Tex hóa SGK 11 Cánh Diều T2, Lương Như Quỳnh]
	Bạn Hoa nói rằng: \lq\lq Nếu hai đường thẳng phân biệt $a$ và $b$ cùng vuông góc với đường thẳng $c$ thì $a$ và $b$ vuông góc với nhau\rq\rq. Bạn Hoa nói đúng hay sai? Vì sao?
	\loigiai{
		Bạn Hoa nói sai vì $a$ và $b$ chưa chắc vuông góc, chúng có thể cắt nhau, chéo nhau hay song song.\\
		\immini{Ví dụ. Hình lập phương $ABCD.A'B'C'D'$ có $ AB $ và $ CD $ cùng vuông góc với $ BC $ nhưng $ AB $ và $ CD $ song song.}
		{
			\begin{tikzpicture}[line cap=round,line join=round,font=\footnotesize,>=stealth,scale=.8]
				\coordinate[label=left:$A$] (A) at (0,0);
				\coordinate[label=below:$B$] (B) at (-1.5,-1.5);
				\coordinate[label=right:$D$] (D) at (4,0);
				\coordinate[label=below:$C$] (C) at ($(B)+(D)-(A)$);
				\coordinate[label=left:$A'$] (A') at	($(A)+(0,3)$);
				\coordinate[label=left:$B'$] (B') at ($(B)+(A')-(A)$);
				\coordinate[label=above:$C'$] (D') at ($(D)+(A')-(A)$);
				\coordinate[label=above:$D'$] (C') at ($(B')+(D')-(A')$);
				\draw (A')--(B')--(C')--(D')--cycle (B)--(C)--(D) (B')--(C')--(C) (D)--(D') (B')--(B);
				\draw[dashed] (B)--(A)--(D) (A)--(A');
			\end{tikzpicture}
		}
	}
\end{bt}

\begin{bt}%[1H8V1-3]
% nguon: Archive/2024/TOAN_11_XE/TOAN_11_XE/data4/11C7/KNTT-7.1.tex #1389

	Cho hình chóp $S.ABCD$ có đáy là hình thoi $A B C D$ cạnh $a$. Cho biết $S A=a \sqrt{3}, S A \perp A B$ và $S A \perp A D$. Tính góc giữa $S B$ và $C D$, $S D$ và $C B$.
	\loigiai{
		\immini{
			Vì $CD\parallel AB$ nên $(SB,CD)=(SB,AB)=\widehat{SBA}$.\\
			$\triangle SBA$ vuông tại $A$, có $\tan \widehat{SBA}=\dfrac{SA}{AB}=\dfrac{a\sqrt{3}}{a}=\sqrt{3}\Rightarrow \widehat{SBA}=60^\circ$.\\
			Tương tự, $CB\parallel AD$ nên $(SD,CB)=(SD,AD)=\widehat{SDA}$. \\
			Do $\triangle SAD=\triangle SAB$ (c.g.c) nên $\widehat{SDA}=\widehat{SBA}=60^\circ$.
		}{
			\begin{tikzpicture}[line join=round, line cap=round,every node/.style={scale=0.8}]
				\path 
				(0,0) coordinate (A)
				(2,0) coordinate (D)
				(-130:1) coordinate (B)
				($(B)+(D)-(A)$) coordinate (C)
				(0,1.5) coordinate (S)
				(intersection of A--C and B--D) coordinate (O);
				\draw (S)--(B)--(C)--(D)--cycle (S)--(C);
				\draw[dashed] (A)--(B) (A)--(D) (A)--(S)
				;
				\foreach \t/\g in {S/90,A/-90,B/-90,C/-90,D/0}{
					\draw[fill=red,draw=black] (\t) circle (1pt) node[shift={(\g:7pt)}]{$ \t $};
				}
				\draw pic[draw, angle radius=2mm]{right angle=D--A--S}; 
				\draw pic[draw, angle radius=2mm]{right angle=S--A--B};
				\tkzMarkSegments[mark=x,size=2pt](A,B A,D C,B C,D) 
			\end{tikzpicture}
		}
	}
\end{bt}

\begin{bt}%[1H8V2-3]
% nguon: Archive/2024/TOAN_11_XE/TOAN_11_XE/data4/11C7/KNTT-7.1.tex #1391

	Cho hình chóp $S . A B C$ có $S A=S B=S C=a$, $\widehat{B S A}=\widehat{C S A}=60^{\circ}$, $\widehat{B S C}=90^{\circ}$. Cho $I$ và $J$ lần lượt là trung điểm của $S A$ và $B C$. Chứng minh rằng $I J \perp S A$ và $I J \perp B C$.
	\loigiai{
		\immini{
			$\triangle SAB$ và $\triangle SAC$ cân tại $S$, có $\widehat{ASB}=\widehat{ASC}=60^\circ$, suy ra $SAB$ và $SAC$ là các tam giác đều. Suy ra $AB=AC=a$.\\
			$\triangle SBC$ vuông cân tại $S$, suy ra $BC=SA\sqrt{2}=a\sqrt{2}$.\\
			Suy ra $\triangle BAC$ vuông cân tại $A$.\\
			$SBC$ và $ABC$ là các tam giác vuông có cùng cạnh huyền $BC$, $J$ là trung điểm $BC\Rightarrow JS=JA \left(=\dfrac{BC}{2}\right)$.\\
			$\triangle JSA$ cân tại $S$ có $JI$ là đường trung tuyến, suy ra $JI\perp SA$.\\
			$IB$ và $IC$ là các đường cao của tam giác đều có cùng cạnh $a$, suy ra $IB=IC$. \\
			$\triangle IBC$ cân tại $I$, có $IJ$ là đường trung tuyến nên $IJ\perp BC$.
			
		}{
			\begin{tikzpicture}[line join=round, line cap=round, every node/.style={scale=0.8}]
				\path 
				(0,0) coordinate (A)
				(3,0) coordinate (C)
				(1,-.7) coordinate (B)
				(1,2.3)coordinate (S)
				($(S)!.5!(A)$) coordinate (I)
				($(B)!.5!(C)$) coordinate (J)
				;
				\draw pic[draw, angle radius=2mm,fill=gray!20]{right angle=B--S--C}; 
				\draw (S)--(A)--(B)--(C)--cycle (J)--(S)--(B)--(I);
				\draw[dashed] (A)--(C) (I)--(C) (I)--(J)--(A);
				\foreach \t/\g in {S/90,A/180,B/-90,C/0,I/180,J/-90}{
					\draw[fill=red,draw=black] (\t) circle (1pt) node[shift={(\g:7pt)}]{$ \t $};
				}
				%\draw pic[draw, angle radius=5mm]{angle=A--S--B};
				%\draw pic[draw, angle radius=2mm]{right angle=A--S--C};  
				\tkzMarkSegments[mark=x,size=2pt](I,S I,A) 
				\tkzMarkSegments[mark=||,size=2pt](J,B J,C)
				\draw pic[draw, angle radius=2mm]{right angle=B--A--C}; 
			\end{tikzpicture}
		}
	}
\end{bt}

\begin{bt}%[1H8V1-3]
% nguon: Archive/2024/TOAN_11_XE/TOAN_11_XE/data4/11C7/KNTT-7.1.tex #1392

	Cho tứ diện đều $A B C D$ cạnh $a$. Gọi $K$ là trung điểm của $C D$. Tính góc giữa hai đường thẳng $A K$ và $B C$.
	\loigiai{
		\immini{
			Gọi $I$ là trung điểm $BD$. Khi đó $IK$ là đường trung bình của $\triangle BCD$ nên $IK\parallel BC$. 
			Do đó $(AK,BC)=(AK,IK)$.\\
			Xét $\triangle AIK$ có: $\heva{&AI=AK=\dfrac{a\sqrt{3}}{2} \text{ (chiều cao của tam giác đều cạnh $a$)}\\&IK=\dfrac{BC}{2}=\dfrac{a}{2}\text{ ($IK$ là đường trung bình của $\triangle BCD$). }}$\\
			$\Rightarrow \cos \widehat{AKI}=\dfrac{KA^2+KI^2-AI^2}{2KI\cdot KA}=\dfrac{\sqrt{3}}{6}\Rightarrow \widehat{AKI}\approx 73^\circ 13'$.
		}{
			\begin{tikzpicture}[line join=round, line cap=round,every node/.style={scale=0.8}] 
				\path 
				(0,0) coordinate (B)
				(2.4,0) coordinate (D)
				(.6,-.7) coordinate (C)
				($(C)!.5!(D)$) coordinate (K)
				($(B)!.5!(K)$) coordinate (G)
				--+(0,2.5) coordinate (A)
				($(B)!.5!(D)$) coordinate (I);
				\draw (A)--(B)--(C)--(D)--cycle (K)--(A)--(C);
				\draw[dashed] (B)--(D) (I)--(A) (I)--(K);
				\foreach \t/\g in {A/90,B/180,C/-90,D/0,K/-90,I/-90}{
					\draw[fill=red,draw=black] (\t) circle (1pt) node[shift={(\g:7pt)}]{$ \t $};
				}
				\tkzMarkSegments[mark=||,size=2pt](K,C K,D)
				\tkzMarkSegments[mark=x,size=2pt](I,B I,D)
			\end{tikzpicture}
		}
	}
\end{bt}

\begin{bt}%[1H8V1-3]
% nguon: Archive/2024/TOAN_11_XE/TOAN_11_XE/data4/11C7/KNTT-7.1.tex #1393

	Cho tứ diện $A B C D$. Gọi $M, N$ lần lượt là trung điểm của $B C$ và $A D$. Biết $A B=C D=2 a$ và $M N=a \sqrt{3}$. Tính góc giữa $A B$ và $C D$.
	\loigiai{
		\immini{
			Gọi $P$ là trung điểm $AC$. Ta có $\heva{&MP\parallel AB \text{ ($MP$ là đường trung bình của $\triangle CAB$) }\\&NP\parallel CD \text{ ($NP$ là đường trung bình của $\triangle ACD$). }}$\\
			$\Rightarrow (AB,CD)=(MP,NP)$.\\
			Xét $\triangle MNP$ có $\heva{&NP=\dfrac{CD}{2}=a\\&MP=\dfrac{AB}{2}=a\\&MN=a\sqrt{3}.}$ \\
			$\Rightarrow \cos\widehat{NPM}=\dfrac{PN^2+PM^2-MN^2}{2PN\cdot PM}=-\dfrac{1}{2}\Rightarrow \widehat{NPM}=120^\circ>90^\circ$.\\
			Suy ra $(MP,NP)=180^\circ-\widehat{NPM}=180^\circ-120^\circ=60^\circ$.\\
			Vậy $(AB,CD)=60^\circ$.
		}{
			\begin{tikzpicture}[line join=round, line cap=round, every node/.style={scale=0.8}]
				\path 
				(0,0) coordinate (D)
				(3,0) coordinate (B)
				(1,-.7) coordinate (C)
				(1.3,2.3)coordinate (A)
				($(B)!.5!(C)$) coordinate (M)
				($(A)!.5!(D)$) coordinate (N)
				($(A)!.5!(C)$) coordinate (P)
				;
				\draw (A)--(D)--(C)--(B)--cycle (A)--(C)
				(N)--(P)--(M);
				\draw[dashed] (B)--(D) (M)--(N);
				\foreach \t/\g in {A/90,D/180,C/-90,B/0,M/-90,N/180,P/40}{
					\draw[fill=red,draw=black] (\t) circle (1pt) node[shift={(\g:7pt)}]{$ \t $};
				}
				\tkzMarkSegments[mark=x,size=2pt](M,B M,C) 
				\tkzMarkSegments[mark=||,size=2pt](N,A N,D)
				\tkzMarkSegments[mark=|,size=2pt](A,B C,D)
			\end{tikzpicture}
		}
	}
\end{bt}

\begin{bt}%[1H8H6-1]
% nguon: Archive/2024/TOAN_11_XE/TOAN_11_XE/data4/11C7/KNTT-7.4.tex #1394

	Cho hình chóp $S.ABC$ có $SA \perp(ABC)$, $SA=a$, $CA=CB=a \sqrt{7}$, $AB=2a$.
	\begin{listEX}[1]
		\item Gọi $\alpha$ là góc giữa $SB$ và $(ABC)$. Tính $\tan \alpha$.
		\item Tính góc giữa $SC$ và $(SAB)$.
	\end{listEX}
	\loigiai{
		\begin{listEX}[1]
			\item  \immini{Do $SA \perp(ABC)$ nên $\alpha=\widehat{SBA}$. Tam giác $SAB$ vuông tại $A$ nên
				$$
				\tan \alpha=\tan \widehat{SBA}=\dfrac{SA}{AB}=\dfrac{a}{2a}=\dfrac{1}{2} .
				$$}{\begin{tikzpicture}[scale=.8,>=stealth, font=\footnotesize, line join=round, line cap=round]
					\path 
					(0,0)coordinate(A) 
					(5,0)coordinate(B) 
					(-60:2.5)coordinate(C) 
					($(A)+(0,4)$) coordinate(S)
					($(A)!0.5!(B)$)coordinate(M)	 
					;
					\draw (S)--(A)--(C)--(B)--cycle  (C)--(S) 
					pic[draw,angle radius=3mm]{right angle=S--A--B}
					pic[draw,angle radius=3mm]{right angle=S--A--C}
					pic[draw,angle radius=3mm]{right angle=C--M--B}
					;
					\path
					(S)--(A)node[midway,left]{$a$}
					(M)--(A)node[midway,above]{$a$}
					(B)--(M)node[midway,above]{$a$}
					(C)--(A)node[midway,left]{$a\sqrt{7}$}
					(C)--(B)node[midway,right]{$a\sqrt{7}$}
					;
					\draw [dashed](S)--(M)--(C) (B)--(A) 
					;	
					\foreach \x/\g in {S/90,A/180,C/-90,M/40,B/0}\fill[black] (\x) circle (1pt) +(\g:.3)node{$\x$};
			\end{tikzpicture}}
			\item Gọi $M$ là trung điểm của $A B$. Tam giác $ABC$ cân tại $C$ nên $CM \perp AB$.\\		
			Mặt khác, từ $SA \perp(ABC)$ ta có $CM \perp SA$. Do đó $CM \perp(SAB)$.
			Vậy góc giữa $SC$ và $(SAB)$ bằng $\widehat{CSM}$.\\
			Tam giác $SAC$ vuông tại $A$ nên $SC=\sqrt{S A^2+A C^2}=\sqrt{a^2+7 a^2}=a \sqrt{8}$.\\
			Ta có $AM=\dfrac{1}{2} AB=a$. Do đó, tam giác $SAM$ vuông cân tại $A$ và $SM=a \sqrt{2}$.\\
			Tam giác $CMS$ vuông tại $M$ và $\cos \widehat{CSM}=\dfrac{SM}{SC}=\dfrac{a \sqrt{2}}{a \sqrt{8}}=\dfrac{1}{2}$.\\
			Vậy $\widehat{CSM}=60^{\circ}$ và do đó góc giữa $SC$ và $(SAB)$ bằng $60^{\circ}$.
		\end{listEX}
	}
\end{bt}

\begin{bt}%[1H8H2-2]
% nguon: Archive/2024/TOAN_11_XE/TOAN_11_XE/data4/11C7/KNTT-7.4.tex #1396
%[HK2 (2016-2017), THPT THỚI LAI, CẦN THƠ]%[Lê Thanh Nin]
	Cho hình chóp $ S.ABCD $ có đáy $ ABCD $ là hình vuông cạnh $ a, SA\perp (ABCD) $ và $ SA=a\sqrt{10} $. Gọi $ M,N $ lần lượt là trung điểm của $ BC $ và $ CD $.
	\begin{tasks}(1)
		\task Chứng minh $ BD\perp (SAC) $.
		\task Tính góc giữa $ SM $ và $ (ABCD) $.
	\end{tasks}
	\loigiai{
		\begin{center}
			\begin{tikzpicture}[line cap=round,line join=round, >=stealth,scale=1]
				\def \xa{-1.3} 
				\def \xb{-1.5}
				\def \y{4}
				\def \z{3}
				\coordinate (A) at (0,0);
				\coordinate (B) at ($(A)+(\xa,\xb)$);
				\coordinate (D) at ($(A)+(\y,0)$);
				\coordinate (C) at ($ (B)+(D)-(A) $);
				\coordinate (S) at ($ (A)+(0,\z) $);
				\tkzDefMidPoint(B,C)    \tkzGetPoint{M}
				\tkzDefMidPoint(D,C)    \tkzGetPoint{N}
				\tkzInterLL(A,C)(B,D) \tkzGetPoint{O}
				\tkzInterLL(A,C)(M,N) \tkzGetPoint{I}
				\coordinate (H) at ($ (S)!3/5!(I) $);
				\draw [dashed] (B)--(A)--(D)--(B) (H)--(A)--(S)--(I) (N)--(M)--(A)--(C);
				\draw (S)--(B)--(C)--(D)--(S)--(C) (N)--(S)--(M);
				\tkzDrawPoints(S,A,B,C,D,M,N,I,H,O)
				\tkzLabelPoints[right](D,N)
				\tkzLabelPoints[above](S,I,O)
				\tkzLabelPoints[left](A,H)
				\tkzLabelPoints[below left](C,B,M)
			\end{tikzpicture}
		\end{center}
		\begin{enumerate}
			\item Ta có $ \heva{&BD\perp AC\\&BD\perp SA}\Rightarrow BD\perp (SAC) $.
			\item Hình chiếu vuông góc của $ SM $ lên mặt phẳng $ (ABCD) $ là $ AM $.\\
			Nên $ \widehat{(SM,(ABCD))}=\widehat{(SM,AM)}=\widehat{SMA} $.\\
			Xét $ \triangle  SAM $ vuông tại $ A $, ta có:\\
			$ \tan \widehat{SMA}=\dfrac{SA}{AM}=\dfrac{a\sqrt{10}}{\dfrac{a\sqrt{5}}{2}}=2\sqrt{2}\Rightarrow \widehat{SMA}\approx 70^\circ31'. $
		\end{enumerate}	
	}
\end{bt}

\begin{bt}%[1H8H2-2]
% nguon: Archive/2024/TOAN_11_XE/TOAN_11_XE/data4/11C7/KNTT-7.4.tex #1397
%[Phan Hoàng Anh]
	Cho tứ diện $OABC$ có các cạnh $OA$, $OB$, $OC$ đôi một vuông góc với nhau. Gọi $H$ là hình chiếu vuông góc của điểm $O$ trên mặt phẳng $(ABC)$.
	\begin{tasks}(1)
		\task Chứng minh rằng $BC\perp (OAH)$, $CA\perp (OBH)$, $AB\perp (OCH)$.
		\task Chứng minh $H$ là trực tâm của tam giác $ABC$.
		\task Chứng minh rằng $\dfrac{1}{OH^2}=\dfrac{1}{OA^2}+\dfrac{1}{OB^2}+\dfrac{1}{OC^2}$.
		\task Chứng minh rằng các góc của tam giác $ABC$ đều nhọn.
	\end{tasks}
	\loigiai{
		\begin{enumerate}[a)]
			\item \immini{Từ giả thiết ta có $OH\perp (ABC)$ suy ra $OH\perp BC$.\hfill$(1)$\\
				Ta có $\left\{\begin{aligned}&OA\perp OB \\&OA\perp OC\end{aligned}\right.\Rightarrow OA\perp (OBC) \Rightarrow OA\perp BC$.\hfill$(2)$\\
				Từ $(1)$ và $(2)$, suy ra $BC\perp (OAH)$. \\
				Chứng minh tương tự ta được $$CA\perp (OBH)\text{ và }AB\perp (OCH).$$}
			{\begin{tikzpicture}[scale=0.6]
					\tkzDefPoints{0/0/O, -2/-1.8/B,5/0/C,0/4/A}
					\coordinate (K) at ($(B)!0.55!(C)$);
					\coordinate (H) at ($(A)!0.6!(K)$);
					\tkzDrawSegments(A,B B,C C,A A,K)
					\tkzDrawSegments[dashed](O,A O,B O,C O,K O,H)
					\tkzDrawPoints(A,B,C,O,H,K)
					\tkzLabelPoints[above left](O)
					\tkzLabelPoints[above right](C,H)
					\tkzLabelPoints[below left](B)
					\tkzLabelPoints[below right](K)
					\tkzLabelPoints[above](A)
					\tkzMarkRightAngle(O,K,B)
					\tkzMarkRightAngle(A,H,O)
			\end{tikzpicture}}
			\item Từ kết quả câu a, ta có $BC\perp (OAH)$ suy ra $BC\perp AH$.\hfill$(3)$\\
			Tương tự, từ $AC\perp (OBH)$ suy ra $AC\perp BH$.\hfill$(4)$\\
			Từ $(3)$ và $(4)$, ta có $\left\{\begin{aligned}
				&AH\perp BC \\
				&BH\perp AC
			\end{aligned}\right. $ suy ra $H$ là trực tâm của $\triangle ABC$.
			\item Giả sử $AH$ cắt $BC$ tại $K$, suy ra $OK\perp BC$. \\
			Trong $\triangle OBC$ vuông tại $O$, ta có $\dfrac{1}{OK^2}=\dfrac{1}{OB^2}+\dfrac{1}{OC^2}$.\hfill$(5)$\\
			Trong $\triangle OAK$ vuông tại $O$, ta có $\dfrac{1}{OH^2}=\dfrac{1}{OA^2}+\dfrac{1}{OK^2}$.\hfill$(6)$\\
			Từ $(5)$ và $(6)$, suy ra $\dfrac{1}{OH^2}=\dfrac{1}{OA^2}+\dfrac{1}{OB^2}+\dfrac{1}{OC^2}$.
			\item Xét $\cos \widehat{BAC}=\dfrac{AB^2+AC^2-BC^2}{2AB\cdot AC}$. \\
			Để chứng minh $\widehat{BAC}$ nhọn ta cần chứng minh $\cos \widehat{BAC}>0 \Leftrightarrow AB^2+AC^2-BC^2>0$. \\
			Thật vậy, ta có $AB^2=OA^2+OB^2;BC^2=OB^2+OC^2;AC^2=OA^2+OC^2$. \\
			Suy ra $AB^2+AC^2-BC^2=2OA^2>0$ (đpcm). \\
			Chứng minh tương tự, ta được các góc $\widehat{BAC}$ và $\widehat{ACB}$ đều nhọn.
	\end{enumerate}}
\end{bt}

\begin{bt}%[1H8H4-5]
% nguon: Archive/2024/TOAN_11_XE/TOAN_11_XE/data4/11C7/KNTT-7.5.tex #1398

	Cho hình chóp $S . ABC$ có đáy là tam giác vuông tại $C$, mặt bên $SAC$ là tam giác đều và nằm trong mặt phẳng vuông góc với $(ABC)$.
	\begin{enumerate}
		\item Chứng minh rằng $(SBC)\perp(SAC)$.
		\item Gọi $I$ là trung điểm của $SC$. Chứng minh rằng $(ABI)\perp(SBC)$.
	\end{enumerate}
	
	\loigiai{
		\immini
		{
			\begin{enumerate}
				\item Gọi $H$ là trung điểm $AC$. Suy ra $SH\perp AC$ ($\triangle  SAC $ đều).\\
				Vì $(SAC)\perp (ABC)$ nên $SH\perp (ABC)$ suy ra $BC\perp SH$.\tagEX{1}
				Mà $BC\perp AC$. \tagEX{2}
				Từ (1), (2) suy ra $BC\perp (SAC)$ mà $BC\subset (SBC)$\\ nên $(SBC)\perp (SAC)$.
				\item $\triangle  SAC $ đều suy ra $AI\perp SC$.\tagEX{3}
				Ta có $BC\perp (SAC)$ suy ra $BC\perp AI$.\tagEX{4}			
				Từ (3),(4) suy ra $AI\perp (SBC)$ mà $AI\subset (ABI)$ \\nên $(ABI)\perp(SBC)$.
			\end{enumerate}
		}{\begin{tikzpicture}[scale=0.8,font=\footnotesize, line join=round, line cap=round, >=stealth]
				\path
				(0,0) coordinate (A)
				(2,-1.5) coordinate (B)
				(5,0) coordinate (C)
				($(A)!.5!(C)$) coordinate (H)
				($(H)+(0,4)$) coordinate (S)
				($(S)!.5!(C)$) coordinate (I)
				;
				\draw (A)--(B)--(C) (S)--(A) (S)--(B) (S)--(C) (B)--(I);
				\draw[dashed] (A)--(C) (S)--(H) (A)--(I);
				\draw pic[draw, angle radius=2mm]{right angle=C--H--S};
				\draw pic[draw, angle radius=2mm]{right angle=A--C--B};
				\foreach \x/\g in {A/180,B/-90,C/0,S/90,H/-90,I/0} \fill[black] (\x) circle (1pt)+(\g:.3) node {$\x$};
			\end{tikzpicture}
		}
		
	}
\end{bt}

\begin{bt}%[1H8H2-4]
% nguon: Archive/2024/TOAN_11_XE/TOAN_11_XE/data4/11C7/KNTT-7.5.tex #1399

	Cho tam giác đều $ABC$ cạnh $a, I$ là trung điểm của $BC, D$ là điểm đối xứng với $A$ qua $I$. Vẽ đoạn thẳng $SD$ có độ dài bằng $\dfrac{a\sqrt{6}}{2}$ và vuông góc với $(ABC)$. Chứng minh rằng
	\begin{enumEX}{2}
		\item $(SBC)\perp(SAD)$
		\item $(SAB)\perp(SAC)$.
	\end{enumEX}
	
	\loigiai{
		\immini
		{
			\begin{enumerate}
				\item Ta có $SD\perp (ABC)$ suy ra $SD\perp AD$.\tagEX{1}
				Vì $\triangle ABC$ đều cạnh $a$ và $D$ đối xứng với $I$ qua $A$ nên $ABCD$ là hình  thoi.\\
				Suy ra $AD\perp BC$.	\tagEX{2}
				Từ (1), (2) suy ra $AD\perp (SBC)$ mà $AD\subset (SAD)$ nên $(SBC)\perp (SAD)$.
				\item Kẻ $BK\perp SA$.\\
				Ta có $\triangle SBA=\triangle SCA$ (c.c.c)
				$\Rightarrow \heva{&BK=CK\\&CK\perp SA.\\}$				
			\end{enumerate}
		}{
			\begin{tikzpicture}[scale=0.8,font=\footnotesize, line join=round, line cap=round, >=stealth]
				\path
				(180:2) coordinate (A)
				(-90:2) coordinate (B)
				(0:2) coordinate (C)
				($(B)!.5!(C)$) coordinate (I)
				($(A)!2!(I)$) coordinate (D)
				($(D)+(0,6)$) coordinate (S)
				($(A)!(B)!(S)$) coordinate (K)
				;
				\draw (A)--(B)--(D) (S)--(D) (S)--(A) (S)--(B) (B)--(K)--(C);
				\draw[dashed] (A)--(C)--(B) (C)--(D) --(A) (S)--(C);
				%\draw pic[draw, angle radius=2mm]{right angle=C--H--S};
				\draw pic[draw, angle radius=2mm]{right angle=A--D--S};
				\draw pic[draw, angle radius=2mm]{right angle=B--K--A};
				\draw pic[draw, angle radius=2mm]{right angle=C--K--S};
				\foreach \x/\g in {A/180,B/-90,C/0,I/-90,D/0,S/90,K/140} \fill[black] (\x) circle (1pt)+(\g:.3) node {$\x$};
			\end{tikzpicture}
		}
		\noindent  Mà $AD=2AI=\sqrt{3}a\Rightarrow SA=\dfrac{3\sqrt{2}}{2}a$; $BD=a\Rightarrow SB=\dfrac{\sqrt{10}}{2}a$.\\
			Suy ra $S_{\triangle SAB}=\sqrt{p(p-SB)(p-BA)(p-SA)}=\dfrac{3}{4}a^2$.\\
			Mặt khác, $BK=\dfrac{2S_{\triangle SAB}}{SA}=\dfrac{\sqrt{2}}{2}a\Rightarrow BK=CK=\dfrac{\sqrt{2}}{2}a$.\\
			Mà $BK^2+CK^2=BC^2\Rightarrow BK\perp CK$.\\
			Suy ra $(SAB)\perp (SAC)$.
	}
\end{bt}

\begin{bt}%[1H8H4-4]
% nguon: Archive/2024/TOAN_11_XE/TOAN_11_XE/data4/11C7/KNTT-7.5.tex #1400
%[Nguyễn Thành Khang, dự án Tex hoá lần 3]
	Cho hình chóp $S.ABCD$ có đáy $ABCD$ là hình vuông và $SA$ vuông góc với mặt đáy. Gọi $(P)$ là mặt phẳng chứa $AD$ và vuông góc với mặt phẳng $(SBC)$. Xác định thiết diện do mặt phẳng $(P)$ cắt hình chóp. 
	\loigiai{
		\immini{
			Kẻ $AH \perp SB\ (H \in SB)$. Ta có $$\heva{&AD \perp SA\ (SA \perp (ABCD)) \\ &AD \perp AB} \Rightarrow AD \perp (SAB) \Rightarrow AD \perp SB.$$
			Ta có $AD \perp SB$ và $AH \perp SB$ nên $SB \perp (ADH)$, suy ra $(SBC) \perp (ADH)$.\\
			Do đó, mặt phẳng $(P)$ chứa $AD$ và vuông góc với mặt phẳng $(SBC)$ là $(ADH)$.\\
			Trong mặt phẳng $(SBC)$ dựng $HK \parallel BC (K \in BC)$, suy ra $HK \parallel AD$ (do cùng song song với $AD$).
		}{
			\begin{tikzpicture}[scale=.7, font=\footnotesize, line join=round, line cap=round,>=stealth]
				\tkzDefPoints{0/0/A, -2/-2/D, 4/-2/C, 6/0/B, 0/4/S}
				\coordinate (H) at ($(S)!0.7!(B)$);
				\coordinate (K) at ($(S)!0.7!(C)$);
				\tkzDrawPoints[fill=black](A,B,C,D,S,H,K)
				\tkzDrawSegments(S,B S,C S,D B,C C,D H,K D,K)
				\tkzDrawSegments[dashed](A,B S,A A,D A,H D,H)
				\tkzMarkRightAngles(S,A,D S,A,B B,A,D)
				\tkzLabelPoints[above left](A)
				\tkzLabelPoints[above left](D)
				\tkzLabelPoints[below right](C)
				\tkzLabelPoints[above right](B,H)
				\tkzLabelPoints[above](S)
				\tkzLabelPoints[below](K)
				\tkzDrawPolygon[dashed,pattern=north east lines,opacity=.3](A,D,K,H)
				\tkzDrawPoints(H,K)
			\end{tikzpicture}
		}
		\noindent Vậy thiết diện là hình thang $ADKH$ có hai đáy là $AD$ và $HK$.
	}
\end{bt}

\begin{bt}%[1H8H2-4]
% nguon: Archive/2024/TOAN_11_XE/TOAN_11_XE/data4/11C7/KNTT-7.5.tex #1401
%[Nguyễn Thành Khang, dự án Tex hoá lần 3]
	Cho tứ diện $OABC$ có các cạnh $OA, OB, OC$ đôi một vuông góc nhau và $OA=OB=OC=a$. Gọi $M$ là điểm trên cạnh $AC$ sao cho $AC=3MC$. Gọi $H$ là hình chiếu của $O$ lên mặt phẳng $(ABC)$. Gọi $(\alpha)$ là mặt phẳng chứa $OM$ và vuông góc với mặt phẳng $(ABC)$. Xác định và tính diện tích thiết diện do mặt phẳng $(\alpha)$ cắt tứ diện.
	\loigiai{
		\immini{
			Ta có $OA \perp OB$ và $OA \perp OC$ nên $OA \perp (OBC)$.\\
			Vì $OH \perp BC$ (do $OH \perp (ABC)$) và $OA \perp BC$ (do $OH \perp (ABC)$) nên $BC \perp (OAH)$, suy ra $BC \perp AH$.\hfill $(1)$\\
			Tương tự ta có $AC \perp BH$.\hfill $(2)$\\
			Từ $(1)$ và $(2)$ suy ra $H$ là trực tâm của $\triangle ABC$.\\
			Ta có $OH  \perp (ABC)$ nên mặt phẳng $(\alpha)$ chứa $OM$ và vuông góc với mặt phẳng $(ABC)$ là $(OHM)$.\\
			Trong mặt phẳng $(ABC)$ gọi $N$ là giao điểm của $HM$ và $AB$, suy ra thiết diện của mặt phẳng $(OHM)$ và tứ diện là $\triangle OMN$.\\
			Do $OA=OB=OC=a$ và $OA, OB, OC$ đôi một vuông góc nhau nên $AB=AC=BC= a\sqrt{2}$, suy ra $\triangle ABC$ đều.\\
			Mà $H$ là trực tâm của $\triangle ABC$ nên $AH$ cắt $BC$ tại trung điểm của đoạn $BC$.
		}{
			\begin{tikzpicture}[scale=.55, font=\footnotesize, line join=round, line cap=round,>=stealth]
				\tkzDefPoints{0/0/O, -3/-2/B, 6/-3/C, 0/9/A, 0.8/1.3/H}
				\coordinate (M) at ($(A)!0.7!(C)$);
				\tkzInterLL(A,H)(B,C) \tkzGetPoint{E}
				\tkzInterLL(H,M)(A,B) \tkzGetPoint{N}
				\tkzDrawSegments[dashed](O,A O,B O,C O,H O,M O,E O,N)
				\tkzDrawSegments(A,B A,C B,C M,N A,E)
				\tkzMarkRightAngles(A,O,B A,O,C C,O,B)
				\tkzLabelPoints[below](O,E)
				\tkzLabelPoints[above left](B,N)
				\tkzLabelPoints[above](A)
				\tkzLabelPoints[above right](C,M,H)
				\tkzDrawPolygon[dashed,pattern=north east lines,opacity=.3](O,M,N)
				\tkzDrawPoints[fill=black](A,B,C,O,H,M,N,E)
			\end{tikzpicture}
		}
		\noindent Gọi $E$ là trung điểm $BC$.\\
		Ta có $H$ là trọng tâm $\triangle ABC$ nên $3HE=AE$. Do đó, $MN \parallel BC$, suy ra $MN= \dfrac{2}{3} BC= \dfrac{2a \sqrt{3}}{3}$.\\
		Vì $\triangle AOE$ vuông tại $A$ nên $\dfrac{1}{OH^2}=\dfrac{1}{OA^2}+\dfrac{1}{OE^2}=\dfrac{3}{4a^2}$, suy ra $AH=\dfrac{a \sqrt{3}}{3}$.\\
		Vậy $S_{\triangle OMN}= \dfrac{1}{2} \cdot MN \cdot OH= \dfrac{2a^2 \sqrt{6}}{9}$.
	}
\end{bt}

\begin{bt}%[1H8H4-5]
% nguon: Archive/2024/TOAN_11_XE/TOAN_11_XE/data4/11C7/KNTT-7.5.tex #1402
%[TeX hóa SGK CTST]%[Truc Danh]
	Cho hình lăng trụ đứng $ABCD\cdot A'B'C'D'$ có đáy $ABCD$ là hình thang vuông tại $A$ và $B$, $AA'=2a, AD=2a, AB=BC=a$.
	\begin{enumerate}
		\item Tính độ dài đoạn thẳng $AC'$.
		\item Tính tổng diện tích các mặt của hình lăng trụ.
	\end{enumerate}
	\loigiai{
		\immini
		{
			\begin{enumerate}
				\item Tam giác $ABC$ vuông tại $B$ có $AC=\sqrt{AB^2+AC^2}=a\sqrt{2}$.\\
				Tam giác $ACC'$ vuông tại $C$ có $$AC'=\sqrt{AC^2+CC'^2}=\sqrt{(a\sqrt{2})^2+(2a)^2}=a\sqrt{6}.$$
				\item Gọi $I$ là trung điểm $AD$, suy ra $ABCI$ là hình vuông nên $\triangle CID$ vuông tại $I$. Ta có
				$$CD=\sqrt{IC^2+ID^2}=a\sqrt{2}.$$
			\end{enumerate}
		}{
			\begin{tikzpicture}[scale=0.8,font=\footnotesize, line join=round, line cap=round, >=stealth]
				\path 
				(0,0) coordinate (B)
				(2.5,0) coordinate (C)
				(4,1) coordinate (D)	
				(-1,1) coordinate (A)
				($(A)+(0,3)$) coordinate (A')
				($(B)+(0,3)$) coordinate (B')
				($(C)+(0,3)$) coordinate (C')
				($(D)+(0,3)$) coordinate (D')
				($(D)!.5!(A)$) coordinate (I)
				;
				\draw (A)--(B)--(C)--(D) (A)--(A') (B)--(B') (C)--(C') (A')--(B')--(C')--(D')--(A') (D)--(D');
				\draw[dashed] (A)--(D)  (A)--(C') (A)--(C)--(I);
				\draw pic[draw, angle radius=2mm]{right angle=D--A--B};
				\draw pic[draw, angle radius=2mm]{right angle=A--B--C};
				\path (A)--(A') node[left,midway]{$2a$};
				\path (A)--(I) node[above,midway]{$a$};
				\path (D)--(I) node[above,midway]{$a$};
				\path (C)--(I) node[left,midway]{$a$};
				\path (A)--(B) node[left,midway]{$a$};
				\path (B)--(C) node[below,midway]{$a$};
				\draw pic[draw, angle radius=2mm]{right angle=C--I--D};
				\foreach \x/\g in {A/180,B/-90,C/-90,D/30,A'/90,B'/70,C'/90,D'/30,I/90} \fill[black] (\x) circle (1pt)+(\g:.3) node {$\x$};
			\end{tikzpicture}
		}
		\begin{enumerate}{\setcounter{enumi}{1}}
			\item 	Diện tích $S_{ABB'A'}=2a\cdot a=2a^2$; 
			Diện tích $S_{ABCD}=S_{A'B'C'D'}=\dfrac{(2a+a)\cdot a}{2}=\dfrac{3}{2}a^2$; 
			Diện tích $S_{BCC'B'}=2a\cdot a=2a^2$; Diện tích $S_{ADD'A'}=2a\cdot 2a=4a^2$.\\
			Diện tích $S_{CDD'C'}=2a\cdot a\sqrt{2}=2a^2\sqrt{2}$.\\
			Vậy tổng diện tích các mặt của hình lăng trụ bằng $$2a^2+3a^2+2a^2+4a^2+2a^2\sqrt{2}=(11+2\sqrt{2})a^2.$$
		\end{enumerate}
	}
\end{bt}

\begin{bt}%[1H8H7-5]
% nguon: Archive/2024/TOAN_11_XE/TOAN_11_XE/data4/11C7/KNTT-7.7.tex #1403
%[TeX hóa SGK CTST]%[Thành Đức Trung]
	\begin{enumerate}
		\item[a)] Tính thể tích của khối hộp chữ nhật có ba kích thước là: $6a ; 4a ; 3a$.
		\item[b)] Tính thể tích khối tứ diện đều $SABC$ cạnh $a$.
		\item[c)] Cho khối lăng trụ $ABCD. A'B'C'D'$  có đáy là hình vuông cạnh $2a$, cạnh bên $AA'=2 a$, hình chiếu của $A'$ trên $(ABCD)$ trùng với giao điểm $O$ của $AC$ và $BD$. Tính thể tích khối lăng trụ đó.
	\end{enumerate}
	\loigiai{
		\begin{enumerate}
			\item[a)] Thể tích khối hộp chữ nhật bằng tích ba kích thước \[V=6a \cdot 4a \cdot 3a= 72 a^3.\]
			\item[b)] Gọi $H$ là hình chiếu vuông góc của $S$ xuống $(ABC)$.\\ Ta có ba tam giác vuông $SHA$, $SHB, SHC$ bằng nhau, suy ra $HA=HB=HC$.\\ Vậy $H$ là tâm của tam giác đều $ABC$.
			\immini{
				Ta có
				$$
				\begin{aligned}
					& AM=\dfrac{a\sqrt{3}}{2}, AH=\dfrac{2}{3} AM=\dfrac{a\sqrt{3}}{3}, \\
					& SH=\sqrt{SA^2-AH^2}=\sqrt{a^2-\dfrac{3a^2}{9}}=\dfrac{a\sqrt{6}}{3} .
				\end{aligned}
				$$
				Khối tứ diện đều $SABC$ có thể tích là
				\[V=\dfrac{1}{3} S_{ABC} \cdot SH=\dfrac{1}{3} \cdot \dfrac{a^2 \sqrt{3}}{4} \cdot \dfrac{a \sqrt{6}}{3}=\dfrac{a^3 \sqrt{2}}{12} .\]}
			{\begin{tikzpicture}[scale=.8,>=stealth, font=\footnotesize, line join=round, line cap=round]
					\path(0,0)coordinate(A) (1.2,-1.6)coordinate(B) (4.5,0)coordinate(C)(barycentric cs:B=1,C=1)coordinate(M)
					(barycentric cs:A=1,B=1,C=1)coordinate(H)+(0,3.5)coordinate(S);
					\draw[dashed](M)--(A)--(C)(S)--(H);
					\draw(S)--(A)--(B)--(S)--(C)--(B) (S)--(M);
					\pic[draw,angle radius=6]{right angle=S--H--M};
					\foreach\p/\g in{A/180,B/-90,C/0,S/90,H/-90,M/-30}\fill(\p)circle(1.2pt)node[shift={(\g:.2)}]{$\p$};		
			\end{tikzpicture}}
			\item[c)] \immini{Chiều cao của khối lăng trụ
				\[h=A'O=\sqrt{A'A^2-A O^2}=\sqrt{4 a^2-2 a^2}=a \sqrt{2}.\]
				Thể tích khối lăng trụ \[V=Sh=4a^2 \cdot a \sqrt{2}=4 a^3 \sqrt{2}.\]}
			{	\begin{tikzpicture}[scale=0.8, font=\footnotesize, line join=round, line cap=round,>=stealth]
					\path (0,3) coordinate (h) (0,0)coordinate (B) (3.5,0) coordinate (C) (5,1) coordinate (D) ($(B)+(D)-(C)$) coordinate (A) ($(B)!0.5!(D)$) coordinate (O) ($(O)+(h)$) coordinate (A') ($(B)!0.7!(O)$) coordinate (H);
					\foreach \p in {B,C,D}\path ($(\p)+(A')-(A)$) coordinate (\p');
					\draw (A')--(B')--(B)--(C)--(D)--(D')--cycle (B')--(C')--(D') (C)--(C');
					\draw[dashed] (B)--(A)--(D)--cycle (C)--(A)--(A')--(O);
					\foreach \p/\g in {A/150,B/240,C/-60,D/0,A'/90,B'/180,C'/-30,D'/45,O/-90} \fill[black] (\p) circle(1pt)+(\g:0.3) node{$\p$};
					\draw pic[draw, angle radius=0.2cm]{right angle= A'--O--A} pic[draw, angle radius=0.2cm]{right angle= B--A--D};
			\end{tikzpicture}}
	\end{enumerate}}
\end{bt}

\begin{bt}%[1H8H7-6]
% nguon: Archive/2024/TOAN_11_XE/TOAN_11_XE/data4/11C7/KNTT-7.7.tex #1404
%[TeX hóa SGK CTST]%[Thành Đức Trung]
	Cắt khối chóp tam giác đều $S.ABC$ với cạnh đáy bằng $a$ và chiều cao $2 a$ bởi một mặt phẳng song song với đáy và đi qua trung điểm các cạnh bên. Tính thể tích khối chóp cụt đều được tạo thành.
	\loigiai{
		\immini{Gọi $ABC. A'B'C'$  là khối chóp cụt đều được tạo thành, $O$ và $O'$ lần lượt là tâm của hai đáy. 
			Chiều cao của khối chóp cụt đều là $h=OO'=\dfrac{SO}{2}=\dfrac{2a}{2}=a$;\\
			Tam giác đều $ABC$ có diện tích $S=\dfrac{AB^2 \sqrt{3}}{4}=\dfrac{a^2 \sqrt{3}}{4}$;\\
			Tam giác đều $A'B'C'$ có cạnh $A'B'=\dfrac{A B}{2}$ nên có diện tích $S'=\dfrac{AB^2 \sqrt{3}}{16}=\dfrac{S}{4}$.\\
			Do đó, thể tích khối chóp cụt đều được tạo thành là}
		{\begin{tikzpicture}[scale=1,>=stealth, font=\footnotesize, line join=round, line cap=round]
				\path(0,0)coordinate(A) (1.2,-1.6)coordinate(B) (4.5,0)coordinate(C)(barycentric cs:B=1,C=1)coordinate(M)
				(barycentric cs:A=1,B=1,C=1)coordinate(O)+(0,3.5)coordinate(S);
				\coordinate (A') at ($(A)!0.5!(S)$);
				\coordinate (B') at ($(B)!0.5!(S)$);
				\coordinate (C') at ($(C)!0.5!(S)$);
				\coordinate (O') at ($(O)!0.5!(S)$);
				\coordinate (N) at ($(B')!0.5!(C')$);
				\draw[dashed](M)--(A)--(C)(S)--(O) (A')--(C') (A')--(N);
				\draw(S)--(A)--(B)--(S)--(C)--(B) (A')--(B') (B')--(C');
				\pic[draw,angle radius=6]{right angle=S--O--M};
				\foreach\p/\g in{A/180,B/-90,C/0,S/90,O/-90,A'/180,B'/180,C'/0,O'/30}\fill(\p)circle(1pt)node[shift={(\g:.3)}]{$\p$};		
		\end{tikzpicture}}
		\[
		V=\dfrac{1}{3}h\left(S+\sqrt{SS'}+S'\right)=\dfrac{1}{3} a\left(S+\dfrac{S}{2}+\dfrac{S}{4}\right)=\dfrac{7aS}{12}=\dfrac{7a}{12} \cdot \dfrac{a^2\sqrt{3}}{4}=\dfrac{7a^3\sqrt{3}}{48}.
		\]}
\end{bt}

\begin{bt}%[1H8H7-9]
% nguon: Archive/2024/TOAN_11_XE/TOAN_11_XE/data4/11C7/KNTT-7.7.tex #1405

	Từ một tấm bìa hình vuông (Hình $92a$), người ta cắt ở bốn góc của tấm bìa đó bốn hình vuông bằng nhau, mỗi hình vuông có cạnh bằng $6$ cm, rồi gập tấm bìa lại để được một chiếc hộp không nắp có dạng hình hộp chữ nhật (Hình $92b$). Tính cạnh của tấm bìa ban đầu, biết rằng thể tích của chiếc hộp bằng $600$ cm$^3$.
	\begin{center}
		\begin{tikzpicture}[scale=1, font=\footnotesize, line join=round, line cap=round,>=stealth]
			\begin{scope}
				\def\a{4};
				\def\k{0.2};
				\pgfmathsetmacro{\m}{\k*\a};
				\path (0,0)coordinate (A)--++(\a,0) coordinate (B)--++(0,\a) coordinate (C)--++(-\a,0) coordinate (D) ($(A)!\k!(B)$) coordinate (a) ($(B)!\k!(C)$) coordinate (b) ($(C)!\k!(D)$) coordinate (c) ($(D)!\k!(A)$) coordinate (d) ($(B)!\k!(A)$) coordinate (a1) ($(C)!\k!(B)$) coordinate (b1) ($(D)!\k!(C)$) coordinate (c1) ($(A)!\k!(D)$) coordinate (d1);
				\draw (A)--(B)--(C)--(D)--cycle;
				\draw[dashed] (a)--(c1) (a1)--(c) (b)--(d1) (b1)--(d);
				\fill[cyan] (A)--++(\m,0)--++(0,\m)--++(-\m,0)--cycle;
				\fill[cyan] (B)--++(0,\m)--++(-\m,0)--++(0,-\m)--cycle;
				\fill[cyan] (C)--++(-\m,0)--++(0,-\m)--++(\m,0)--cycle;
				\fill[cyan] (D)--++(0,-\m)--++(\m,0)--++(0,\m)--cycle;
				\draw ($(A)!0.5!(B)$)++(-90:0.5)node{Hình $92a$};
			\end{scope}
			\begin{scope}[xshift=8cm]
				\path (0,1) coordinate (h) (0,0)coordinate (B) (2.7,0) coordinate (C) (4,0.7) coordinate (D) ($(B)+(D)-(C)$) coordinate (A);
				\foreach \p in {A,B,C,D}\path ($(\p)+(h)$) coordinate (\p');
				\draw (A')--(B')--(B)--(C)--(D)--(D')--cycle (B')--(C')--(D') (C)--(C');
				\draw[dashed] (B)--(A)--(D) (A)--(A');
				\draw ($(B)!0.5!(D)$)++(-90:0.9)node{Hình $92b$};
			\end{scope}
		\end{tikzpicture}
	\end{center}
	\loigiai
	{Giả sử tấm bìa ban đầu có cạnh $x$ (cm) ($x>12$). Khi đó,
		đáy của chiếc hộp là hình vuông cạnh $x-12$ (cm) nên diện tích của đáy chiếc hộp là $(x-12)^2$ $\left(\mathrm{cm}^2\right)$.\\
		Mà chiều cao của chiếc hộp là $6$ cm, suy ra thể tích của chiếc hộp bằng $(x-12)^2 \cdot 6$ $\left(\mathrm{cm}^3\right)$.\\
		Theo đề bài, thể tích của chiếc hộp bằng $600$ cm$^3$ nên
		$$(x-12)^2\cdot 6=600\Leftrightarrow(x-12)^2=100.$$
		Với $x>12$ ta có $x-12=10 \Leftrightarrow x=22$ (cm).
		Vậy độ dài cạnh của tấm bìa ban đầu là $22$ cm.
	}
\end{bt}

\begin{bt}%[1H8H7-3]
% nguon: Archive/2024/TOAN_11_XE/TOAN_11_XE/data4/11C7/KNTT-7.7.tex #1406

	\immini{Tính thể tích của khối chóp $S.ABCD$. Biết đáy $ABCD$ là hình vuông cạnh $a$, $SA \perp (ABCD)$, góc giữa đường thẳng $SB$ và mặt phẳng $(ABCD)$ bằng $60^{\circ}$ (Hình $94$).}{
		\begin{tikzpicture}[scale=1, font=\footnotesize, line join=round, line cap=round,>=stealth]
			\path (0,0)coordinate (B) (3,0) coordinate (C) (4.5,1) coordinate (D) ($(B)+(D)-(C)$) coordinate (A) (0,2.5) coordinate (h) ($(A)+(h)$) coordinate (S);
			\draw (S)--(B)--(C)--(D)--cycle (S)--(C);
			\draw[dashed] (S)--(A)--(B) (A)--(D);
			\foreach \p/\g in {A/150,B/240,C/-60, D/0, S/90} \fill[black] (\p) circle(1pt)+(\g:0.3) node{$\p$};
			\draw pic[draw,"$60^\circ$", angle eccentricity=1.5, angle radius=0.7cm]{angle=A--B--S};
			\draw ($(B)!0.5!(D)$)++(-90:1.3)node{Hình $94$};
	\end{tikzpicture}}
	\loigiai
	{Do $SA \perp (ABCD)$ và $BA$ chứa trong $(ABCD)$ nên $SA \perp AB$, suy ra $S A=A B \cdot \tan \widehat{SBA}$.\\
		Vì $SA \perp (ABCD)$ nên góc giữa đường thẳng $SB$ và mặt phẳng $(ABCD)$ bằng $\widehat{SBA}$. Suy ra $\widehat{SBA}=60^{\circ}$.\\
		Từ đó, ta có $SA=a\cdot\tan 60^{\circ}=a\sqrt{3}$. Mặt khác, do diện tích hình vuông $ABCD$ là $S_{ABCD}=a^2$ nên thể tích của khối chóp $S.ABCD$ là
		$$V_{S.ABCD}=\dfrac{1}{3}S_{ABCD}\cdot SA=\dfrac{1}{3}\cdot a^2\cdot a\sqrt{3}=\dfrac{a^3\sqrt{3}}{3}.$$
	}
\end{bt}

\begin{bt}%[1H8V7-6]
% nguon: Archive/2024/TOAN_11_XE/TOAN_11_XE/data4/11C7/KNTT-7.7.tex #1407

	Tính thể tích của khối chóp cụt lục giác đều $ABCDEF \cdot A'B'C'D'E'F'$ với $O$ và $O'$ là tâm hai đáy, cạnh đáy lớn và đáy nhỏ lần lượt là $a$ và $\dfrac{a}{2}$, $OO'=a$.
	\loigiai{
		\immini{Diện tích đáy lớn $S=6\cdot a^2\dfrac{\sqrt{3}}{4}=\dfrac{3\sqrt{3}}{2}a^2$.\\
			Diện tích đáy lớn $S'=6\cdot \left(\dfrac{a}{2}\right)^2\dfrac{\sqrt{3}}{4}=\dfrac{3\sqrt{3}}{8}a^2$.\\	
			Thể tích của khối chóp cụt lục giác đều $ABCDEF \cdot A'B'C'D'E'F'$ là 
			\allowdisplaybreaks
			\begin{eqnarray*}
				V&=&\dfrac{1}{3}OO'\cdot \left(S+\sqrt{S\cdot S'}+S'\right)\\
				&=&\dfrac{1}{3}\cdot a \cdot \left(\dfrac{3\sqrt{3}}{2}a^2+\sqrt{\dfrac{3\sqrt{3}}{2}a^2\cdot \dfrac{3\sqrt{3}}{8}a^2}+\dfrac{3\sqrt{3}}{8}a^2\right)\\
				&=&	\dfrac{7\sqrt{3}}{8}a^3.
		\end{eqnarray*}}
		{\begin{tikzpicture}[scale=1, font=\footnotesize,line join=round, line cap=round, >=stealth];
				\clip (-3,-1.5) rectangle (5,3.8);
				\def\r{2.25};
				\def\rr{1.5}
				\def\a{45}
				\def\h{7}
				\path 
				(0,0) coordinate (O)
				(0,\h) coordinate (S)
				(-\r,0) coordinate (A)
				(\r,0) coordinate (D)
				(\a:\rr) coordinate (E)
				(-\a:\rr) coordinate (C)
				($(C)!2!(O)$) coordinate (F)
				($(E)!2!(O)$) coordinate (B)
				($(S)!0.6!(O)$) coordinate (O')
				($(S)!0.6!(A)$) coordinate (A')
				($(S)!0.6!(B)$) coordinate (B')
				($(S)!0.6!(C)$) coordinate (C')
				($(S)!0.6!(D)$) coordinate (D')
				($(S)!0.6!(E)$) coordinate (E')
				($(S)!0.6!(F)$) coordinate (F')
				($(A')!0.5!(B')$) coordinate (I')
				($(A)!0.5!(B)$) coordinate (I)
				;
				\draw 
				(A)--(B)--(C)--(D) 
				(A')--(B')--(C')--(D')--(E')--(F')--cycle
				(A)--(A') (B)--(B') (C)--(C') (D)--(D')
				;
				\draw[dashed] 
				(A)--(F)--(E)--(D)
				(E)--(E') (F)--(F')
				(O)--(O')
				;
				\foreach \l/\g in {A/180,B/-135,C/-45,D/0,E/45,F/135,O/-90,O'/90,A'/180,B'/45,C'/135,D'/0,E'/45,F'/135}
				\draw[fill=black] (\l) circle (1pt) +(\g:.3) node{$\l$};
		\end{tikzpicture}}
	}
\end{bt}

\begin{bt}%[1H8H7-6]
% nguon: Archive/2024/TOAN_11_XE/TOAN_11_XE/data4/11C7/KNTT-7.7.tex #1408

	Cho khối chóp cụt tam giác đều $ABC.A'B'C'$ có chiều cao bằng $3a$, $AB=4a$, $A'B'=a$ (Hình $95$). Tính thể tích của khối chóp cụt đều $ABC.A'B'C'$.
	\loigiai
	{\immini{Diện tích tam giác đều $A B C$ là
			$$S_1=\dfrac{1}{2}AB\cdot AC\cdot\sin \widehat{BAC}=\dfrac{1}{2}\cdot 4a\cdot 4a\cdot\sin 60^{\circ}=4\sqrt{3}a^2.$$
			Diện tích tam giác đều $A'B'C'$ là
			$$S_2=\dfrac{1}{2}A'B'\cdot A'C'\cdot\sin \widehat{B'A'C'}=\dfrac{1}{2}\cdot a\cdot a\cdot\sin 60^{\circ}=\dfrac{\sqrt{3}a^2}{4}.$$
			Thể tích khối chóp cụt đều $ABC.A'B'C'$ là
			$$V=\dfrac{1}{3}\cdot 3a\cdot\left(4\sqrt{3}a^2+\sqrt{4\sqrt{3}a^2\cdot\dfrac{\sqrt{3}a^2}{4}}+\dfrac{\sqrt{3}a^2}{4}\right)
			=\dfrac{21\sqrt{3}a^3}{4}.$$}{
			\begin{tikzpicture}[scale=1, font=\footnotesize, line join=round, line cap=round,>=stealth]
				\def \k{0.6};
				\path (0,0)coordinate (A) (3.5,0) coordinate (C) (2.5,-1) coordinate (B) ($(B)!0.5!(C)$) coordinate (M) ($(A)!2/3!(M)$) coordinate (O) (0,3) coordinate (h) ($(O)+(h)$) coordinate (S);
				\foreach \p in {A,B,C}{\path ($(\p)!\k!(S)$) coordinate (\p');}
				\draw (A')--(A)--(B)--(C)--(C')--cycle (B')--(C') (A')--(B')--(B);
				\draw[dashed] (A)--(C);
				\foreach \p/\g in {A/150,B/-60,C/30, A'/150, B'/-30, C'/30} \fill[black] (\p) circle(1pt)+(\g:0.3) node{$\p$};
				\draw ($(A)!0.5!(C)$)++(-90:1.6)node{Hình $95$};
		\end{tikzpicture}}
	}
\end{bt}

\begin{bt}%[1H8H5-5]
% nguon: Archive/2024/TOAN_11_XE/TOAN_11_XE/data4/11C7/KNTT-7.6.tex #1409
%[Nguyễn Thành Khang, dự án Tex hoá lần 3]%[Đinh Thanh Hoàng]
	Cho tứ diện $OABC$ có $OA$, $OB$, $OC$ vuông góc với nhau đôi một và $OA = OB = OC = a$. Gọi $I$ là trung điểm của $BC$. Hãy dựng và tính độ dài đoạn vuông góc chung của
	\begin{enumEX}{2}
		\item $OA$ và $BC$.
		\item $AI$ và $OC$.
	\end{enumEX}
	\loigiai{
		\immini{
			\begin{enumerate}
				\item Ta có $\heva{&OA \perp OI\\ &BC \perp OI}$ $\Rightarrow OI$ là đoạn vuông góc chung của $OA$ và $BC$.\\
				Tam giác $OBC$ vuông cân tại $O$ nên $OI = \dfrac{BC}{2} = \dfrac{a\sqrt{2}}{2}$.\\
				Vậy $\mathrm{d}(OA, BC) = \dfrac{a\sqrt{2}}{2}$.
				\item Gọi $K$ là trung điểm của $OB$, ta có $IK \parallel OC$ $\Rightarrow OC \parallel (AIK)$.\\
				Ta có $\heva{&IK \perp OB\\ &IK \perp OA}$ $\Rightarrow IK \perp (OAB)$ $\Rightarrow (AIK) \perp (OAB)$ theo giao tuyến $AK$.
			\end{enumerate}
		}{
			\begin{tikzpicture}[scale=.6, font=\footnotesize, line join=round, line cap=round,>=stealth]
				\tkzDefPoints{0/0/O,0/5/A,5/-3/B,7/0/C}
				\tkzDefMidPoint(B,C) \tkzGetPoint{I}
				\tkzDefMidPoint(O,B) \tkzGetPoint{K}
				\tkzDefPointBy[homothety = center A ratio 0.6](K) \tkzGetPoint{H}
				\tkzDefPointBy[homothety = center A ratio 0.6](I) \tkzGetPoint{E}
				\tkzDefPointBy[translation = from H to O](E) \tkzGetPoint{F}
				\tkzDrawSegments(O,A O,B B,C A,K A,B A,I A,C O,H)
				\tkzDrawSegments[dashed](O,C K,I H,E E,F O,I)
				\tkzMarkRightAngles[size=.3](A,O,B A,O,C B,O,C O,H,K F,E,I E,F,C I,K,B)
				\tkzDrawPoints[fill=black](A,B,C,O,I,K,H,E,F)
				\tkzLabelPoints[above](A)
				\tkzLabelPoints[below](B)    		
				\tkzLabelPoints[right](E,I,C)
				\tkzLabelPoints[left](O,H)
				\tkzLabelPoints[below left](K)
				\tkzLabelPoints[below right](F)
			\end{tikzpicture}
		}
		\begin{enumerate}
			\item[] Trong mặt phẳng $(OAB)$, kẻ $OH \perp AK$ tại $H$, ta có: $OH \perp (AIK)$ $\Rightarrow OH \perp AI$.\\
			Trong mặt phẳng $(AIK)$, kẻ $HE \parallel IK$ $(E \in SI)$.\\
			Trong mặt phẳng $(HE, OC)$, kẻ $EF \parallel OH$ $(F \in OC)$ $\Rightarrow EF \perp AI$.\\
			Lại có $OC \perp (OAB) \Rightarrow OC \perp AH \Rightarrow OC \perp EF$.\\
			Do đó $EF$ là đoạn vuông góc chung của $OC$ và $AI$.\\
			Trong tam giác vuông $OAK$ ta có: $\dfrac{1}{OH^2} = \dfrac{1}{OA^2} + \dfrac{1}{OK^2} = \dfrac{1}{a^2} + \dfrac{4}{a^2} = \dfrac{5}{a^2}$.\\
			Suy ra $EF = OH = \dfrac{a\sqrt{5}}{5}$.\\
			Vậy $\mathrm{d}(AI, OC) = \dfrac{a\sqrt{5}}{5}$.
		\end{enumerate}
	}
\end{bt}

\begin{bt}%[1H8H5-4]
% nguon: Archive/2024/TOAN_11_XE/TOAN_11_XE/data4/11C7/KNTT-7.6.tex #1410
%[Nguyễn Thành Khang, dự án Tex hoá lần 3]%[Đinh Thanh Hoàng]
	Cho hình chóp $S.ABCD$ có đáy là hình vuông cạnh bằng $a$, $SA$ vuông góc với đáy và $SA = a$. $M$ là trung điểm của $SB$. Tính khoảng cách giữa các đường thẳng
	\begin{enumEX}{3}		
		\item $SC$ và $BD$.
		\item $AC$ và $SD$.
		\item $SD$ và $AM$.
	\end{enumEX}
	\loigiai{
		\begin{center}
			\begin{tikzpicture}[line join=round,scale=0.7]
				\tkzDefPoints{0/0/A,-3/-3/D,3/-3/C,6/0/B,0/5/S}
				\tkzDefMidPoint(S,B) \tkzGetPoint{M}
				\tkzInterLL(A,C)(B,D) \tkzGetPoint{O}
				\tkzDefPointBy[translation = from C to D](A) \tkzGetPoint{E}		
				\tkzDefPointBy[translation = from O to D](A) \tkzGetPoint{I}
				\tkzDefPointBy[homothety = center S ratio 0.7](C) \tkzGetPoint{H}
				\tkzDefPointBy[homothety = center S ratio 0.6](I) \tkzGetPoint{K}
				\tkzDrawSegments(S,E S,D S,C S,B C,D B,C M,C D,E S,I)
				\tkzDrawSegments[dashed](S,A A,D A,B A,K A,E A,M O,M O,H A,C B,D A,I)
				\tkzMarkRightAngles[size=.3](S,A,B S,A,D B,A,D A,I,D A,K,I O,H,C D,O,H)
				\tkzDrawPoints[fill=black](A,B,C,D,E,S,I,O,M,H,K) 
				\tkzLabelPoints(C)
				\tkzLabelPoints[above](S)
				\tkzLabelPoints[below](O)    		
				\tkzLabelPoints[right](H,B)
				\tkzLabelPoints[left](E,K,I)
				\tkzLabelPoints[below left](D)
				\tkzLabelPoints[above left](A)
				\tkzLabelPoints[above right](M)
			\end{tikzpicture}
		\end{center}
		\begin{enumerate}
			\item Gọi $O$ là giao điểm của $AC$ và $BD$. Ta có: $\heva{&BD \perp SA\\ &BD \perp AC}$ $\Rightarrow BD \perp (SAC)$ tại $O$.\\
			Trong mặt phẳng $(SAC)$, kẻ $OH \perp SC$ tại $H$, ta có $OH \perp SC$ và $OH \perp BD$ (vì $BD \perp (SAC)$).\\
			Vậy $OH$ là đoạn vuông góc chung của $BD$ và $SC$.\\
			Ta có $\dfrac{OH}{OC} = \dfrac{SA}{SC} = \sin \widehat{ACS}$ $\Rightarrow OH = \dfrac{OC.SA}{SC} = \dfrac{\dfrac{a\sqrt{2}}{2}\cdot a}{a\sqrt{3}} = \dfrac{a\sqrt{6}}{6}$.\\
			Vậy $\mathrm{d}(SC, BD) = OH = \dfrac{a\sqrt{6}}{6}$.
			\item Dựng hình bình hành $ACDE$, ta có: $AC \parallel DE$ $\Rightarrow AC \parallel (SDE)$.\\
			$\Rightarrow \mathrm{d}(AC, SD) = \mathrm{d}(AC, (SDE)) = \mathrm{d}(A, (SDE))$.\\
			Trong mặt phẳng $(ABCD)$, kẻ $AI \perp DE$ tại $I$, ta có $\heva{&DE \perp AI\\ &DE \perp SA}$ $\Rightarrow DE \perp (SAI)$.\\
			$\Rightarrow (SDE) \perp (SAI)$ theo giao tuyến $SI$.\\
			Trong mặt phẳng $(SAI)$, kẻ $AK \perp SI$ tại $K$, ta có: $AK \perp (SDE)$ $\Rightarrow AK = \mathrm{d}(A, (SDE))$.\\
			Ta có $AIDO$ là hình bình hành nên $AI = OD = \dfrac{a\sqrt{2}}{2}$.\\
			Trong tam giác vuông $SAI$ ta có: $\dfrac{1}{AK^2} = \dfrac{1}{AI^2} + \dfrac{1}{SA^2} = \dfrac{2}{a^2} + \dfrac{1}{a^2} = \dfrac{3}{a^2}$.\\
			Suy ra $AK = \dfrac{a\sqrt{3}}{3}$.\\
			Vậy $\mathrm{d}(AC, SD) = \mathrm{d}(A, (SDE)) = AK = \dfrac{a\sqrt{3}}{3}$.
			\item Ta có $OM \parallel SD$ và $AC \parallel DE$ nên $(AMC) \parallel (SDE)$.\\
			Suy ra $\mathrm{d}(SD, AM) = \mathrm{d}((AMC), (SDE)) = \mathrm{d}(A, (SDE)) = AK = \dfrac{a\sqrt{3}}{3}$.
		\end{enumerate}
	}
\end{bt}

\begin{bt}%[1H8V7-9]
% nguon: Archive/2024/TaiLieuGiangDay/5.HH Oxyz/VDC_TDM.tex #2416

	Cho một hình lăng trụ tam giác đều $A B C .A'B'C'$ có độ dài cạnh đáy bằng $12$ và cạnh bên bằng $16$. Dùng một sợi dây có chiều dài $L$ quấn quanh hình lăng trụ từ điểm $A'$ đến điểm $A$ như hình vẽ, sao cho sợi dây luôn áp sát vào các mặt bên. Hãy xác định chiều dài ngắn nhất của sợ dây?
	\loigiai{
		\begin{center}
			\begin{tikzpicture}[scale=1, font=\footnotesize,>=stealth]
				%Gán số liệu.
				\def\canhA'C'{4};\def\canhB'A'{2};\def\gocB'A'C'{-50};\def\h{3};\def\xdinhA{0};
				%Gán tọa độ.
				\coordinate (A') at (0,0);
				\coordinate (B') at ($(A') + (\gocB'A'C':\canhB'A')$);
				\coordinate (C') at ($(A') + (0:\canhA'C')$);
				\coordinate (A) at ($(A') + (\xdinhA,\h)$);
				\coordinate (B) at ($(B') + (\xdinhA,\h)$);
				\coordinate (C) at ($(C') + (\xdinhA,\h)$);
				\coordinate (M) at ($(B)!0.7!(B')$);
				\coordinate (N) at ($(C)!0.3!(C')$);
				%Vẽ khối lăng trụ A'B'C'.ABC.
				\draw (A)--(B)--(C)--cycle (B')--(B) (A)--(A')--(B') (C)--(C')--(B') (A')--(M)--(N);
				\draw[dashed] (A')--(C') (A)--(N);
				%Gán nhãn.
				\foreach \x/\y in {A'/180, B'/-90, C'/0, A/180, B/90, C/0, M/-45, N/0}{\fill (\x) circle(1pt) ($(\x)+(\y:0.3cm)$) node{$\x$};}
			\end{tikzpicture}	
		\end{center}
Ta thực hiện việc trải phẳng các mặt bên của hình lăng trụ, ta được hình sau		
		\begin{center}
\begin{tikzpicture}[scale=1, font=\footnotesize,>=stealth]
	\coordinate (A_1') at (0,0);
	\coordinate (A_1) at (0,4);
	\coordinate (B') at (2,0);
	\coordinate (C') at (4,0);
	\coordinate (A_2') at (6,0);
	\coordinate (B) at (2,4);
	\coordinate (C) at (4,4);
	\coordinate (A_2) at (6,4);
	\coordinate (M) at ($(B)!0.3!(B')$);
	\coordinate (N) at ($(C)!0.7!(C')$);;
	\draw (A_1') rectangle (A_2) (B)--(B') (C)--(C') (A_1')--(A_2) (A_1')--node[below]{$12$}(B') --node[below]{$12$}(C')--node[below]{$12$} (A_2')--node[right]{$16$}(A_2);
	\draw [dashed] (A_1')--(M)--(N)--(A_2);
	%Gán nhãn.
	\foreach \x/\y in {A_1'/-90, B'/-90,C'/-90, A_2'/-90, A_1/90, B/90, C/90,A_2/90, M/180, N/0}{\fill (\x) circle(1pt) ($(\x)+(\y:0.3cm)$) node{$\x$};}
\end{tikzpicture}
		\end{center}
		Ta có $A_1'A_2=\sqrt{(12+12+12)^2+16^2}=4\sqrt{97}$.\\
		Ta có $L=A'M+MN+NA_2=A_1'M+MN+NA_2\geq A_1'A_2=4\sqrt{97}$.\\
		Dấu \lq\lq$=$\rq\rq\, xảy ra khi $M=A_1'A_2\cap BB'$, $N=A_1'A_2\cap CC'$.
	}
\end{bt}

\begin{bt}%[1H8H4-5]
% nguon: Archive/SP_DA_SauLuoi/1-TL-TN-THPT-NgoQuyen-ThaiNguyen-HKII-NH23-24.tex #2430
%[Dự án đề kiểm tra Toán 11 HKII NH23-24- Nguyễn Cường]%[THPT Ngô Quyền - Thái Nguyên]
Cho hình chóp $S.ABC$ đáy là tam giác vuông tại $B$, $SA$ vuông góc với đáy. Chứng minh mặt phẳng $(SBA)$ vuông góc với mặt phẳng $(SBC)$.
\loigiai{
\begin{center}
	\begin{tikzpicture}[scale=1, font=\footnotesize, line join=round, line cap=round, >=stealth]
			\path
			(0,0) coordinate (A)++(0:4)coordinate(C)++(230:2)coordinate(B)
			(A)++(90:3)coordinate (S)
			;
			\draw (B)--(S)--(A)--(B)--(C)--(S);
			\draw [dashed] (A)--(C);
			\foreach \x/\y/\z in {A/B/C}{
				\draw[thin] pic[draw, angle radius = 6pt]{right angle = \x--\y--\z};
			}
			\foreach \p/\g in {A/180,B/-90,C/0,S/90}\fill[black] (\p) circle (1pt) +(\g:.3)node{$\p$};
	\end{tikzpicture}
\end{center}
Ta có $\heva{&BC\perp BA\\&BC\perp SA}\Rightarrow BC\perp (SAB)$.\\
Mà $BC\subset (SBC)$ nên $(SBC)\perp (SBA)$.
}
\end{bt}

\begin{bt}%[1H8V5-3]
% nguon: Archive/SP_DA_SauLuoi/1-TL-TN-THPT-NgoQuyen-ThaiNguyen-HKII-NH23-24.tex #2431
%[Dự án đề kiểm tra Toán 11 HKII NH23-24- Nguyễn Cường]%[THPT Ngô Quyền - Thái Nguyên]
Cho hình chóp $S.ABCD$ có đáy $ABCD$ là hình chữ nhật $AB=2a$ và $AD=a$, cạnh bên $SA$ vuông góc với đáy. Tính khoảng cách từ điểm $C$ đến mặt phẳng $(SBD)$, biết số đo góc $[S,BD,A]$ bằng $60^\circ$.
\loigiai{
\begin{center}
	\begin{tikzpicture}[scale=1, font=\footnotesize, line join=round, line cap=round, >=stealth]
		\path
		(0,0) coordinate (A)++(0:4)coordinate(D)++(230:2)coordinate(C)++(180:4)coordinate(B)
		(A)++(90:3)coordinate (S)
		($(B)!.4!(D)$)coordinate(H)
		($(B)!.5!(D)$)coordinate(O)
		($(S)!.6!(H)$)coordinate(I)
		;
		\draw (S)--(B)--(C)--(D)--cycle (S)--(C);
		\foreach \x/\y/\z in {A/H/B,S/H/D,A/I/H}{
			\draw[thin] pic[draw, angle radius = 6pt]{right angle = \x--\y--\z};
		}
		\draw [dashed] (B)--(A)--(D)--(B) (S)--(A)--(H)--(S) (C)--(A)--(I);
		\foreach \p/\g in {A/180,B/-90,C/-90,S/90,D/0,H/-90,I/0,O/-90}\fill[black] (\p) circle (1pt) +(\g:.3)node{$\p$};
	\end{tikzpicture}
\end{center}
Gọi $AH$ và $AI$ lần lượt là đường cao của tam giác $ABD$ và $SAH$.\\
Ta có $\heva{&BD\perp AH\\&BD\perp SA}\Rightarrow BD\perp (SAH)$.\\
Mà $SH\subset (SAH)\Rightarrow BD\perp SH$.\\
Do đó, góc $[S,BD,A]=\widehat{SHA}=60^\circ$.\\
Ta có $\heva{&AH\perp SH\\&AI\perp BD}\Rightarrow AH\perp (SBD)$, suy ra $\mathrm{d}(A,(SBD))=AI$.\\
$\triangle ABD$ có $\dfrac{1}{AH^2}=\dfrac{1}{AB^2}+\dfrac{1}{AD^2}=\dfrac{1}{4a^2}+\dfrac{1}{a^2}=\dfrac{5}{4a^2}$, suy ra $AH=\dfrac{2a\sqrt{5}}{5}$.\\
$\triangle AIH$ vuông tại $I$ có $\sin 60^\circ=\dfrac{AI}{AH}\Rightarrow AI=AH\cdot\sin 60^\circ=\dfrac{2a\sqrt{5}}{5}\cdot \dfrac{\sqrt{3}}{2}=\dfrac{a\sqrt{15}}{5}$.\\
Gọi $O$ là giao điểm $AC$ và $BD$, ta có $\dfrac{\mathrm{d}(C,(SBD))}{\mathrm{d}(A,(SBD))}=\dfrac{CO}{AO}=1$.\\
Vậy $\mathrm{d}(C,(SBD))=\mathrm{d}(A,(SBD))=AI=\dfrac{a\sqrt{15}}{5}$.
}
\end{bt}

\begin{bt}%[1H8C5-3]
% nguon: Archive/SP_DA_SauLuoi/1-TL-TN-THPT-NgoQuyen-ThaiNguyen-HKII-NH23-24.tex #2432
%[Dự án đề kiểm tra Toán 11 HKII NH23-24- Nguyễn Cường]%[THPT Ngô Quyền - Thái Nguyên]
Cho hình chóp $S.ABCD$ có đáy $ABCD$ là hình thang vuông tại $A$ và $D$; $AB=AD=2a$ và $DC=a$. Điểm $I$ là trung điểm đoạn thẳng $AD$, mặt phẳng $(SIB)$ và $(SIC)$ cùng vuông góc với mặt phẳng $(ABCD)$. Mặt phẳng $(SBC)$ tạo với mặt phẳng $(ABCD)$ một góc bằng $60^\circ$. Tính khoảng cách từ điểm $D$ đến $(SBC)$ theo $a$.
\loigiai{
\begin{center}
	\begin{tikzpicture}[scale=1, font=\footnotesize, line join=round, line cap=round, >=stealth]
		\path
		(0,0) coordinate (A)++(0:4)coordinate(B)++(200:4)coordinate(C)++(180:2)coordinate(D)
		($(A)!.5!(D)$)coordinate(I)
		(I)++(90:4)coordinate (S)
		($(A)!2!(D)$)coordinate(E)
		($(B)!.7!(C)$)coordinate(H)
%		(intersection of I--H and D--E)coordinate(F)
		($(S)!.6!(H)$)coordinate(K)
		;
		\draw (S)--(B)--(C)--(D)--cycle (H)--(S)--(C) (D)--(E)--(C);
		\draw [dashed] (B)--(A)--(D) (B)--(I)--(C) (I)--(S)--(A) (I)--(H) (D)--(E) (I)--(K);
		\foreach \x/\y/\z in {I/H/C,S/H/B,I/K/H}{
			\draw[thin] pic[draw, angle radius = 6pt]{right angle = \x--\y--\z};
		}
		\foreach \p/\g in {A/-45,B/0,C/-90,S/90,D/-90,H/-90,I/150,K/30,E/90}\fill[black] (\p) circle (1pt) +(\g:.3)node{$\p$};
	\end{tikzpicture}
\end{center}
Ta có $\heva{&(SIB)\perp (ABCD)\\&(SIC)\perp (ABCD)}\Rightarrow (SBI)\cap (SCI)=SI\perp (ABCD)$.\\
Gọi $IH$ và $IK$ lần lượt là đường cao của tam giác $IBC$ và $SIH$.\\
Ta có $\heva{&BC\perp IH\\&BC\perp SI}\Rightarrow BC\perp (SIH)$.\\
Mà $SH\subset (SIH)\Rightarrow BC\perp SH$.\\
Do đó, góc giữa hai mặt phẳng $(SBC)$ và $(ABCD)$ là $\widehat{SHI}=60^\circ$.\\
Ta có $\heva{&IK\perp SH\\&IK\perp BC}\Rightarrow IK\perp (SBC)$, suy ra $\mathrm{d}(I,(SBC))=IK$.\\
Gọi $E$ là giao điểm của $BC$ và $AD$.\\
Do $CD=\dfrac{AB}{2}$ nên $D$, $C$ lần lượt là trung điểm của $AE$ và $BE$.\\
Khi đó, $BC=CE=\sqrt{CD^2+DE^2}=\sqrt{a^2+4a^2}=a\sqrt{5}$.
\begin{itemize}
	\item $S_{ABCD}=\dfrac{(AB+CD)AD}{2}=\dfrac{(2a+a)\cdot 2a}{2}=3a^2$.
	\item $S_{ICD}=\dfrac{CD\cdot DI}{2}=\dfrac{a^2}{2}$.
	\item $S_{IAB}=\dfrac{AB\cdot AI}{2}=\dfrac{2a^2}{2}=a^2$.
\end{itemize}
Suy ra, $S_{IBC}=S_{ABCD}-S_{ICD}-S_{IAB}=3a^2-a^2-\dfrac{a^2}{2}=\dfrac{3a^2}{2}$.\\
Mà $S_{IBC}=\dfrac{BC\cdot IH}{2}\Rightarrow IH=\dfrac{2S_{IBC}}{BC}=\dfrac{3a^2}{a\sqrt{5}}=\dfrac{3a\sqrt{5}}{5}$.\\
$\triangle IHK$ vuông tại $K$ có $\sin 60^\circ=\dfrac{IK}{IH}\Rightarrow IK=IH\cdot \sin 60^\circ=\dfrac{3a\sqrt{5}}{5}\cdot\dfrac{\sqrt{3}}{2}=\dfrac{3a\sqrt{15}}{10}$.\\
Ta có $\dfrac{\mathrm{d}(D,(SBC))}{\mathrm{d}(I,(SBC))}=\dfrac{DE}{IE}=\dfrac{2}{3}$.\\
Do đó, $\mathrm{d}(D,(SBC))=\dfrac{2IK}{3}=\dfrac{a\sqrt{15}}{5}$.
}
\end{bt}

\begin{bt}%[1H8V5-3]
% nguon: Archive/SP_DA_SauLuoi/1-TN-TL-THPT-LeLoi-QuangTri-HKII-NH23-24.tex #2439
%[Dự án đề kiểm tra Toán 11 CHKII NH23-24-Nguuyễn Văn Sang]%[THPT Lê Lợi -Quảng Trị]
	Cho hình chóp $S.ABCD$ có cạnh bên $SA$ vuông góc với mặt phẳng đáy, $A B C D$ là hình vuông cạnh $a$. Biết góc giữa $S B$ và mặt đáy bằng $60^{\circ}$. Gọi $M$ là điểm nằm trên cạnh $CD$ sao cho $DM=2MC$. Tính khoảng cách từ điểm $M$ đến mặt phẳng $(SBD)$.
	\loigiai{
		\immini{
			Do $SA \perp(ABCD)$ nên $$\widehat{(SB,(ABCD))}=\widehat{SBA}=60^{\circ}.$$
			Tam giác $S A B$ vuông tại $A$ nên $$SA=AB \cdot \tan \widehat{SBA}=a \sqrt{3}.$$
			Gọi $I$ là giao điểm của $AM$ và $BD$. Ta có
			$$
			\dfrac{IM}{IA}=\dfrac{DM}{AB}=\dfrac{2}{3}=\dfrac{\mathrm{d}(M,(SBD))}{\mathrm{d}(A,(SBD))}.
			$$
			Suy ra $\mathrm{d}(M,(SBD))=\dfrac{2}{3} \mathrm{d}(A,(SBD))$.
		}
		{
			\begin{tikzpicture}[scale=1, font=\footnotesize, line join=round, line cap=round, >=stealth]
				\def\bc{4} \def\ba{2.25} \def\h{4} \def\gocB{30} 
				\coordinate[label=below left:$D$] (B) at (0,0);
				\coordinate[label=above left:$A$] (A) at (\gocB:\ba);
				\coordinate[label=below:$C$] (C) at (\bc,0);
				\coordinate[label=right:$B$] (D) at ($(C)-(B)+(A)$);
				\coordinate[label=above:$S$] (S) at ($(A)+(90:\h)$);
				\draw (B)--(C)--(D)--(S)--cycle (S)--(C);
				\draw[dashed] (A)--(D) (S)--(A)--(B);
				\foreach \diem in {A,B,C,D,S}	
				\fill (\diem)circle(1pt);
				\draw (A)--++(90:0.2)--++(0:0.2)--++(-90:0.2);
				\coordinate (M) at ($(C)!1/3!(B)$);
				\coordinate (O) at ($(C)!1/2!(A)$);
				\coordinate (H) at ($(O)!1/3.5!(S)$);
				\fill (M) circle (1pt) node[below]{$M$};
				\fill (O) circle (1pt) node[below]{$O$};
				\fill (H) circle (1pt) node[right]{$H$};
				\draw[dashed] (C)--(A)--(H) (S)--(O);
				% Giao điểm của AM và BD
				\draw[dashed,name path=AM] (A) -- (M);
				\draw[dashed,name path=BD] (B) -- (D);
				\path[name intersections={of=AM and BD, by=I}];
				\fill (I) circle (1pt) node[below]{$I$};
				\clip (A) -- (D) -- (S);
				\draw (D) circle (0.5cm);
				\draw ($(D)+(-0.25,0.25)$) node[left]{$60^\circ$};
			\end{tikzpicture}
		}\noindent
		Gọi $O$ là tâm đáy. Kẻ đường cao $A H$ của tam giác $SAO$.\\
		Khi đó $AH \perp(SBD)$.\\
		Suy ra $\mathrm{d}(A,(SBD))=AH=\dfrac{SA \cdot AO}{\sqrt{SA^2+AO^2}}=\dfrac{a\sqrt{3} \cdot \dfrac{a \sqrt{2}}{2}}{\sqrt{(a \sqrt{3})^2+\left(\dfrac{a \sqrt{2}}{2}\right)^2}}=\dfrac{a \sqrt{21}}{7}$.\\
		Do đó, $\mathrm{d}(M,(SBD))=\dfrac{2}{3} \mathrm{d}(A,(SBD))=\dfrac{2}{3} \cdot \dfrac{\sqrt{21}}{7}=\dfrac{2 a \sqrt{21}}{21}$.
	}
\end{bt}

\begin{bt}%[1H8V7-8]
% nguon: Archive/SP_DA_SauLuoi/1-TL-TN-DS-THPT-DamDoi-CaMau-HKII-NH23-24.tex #2443
%[Dự án đề kiểm tra Toán 11 GHKI NH23-24- Tacgia]%[THPT - Tp HCM]
	Cho hình chóp $S.ABCD$ có đáy $ABCD$ là hình vuông cạnh $a$, $SA \perp (ABCD)$, $SB=a\sqrt{5}$.
	\begin{enumerate}
		\item Tính thể tích khối chóp $S.ABCD$.
		\item Gọi $M$, $N$ lần lượt là trung điểm của $SB$, $SD$. Tính khoảng cách của hai đường thẳng chéo nhau $MD$ và $CN$.
	\end{enumerate}
	\loigiai{
		\begin{center}
			\begin{tikzpicture}[scale=1, font=\footnotesize,line join=round, line cap=round, >=stealth]
			\path 
			(0,0) coordinate (A)
			++(-140:2) coordinate (B)
			++(0:3.5) coordinate (C)
			($(A)+(C)-(B)$) coordinate (D)
			($(A)+(0,3)$) coordinate (S)
			($(S)!.5!(B)$) coordinate (M)
			($(S)!.5!(D)$) coordinate (N)
			;
			\foreach \i in{B,C,D}{\draw (S)--(\i);};
			\draw (B)--(C)--(D) (C)--(N);
			\draw[dashed] (S)--(A)--(B) (A)--(D) (M)--(D)  (M)--(N);
			\pic[draw,angle eccentricity=1.8,angle radius=2mm]{right angle=S--A--D};
			\foreach \i/\g in {A/-90,B/-90,C/-90,D/-90,S/90,M/180,N/90}
			\fill[black] (\i) circle(1pt)+(\g:3mm)node[scale=1]{$\i$};
		\end{tikzpicture}
		\end{center}
	\begin{enumerate}
	\item Xét $\triangle SAB$ vuông tại $A$, theo định lý Pi-ta-go ta có\\ $SA=\sqrt{SB^{2}-AB^{2}}=\sqrt{5a^{2}-a^2}=2a$.\\
	Thể tích khối chóp $S.ABCD$ là $V_{S.ABCD}=\dfrac{1}{3}\cdot SA \cdot S_{ABCD}=\dfrac{1}{3}\cdot 2a  \cdot a^{2}=\dfrac{2}{3}a^{3}$.
	\item Trước hết, ta xác định toạ độ các điểm $A(0;0;0)$, $B(a;0;0)$, $C(a;a;0)$, $D(0;a;0)$, $S(0;0;2a)$.\\
	Gọi $M$ là trung điểm của $SB$, ta có
	\begin{eqnarray*}
		M\left(\dfrac{0+a}{2};\dfrac{0+0}{2};\dfrac{2a+0}{2} \right)=\left(\dfrac{a}{2};0;a \right). 
	\end{eqnarray*}
	Gọi $N$ là trung điểm của $SD$, ta có
	\begin{eqnarray*}
		N\left(\dfrac{0+0}{2};\dfrac{0+a}{2};\dfrac{2a+0}{2} \right)=\left(0;\dfrac{a}{2};a \right).
	\end{eqnarray*}
	Phương trình đường thẳng $MD$ qua $M\left(\dfrac{a}{2};0;a \right) $ và $D(0;a;0)$ có dạng tham số 
	\begin{eqnarray*}
		\heva{&x=\dfrac{a}{2}-\dfrac{a}{2}t\\&y=at\\&z=a-at} \left(t\in\mathbb{R} \right).
	\end{eqnarray*}
		Phương trình đường thẳng $CN$ qua $C\left(a;a;0 \right)$ và $N\left(0;\dfrac{a}{2};a\right) $ có dạng tham số 
	\begin{eqnarray*}
		\heva{&x=a-at\\&y=a-\dfrac{a}{2}t\\&z=at} \left(t\in\mathbb{R} \right). 
	\end{eqnarray*}
	Để tính khoảng cách giữa hai đường thẳng chéo nhau $MD$ và $CN$, ta dùng công thức
	\begin{eqnarray*}
		d=\dfrac{\left|\overrightarrow{AB}\cdot\left(\overrightarrow{u}\times\overrightarrow{v} \right)  \right|}{\left|\overrightarrow{u}\times\overrightarrow{v} \right| }
	\end{eqnarray*}
	Véc-tơ chỉ phương của $MD$ là $\overrightarrow{u}=\left(-\dfrac{a}{2};a;-a\right)$.\\
	Véc-tơ chỉ phương của $CN$ là $\overrightarrow{v}=\left(-a;-\dfrac{a}{2};a\right)$.\\
	Tích có hướng $\overrightarrow{u}\times\overrightarrow{v}=\left(a^{2};a^{2};\dfrac{3a^{2}}{4} \right) $.\\
	Chọn điểm $A\left(\dfrac{a}{2};0;a \right) $ trên $MD$ và $B\left(a;a;0 \right)$ trên $CN$, ta có $\overrightarrow{AB}=\left(\dfrac{a}{2};a;-a\right) $.\\
	Khi đó, $\overrightarrow{AB}\cdot\left(\overrightarrow{u}\times\overrightarrow{v} \right)=\dfrac{a}{2}\times a^{2}+a\times a^{2}+(-a)\times \dfrac{3a^{2}}{4}=\dfrac{a^{3}}{2}+a^{3}-\dfrac{3a^{3}}{4}=\dfrac{3a^{3}}{4}$.\\
	Độ dài của $\overrightarrow{u}\times\overrightarrow{v}$ là
	
	$\left| \overrightarrow{u}\times\overrightarrow{v}\right|=\sqrt{a^{4}+a^{4}+\left(\dfrac{3a^{2}}{4} \right)^{2} }=\sqrt{2a^{4}+\dfrac{9a^{4}}{16}}=\sqrt{\dfrac{41a^{4}}{16}}=\dfrac{\sqrt{41}a^{2}}{4} $.\\
	Do đó, khoảng cách giữa hai đường thẳng là $d=\dfrac{\left|\dfrac{3a^{3}}{4} \right| }{\dfrac{\sqrt{41}a^{2}}{4}}=\dfrac{3a}{\sqrt{41}}$.
	\end{enumerate}
	}
\end{bt}

\begin{bt}%[1H8V5-3]
% nguon: Archive/SP_DA_SauLuoi/1-TL-THPT-NguyenThiMinhKhai-TpHCM-HKII-NH23-24.tex #2476
%[Dự án đề kiểm tra Toán 11 HKII NH23-24- Ngô Tất Thành]%[THPT Nguyễn Thị Minh Khai]
	Cho hình chóp $S.ABC$ có đáy $ABC$ là tam giác cân tại $B$, $SA\perp(ABC)$, $AB=a$, $\widehat {ABC}=120^{\circ}$, $SA=\dfrac{3a}{2}$. Gọi $I$ là trung điểm của $AC$.
	\begin{enumerate}
		\item Chứng minh rằng $BI\perp(SAC)$.
		\item Tính thể tích khối chóp $S.ABC$.
		\item Tính góc giữa hai mặt phẳng $(SBI)$ và $(ABC)$.
		\item Trong $(ABC)$, lấy điểm $K$ sao cho $\overrightarrow{CK}=-\dfrac{3}{2}\overrightarrow{AB}$. Gọi $H$ là hình chiếu vuông góc của $A$ lên $SK$, $L$ là điểm thuộc cạnh $SC$ thỏa mãn $SL=\dfrac{3}{4}SC$. Tính khoảng cách từ điểm $C$ đến $(AHL)$
	\end{enumerate}
	\loigiai{
\begin{center}
		\begin{tikzpicture}[scale=1,line join=round,line cap=round,font=\footnotesize,>=stealth]
			\path (0,0) coordinate (A) % theo phép tịnh tiến theo 1 góc.
			++(0:4) coordinate (K) 
			++(135:5.3) coordinate (S) 
			++(-110:6) coordinate (B)
			++(-15:4.3) coordinate (C) 
			($(S)!0.7!(K)$) coordinate (H) 
			($(S)!0.8!(C)$) coordinate (L)
			($(A)!0.5!(C)$) coordinate (I);
			\draw (S)--(B)--(C)--(K)--(S)--cycle (S)--(C) (H)--(L);
			\draw[dashed] (B)--(A)--(K) (S)--(A)--(H) (B)--(I)--(S) (L)--(A)--(H) (A)--(C); 
			\foreach \x/\y in {A/-180,B/-180,C/0,K/0,S/90,H/45,I/45,L/0 %hiện điểm
			}{\draw[fill=black] (\x) circle(1pt)++(\y:0.35)node{$\x$};}%nút tròn ngay điểm
			\draw pic[draw, angle radius=2.5mm]{right angle=S--A--K}; %góc vuông
			\draw pic[draw, angle radius=2.5mm]{right angle=A--K--C};
			\draw pic[draw, angle radius=2.5mm]{right angle=B--I--A};
			\draw pic[draw, angle radius=2.5mm]{right angle=A--H--S};
			\draw pic[draw, angle radius=2.5mm]{right angle=L--H--K};
			\draw pic[draw, angle radius=2.5mm]{right angle=S--A--B};
			\draw[thick] (0.38,-0.65) -- (0.65,-0.65); 
			\draw[thick] (1.43,-2) -- (1.7,-2);
		\end{tikzpicture}
\end{center}
		\begin{enumerate}
			\item Chứng minh rằng $BI\perp(SAC)$.\\
			Tam giác $ABC$ cân tại $B$ có $I$ là trung điểm của $AC$ $\Rightarrow BI\perp AC\quad(1)$\\
			$SA\perp(ABC)\Rightarrow BI\perp SA\quad(2)$\\
			Từ $(1)$ và $(2)$ suy ra $BI\perp(SAC)$.
			\item Tính thể tích khối chóp $S.ABC$.\\
			$S_{ABC}=\dfrac{1}{2}BA\cdot BC\cdot\sin\widehat{ABC}=\dfrac{1}{2}\cdot a\cdot a\cdot\sin 120^{\circ}=\dfrac{a^2\sqrt{3}}{4}$.\\
			$V_{SABC}=\dfrac{1}{3} S_{ABC}\cdot SA=\dfrac{1}{3}\cdot\dfrac{a^2\sqrt{3}}{4}\cdot\dfrac{3a}{2}=\dfrac{a^3\sqrt{3}}{8}$.
			\item Tính góc giữa hai mặt phẳng $(SBI)$ và $(ABC)$.\\
			$\heva{&(SBI)\cap(ABC)=BI\\
				&BI\perp(SAC)\\
				&(SAC)\cap(SBI)=SI\\
				&(SAC)\cap(ABC)=AC}$\\
			$\Rightarrow((SBI),(ABC)) =(SI,AC)=\widehat{SIA}$ (do $SA\perp(ABC)\Rightarrow SA\perp AI$).\\
			$AI=AB\cdot\sin 60^{\circ}=\dfrac{a\sqrt{3}}{2}$.\\
			$\tan\widehat{SIA}=\dfrac{SA}{AI}=\dfrac{3a}{\dfrac{a\sqrt{3}}{2}}=\sqrt{3}\Rightarrow\widehat{SIA}=60^{\circ}$. 
			Vậy $\left((SBI),(ABC)\right)=60^{\circ}$.
			\item Trong $(ABC)$, lấy điểm $K$ sao cho $\overrightarrow{CK}=-\dfrac{3}{2}\overrightarrow{AB}$. Gọi $H$ là hình chiếu vuông góc của $A$ lên $SK$, $L$ là điểm thuộc cạnh $SC$ thỏa mãn $SL=\dfrac{3}{4}SC$. Tính khoảng cách từ điểm $C$ đến $(AHL)$.\\
			Ta có
			$CK=\dfrac{3}{2}a$; $AC=2AI=a\sqrt{3}$, $AK=\sqrt{AC^2+CK^2 - 2AC\cdot CK\cdot\cos 30^\circ}=a\dfrac{\sqrt{3}}{2}$;\\
			$AK^2+CK^2=AC^2\Rightarrow AK\perp KC$.\\
			Mà  $SA\perp CK(SA\perp(ABC)) $ nên $ CK\perp(SAK)\Rightarrow CK\perp KH$.\\
			$SH\cdot SK=SA^2\Rightarrow\dfrac{SH}{SK}=\dfrac{SA^2}{SK^2}=\dfrac{SA^2}{SA^2+AK^2}=\dfrac{3}{4}\Rightarrow\dfrac{SL}{SC}=\dfrac{3}{4}\Rightarrow HL\parallel KC$.\\
			$\Rightarrow CK\parallel(AHL)\Rightarrow d(C,(AHL))=d(K,(AHL))$.\\
			$\heva{&KH\perp AH\\& KH\perp HL(CK\perp HK)}\Rightarrow HK\perp(AHL)$.\\
			$\Rightarrow d(K,(AHL))=KH=\dfrac{1}{4}\cdot SK=\dfrac{1}{4}\sqrt{SA^2+AK^2}=a\dfrac{\sqrt{3}}{4}$.\\
		\end{enumerate}}
	\end{bt}

\begin{bt}%[1H8V5-3]
% nguon: Archive/SP_DA_SauLuoi/1-TL-TN-THPT-AnThoi-BenTre-HKII-NH23-24.tex #2489
%[Dự án đề kiểm tra Toán 11 GHKI NH23-24- Nguyễn Mạnh Trung Anh]%[THPT AN THỚI- BẾN TRE]
	Cho hình chóp $S.ABCD$ có đáy $ABCD$ là hình vuông cạnh $a$, cạnh bến $SA$ vuông góc với mặt phẳng đáy và $SA=a\sqrt{3}$.
	\begin{enumerate}
		\item Chứng minh $BD\perp(SAC)$.
		\item Tính khoảng cách từ $A$ đến mặt phẳng $(SCD)$.
	\end{enumerate}
	\loigiai{
		\begin{center}
			\begin{tikzpicture}[>=stealth,line join=round,line cap=round,font=\footnotesize,scale=0.65]
				\path (0,0) coordinate (A)
				(-2.5,-1.4) coordinate (B)
				(6,0) coordinate (D)
				($(B)+(D)-(A)$) coordinate (C)
				($(A)!0.5!(C)$) coordinate (O)
				($(A)+(0,6)$)coordinate (S)
				($(S)!(A)!(D)$) coordinate (H)
				;
				\draw (S)--(B)--(C)--(S)--(D)--(C);
				\draw[dashed] (S)--(A)--(B)--(D)--(A)--(C) (A)--(H);
				\foreach \p / \r in {A/45,B/-135,C/-45,S/90,O/-90,D/45, H/45}
				\fill (\p) circle (1.2pt) node[shift={(\r:2mm)}]{$\p$};
				\draw pic[thin,draw,angle radius=2mm] {right angle= B--A--S};
				\draw pic[thin,draw,angle radius=2mm] {right angle= D--H--A};
			\end{tikzpicture}
		\end{center}
		Kẻ $AH\perp SD (H\in SD)$.
		\begin{enumerate}
			\item Ta có $\heva{&BD\perp AC, BD\perp SA\\&SA, AC\subset(SAC)\\&SA\cap AC=\{A\}}\Rightarrow BD\perp (SAC)$.
			\item Ta có $\heva{&CD\perp AD, CD\perp SA\\& SA, AD\subset(SAD)\\&SA\cap AD=\{A\}}\Rightarrow CD\perp(SAD)\Rightarrow CD\perp AH$.\\
			Lại có $\heva{&AH\perp SD, AH\perp CD\\& SD, CD\subset(SCD)\\&SD\cap CD=\{D\}}\Rightarrow AH\perp (SCD)\Rightarrow \mathrm{d}(A, (SCD))=AH$.\\
			Xét $\triangle{SAD}$ vuông tại $A$ có
			\begin{align*}
				&\dfrac{1}{AH^2}=\dfrac{1}{SA^2}+\dfrac{1}{AD^2}\\
				\Rightarrow&\dfrac{1}{AH^2}=\dfrac{1}{\left(a\sqrt{3}\right)^2}+\dfrac{1}{a^2}=\dfrac{4}{3a^2}\\
				\Rightarrow&AH=\dfrac{a\sqrt{3}}{2}.
			\end{align*}
			Vậy khoảng cách từ $A$ đến $(SCD)$ bằng $\dfrac{a\sqrt{3}}{2}$.
		\end{enumerate}
	}
\end{bt}

\begin{bt}%[1H8V5-4]
% nguon: Archive/SP_DA_SauLuoi/1-TL-TN-DS-TLN-THPT-NguyenKhuyen-TpHCM-HKII-NH23-24.tex #2505
%[Dự án đề kiểm tra Toán 11 HKII NH23-24- Quan Ón]%[THPT Nguyễn Khuyến - Tp HCM]
	Cho hình chóp $S.ABCD$ có đáy là hình bình hành $ABCD$ tâm $O$, tam giác $ABC$ vuông cân tại $A$ có $AB = AC = 1$, $SA \perp (ABCD)$. Đường thẳng $SD$ tạo với mặt đáy một góc $45^\circ$. Tính khoảng cách giữa hai đường thẳng $AD$ và $SB$ (kết quả gần đúng làm tròn đến hàng phần trăm).
	\loigiai{
		\immini{
			Gọi $M$ là trung điểm $BC$, $H$ là hình chiếu của $A$ lên $SM$.\\
			Khi đó $AM$ là đường trung tuyến cũng là đường cao trong $\triangle ABC$ vuông cân tại $A$. Do đó $AM \perp BC$.\\
			Ta có $\heva{&BC \perp SA \,\,(SA \perp (ABCD), BC \subset (ABCD))\\&BC \perp AM \textrm{ (cmt)}}$\\
			$\Rightarrow BC \perp (SAM)$, mà $AH \subset (SAM)$ nên $BC \perp AH$.\\
			Do đó $\heva{&AH \perp BC \textrm{ (cmt)}\\&AH \perp SM} \Rightarrow AH \perp (SBC)$\\
			$\Rightarrow \mathrm{d}(A,(SBC)) = AH$.
		}{
			\begin{tikzpicture}[>=stealth,line join=round,line cap=round,font=\footnotesize,scale=.8]
				\path 
				(0,0) coordinate (A)
				(4,0) coordinate (B)
				(-1.75,-1.5) coordinate (D)
				($(B)+(D)-(A)$) coordinate (C)
				(0,2.6) coordinate (S)
				($(B)!0.5!(C)$) coordinate (M)
				($(S)!(A)!(M)$) coordinate (H)
				;
				\draw 
				(C)--(D)--(S)--(C)--(B)--(S)--(M)
				;
				\draw[dashed]
				(D)--(A)--(S) (C)--(A)--(B) (M)--(A)--(H)
				;
				\foreach \p/\g in {A/180, B/0, C/-90, D/-90, S/90, M/-45, H/0}
				\draw[fill=black] (\p) circle (1pt) node[shift=(\g:3mm)] {$\p$};
				\pic[draw,angle radius=2mm,angle eccentricity=1.5] {right angle = A--M--C};
				\pic[draw,angle radius=2mm,angle eccentricity=1.5] {right angle = A--H--M};
				\pic[draw,angle radius=3.5mm,angle eccentricity=1.5] {right angle = B--A--C};
			\end{tikzpicture}
		}
		\noindent
		Vì $AD \parallel (SBC)$ chứa $BC$ nên 
		$$ \mathrm{d}(SB,AD) = \mathrm{d}(AD,(SBC)) = \mathrm{d}(A,(SBC)) = AH.$$
		Mặt khác, ta có $SA \perp AD$ nên $\left(\widehat{SD,(ABCD)}\right) = \widehat{SDA} = 45^\circ$.\\
		Xét $\triangle ABC$ vuông cân tại $A$, ta có
		$$ BC = \sqrt{AB^2 + AC^2} = \sqrt{1^2 + 1^2} = \sqrt{2}. $$
		Suy ra $AD = BC = \sqrt{2}$ và $AM = \dfrac{1}{2}BC = \dfrac{\sqrt{2}}{2}$.\\
		Xét $\triangle SAD$ vuông tại $A$ có $\widehat{SDA} = 45^\circ$ nên $\triangle SAD$ vuông cân tại $A$, do đó $SA = \sqrt{2}$.\\
		Xét $\triangle SAM$ vuông tại $A$, có đường cao $AH$ nên
		$$ AH = \dfrac{SA\cdot AM}{\sqrt{SA^2 + AM^2}} = \dfrac{\sqrt{2}\cdot \dfrac{\sqrt{2}}{2}}{\sqrt{\left(\sqrt{2}\right)^2 + \left(\dfrac{\sqrt{2}}{2}\right)^2}} = \dfrac{\sqrt{10}}{5} \approx 0{,}63. $$
		Vậy $\mathrm{d}(SB,AD) \approx 0{,}63$.
	}
\end{bt}

\begin{bt}%[1H8H6-3]
% nguon: Archive/SP_DA_SauLuoi/1-TL-TN-DS-TLN-THPT-NguyenKhuyen-TpHCM-HKII-NH23-24.tex #2508
%[Dự án đề kiểm tra Toán 11 HKII NH23-24- Quan Ón]%[THPT Nguyễn Khuyến - Tp HCM]
	Cho hình chóp tam giác $S.ABC$ có $SA$ vuông góc với mặt đáy $(ABC)$, tam giác $ABC$ vuông cân tại $B$ và $SA = AB = a$. Gọi $M$ là trung điểm của $AC$.
	\begin{enumerate}
		\item Chứng minh $BM \perp (SAC)$.
		\item Tính góc giữa hai mặt phẳng $(SBC)$ và $(ABC)$.
	\end{enumerate}
	\loigiai{
		\begin{center}
			\begin{tikzpicture}[>=stealth,line join=round,line cap=round,font=\footnotesize,scale=1]
				\path 
				(0,0) coordinate (A)
				(2.6,-1.8) coordinate (B)
				(4,0) coordinate (C)
				(0,2.6) coordinate (S)
				($(A)!0.5!(C)$) coordinate (M)
				;
				\draw 
				(A)--(S)--(B)--(C)--(S) (A)--(B)
				;
				\draw[dashed]
				(A)--(C) (B)--(M)
				;
				\foreach \p/\g in {A/-135, B/-90, C/-45, S/90, M/90}
				\draw[fill=black] (\p) circle (1pt) node[shift=(\g:3mm)] {$\p$};
				\pic[draw,angle radius=2mm,angle eccentricity=1.5] {right angle = B--M--C};
				\pic[draw,angle radius=2mm,angle eccentricity=1.5] {right angle = A--B--C};
				\pic[draw,angle radius=2mm,angle eccentricity=1.5] {right angle = S--A--B};
			\end{tikzpicture}
		\end{center}
		\begin{enumerate}
			\item Xét $\triangle ABC$ vuông cân tại $B$ nên $BM$ vừa là đường trung tuyến vừa là đường cao trong $\triangle ABC$.\\
			Do đó $BM \perp AC$.\\
			Ta có $\heva{&BM \perp AC \textrm{ (cmt)}\\&BM \perp SA\,\, (SA \perp (ABC), BM \subset (ABC))} \Rightarrow BM \perp (SAC)$.
			\item Ta có $\heva{&BC \perp AB \textrm{ ($\triangle ABC$ vuông tại $B$)}\\&BC \perp SA \,\, (SA \perp (ABC), BC \subset (ABC))} \Rightarrow BC \perp (SAB)$.\\
			Mà $SB \subset (SAB) \Rightarrow BC \perp SB$.\\
			Ta có $\heva{&AB \perp BC, AB \subset (ABC)\\&SB \perp BC, SB \subset (SBC)\\&(ABC) \perp (SBC) = BC} \Rightarrow \left(\widehat{(ABC),(SBC)}\right) = \left(\widehat{AB,SB}\right) = \widehat{SBA}$.\\
			Xét $\triangle SAB$ vuông tại $A$, có $SA = AB = a$ (giả thiết) nên $\triangle SAB$ vuông cân tại $A$, do đó $\widehat{SBA} = 45^\circ$.\\
			Vậy $\left(\widehat{(ABC),(SBC)}\right) = 45^\circ$.
		\end{enumerate}
	}
\end{bt}

\begin{bt}%[1H8H7-3]
% nguon: Archive/SP_DA_SauLuoi/1-TN-DS-TLN-THPT-LuongTheVinh-HaNoi-HKII-NH23-24.tex #2516
%[Dự án đề kiểm tra Toán 11 GHKI NH23-24- Don Lee]%[THPT Lương Thế Vinh - Tp Hà Nội]
	Cho hình chóp $S.ABCD$ có đáy $ABCD$ là hình chữ nhật cạnh $BC=2$; $CD=4$. Biết $SA$ vuông góc với $(ABCD)$, góc giữa $SD$ và $(ABCD)$ bằng $45^\circ$, thể tích khối chóp $S.ABCD$ là $\dfrac{a}{b}$ với $a$, $b$ là các số tự nhiên và $\dfrac{a}{b}$ là phân số tối giản. Giá trị của $a+b$ bằng bao nhiêu?
	\loigiai{
		\immini
		{Từ giả thiết ta có $SA \perp (ABCD)$ và\\
			$\left(\widehat{SD,(ABCD)}\right)=\widehat{SDA}=45^\circ$.\\		
			Suy ra tam giác $SDA$ vuông cân tại $A$ nên $SA=DA=4$.\\		
			Thể tích khối chóp $S.ABCD$ là		
			$$V_{S.ABCD}=\dfrac{1}{3}\cdot SA\cdot S_{ABCD}=\dfrac{1}{3}\cdot 4\cdot 2\cdot 4= \dfrac{32}{3}.$$
			Vậy $a+b=32+3=35$.}
		{\begin{tikzpicture}[>=stealth,line join=round,line cap=round]
				\draw[dashed,black,scale=2.5] (0,0)coordinate(A)--(-130:0.5)coordinate(B) (1,0)coordinate(D)--(A)--($(A)+(90:1)$)coordinate(S);
				\draw[black,scale=2.5] (B)--($(D)-(A)+(B)$)coordinate(C)--(D)--(S)--(B) (S)--(C);
				\path pic[angle radius=2mm,draw=blue] {right angle = D--A--S};
				\path pic[angle radius=3mm,draw=blue] {angle = S--D--A};
				\foreach \diem/\vitrin in {S/above,A/left,B/below left,C/below right,D/right} \fill (\diem)circle(1.5pt)node[\vitrin]{$\diem$};
		\end{tikzpicture}}
	}
\end{bt}

\begin{bt}%[1H8H2-3]
% nguon: Archive/SP_DA_SauLuoi/1-TL-TN-THPT-Chuyen-LeQuyDon-NinhThuan-HKII-NH23-24.tex #2532
%[Dự án đề kiểm tra Toán 11 CKII NH23-24-Đợt 16-Lê Hữu Kiệt]%[THPT Chuyên Lê Quý Đôn- Ninh Thuận]
	Cho hình chóp $S.ABCD$ có đáy $ABCD$ là hình vuông và $SA$ vuông góc với đáy. Chứng minh rằng:
	\begin{enumerate}
		\item $SA$ vuông góc với $BC$;
		\item $BC$ vuông góc với mặt phẳng $\left( SAB \right)$.
	\end{enumerate}
	\loigiai{
	\immini
	{\begin{enumerate}
		\item Theo giả thiết ta có $SA \perp (ABCD)$, $BC\subset (ABCD)$.\\
		Vậy $SA\perp BC$.
		\item Ta có
		\begin{center}
			$\heva{&BC\perp SA \\ &BC\perp AB \, (ABCD\text{\, là hình vuông}) \\ &SA\cap AB=A \\ &AB,\, SA\in(SAB).}$
		\end{center}
		Vậy $BC\perp (SAB)$.
	\end{enumerate}}
	{\begin{tikzpicture}[>=stealth, line join=round, line cap=round, font=\footnotesize]
	\path
		(0,0) coordinate (A) (3,0) coordinate (B) (-150:1.8) coordinate (D)++(3,0) coordinate (C)
		(0,2) coordinate (S);
	\draw[dashed] (D)--(A)--(B) (S)--(A);
	\draw (S)--(D)--(C)--(B)--cycle (S)--(C);
	\foreach \x/\g in {S/90, C/0, B/0, A/180, D/180}\fill[black](\x)circle(1pt)+(\g:.3)node{$\x$};
	\end{tikzpicture}}
	}
\end{bt}

\begin{bt}%[1H8V5-4]
% nguon: Archive/SP_DA_SauLuoi/1-TL-TN-THPT-HungVuong-HKII-TpHCM-NH23-24.tex #2550
%[Dự án đề kiểm tra Toán 11 HKII NH23-24- Nguyễn Tấn Linh]%[THPT Hùng Vương - Tp HCM]
	Cho hình chóp $S. ABCD$ có đáy $ABCD$ là hình vuông cạnh $a$, $SA\perp(ABCD)$. Biết $SA=a\sqrt{2}$.
	\begin{enumerate}
		\item Chứng minh $CD\perp(SAD)$.
		\item Tính góc giữa $SC$ và mặt phẳng $(ABCD)$.
		\item Gọi $M,N$ lần lượt là trung điểm của $BC$, $CD$. Tính khoảng cách giữa hai đường thẳng chéo nhau $SD$ và $MN$.
	\end{enumerate}
	\loigiai
	{
		\begin{center}
			\begin{tikzpicture}[>=stealth,line join=round,line cap=round,font=\footnotesize,scale=1]
				\coordinate (B) at (0,0);
				\coordinate (A) at (1,1.2);
				\coordinate (C) at (3.5,0);
				\coordinate (D) at ($(C)-(B)+(A)$);
				\coordinate (S) at ($(A)+(90:3.5)$);
				\coordinate (M) at ($(B)!0.5!(C)$);
				\coordinate (N) at ($(C)!0.5!(D)$);
				\coordinate (O) at ($(A)!0.5!(C)$);
				\coordinate (H) at ($(S)!0.4!(O)$);
				\coordinate (I) at ($(O)!0.5!(C)$);
				\draw (B)--(C)--(D)--(S)--cycle (S)--(C);
				\draw[dashed] (A)--(D) (S)--(A)--(B) (M)--(N) (B)--(D) (A)--(C) (A)--(H) (S)--(O);
				\def \dolongoc{5}
				\foreach \x/\dinh/\y in {S/A/D,A/H/O} \draw ($(\dinh)!\dolongoc pt!(\x)$)--($(\dinh)!\dolongoc pt!(\x)+(\dinh)!\dolongoc pt!(\y)-(\dinh)$)--($(\dinh)!\dolongoc pt!(\y)$); 
				\foreach \diem/\goc in {A/180,B/-90,C/-90,D/0,S/90,M/-90,N/0,O/-90,H/-30,I/90} \fill[black](\diem) circle (1pt) ($(\diem)+(\goc:3mm)$) node{$\diem$};
			\end{tikzpicture}
		\end{center}
		\begin{enumerate}
			\item Ta có $\heva{& CD\perp AD&&\text{($ABCD$ là hình vuông)} \\ & CD\perp SA&&\left(SA\perp (ABCD)\right)\\ & AD\cap SA=A\text{ trong }(SAD)}\Rightarrow CD\perp (SAD)$.
			\item Ta có $SA\perp (ABCD)\Rightarrow A$ là hình chiếu của $S$ lên mặt phẳng $(ABCD)$.\\
			Lại có $SC\cap (ABCD)=C$ nên góc giữa đường thẳng $SC$ và mặt phẳng $(ABCD)$ là góc giữa $SC$ và $AC$ và bằng $\widehat{SCA}$.\\
			Ta có $AC=a\sqrt{2}$ nên tam giác $SAC$ vuông cân tại $A$, do đó $\widehat{SCA}=45^\circ$.
			\item Ta có $\heva{&BD\perp AC\\&BD\perp SA\\& SA\cap AC=A\text{ trong }(SAC)}\Rightarrow BD \perp (SAC)$.\\ 
			Trong $(SAC)$, vẽ $AH \perp SO$.\\
			Ta có $\heva{&AH\perp SO\\&AH\perp BD \text{ vì } BD\perp (SAC)\\& BD\cap SO=O\text{ trong }(SBD)}\Rightarrow AH \perp (SBD)$.\\
			Vậy $\mathrm{d}(A,(SBD))=AH$.\\
			Vì $ABCD$ là hình vuông nên $AC=\sqrt{AB^2+BC^2}=a\sqrt{2}$.\\
			 Do đó, $AO=\dfrac{AC}{2}=\dfrac{a\sqrt{2}}{2}$.\\
			 Ta có $\dfrac{1}{AH^2}=\dfrac{1}{SA^2}+\dfrac{1}{AO^2}=\dfrac{1}{(a\sqrt{2})^2}+\dfrac{1}{(\dfrac{a\sqrt{2}}{2})^2}=\dfrac{5}{2a^2}\\
			 \Rightarrow AH=\dfrac{a\sqrt{10}}{5}$. Vậy $\mathrm{d}(A,(SBD))=AH=\dfrac{a\sqrt{10}}{5}$.\\
			Trong $(ABCD)$, gọi $I=MN\cap OC$.\\
			Vì $M,N$ lần lượt là trung điểm của $BC,CD$ nên ta có $MN \parallel BD$ và $MN=\dfrac{BD}{2}$.\\
			Vì $MN \parallel BD$ nên $\dfrac{CI}{OC}=\dfrac{CM}{CB}=\dfrac{1}{2}$. Do đó $\dfrac{IO}{OC}=\dfrac{1}{2}$.\\
			Ta có $\heva{&MN \parallel BD\\ & BD \subset (SBD)}\Rightarrow MN \parallel (SBD)$.\\
			Ta có $\heva{&MN \parallel (SBD)\\& BD \subset (SBD)\\& I \in MN}\Rightarrow \mathrm{d}(MN,SD)=\mathrm{d}(MN,(SBD))=\mathrm{d}(I,(SBD))$.\\
			Ta có $\dfrac{\mathrm{d}(I,(SBD))}{\mathrm{d}(A,(SBD))}=\dfrac{IO}{AO}=\dfrac{IO}{OC}=\dfrac{1}{2} \Leftrightarrow \mathrm{d}(I,(SBD))=\dfrac{\mathrm{d}(A,(SBD))}{2}=\dfrac{a\sqrt{10}}{10}$.\\
		\end{enumerate}
	}
\end{bt}

\begin{bt}%[1H8H4-5]
% nguon: Archive/SP_DA_SauLuoi/1-TL-THPT-NguyenTatThanh-HCM-HKII-NH23-24.tex #2558
%[Dự án đề kiểm tra Toán 11 HKII NH23-24- Trần Xuan Hòa]%[THPT NguyenTatTHanh]
Cho hình chóp $S.ABCD$ có đáy là hình chữ nhật, $SA\perp (ABCD), AD=SA=a$, $AB=\sqrt{3}a; H, K$ lần lượt là hình chiếu của $A$ trên $SB$ và $SD$.
\begin{enumerate}
	\item Chứng minh $(SCD) \perp (SAD), (AHK) \perp (SAC)$.
	\item Tính thể tích khối chóp $S.ABC$.
	\item Tính khoảng cách giữa hai đường thẳng $HK$ và $AC$.
\end{enumerate}	
	\loigiai{
	\begin{center}
		\begin{tikzpicture}[scale=1.1, font=\footnotesize,line join=round, line cap=round, >=stealth]
			\path 
			(0,0) coordinate (A)
			++(-140:2) coordinate (B)
			++(0:3.5) coordinate (C)
			($(A)+(C)-(B)$) coordinate (D)
			($(A)+(0,3)$) coordinate (S)
			;
			\coordinate (H) at ($(S)!.45!(B)$);
			\coordinate (K) at ($(S)!.5!(D)$);
			\coordinate (O) at ($(B)!.5!(D)$);
			\coordinate (N) at ($(S)!.4!(A)$);
			\coordinate (M) at (intersection of S--O and H--K);
			\foreach \i in{B,C,D}{\draw (S)--(\i);};
			\draw (B)--(C)--(D);
			\draw[dashed] (S)--(A)--(B) (C)--(A)--(D) (B)--(D) (S)--(O) (H)--(K) (K)--(A)--(H) (M)--(N)--(K);
			\pic[draw,angle eccentricity=1.8,angle radius=2mm]{right angle=S--A--D};
			\pic[draw,angle eccentricity=1.8,angle radius=2mm]{right angle=D--K--A};
			\pic[draw,angle eccentricity=1.8,angle radius=2mm]{right angle=B--H--A};
			\foreach \i/\g in {A/-90,B/-90,C/-90,D/-90,S/90,H/190,K/60,O/-90,M/-140,N/180}
			\fill[black] (\i) circle(1pt)+(\g:3mm)node[scale=1]{$\i$};
		\end{tikzpicture}
	\end{center}
	\begin{enumerate}
		\item 
		\begin{itemize}
			\item Ta có $C D \perp A D, C D \perp S A \Rightarrow C D \perp(S A D) \Rightarrow(S C D) \perp(S A D)$
			\item Ta có $C D \perp A K, S D \perp A K \Rightarrow A K \perp(S C D) \Rightarrow A K \perp S C$.
		\end{itemize}
		Tương tự $A H \perp S C; SC \perp (AHK) \Rightarrow(SAC) \perp(AHK)$.
		\item $V_{\text{SABC}}=\dfrac{1}{3} S A \cdot S_{\text{ABC}}=\dfrac{\sqrt{3}}{6} a^3$.\\
		\item Gọi $O=B D \cap A C, M=S O \cap H K$.\\
		Trong $\triangle SAO$ kẻ $MN \parallel AO$, nên $AC\parallel (MNK)$.\\
		Ta có $\mathrm{d}(AC,HK)=\mathrm{d}(AC,(MNK))=\mathrm{d}(A,(MNK))$.\\
		Ta có $\mathrm{d}(A,(MNK))=2\mathrm{d}(S,(MNK))$.\\
		Tính được 
		\begin{itemize}
			\item $V_{S.MNK}=\dfrac{SN}{SA}\cdot \dfrac{SM}{SO}\cdot \dfrac{SK}{SD}\cdot V_{S.OAD}=\dfrac{1}{18}\cdot \dfrac{1}{2}V_{S.ABC}=\dfrac{1}{36}\cdot \dfrac{\sqrt{3}}{6} a^3=\dfrac{\sqrt{3}}{216}a^3$.
			\item $MN=\dfrac{OA}{3}=\dfrac{a}{3}$; $NK^2=SN^2+SK^2-2SN\cdot SK\cos 45^\circ=\dfrac{5}{18}$.
			\item $MK^2=\dfrac{2}{9}$, trong $\triangle MNK$ có $\cos \widehat{M}=\dfrac{MK^2+MN^2-NK^2}{2MK\cdot MN}=\dfrac{1}{4\sqrt{2}}$.\\ Do đó $\sin \widehat{M}=\dfrac{\sqrt{31}}{4\sqrt{2}}$. \\Vậy diện tích $\triangle MNK$ là $S=\dfrac{MN\cdot MK\cdot \sin \widehat{M}}{2}=\dfrac{\sqrt{31}}{72}$.
			\item Ta có $V_{S.MNK}=\dfrac{\mathrm{d}(S,(MNK))\cdot S}{3}$ nên $\mathrm{d}(S,(MNK))=\dfrac{\sqrt{93}}{31}$.
		\end{itemize}
		Vậy $\mathrm{d}(A,(MNK))=2\mathrm{d}(S,(MNK))=\dfrac{2 \sqrt{93}}{31}a$.
		\end{enumerate}
	}
\end{bt}

\begin{bt}%[1H8V5-4]
% nguon: Archive/SP_DA_SauLuoi/1-TL-TN-DS-THPT-ChuyenTranPhu-HaiPhong-HKII-NH23-24.tex #2566
%[Dự án đề kiểm tra Toán 11 HK2 NH23-24 - Phan Anh]%[THPT Chuyên Trần Phú - Hải Phòng]
Cho hình chóp $S.ABCD$ có đáy là hình vuông cạnh bằng $2$, $SA$ vuông góc với mặt phẳng đáy và $SA=\sqrt{5}$. Gọi $M$, $N$ là trung điểm của $SA$ và $CD$. Tính khoảng cách giữa hai đường thẳng $MN$ và $SC$, (làm tròn đến hàng phần trăm).
\loigiai{\immini{Gọi $O$ là tâm của hình vuông $ABCD$.\\
	Ta có $OM$ là đường trung bình của tam giác $SAC$ nên $OM\parallel SC$.\\
	Mặt khác ta có $\heva{&OM\subset(OMN)\\&SC\not\subset(OMN)}\Rightarrow SC\parallel(OMN)$.\\
	Mà $MN\subset(OMN)$ nên
	\[\mathrm{d}(MN,SC)=\mathrm{d}(SC,(OMN))=\mathrm{d}(C,(OMN)).\]
	Do $AC$ cắt $(OMN)$ tại $O$ nên}
{\begin{tikzpicture}[>=stealth,line join=round,line cap=round,font=\footnotesize,scale=0.5]
		\path (0,0) coordinate (A)
		(-2.5,-1.4) coordinate (B)
		(6,0) coordinate (D)
		($(B)+(D)-(A)$) coordinate (C)
		($(A)!0.5!(C)$) coordinate (O)
		($(A)+(0,6)$)coordinate (S)
		($(A)!0.5!(S)$) coordinate (M)
		($(A)!0.5!(B)$) coordinate (P)
		($(C)!0.5!(D)$) coordinate (N)
		($(M)!0.65!(P)$) coordinate (H)
		;
		\draw (S)--(B)--(C)--(S)--(D)--(C);
		\draw[dashed] (S)--(A)--(B)--(D)--(A)--(C) (A)--(H) (M)--(P)--(N)--(M)--(O);
		\foreach \p / \r in {A/45,B/-135,C/-45,S/90,M/135,O/-90,D/45,N/-45,H/135,P/135}
		\fill (\p) circle (1.2pt) node[shift={(\r:2.5mm)}]{$\p$};
\end{tikzpicture}}
\[\dfrac{\mathrm{d}(C,(OMN))}{\mathrm{d}(A,(OMN))}=\dfrac{OC}{OA}=1\Rightarrow\mathrm{d}(C,(OMN))=\mathrm{d}(A,(OMN)).\]
Gọi $P$ là trung điểm của $AB$, ta có $(OMN)$ cũng chính là $(MNP)$.\\
Ta có $\heva{&NP\perp AB\\&NP\perp SA}\Rightarrow NP\perp(SAB)$.\\
Do $NP\subset(MNP)$ nên $(MNP)\perp(SAB)$ theo giao tuyến $MP$.\\
Gọi $H$ là hình chiếu vuông góc của $A$ lên $MP$ trong $(SAB)$, suy ra $AH\perp(MNP)$ tại $H$.\\
Tam giác $AMP$ có $\dfrac{1}{AH^2}=\dfrac{1}{AP^2}+\dfrac{1}{AM^2}=\dfrac{9}{5}\Rightarrow AH=\dfrac{\sqrt{5}}{3}$.\\
Vậy $\mathrm{d}(MN,SC)=\mathrm{d}(A,(MNP))=AH=\dfrac{\sqrt{5}}{3}$.}
\end{bt}

\begin{bt}%[1H8H7-4]
% nguon: Archive/SP_DA_SauLuoi/1-TL-TN-DS-THPT-ChuyenTranPhu-HaiPhong-HKII-NH23-24.tex #2567
%[Dự án đề kiểm tra Toán 11 HK2 NH23-24 - Phan Anh]%[THPT Chuyên Trần Phú - Hải Phòng]
Cho hình lăng trụ đứng $ABC.A'B'C'$ có đáy $ABC$ là tam giác vuông cân tại $A$, $BC=a\sqrt{2}$, $A'B$ tạo với đáy một góc bằng $60^{\circ}$. Gọi $V$ là thể tích của khối lăng trụ $ABC.A'B'C'$. Tính tỷ số $\dfrac{V}{a^3}$, (làm tròn đến hàng phần trăm).
\loigiai{\immini{Ta có $A'A\perp(ABC)$ nên $AB$ là hình chiếu của $A'B$ lên đáy $(ABC)$.\\
	Suy ra $(A'B,(ABC))=(A'B,AB)=\widehat{A'BA}=60^\circ$.\\
	Tam giác $ABC$ vuông cân tại $A$ và có $BC=a\sqrt{2}$ nên $AB=AC=a$.\\
	Tam giác $A'BA$ có $\tan60^\circ=\dfrac{A'A}{AB}\Rightarrow A'A=AB\tan60^\circ=a\sqrt{3}$.\\
	Thể tích của khối lăng trụ $ABC.A'B'C'$ là 
	\[V=A'A\cdot S_{ABC}=A'A\cdot\dfrac{1}{2}\cdot AB\cdot AC=\dfrac{a^3\sqrt{3}}{2}.\]
	Suy ra $\dfrac{V}{a^3}=\dfrac{\sqrt{3}}{2}\approx0{,}87$.
}
{\begin{tikzpicture}[>=stealth,line join=round,line cap=round,font=\footnotesize,scale=0.5]
		\path (0,0) coordinate (A)
		(2,-1.4) coordinate (B)
		(5,0) coordinate (C)
		($(A)+(0,5)$)coordinate (A')
		($(B)+(0,5)$)coordinate (B')
		($(C)+(0,5)$)coordinate (C')
		;
		\draw (B')--(B)--(A')--(A)--(B)--(C)--(C')--(A')--(B')--(C')--(A');
		\draw[dashed] (A)--(C);
		\foreach \p / \r in {A/135,B/-90,A'/135,B'/90,C/45,C'/45}
		\fill (\p) circle (1.2pt) node[shift={(\r:3mm)}]{$\p$};
\end{tikzpicture}}}
\end{bt}

\begin{bt}%[1H8V6-1]
% nguon: Archive/SP_DA_SauLuoi/1-TL-THPT-PTNK-TpHCM-HKII-NH23-24.tex #2581
%[Dự án A - đợt 17, Nguyễn Quang Hiệp]%[Đề kiểm tra học kỳ 2, THPT PTNK, TP HCM, năm học 2023 2024]
	Cho hình chóp $S.ABC$ có đáy $ABC$ là tam giác đều có cạnh bằng $\sqrt{6}$, $ \widehat{SBA} = 45^\circ$ và $SA \perp (ABC)$. Gọi $I$ là trung điểm $BC$.
	\begin{enumerate}
		\item  Chứng minh $BC \perp (SIA)$.
		\item  Gọi $\alpha = ((SBC); (ABC))$. Tính $\tan \alpha$.
		\item  Gọi $\beta = (SA, (SBC))$. Tính $\sin \beta$ và tính khoảng cách từ $I$ đến $(SAB)$.		
	\end{enumerate}
	\loigiai{
		\begin{enumerate}
			\item Ta có $\triangle ABC$ là tam giác đều, có $I$ là trung điểm của $BC$.
			\immini{
				Do đó $AI\perp BC$.\\
				Lại có $SA\perp (ABC)$, nên $SA\perp BC$.\\
				Từ đó suy ra $BC\perp (SIA)$.
				\item Ta có $BC\perp (SIA)$, nên $BC\perp SI$.\\
				Lại có $BC\perp AI$, nên $((SBC); (ABC))=\widehat{SIA}$.\\
				Ta có $\triangle ABC$ đều cạnh bằng $\sqrt{6}$, có $AI$ là đường trung tuyến nên $AI=\dfrac{3\sqrt{2}}{2}$.\\
				Mặt khác, ta có $\triangle SAB$ vuông tại $A$, có $\widehat{SBA}=45^{\circ}$, suy ra $\triangle SAB$ vuông cân tại $A$.\\
				Do đó $SA=AB=\sqrt{6}$.\\
				Ta có $\triangle SAI$ vuông tại $A$, nên \allowdisplaybreaks
				\begin{eqnarray*}
					\tan \widehat{SIA}=\dfrac{SA}{AI}=\dfrac{\sqrt{6}}{\dfrac{3\sqrt{2}}{2}}=\dfrac{2\sqrt{3}}{3}.
				\end{eqnarray*}	
				Vậy $\tan \alpha=\dfrac{2\sqrt{3}}{3}$.	
				
			}{
				\begin{tikzpicture}[>=stealth,line join=round,line cap=round,font=\footnotesize,scale=1]
					\tikzset{
						pics/chopTamGiac/.style n args={4}{
							code={
								\tikzset{declare function={a=4;b=3;h=3;goc=-60;}}
								\path 
								(0,0)coordinate (#1)+(0:a)coordinate (#3)+
								(goc:b)coordinate (#2)+(90:h)coordinate (#4)
								
								;
								\foreach \pointo/\pointt in {#4/#1,#1/#2,#2/#3,#3/#4,#4/#2}{
									\draw[fill=black](\pointo)--(\pointt);
								}
								\foreach \pointo/\pointt in {#1/#3}{
									\draw[fill=black,dashed](\pointo)--(\pointt);
								}
								
							}
					}}
					\path 
					(0,0)coordinate (a)pic{chopTamGiac={A}{B}{C}{S}}
					($(B)!.5!(C)$)coordinate (I)
					($(S)!.3!(I)$)coordinate (H)
					pic[draw,angle radius=2mm]{right angle=C--A--S}
					pic[draw,angle radius=2mm]{right angle=A--H--S}
					;
					\foreach \pointo/\pointt in {A/I,A/H}{
						\draw[fill=black,dashed](\pointo)--(\pointt);
					}
					\foreach \pointo/\pointt in {S/I}{
						\draw[fill=black](\pointo)--(\pointt);
					}
					\foreach \point/\goc in {S/90,A/135,B/-90,C/-10,I/-45,H/45}{
						\draw[fill=black](\point)circle(.8pt)+(\goc:2mm)node[scale=.8]{$\point$};
					}
				\end{tikzpicture}
			}
			\item Gọi $H$ là hình chiếu của $A$ trên $SI$, khi đó $AH\perp SI$.\\
			Lại có $BC\perp AH$ và $AH\subset (SIA)$.\\
			Từ đó suy ra $AH\perp (SBC)$.\\
			Do đó $\beta= (SA,(SBC))=\widehat{ASI}$.\\
			Ta có $\triangle SIA$ vuông tại $A$, suy ra $SI=\sqrt{SA^2+AI^2}=\dfrac{\sqrt{42}}{2}$.\\
			Suy ra $\sin \beta =\sin \widehat{ASI}=\dfrac{AI}{SI}=\dfrac{\sqrt{21}}{7}$.\\
			Gọi $K$ là hình chiếu của $I$ trên $AB$.\\
			Khi đó $IK\perp  AB$ và $AI\perp SA$ $(SA \perp (ABC)$ và  $IK\subset(ABC))$, do đó $IK\perp (SAB)$.\\
			Nên $\mathrm{d}\left(I,(SAB)\right)=IK$.\\
			Ta có $\triangle AIB$ vuông tại $I$, có $IK$ là đường cao nên \allowdisplaybreaks
			\begin{eqnarray*}
				IK=\dfrac{AI\cdot IB}{AB}=\dfrac{\dfrac{3\sqrt{2}}{2}\cdot \dfrac{\sqrt{6}}{2}}{\sqrt{6}}=\dfrac{3\sqrt{2}}{4}.
			\end{eqnarray*}
		\end{enumerate}
	}
\end{bt}

\begin{bt}%[1H8H7-3]
% nguon: Archive/SP_DA_SauLuoi/1-TL-TN-DS-THPT-BaoLoc-LamDong-HKII-NH23-24.tex #2601
%[Dự án đề kiểm tra Toán 11 GHKI NH23-24- Phan Trung Hiếu]%[THPT Bảo Lộc- Lâm Đồng]
	Cho hình chóp $S.ABCD$ có đáy $ABCD$ là hình vuông cạnh $a$, $SA\perp (ABCD)$, góc giữa $SC$ và mặt phẳng $(ABCD)$ bằng $45^\circ$.
	\begin{enumerate}
		\item Tính theo $a$ thể tích khối chóp $S.ABCD$.
		\item Tính khoảng cách từ $A$ đến mặt phẳng $(SBC)$.
	\end{enumerate}
	\loigiai{
		\begin{center}
			\begin{tikzpicture}[thick, font=\footnotesize]
				\def\h{3}
				\path 
				(0,0) coordinate (A)
				($(A)+(90:\h)$) coordinate (S)
				(-1.5,-1.5) coordinate (D)
				(3,-1.5) coordinate (C)
				($(A)+(C)-(D)$) coordinate (B)
				($(S)!(A)!(B)$) coordinate (H)
				;
				\draw (S)--(D)--(C)--(B)--cycle (S)--(C);
				\draw[dashed] (A)--(S) (A)--(B) (A)--(D) (A)--(C) (A)--(H);
				\foreach \x/\g in {A/130,D/180,C/-45,B/0,S/90,H/60}
				\fill[black] 	(\x) circle (1pt)
				($(\g:3mm)+(\x)$) node {$\x$};
			\end{tikzpicture}
		\end{center}
		\begin{enumerate}
			\item Ta có $SA=AC\tan45^\circ =a\sqrt{2}$.\\
			Thể tích khối chóp $S.ABCD$ là $V=\dfrac{1}{3}S_{ABCD}\cdot SA=\dfrac{1}{3}\cdot a^2\cdot a\sqrt{2}=\dfrac{a^3\sqrt{2}}{3}$.
			\item Trong $(SAB)$, kẻ $AH\perp SB$.\\
			Ta có $\heva{&BC\perp AB\\&BC\perp SA}\Rightarrow BC\perp(SAB)$.\\
			Mặt khác, $\heva{&AH\perp SB\\&AH\perp SB}\Rightarrow AH\perp(SBC)\Rightarrow\mathrm{d}(A,(SBC))=AH$.\\
			Ta có $AH=\dfrac{SA\cdot AB}{\sqrt{SA^2+AB^2}}=\dfrac{a\sqrt{2}\cdot a}{\sqrt{(a\sqrt{2})^2+a^2}}=\dfrac{\sqrt{6}}{3}$.
		\end{enumerate}
	}
	
\end{bt}

\begin{bt}%[1H8H7-2]
% nguon: Archive/SP_DA_SauLuoi/1-TN-DS-TLN-THPT-TruongVinhKy-BenTre-HKII-2023-2024.tex #2611
%[Dự án đề kiểm tra Toán 11 HKII NH23-24- Nhật Thiện]%[THPT Trương Vĩnh Ký - Bến Tre]
	Cho hình chóp $S.ABC$ có đáy là tam giác đều cạnh bằng $2a$, $SA$ vuông góc với đáy. Góc giữa $(SBC)$ với đáy bằng $30^\circ$.
	\begin{enumerate}
		\item Tính thể tích khối chóp $S.ABC$ theo $a$.
		\item Xác định và tính khoảng cách từ $A$ đến mặt phẳng $(SBC)$.
	\end{enumerate}
\loigiai{
	\begin{center}
		\begin{tikzpicture}[scale=1, font=\footnotesize, line join=round, line cap=round, >=stealth]
\path (0,0) coordinate (A)	(1,-1) coordinate (B) (4,0) coordinate (C) (A)+(0,3) coordinate (S) ($(B)!.5!(C)$) coordinate (M) ($(S)!2/3!(M)$) coordinate (H);
\draw (S)--(A)--(B)--(C)--cycle (S)--(B) (S)--(M);
\draw[dashed] (A)--(C) (A)--(M) (A)--(H);
\foreach \x/\dinh/\y in {A/H/M,S/A/B} \draw[fill = gray!50] ($(\dinh)!3pt!(\x)$)--($(\dinh)!3pt!(\x)+(\dinh)!3pt!(\y)-(\dinh)$)--($(\dinh)!3pt!(\y)$)--(\dinh)--cycle; 
\foreach \p/\r in {A/120,S/90,B/-120,C/-60,M/-30,H/60}
\fill (\p) circle (1pt) node[shift={(\r:3mm)}]{$\p$};
\end{tikzpicture}
	\end{center}
	\begin{enumerate}
		\item Gọi $M$ là trung điểm cạnh $BC$.\\
		Ta có $\heva{&AM\perp BC\\&SA\perp BC}\Rightarrow BC\perp SM$.\\
		Khi đó 
		$\heva{&SM\perp BC,\ SM\subset (SBC)\\&AM\perp BC,\ AM\subset (ABC)\\&(SBC)\cap (ABC)=BC.}$\\
		Suy ra góc giữa $(SBC)$ và mặt đáy là góc giữa $SM$ và $AM$.\\
		Hay $\widehat{SMA}=30^\circ$.\\
		Ta có $AM=\dfrac{2a\sqrt{3}}{2}=a\sqrt{3}$.\\
		Suy ra $SA=AM\cdot \tan \widehat{SMA}=a$.\\
		Khi đó thể tích của khối chóp là $V=\dfrac{1}{3}\cdot SA\cdot S_{\triangle ABC}=\dfrac{1}{3}\cdot a\cdot \dfrac{a^2\sqrt{3}}{4}=\dfrac{a^3\sqrt{3}}{12}$.
		\item Gọi $H$ là hình chiếu của $A$ lên $SM$.\\
		Ta có $BC\perp (SAM)\Rightarrow BC\perp AH$.\\
		Khi đó $\heva{&AH\perp SM\\&AH\perp BC}\Rightarrow \mathrm{d}(A,(SBC))=AH$.\\
		Xét tam giác $SAM$ vuông tại $A$ có
		$$AH=\dfrac{AM\cdot SA}{\sqrt{AM^2+SA^2}}=\dfrac{a\sqrt{3}\cdot a}{\sqrt{(a\sqrt{3})^2+a^2}}=\dfrac{a\sqrt{3}}{2}.$$
		Vậy khoảng cách từ $A$ đến $(SBC)$ bằng $\dfrac{a\sqrt{3}}{2}$.
	\end{enumerate}
}
\end{bt}

\begin{bt}%[1H8V5-3]
% nguon: Archive/SP_DA_SauLuoi/1-TL-TN-THPT-ChauVanLiem-HKII-NH23-24.tex #2615
%[Dự án đề kiểm tra Toán 11 HKII NH23-24- Phạm Văn Long]%[THPT Châu Văn Liêm - Cần Thơ]
	Cho hình chóp tam giác đều $S. ABC$ có cạnh đáy bằng $a$, chiều cao bằng $2a$. Gọi $I$ là trung điểm của cạnh $BC$, $O$ là trọng tâm của tam giác $ABC$. Biết $SO$ vuông góc với mặt đáy.
	\begin{enumerate}
		\item Chứng minh rằng $(SOA) \perp(ABC)$.
		\item Tính khoảng cách từ điểm $O$ đến mặt phẳng $(SBC)$.
	\end{enumerate}
	\loigiai{
	\begin{center}
		\begin{tikzpicture}[=>stealth,line join=round,line cap=round, font=\footnotesize, scale=.8]
			\def\a{6}
			\def\goc{230}
			\def\b{3}
			\def\h{4.5}
			\path
			(0,0)coordinate (A)++(0:\a)coordinate (C)++(\goc:\b)coordinate (B)
			($(B)!.5!(C)$)coordinate (I)
			($(A)!2/3!(I)$)coordinate (O)
			(O)++(90:\h)coordinate (S)
			($(S)!2/3!(I)$)coordinate (H)			
			;
			\draw pic[draw,angle radius=2mm] {right angle = S--O--A};
			\draw pic[draw,angle radius=2mm] {right angle = A--I--B};
			\draw pic[draw,angle radius=2mm] {right angle = O--H--I};
			\draw (S)--(A)--(B)--(C)--cycle (S)--(B) (S)--(I);
			\draw[dashed] (C)--(A) (A)--(I) (S)--(O) (O)--(H);
			\foreach \x/\g in{A/180,B/-90,C/0,S/90,I/10,O/-90,H/10}
			\fill[black](\x)circle(1pt) ($(\x)+(\g:3mm)$)node{$\x$}
			;
		\end{tikzpicture}
	\end{center}
	\begin{enumerate}
		\item Ta có $\heva{&SO\perp (ABC)\\&SO\subset (SOA)}\Rightarrow (SOA) \perp(ABC)$.
		\item Kẻ $OH \perp SI$ tại $H$.\\
		Ta có $\heva{&BC\perp AI\\&BC  \perp SO\\&SO \cap AI=\{O\}\\&AI,SO \subset (SAI)}\Rightarrow BC \perp (SAI)\Rightarrow CB\perp OH$.\\
		Khi đó $\heva{&OH\perp SI\\&OH  \perp BC\\&SI \cap BC=\{I\}\\&SI,BC \subset (SBC)}\Rightarrow OH \perp (SBC)\Rightarrow \mathrm{d}(O,(SBC))=OH$.\\
		Ta có $AI=\dfrac{a\sqrt{3}}{2}\Rightarrow OI=\dfrac{a\sqrt{3}}{6}$.\\
		Khi đó $OH=\dfrac{SO\cdot OI}{\sqrt{SO^2+OI^2}}=\dfrac{2a\cdot \dfrac{a\sqrt{3}}{6}}{\sqrt{(2a)^2+\left(\dfrac{a\sqrt{3}}{6}\right)^2}}=\dfrac{2a}{7}$.\\
		Vậy $\mathrm{d}(O,(SBC))=\dfrac{2a}{7}$.
	\end{enumerate}
	}
\end{bt}

\begin{bt}%[1H8H5-3]
% nguon: Archive/SP_DA_SauLuoi/1-TN-DS-TLN-THPT-TruongVinhKy-TpHCM-HKII-NH23-24.tex #2623
%[Dự án đề kiểm tra Toán 11 CKII NH23-24- TRƯƠNG ĐĂNG KHOA]%[THPT TRƯƠNG VĨNH KÝ]
	Cho hình chóp $S.ABCD$ có đáy $ABCD$ là hình vuông cạnh bằng $1$, cạnh $SA$ vuông góc với đáy và $SA=\sqrt{3}$. Khoảng cách từ điểm $D$ đến $(SBC)$ bằng bao nhiêu?
	\dapso{$\dfrac{\sqrt{3}}{2}$}
	\loigiai{
		\immini{
			Ta có $\heva{
				&(SAB)\perp (SBC)\\
				&(SAB)\cap (SBC)=SB\\
				&\text{Trong } (SAB), \text{ kẻ } AH\perp SB \text{ tại } H 
			}\Rightarrow AH\perp (SBC)$.\\
			Ta có $AD\parallel BC\subset (SBC)$ nên $AD \parallel (SBC)$.\\
			Khi đó $$\mathrm{d}(D,(SBC))=\mathrm{d}(A,(SBC))=AH=\dfrac{SA\cdot AB}{\sqrt{SA^2+AB^2}}=\dfrac{\sqrt{3}\cdot 1}{\sqrt{3+1}}=\dfrac{\sqrt{3}}{2}.$$
		}{\begin{tikzpicture}[scale=0.8, font=\footnotesize,>=stealth]
				%Gán số liệu.
				\def\canhAD{3};\def\canhBA{1};\def\gocBAD{-130};\def\h{2};\def\xdinhS{0};
				%Gán tọa độ.
				\coordinate (A) at (0,0);
				\coordinate (B) at ($(A)+(\gocBAD:\canhBA)$);
				\coordinate (C) at ($(B)+(0:\canhAD)$);
				\coordinate (D) at ($(A)+(0:\canhAD)$);
				\coordinate (S) at ($(A)+(\xdinhS,\h)$);
				\coordinate (H) at ($(S)!0.6!(B)$);
				%Vẽ khối chóp S.ABCD.
				\draw (B)--(S)--(C)--cycle (S)--(D)--(C);
				\draw[dashed] (A)--(D) (S)--(A)--(B) (A)--(H);
				%Gán nhãn.
				\foreach \x/\y in {A/60,B/-90,C/-90,D/0,S/90,H/120}{\fill (\x) circle(1pt) ($(\x)+(\y:0.3cm)$) node{$\x$};}
		\end{tikzpicture}}
	}
\end{bt}

\begin{bt}%[1H8H7-9]
% nguon: Archive/SP_DA_SauLuoi/1-TN-DS-TLN-THPT-TruongVinhKy-TpHCM-HKII-NH23-24.tex #2624
%[Dự án đề kiểm tra Toán 11 CKII NH23-24- TRƯƠNG ĐĂNG KHOA]%[THPT TRƯƠNG VĨNH KÝ]
	Một tấm bìa hình vuông cắt bỏ ở mỗi góc tấm bìa một hình vuông có cạnh bằng $12$ cm rồi gấp lại thành hình hộp chữ nhật không nắp. Nếu dung tích cái hộp $4800$ cm$^3$ thì độ dài cạnh tấm bìa bằng bao nhiêu?
	\dapso{$44$ cm}
	\loigiai{
		Đặt độ dài cạnh của tấm bìa hình vuông là $x$ (cm), $(x>0)$.\\
		Độ dài cạnh của tấm bìa hình vuông sau khi cắt bỏ mỗi góc của tấm bìa một hình vuông có cạnh bằng $12$ cm là $x-24$ (cm).\\
		Theo đề bài có $V=12(x-24)^2=4800\Leftrightarrow x=44$ cm.
	}
\end{bt}

\begin{bt}%[1H4C1-6]
% nguon: Archive/SP_KhoLatex/ID6_old/1H3G1-5.tex #2661
%[Hồ Hà Đặng - dự án Latex hóa 11]
	Cho tứ diện $SABC$. Gọi $M,N,P$ là ba điểm lần lượt nằm trên ba cạnh $SA,SB,SC$ và $I$ là giao điểm của ba mặt phẳng $(BCM)$, $(CAN)$, $(ABP)$; $J$ là giao điểm của ba mặt phẳng $(ANP)$, $(BPM)$, $(CMN)$. Chứng minh $S,I,J$ thẳng hàng và $\dfrac{MS}{MA}+\dfrac{NS}{NB}+\dfrac{PS}{PC}+1=\dfrac{JS}{JI}$. 
	\loigiai{
		\begin{itemize}
			\item  Chỉ ra sự tồn tại của điểm $I$ và $J$.  
			Gọi $E=BP\cap CN,F=CM\cap AP$ và $T=AN\cap BM$.  
			Trong mặt phẳng $(BCM)$ có $BF\cap CT=I$, trong mặt phẳng $(ANP)$ có $NF\cap PT=J$.
			\item   \immini{Chứng minh $S,I,J$ thẳng hàng và $$\dfrac{MS}{MA}+\dfrac{NS}{NB}+\dfrac{PS}{PC}+1=\dfrac{JS}{JI}.$$
				Đặt $\overrightarrow{SA}=\overrightarrow{a},\overrightarrow{SB}=\overrightarrow{b},\overrightarrow{SC}=\overrightarrow{c}$ và $\overrightarrow{SM}=x\overrightarrow{MA},\overrightarrow{SN}=y\overrightarrow{NB},\overrightarrow{Sp}=z\overrightarrow{PC}$ với $x,y,z>0$. \\
				Ta có $\overrightarrow{SM}=x\overrightarrow{MA}=x\left(\overrightarrow{MS}+\overrightarrow{SA}\right)$, suy ra $(x+1)\overrightarrow{SM}=x\overrightarrow{SA}$ hay $\overrightarrow{SM}=\dfrac{x}{x+1}\overrightarrow{a}$. \\
				Tương tự, ta cũng có $\overrightarrow{SN}=\dfrac{y}{y+1}\overrightarrow{b}$ và $\overrightarrow{SP}=\dfrac{z}{z+1}\overrightarrow{c}$. \\ 
			}{
			\begin{tikzpicture}[scale=1, line join=round, line cap=round]
			\tkzDefPoints{0/0/A,3/-2/B,7/1/C,1.5/5/S}
			\coordinate (M) at ($(S)!.3!(A)$);
			\coordinate (N) at ($(S)!.6!(B)$);
			\coordinate (P) at ($(S)!.6!(C)$);
			\tkzInterLL(M,B)(A,N)\tkzGetPoint{T}
			\tkzInterLL(P,B)(C,N)\tkzGetPoint{E}
			\tkzInterLL(M,C)(A,P)\tkzGetPoint{F}
			\tkzInterLL(T,C)(B,F)\tkzGetPoint{I}
			\tkzInterLL(F,N)(T,P)\tkzGetPoint{J}
			\tkzDrawSegments(A,B B,C S,A S,B S,C B,M B,P A,N C,N N,P N,M)
			\tkzDrawSegments[dashed](A,C M,P C,M A,P F,N P,T B,F C,T S,I)
			\tkzDrawPoints[fill=black](A,B,C,S,M,N,P,E,F,J,J,T,I)
			\tkzLabelPoints[above](S,P,F,N)
			\tkzLabelPoints[below](B,I)
			\tkzLabelPoints[left](A,M,J)
			\tkzLabelPoints[right](C)
			\tkzLabelPoints[above right](E)
			\tkzLabelPoints[below left](T)
			%\tkzLabelPoints[right](C)
			\end{tikzpicture}}
		Do $T=AN\cap BM$ nên $\left\{\begin{aligned}
		&T\in BM \\
		&T\in AN
		\end{aligned}\right. \Rightarrow \left\{\begin{aligned}
		&\overrightarrow{ST}=\alpha \overrightarrow{SM}+\left(1-\alpha \right)\overrightarrow{SB} \\
		&\overrightarrow{ST}=\beta \overrightarrow{SN}+\left(1-\beta \right)\overrightarrow{SA}
		\end{aligned}\right. $  
		suy ra $$\alpha \overrightarrow{SM}+\left(1-\alpha \right)\overrightarrow{SB}=\beta \overrightarrow{SN}+\left(1-\beta \right)\overrightarrow{SA}$$  
		hay $\dfrac{\alpha x}{x+1}\overrightarrow{a}+\left(1-\alpha \right)\overrightarrow{b}=\dfrac{\beta y}{y+1}\overrightarrow{b}+\left(1-\beta \right)\overrightarrow{a}$.
				Vì $\overrightarrow{a},\overrightarrow{b}$ không cùng phương nên ta có  
				$$\left\{\begin{aligned}
				&\dfrac{\alpha x}{x+1}=1-\beta \\
				&\dfrac{\beta y}{y+1}=1-\alpha
				\end{aligned}\right. \Leftrightarrow \left\{\begin{aligned}
				&\alpha =\dfrac{x+1}{x+y+1} \\
				&\beta =\dfrac{y+1}{x+y+1}
				\end{aligned}\right. .$$ \\
				Suy ra $\overrightarrow{ST}=\dfrac{x}{x+y+1}\overrightarrow{a}+\dfrac{y}{x+y+1}\overrightarrow{b}$. \\
				Tương tự, ta có $\overrightarrow{SE}=\dfrac{y}{y+z+1}\overrightarrow{b}+\dfrac{z}{y+z+1}\overrightarrow{c}$, $\overrightarrow{SF}=\dfrac{z}{z+x+1}\overrightarrow{c}+\dfrac{x}{z+x+1}\overrightarrow{a}$. \\
				Làm tương tự như trên đối với hai giao điểm $I=BF\cap CT$ và $J=NF\cap PT$. Ta được  
				$$\overrightarrow{SI}=\dfrac{1}{x+y+z+1}\left(x\overrightarrow{a}+y\overrightarrow{b}+z\overrightarrow{c}\right)\, \text{và}\,\overrightarrow{SJ}=\dfrac{1}{x+y+z+2}\left(x\overrightarrow{a}+y\overrightarrow{b}+z\overrightarrow{c}\right).$$
				Suy ra $\overrightarrow{SJ}=\dfrac{x+y+z+1}{x+y+z+2}\overrightarrow{SI}$. Điều này chứng tỏ $S,I,J$ thẳng hàng. \\
				Mặt khác, ta có \begin{eqnarray*}
					\overrightarrow{SJ}=\dfrac{x+y+z+1}{x+y+z+2}\overrightarrow{SI} &\Leftrightarrow& (x+y+z+2)\overrightarrow{SJ}=(x+y+z+1)\overrightarrow{SI} \\&\Leftrightarrow& (x+y+z+2)\overrightarrow{SJ}=(x+y+z+1)\left(\overrightarrow{SJ}+\overrightarrow{JI}\right).
				\end{eqnarray*}
				Suy ra $$\overrightarrow{SJ}=(x+y+z+1)\overrightarrow{JI}$$ nên $$\dfrac{SJ}{JI}=x+y+z+1=\dfrac{SM}{MA}+\dfrac{SN}{NB}+\dfrac{SP}{PC}+1.$$ 
		\end{itemize}
		
}\end{bt}

\begin{bt}%[1H8C2-3]
% nguon: Archive/SP_KhoLatex/ID6_old/1H3G5-5.tex #2820
%[Nguyễn Đắc Giáp]
	Cho hình thang $ABCD$ có $\widehat{A}$ và $\widehat{B}$ là góc vuông, $AD=2a$, $AB=BC=a$. $S$ là đỉnh nằm trên $Ax$ vuông góc với mặt phẳng $(ABCD)$. Gọi $C'$ và $D'$ lần lượt là hình chiếu của $A$ trên $SC$ và $SD$.
	\begin{enumerate}
		\item Chứng minh $\widehat{SBC}=\widehat{SCD}=90^{\circ}$.
		\item Chứng minh $AD';AC'$ và $AB$ cùng nằm trong một mặt phẳng. Từ đó chứng minh rằng $C'D'$ đi qua một điểm cố định, khi $S$ di động trên $Ax$.
		\item Cho $AS=a\sqrt{2}$. Tính diện tích tứ giác $ABC'D'$.
		\item Dựng và tính độ dài đoạn vuông góc chung của hai đường thẳng $AB$ và $SC$ biết $AS=a\sqrt{2}$.
	\end{enumerate}
	\loigiai{
		\begin{enumerate}
			\item $\left.\begin{aligned}
			&SA\perp (ABCD) \\
			&AB\perp BC
			\end{aligned}\right\} \Rightarrow SB\perp BC$ (định lí 3 đường vuông góc) $ \Rightarrow \widehat{SBC}=90^{\circ}$. \\
			Nối $AC$, dễ thấy $AC\perp CD$. Do đó ta có $\left.\begin{aligned}
			&SA\perp (ABCD) \\
			&AC\perp CD
			\end{aligned}\right\} \Rightarrow SC\perp CD$ (định lí 3 đường vuông góc) $ \Rightarrow \widehat{SCD}=90^{\circ}$.
		\end{enumerate}
		\immini{
			\begin{enumerate}\setcounter{enumi}{1}
		
			\item Dễ thấy $AB\perp (SAD) \Rightarrow AB\perp SD$ (1) \\
			Theo giả thiết $AD'\perp SD$ (2) \\
			Ta đã chứng minh được $\left.\begin{aligned}
			&CD\perp AC \\
			&CD\perp SC
			\end{aligned}\right\} \Rightarrow CD\perp (SAD) \Rightarrow CD\perp AC'$.\\
			Giả thiết $AC'\perp SC$, vậy $AC'\perp (SCD) \Rightarrow AC'\perp SD'$ (3) \\
			Từ (1), (2), (3) suy ra $AB,AC',AD'$ nằm trong một mặt \\
			phẳng $(\alpha) $ qua $AB$ và vuông góc với $SD$. \\
			Gọi $I$ là giao điểm của $AB$ và $CD$ \\
			$\left.\begin{aligned}
			&I\in AB \Rightarrow I\in (\alpha) \\
			&I\in CD \Rightarrow I\in (SCD)
			\end{aligned}\right\} \Rightarrow I\in C'D'$ vì $C'D'=(\alpha) \cap (SCD)$ \\
			Vậy $C'D'$ đi qua điểm cố định $I$.
		\end{enumerate} 
		}
		{
			\begin{tikzpicture}[scale=0.8, line join=round, line cap=round]
			\tkzDefPoints{0/0/A,-1/-2/B,3/-2/C}
			\tkzDefPointBy[translation=from B to C](A)\tkzGetPoint{E}
			\tkzDefPointBy[translation=from B to C](E)\tkzGetPoint{D}
			\coordinate (S) at ($(A)+(0,5)$);
			\coordinate (C') at ($(S)!0.5!(C)$);
			\tkzInterLL(C,D)(A,B)\tkzGetPoint{I}
			\tkzInterLL(S,D)(I,C')\tkzGetPoint{D'}
			\coordinate (F) at ($(S)!0.66!(E)$);
			\tkzDefLine[parallel=through F](C,E)\tkzGetPoint{f}
			\tkzInterLL(F,f)(C,S)\tkzGetPoint{P}
			\tkzDefLine[parallel=through P](A,F)\tkzGetPoint{p}
			\tkzInterLL(P,p)(A,B)\tkzGetPoint{Q}
			\coordinate (C_1) at ($(I)!0.5!(D')$);
			\tkzDrawPolygon(S,I,D)
			\tkzDrawSegments(S,C I,D' C,C_1)
			\tkzDrawSegments[dashed](S,B B,C A,S A,I A,D C,E A,C P,Q P,F A,F S,E A,C' A,D')
			\tkzDrawPoints[fill=black](D,C,A,B,S,I,P,Q,C',D',E,F,C_1)
			\tkzLabelPoints[above](S,Q)
			\tkzLabelPoints[above right](E,F,D')
			\tkzLabelPoints[below](I,C)
			\tkzLabelPoints[left](A,B,C_1)
			\tkzLabelPoints[right](D)
			\tkzLabelPoints[right](C',P)
			\tkzMarkRightAngles[size=0.15](S,D',A S,A,B)
			\end{tikzpicture}	
		}
	
		\begin{enumerate}\setcounter{enumi}{2}
		\item Ta có: ${S}_{ABC'D'}={S}_{AID'} -{S}_{BIC'}$.\\
		Dễ thấy $BC\parallel AD$ và $BC=\dfrac{1}{2}AD$ nên $AI=2AB=2a$. \\
		Trong tam giác vuông $AID',AD'$ là đường cao nên:
		$$\dfrac{1}{AD'^2}=\dfrac{1}{SA^2}+\dfrac{1}{AD^2}=\dfrac{1}{{\left(a\sqrt{2}\right)}^2}+\dfrac{1}{(2a)^2}.$$
		Suy ra $AD'^2=\dfrac{4a^2}{3} \Rightarrow AD'=\dfrac{2a}{\sqrt{3}}$ nên $S_{AID'}=\dfrac{1}{2}AI\cdot AD'=\dfrac{1}{2}\cdot 2a\cdot \dfrac{2a}{\sqrt{3}}=\dfrac{2a^2\sqrt{3}}{3}$. \\
		Ta cũng có, do $\widehat{SAC}=90^{\circ} \Rightarrow $ tam giác $SAC$ vuông. \\
		Ta cũng có $AC=a\sqrt{2} \Rightarrow AS=AC \Rightarrow $ tam giác $SAC$ vuông, cân mà $AC'\perp SC \Rightarrow C'$ là trung điểm của $SC$. \\
		Từ $C$ trong mặt phẳng $(CID)$ kẻ $CC_1\parallel SD$. Dễ thấy $ \Delta SC'D'= \Delta CC'C_1 \Rightarrow C_1C'=C'D'$. \\
		Trong tam giác $AID$, thì $C$ là trung điểm của $ID$, suy ra $C_1$ là trung điểm của $ID'$. Do vậy $C'D'=C_1C'$ mà $C_1D'=\dfrac{1}{2}ID' \Rightarrow \dfrac{IC'}{ID'}=\dfrac{3}{4}$. Nếu từ $C'$ kẻ đường cao $AH$ của tam giác $BIC'$ thì $H\in AB$ và $C'H\parallel AD'$. Từ đó ta có:
		$$\dfrac{C'H}{AD'}=\dfrac{IC'}{ID'}+\dfrac{3}{4} \Rightarrow C'H=\dfrac{3a}{2\sqrt{3}}.$$
		Vậy $$S_{BIC'}=\dfrac{1}{2}BI\cdot C'H=\dfrac{1}{2}\cdot a\cdot \dfrac{3a}{2\sqrt{3}}=\dfrac{3a^2\sqrt{3}}{12}$$ và $$S_{ABC'D'}=\dfrac{2a^2\sqrt{3}}{3}-\dfrac{3a^2\sqrt{3}}{12}=5\dfrac{3a^2\sqrt{3}}{12}.$$
		\item Gọi $E$ là trung điểm của $AD$. \\
		Từ $A$ kẻ $AF\perp SE$. Qua $F$ kẻ đường thẳng song song với $CE$, cắt $SC$ ở $P$.\\
		Từ $P$ kẻ $PQ\parallel AF$ cắt $AB$ ở $Q$. $PQ$ là đoạn vuông góc chung của $SC$ và $AB$. \\
		Thật vậy $CE\perp AD \Rightarrow CE\perp (SAD) \Rightarrow (SEC)\perp (SAD);AF\in (SAD)$ và do $AF\perp SE$  \\
		$\Rightarrow \left.\begin{aligned}
		&AF\perp (SEC) \\
		&PQ\parallel AF
		\end{aligned}\right\} \Rightarrow PQ\perp (SEC);SC\in (SEC)$. \\
		Vậy $PQ\perp SC$. Mặt khác $\left.\begin{aligned}
		&AB\perp (SAD) \Rightarrow AB\perp AF \\
		&PQ\parallel AF
		\end{aligned}\right\} \Rightarrow PQ\perp AB$ \\
		Ta có: $PQ=AF$ mà $\dfrac{1}{AF^2}=\dfrac{1}{SA^2}+\dfrac{1}{AE^2}=\dfrac{1}{2a^2}+\dfrac{1}{a^2} \Rightarrow AF^2=\dfrac{2a^2}{3} \Rightarrow AF=\dfrac{a\sqrt{6}}{3}$.
	\end{enumerate} 
	}
\end{bt}

\begin{bt}%[1H8C4-5]
% nguon: Archive/SP_KhoLatex/ID6_old/1H3G5-5.tex #2821
%[Nguyễn Đắc Giáp]
	Cho tam giác đều $SAD$ và hình vuông $ABCD$ cạnh $a$ nằm trong hai mặt phẳng vuông góc. Gọi $I$ là trung điểm của $AD,M$ là trung điểm của $AB,F$ là trung điểm của $SB$ và $K$ là giao điểm của $CM$ và $BI$.
	\begin{enumerate}
		\item Chứng minh $(CMF)\perp (SIB)$.
		\item Tính $BK$ và $KF$ và suy ra rằng tam giác $BKF$ cân.
		\item Dựng và tính độ dài đoạn vuông góc chung của hai đường thẳng $AB$ và $SD$.
		\item Tính khoảng cách giữa hai đường thẳng $CM$ và $SA$.
	\end{enumerate}
	\loigiai{
			\immini{
			\begin{enumerate}
			\item Ta có $\left.\begin{aligned}
			&SI\perp AD \\
			&(SAD)\perp (ABCD)
			\end{aligned}\right\} \Rightarrow SI\perp (ABCD) \Rightarrow SI\perp CM$  (1). \\
			$\triangle BCM=\triangle ABI \Rightarrow \widehat{B_1}={\widehat{C}}_1$ mà \\
			$\widehat{C_1}+{\widehat{M}}_1=90^{\circ} \Rightarrow \widehat{B_1}+{\widehat{M}}_1=90^{\circ} \Rightarrow M\widehat{K}B=90^{\circ}$. \\
			Hay $BI\perp CM$  (2). \\
			Từ (1) và (2) suy ra $\left.\begin{aligned}
			&CM\perp (SIB) \\
			&CM\subset (CMF)
			\end{aligned}\right\} \Rightarrow (CMF)\perp (SIB)$.
		\end{enumerate}
		}
	{
	\begin{tikzpicture}[scale=1, line join=round, line cap=round]
	\tkzDefPoints{0/0/D,-2.2/-1.6/A,2.5/-1.6/B}
	\coordinate (C) at ($(B)+(D)-(A)$);
	\tkzDefMidPoint(D,A)\tkzGetPoint{I}
	\coordinate (S) at ($(I)+(0,5)$);
	\tkzDefMidPoint(A,B)\tkzGetPoint{M}
	\tkzDefMidPoint(S,B)\tkzGetPoint{F}
	\tkzDefMidPoint(S,D)\tkzGetPoint{J}
	\tkzInterLL(C,M)(B,I)\tkzGetPoint{K}
	\coordinate (H) at ($(F)!-1.2!(K)$);
	\tkzDrawPolygon(S,A,B,C)
	\tkzDrawSegments(S,B M,F F,C S,H F,H)
	\tkzDrawSegments[dashed](S,D S,I A,D D,C B,I M,C F,K A,J)
	\tkzDrawPoints[fill=black](S,A,B,C,D,I,J,M,K,H,F)
	\tkzLabelPoints[above](S)
	\tkzLabelPoints[below](A,B,M)
	\tkzLabelPoints[left](D,I)
	\tkzLabelPoints[above right](H,F,K)
	\tkzLabelPoints[right](J,C)
	\tkzMarkRightAngles[size=0.15,fill=gray!20](S,H,F S,I,D)
	\end{tikzpicture}	
}
		
	\begin{enumerate}\setcounter{enumi}{1}
		\item Trong tam giác vuông $MBC,KB$ là đường cao nên 
		$$\dfrac{1}{KB^2}=\dfrac{1}{MB^2}+\dfrac{1}{CB^2}=\dfrac{1}{\left(\dfrac{a}{2}\right)^2}-\dfrac{1}{a^2} \Rightarrow KB^2=\dfrac{a^2}{5} \Rightarrow KB=\dfrac{a\sqrt{5}}{5}.$$
		Ta dễ dàng chứng minh được: $SB=SC;SC^2=SD^2+DC^2$ (vì $SD\perp DC$) $ \Rightarrow SB=SC=a\sqrt{2}$. \\
		Trong tam giác $SBC,CF$ là trung tuyến nên ta có: $$SC^2+BC^2=2\cdot CF^2+\dfrac{SB^2}{2} \Rightarrow 2a^2+a^2=2\cdot CF^2+a^2 \Rightarrow CF^2=a^2.$$
		Vì $CK\perp (SIB) \Rightarrow CK\perp FK \Rightarrow \triangle CKF$ vuông tại $K$ nên $FK^2=CF^2-KC^2$. \\
		Ta lại có: $$KC^2=BC^2-KB^2=a^2-\dfrac{a^2}{5}=\dfrac{4a^2}{5}.$$
		Do đó $$KF^2=a^2-\dfrac{4a^2}{5}=\dfrac{a^2}{5} \Rightarrow KF=\dfrac{a\sqrt{5}}{5}.$$
		Vậy $KB=KF \Rightarrow \triangle BKF$ cân đỉnh $K$. \\
		\item Ta có $SI\perp (ABCD) \Rightarrow \left.\begin{aligned}
		&SI\perp AB \\
		&AD\perp AB
		\end{aligned}\right\} \Rightarrow AB\perp (SAD)$. \\
		Ta kẻ đường cao $AJ$ của tam giác $SAD$ thì $AB\perp (SAD) \Rightarrow AB\perp AJ$. \\
		$SD\perp IJ \Rightarrow AJ$ là đoạn vuông góc chung của $AB$ và $SD$. \\
		Ta có ngay: $AJ=\dfrac{a\sqrt{3}}{2}$ (vì $\triangle SAD$ đều cạnh $a$). \\
		\item Dễ thấy $SA\parallel FM \Rightarrow SA\parallel (CMF)$ mà $(CMF)\supset CM$. Do vậy khoảng cách giữa $SA$ và $CM$ cũng là khoảng cách từ một điểm trên $SA$ đến $(CMF)$. Ta chọn điểm $S$ chẳng hạn. Từ $S$ kẻ $SH\perp (CMF);SH$ là khoảng cách giữa $SA$ và $CM$. Ta đã có 
		$$\left.\begin{aligned}
		&(SIB)\perp (CMF) \\
		&(SIB)\cap (CMF)=KF \\
		&SH\perp (SIB),SH\perp (CMF)
		\end{aligned}\right\} \Rightarrow H\in KF.$$
		Ta có: $SH=SF\cdot \sin \widehat{SFH}$ vì $\widehat{SFH}=\widehat{BFK}$ (đối đỉnh) mà $\widehat{BFK}=\widehat{FBK}$ ($\triangle FBK$ cân) \\
		$ \Rightarrow \widehat{SFK}=\widehat{FBK} \Rightarrow \sin \widehat{SFH}=\sin \widehat{FBK}$. Do đó\\
		$$\sin \widehat{FPK}=\dfrac{SI}{SB}=\dfrac{\dfrac{a\sqrt{3}}{2}}{a\sqrt{2}}=\dfrac{\sqrt{3}}{2\sqrt{2}} \Rightarrow SH=\dfrac{a\sqrt{2}}{2}\cdot \dfrac{\sqrt{3}}{2\sqrt{2}}=\dfrac{a\sqrt{3}}{4}.$$ 
	\end{enumerate}
		
	}
\end{bt}

\begin{bt}%[1H8C2-3]
% nguon: Archive/SP_KhoLatex/ID6_old/1H3G5-5.tex #2822
%[Nguyễn Đắc Giáp]
	Cho tứ diện $OABC$ có $OA = OB = OC = a$, $\widehat{AOB}=\widehat{AOC}=60^{\circ},\widehat{BOC}=90^{\circ}$. 
	\begin{enumerate}
		\item Chứng minh rằng $ABC$ là tam giác vuông.
		\item Chứng minh $OA$ vuông góc $BC$.
		\item Gọi $I, J$ là trung điểm $OA$ và $BC$. Chứng minh $IJ$ là đoạn vuông góc chung $OA$ và $BC$.
	\end{enumerate}
	\loigiai{
		\immini{
			\begin{enumerate}
			\item Ta có $OA = OB = OC = a$, $\widehat{AOB}=\widehat{AOC}=60^{\circ}$ nên $AOB$ và $AOC$ đều cạnh $a$. (1) \\
			Lại có $\widehat{BOC}=90^{\circ}$ $\Rightarrow  \triangle BOC$ vuông tại $O$ và $BC=a\sqrt{2}$. (2) \\
			Trong $\triangle ABC$ có $$AB^2+AC^2=a^2+a^2=2a^2={\left(a\sqrt{2}\right)}^2=BC^2.$$
			Suy ra tam giác $ABC$ vuông tại $A$.
			\item Do $J$ là trung điểm $BC$, $\triangle ABC$ vuông cân tại $A$ nên $AJ\perp BC$. Tam giác $OBC$ vuông cân tại $O$ nên $OJ\perp BC$. \\
			Suy ra $BC\perp (OAJ) \Rightarrow OA\perp BC$.
		\end{enumerate}
		}
	{
	\begin{tikzpicture}[scale=1,line join=round,line cap=round]
	\tkzDefPoints{0/0/B,1/-1.5/A,4/0/C}
	\coordinate (J) at ($(B)!0.5!(C)$);
	\coordinate (O) at ($(J)+(0,3)$);
	\coordinate (I) at ($(O)!0.5!(A)$);
	\tkzDrawPolygon(O,B,A,C)
	\tkzDrawSegments(O,A)
	\tkzDrawSegments[dashed](B,C O,J I,J A,J)
	\tkzDrawPoints[fill=black](A,B,C,O,J)
	\tkzLabelPoints[above](O)
	\tkzLabelPoints[below](J,A)
	\tkzLabelPoints[left](B,I)
	\tkzLabelPoints[right](C)
	\tkzMarkRightAngles[size=0.15](O,J,C B,A,C B,O,C)
	\tkzMarkSegments[mark=|](J,B J,C)
	\tkzMarkSegments[mark=||](A,I I,O)
	\end{tikzpicture}	
}
		
		\begin{enumerate}\setcounter{enumi}{2}
			\item Từ câu 2) ta có $ IJ\perp BC$. \\
			Ta có $\triangle ABC=\triangle OBC(c.c.c) \Rightarrow AJ= OJ$. (3) \\
			Từ (3) ta có tam giác $JOA$ cân tại $J$, $IA = IO$ nên $IJ \perp OA$. (4) \\
			Từ (3) và (4) ta có $IJ$ là đoạn vuông góc chung của $OA$ và $BC$.
		\end{enumerate}
	}
\end{bt}

\begin{bt}%[1H8H5-3]
% nguon: Archive/SP_KhoLatex/ID6_old/19-OT-KhoangCach.tex #2843
 
		Cho hình chóp $S.ABC$ có đáy là tam giác vuông đỉnh $A$, $AC=a$, $AB=a\sqrt{3}$, $SC$ vuông góc với mặt phẳng đáy. Tính khoảng cách từ $A$ đến mặt phẳng $(SBC)$. 
	\loigiai{
		\immini
		{
			Dựng $AH\perp BC$ tại $H$.\\
			Từ $\heva{& AH\perp BC \\ & AH\perp SC}\Rightarrow AH\perp (SBC)$ tại $H$. \\
			Suy ra $\mathrm{d}(A,(SBC))=AH$. \\ 
			Áp dụng hệ thức lượng cho $\triangle ABC$ vuông tại $A$, đường cao $AH$, ta có
			\begin{align*}
				\dfrac{1}{AH^2}=\dfrac{1}{AB^2}+\dfrac{1}{AC^2}=\dfrac{1}{a^2}+\dfrac{1}{3a^2}=\dfrac{4}{3a^2}\Rightarrow AH=\dfrac{a\sqrt{3}}{2}.
			\end{align*}
			Vậy $\mathrm{d}(A,(SBC))=AH=\dfrac{a\sqrt{3}}{2}$.
		}
		{
			\begin{tikzpicture}[line join = round, line cap = round,>=stealth,font=\footnotesize,scale=0.7]
				\tkzDefPoints{0/0/C,1.5/-2/A,5/0/B}
				\coordinate (S) at ($(C)+(0,4)$);
				\coordinate (H) at ($(B)!0.6!(C)$);
				\pgfresetboundingbox
				\tkzDrawSegments(S,A S,B S,C A,B A,C)
				\tkzDrawSegments[dashed](B,C A,H)
				\tkzMarkRightAngles[size=0.25](S,C,B S,C,A A,H,B)
				\foreach \x/\g in {S/90,A/-90,B/0,C/180,H/90} \fill[black](\x) circle (1pt)+(\g:3mm) node{$\x$};
			
\end{tikzpicture}
		}
	}
\end{bt}

\begin{bt}%[1H8H5-3]
% nguon: Archive/SP_KhoLatex/ID6_old/19-OT-KhoangCach.tex #2844
%[Đề thi thử sở GD-ĐT Hòa Bình]%[Đỗ Viết Lân - 12 EX10 2021]
	Cho hình hộp chữ nhật $ABCD.A'B'C'D'$ có $AB=a \sqrt{3}$, $AD=a$. Khoảng cách từ $A$ đến mặt phẳng $\left(BDD'B'\right)$ bằng
	\loigiai{
		\immini{Dựng $AO \perp BD$, $O \in BD$.\\
			Ta có $\heva{&AO \perp BD \\ &AO \perp BB'} \Rightarrow AO \perp\left(BDD'B'\right)$\\
			$\Rightarrow \mathrm{d}\left(A,\left(BDD'B'\right)\right)=AO$.\\
			Trong tam giác vuông $ABD$ ta có\\
			$\dfrac{1}{AO^{2}}=\dfrac{1}{AB^{2}}+\dfrac{1}{AD^{2}}=\dfrac{1}{a^{2}}+\dfrac{1}{3a^{2}} \Rightarrow AO=\dfrac{a\sqrt{3}}{2}$.\\
			Vậy $\mathrm{d}\left(A,\left(BDD'B'\right)\right)=\dfrac{a \sqrt{3}}{2}$.}{
			\begin{tikzpicture}[font=\footnotesize, line join=round, line cap=round, >=stealth]
				\def\r{3.5}
				\path
				(0:0) coordinate (A)
				++(0:\r) coordinate (B)
				++(30:0.5*\r) coordinate (C)
				($(A)+(C)-(B)$) coordinate (D)
				($(B)!0.7!(D)$) coordinate (O)
				\foreach \x in {A,B,C,D}{
					(\x)++(90:\r) coordinate (\x')
				}
				;
				\draw 
				(A)--(B)--(B')--(A')--cycle
				(A')--(D')--(C')--(B')--(D')
				(B)--(C)--(C')
				;
				\draw[dashed]
				(C)--(A)--(D)--cycle
				(D)--(D') (A)--(O)
				(D)--(B)
				;
				\foreach \x/\g in 
				{A/-135,B/-45,C/0,D/160,A'/150,B'/-30,C'/30,D'/135}
				\draw [fill=black] (\x) circle (.03)+ (\g:.3) 
				node {$\x$};
				\draw [fill=black] (O) circle (.03)+ (10:.4) 
				node {$O$};
				\tkzMarkRightAngles(A,O,B)
		
\end{tikzpicture}}
	}
\end{bt}

\begin{bt}%[1H8H5-3]
% nguon: Archive/SP_KhoLatex/ID6_old/19-OT-KhoangCach.tex #2845
 
		Cho hình chóp $S.ABC$ có đáy là tam giác vuông đỉnh $A$, $AC=a$, $SC$ vuông góc với mặt phẳng đáy và $SC=a$. Tính khoảng cách từ $C$ đến mặt phẳng $(SAB)$.
	\loigiai{
		\immini
		{
			Dựng $CH\perp SA$ tại $H$.\\
			Từ $\heva{& AB\perp AC \\ & AB\perp SC}\Rightarrow AB\perp (SAC)\Rightarrow AB\perp CH$.\\
			Từ $\heva{& CH\perp SA \\ & CH\perp AB}\Rightarrow CH\perp (SAB)$ tại $H$.\\
			Suy ra $\mathrm{d}(C,(SAB))=CH$.\\
			Áp dụng hệ thức lượng cho $\triangle SCA$ (vuông tại $C$) có
			\begin{align*}
				\dfrac{1}{CH^2}=\dfrac{1}{CS^2}+\dfrac{1}{CA^2}=\dfrac{1}{a^2}+\dfrac{1}{a^2}=\dfrac{2}{a^2}\Rightarrow CH=\dfrac{a\sqrt{2}}{2}.
			\end{align*}
			Suy ra $\mathrm{d}(C,(SAB))=\dfrac{a\sqrt{2}}{2}$.
		}
		{
			\begin{tikzpicture}[line join = round, line cap = round,>=stealth,font=\footnotesize,scale=0.7]
				\tkzDefPoints{0/0/C,1.5/-2/A,5/0/B}
				\coordinate (S) at ($(C)+(0,4)$);
				\coordinate (H) at ($(S)!0.6!(A)$);
				\pgfresetboundingbox
				\tkzDrawSegments(S,A S,B S,C A,B A,C C,H)
				\tkzDrawSegments[dashed](B,C)
				\tkzMarkRightAngles[size=0.25](S,C,B S,C,A C,H,S)
				\foreach \x/\g in {S/90,A/-90,B/0,C/180,H/0} \fill[black](\x) circle (1pt)+(\g:3mm) node{$\x$};
			
\end{tikzpicture}
		}
	}
\end{bt}

\begin{bt}%[1H8H5-3]
% nguon: Archive/SP_KhoLatex/ID6_old/19-OT-KhoangCach.tex #2846
%[Thi thử lần 4, THPT Triệu Sơn 3, Thanh Hóa, 2021]%[Phan Duy Anh - Huỳnh Đức Vũ, dự án EX12]
	Cho hình lăng trụ đứng $ABC.A'B'C'$ có độ dài cạnh bên bằng $3 $, đáy $ABC$ là tam giác vuông tại $B$ và $AB=2$. Tình khoảng cách từ $A$ đến mặt phẳng $\left(A'BC\right)$.
	\loigiai{
		\immini{
			Dựng $AH\perp A'B$. Ta có
			\[\heva{&BC\perp AB\\&BC\perp AA'} \Rightarrow BC\perp (AA'B)\Rightarrow BC\perp AH.\]
			Ta có
			\[\heva{&AH\perp BC\\&AH\perp A'B}\Rightarrow AH\perp (A'BC).\]
			Suy ra khoảng cách	từ $A$ đến mặt phẳng $\left(A'BC\right)$ là $\mathrm{\,d} =AH$.
		}
		{
			\begin{tikzpicture}[scale=1,>=stealth, font=\footnotesize, line join=round, line cap=round]
				\tkzDefPoints{0/0/A,1.1/-1.5/B,3.5/0/C}
				\coordinate (A') at ($(A)+(0,3.2)$);
				\coordinate (H) at ($(B)!(A)!(A')$);
				\tkzDefPointsBy[translation=from A to A'](B,C){B'}{C'}
				\tkzDrawPolygon(A,B,C,C',B',A')
				\tkzDrawSegments(A',C' B',B A',B A,H)
				\tkzDrawSegments[dashed](A,C A',C)
				\tkzDrawPoints[fill=black,size=1](A,C,B,A',B',C',H)
				\tkzMarkRightAngles(A,B,C A,H,B)
				\foreach \x/\g in
				{A/180,A'/180,B/-90,B'/-40,C/0,C'/0,H/80}
				\fill[black](\x) circle (1pt)
				($(\x)+(\g:5mm)$) node{$\x$};
		
\end{tikzpicture}}
		\noindent 
		Theo hệ thức lượng trong tam giác vuông, ta có 
		\allowdisplaybreaks
		\begin{eqnarray*}
			&&\dfrac{1}{AH^2}=\dfrac{1}{AB^2}+\dfrac{1}{AA'^2}=\dfrac{1}{4}+\dfrac{1}{9}=\dfrac{13}{36}\\
			&\Rightarrow & AH=\dfrac{6\sqrt{13}}{13}.
		\end{eqnarray*}
		Vậy khoảng cách	từ $A$ đến mặt phẳng $\left(A'BC\right)$ là $\dfrac{6\sqrt{13}}{13}$.
	}
\end{bt}

\begin{bt}%[1H8V5-3]
% nguon: Archive/SP_KhoLatex/ID6_old/19-OT-KhoangCach.tex #2847
%[Thi thử, THPT Hải Lăng (Quảng Trị) 2021]%[Lương Như Quỳnh, 12-EX-11]
	Cho hình chóp tam giác đều $S.ABC$ có đáy $ABC$ với cạnh đáy bằng $a $. Gọi $G$ là trọng tâm $ABC$. Góc giữa mặt bên với mặt đáy bằng $60^{\circ}$. Tính khoảng cách từ điểm $G$ đến mặt phẳng $(SBC)$.
	\loigiai{
		\immini{
			Gọi $ I $ là trung điểm $ BC $.\\
			Trong $ (SAI) $, kẻ $ GH\perp SI $.\\
			Ta có $ \heva{&BC \perp AI \\ &BC \perp SI} $ \\
			suy ra $ BC \perp(SAI) \Rightarrow BC \perp GH$.\\
			Mà $SI \perp GH \Rightarrow GH \perp(SBC) \Rightarrow GH=\mathrm{d}(G,(SBC))$.\\
			Ta có $ \heva{&(SBC)\cap (ABC)=BC \\ &SI \subset (SBC),\,SI \perp BC\\&AI \subset (ABC),\,AI \perp BC} $ suy ra \\
			$((SBC),(ABC))=(SI,AI)=\widehat{SIA}=\widehat{SIG}=60^\circ$.
		}
		{
			\begin{tikzpicture}[scale=.7, font=\footnotesize, line join=round, line cap=round, >=stealth]
				\path
				(-1,0)coordinate(A)
				(2,-2)coordinate(B)
				(7,0)coordinate(C)
				($(C)!1/2!(B)$)coordinate(I)
				($(C)!1/2!(B)$)coordinate(N)
				($(A)!2/3!(I)$)coordinate(G)
				($(G)+(0,6)$)coordinate(S)
				($(S)!2/3!(I)$)coordinate(H)
				;
				\draw[dashed](C)--(A)--(I) (C)--(N) (S)--(G)--(H);
				\draw (I)--(S)--(B) (S)--(A)--(B)--(C)--cycle;
				\foreach \p/\g in {A/180,C/0,B/-90,S/90,G/-45,I/-45,H/0}
				\draw[fill=black] (\p) circle(.5pt) +(\g:.3) node{$\p$};
				\foreach \i/\j/\k/\t in {B/A/C/6,G/S/C/6,H/G/I/6}{
					\def\dgiua{\i}\def\dmot{\j}\def\dhai{\k}\def\tyso{\t pt}
					\draw ($(\dgiua)!\tyso!(\dmot)$)--($(\dgiua)!2!($($(\dgiua)!\tyso!(\dmot)$)!.5!($(\dgiua)!\tyso!(\dhai)$)$)$)--($(\dgiua)!\tyso!(\dhai)$);}
			
\end{tikzpicture}
		}
		\noindent 
		Ta có $ GI=\dfrac{1}{3}AI =\dfrac{a\sqrt{3}}{6}\Rightarrow GH=GI\cdot \sin 60^\circ =\dfrac{a\sqrt{3}}{6}\cdot \dfrac{\sqrt{3}}{2}=\dfrac{a}{4} $.\\
		Vậy $\mathrm{d}(G,(SBC)) = \dfrac{a}{4}$.
	}
\end{bt}

\begin{bt}%[1H8V5-3]
% nguon: Archive/SP_KhoLatex/ID6_old/19-OT-KhoangCach.tex #2848
%[Thi thử, Sở GD và ĐT - Kiên Giang, 2021]%[Lại Thị Hảo, 12EX11]
	Cho hình chóp $S.ABCD$ có $SA \perp (ABCD)$, đáy $ABCD$ là hình vuông, biết $AB=1$, $SA=2$. Tính khoảng cách từ $C$ đến mặt phẳng $(SBD)$.
	\loigiai{
		\immini
		{
			Gọi $O$ là tâm hình vuông $ABCD$, khi đó $AO=CO$ và $AC \cap(SBD)=O$. \\
			Ta có $\heva{&SA \perp (ABCD)\\& BD \subset (ABCD)} \Rightarrow SA \perp BD$.\\
			Mặt khác $ABCD$ là hình vuông nên $AC \perp BD$ hay $AO \perp BD$.\\
			Khi đó $\heva{&BD \perp SA \\& BD \perp AO \\&SA, AO \subset (SAO)}\Rightarrow \heva{&(SAO) \perp(SBD). \\
				&(SAO) \cap(SBD)=SO.}$\\
			Dựng $AH\perp SO$ tại $H$. Khi đó, 
			$\heva{&AH \perp SO \\& AH \perp BD} \Rightarrow AH\perp (SBD)$ tại $H$.  
		}
		{
			\begin{tikzpicture}[font=\footnotesize, line join=round, line cap=round, >=stealth,scale=0.7]
				\path
				(0,0) coordinate (A) 
				(-2,-2) coordinate (B)
				(3,-2) coordinate (C)
				(5,0) coordinate (D)
				(0,4) coordinate (S)
				($(A)!.5!(C)$) coordinate (O)
				($(O)!.3!(S)$) coordinate (H)
				;
				\draw 
				(S)--(B)--(C)--(D)--cycle
				(S)--(C);
				\tkzMarkRightAngles(S,A,D);
				\tkzMarkRightAngles(S,H,A);
				\draw [dashed] (S)--(A)--(D)--(B)--(A)--(C) (S)--(O) (A)--(H); 
				\foreach \x/\g in
				{A/150,B/180,C/0,D/0,S/90,O/-90,H/30}
				\fill[black](\x) circle (1.5pt)
				($(\x)+(\g:3mm)$) node{$\x$};
			
\end{tikzpicture}
		}\noindent
		Suy ra $\mathrm{d}(A,(SBD))=AH.$\\
		Ta có $\dfrac{\mathrm{d}(C, (SBD))}{\mathrm{d}(A, (SBD))}=\dfrac{CO}{AO}=1.$ Suy ra $\mathrm{d}(C,(SBD))=\mathrm{d}(A,(SBD))=AH.$\\
		Ta có $AC=\sqrt{1^2+1^2}=\sqrt{2}$.\\
		Trong tam giác $SAO$ vuông tại $A$, ta có: 
		$$\dfrac{1}{AH^{2}}=\dfrac{1}{AO^{2}}+\dfrac{1}{SA^{2}}
		=\dfrac{1}{\left(\dfrac{AC}{2}\right)^{2}}+\dfrac{1}{2^{2}}=\dfrac{1}{\left(\dfrac{\sqrt{2}}{2}\right)^{2}}+\dfrac{1}{2^{2}}=\dfrac{9}{4}.$$
		$\Rightarrow A H=\dfrac{2}{3} \Rightarrow d(C,(S B D))=\dfrac{2}{3}$.
	} 
\end{bt}

\begin{bt}%[1H8V5-3]
% nguon: Archive/SP_KhoLatex/ID6_old/19-OT-KhoangCach.tex #2849
%[Đề GHK2 THPT Nguyễn Trung Thiên - Hà Tĩnh - 21]%[EX7-21-Đỗ Đường Hiếu]
	Cho hình chóp đều $S.ABCD$ có đáy $ABCD$ là hình vuông tâm $O$ cạnh $a$, cạnh bên tạo với đáy một góc $60^\circ$. Gọi $M$ là điểm thuộc cạnh $SB$ sao cho $SM=\dfrac{2}{3}SB$. Tính khoảng cách từ $M$ đến mặt phẳng $(SCD)$.	
	\loigiai{
		\immini{Hình chóp $S.ABCD$ đều và $O$ là tâm hình vuông $ABCD$ nên suy ra $SO\perp (ABCD)$.\\
			Góc giữa cạnh bên với đáy của hình chóp bằng $60^\circ$ nên suy ra $\widehat{SAO}=60^\circ$.\\
			Ta có $SO= OA\cdot \tan SAO =\dfrac{a\sqrt{2}}{2}\cdot \tan 60^\circ =\dfrac{a\sqrt{6}}{2}$.\\
			Gọi $N$ là trung điểm $CD$, khi đó $CD\perp (SON)$.\\
			Dựng $OH\perp SN$ tại $H$. Từ $CD\perp (SON)\Rightarrow OH\perp CD$.\\
			Do đó $OH\perp (SCD)$, nên $OH=\mathrm{d}(O,(SCD))$.
		}
		{
			\begin{tikzpicture}[scale=0.8, font=\footnotesize,line join=round, line cap=round,>=stealth]
				\def\a{5}
				\path 
				(0,0) coordinate (A)
				(-2,-2) coordinate (B)
				(A)+(\a,0) coordinate (D)
				(B)+(\a,0) coordinate (C) 
				($(S)!2/3!(B)$) coordinate (M)
				($(A)!1/2!(C)$) coordinate (O)
				(O)+(0,4) coordinate (S)
				($(D)!1/2!(C)$) coordinate (N)
				($(S)!0.6!(N)$) coordinate (H)
				;
				\draw 
				(B)--(C)--(D)
				(S)--(B) (S)--(C) (S)--(D) (S)--(N);
				\draw[dashed]
				(B)--(A)--(D)
				(S)--(A) (A)--(C) (B)--(D)
				(S)--(O)--(N)
				(O)--(H);
				\foreach \p/\g in {A/-90,B/180,C/0,D/0,S/90,M/150,O/-90,N/-70,H/60}
				\fill[black] (\p) circle (1.2pt) 
				($(\p)+(\g:3mm)$) node{$\p$};
				\tkzMarkRightAngles(S,O,N O,N,C O,H,S)
		
\end{tikzpicture}}
		\noindent
		Xét tam giác $SON$ vuông tại $O$, ta có
		\[OH=\dfrac{SO\cdot ON}{\sqrt{SO^2+ON^2}}=\dfrac{\dfrac{a\sqrt{6}}{2}\cdot \dfrac{a}{2}}{\sqrt{\left(\dfrac{a\sqrt{6}}{2}\right)^2+\left(\dfrac{a}{2}\right)^2}}=\dfrac{a\sqrt{42}}{14}.\]
		Vì $\dfrac{MS}{BA}=\dfrac{2}{3}$ và $\dfrac{BD}{OD}=2$ nên
		\[\mathrm{d}\left(M,(SCD)\right)=\dfrac{2}{3}\mathrm{d}\left(B,(SCD)\right)=\dfrac{4}{3}\mathrm{d}\left(O,(SCD)\right)=\dfrac{4}{3}\cdot \dfrac{a\sqrt{42}}{14}=\dfrac{2a\sqrt{42}}{21}.\]
	}
\end{bt}

\begin{bt}%[1H8V5-3]
% nguon: Archive/SP_KhoLatex/ID6_old/19-OT-KhoangCach.tex #2850
%[Thi thử L1, THPT Hồng Lĩnh Hà Tĩnh, Hà Tĩnh, 2021]%[Thành Đức Trung, dự án 12-EX-5] 
		Cho hình lăng trụ đứng $ABC.A'B'C'$ có đáy là tam giác cân có $AC=BC=3a$. Đường thẳng $A'C$ tạo với đáy một góc $60^\circ$. Trên cạnh $A'C$ lấy điểm $M$ sao cho $A'M=2MC$. Biết rằng $A'B=a\sqrt{31}$. Tính khoảng cách từ $M$ đến mặt phẳng $(ABB'A')$.
	\loigiai
	{
		\immini
		{
			Ta có $MC\cap(ABB'A')=A' \Rightarrow \dfrac{\mathrm{d}(M,(ABB'A'))}{\mathrm{d}(C,(ABB'A'))}=\dfrac{A'M}{A'C}=\dfrac{2}{3}$ \\
			$\Rightarrow \mathrm{d}(M,(ABB'A'))=\dfrac{2}{3}\mathrm{d}(C,(ABB'A'))$. \\
			Gọi $H$ là trung điểm $AB$, suy ra $CH \perp AB$ mà $CH \perp AA'$ nên $CH \perp (ABB'A')$, suy ra $\mathrm{d}(C,(ABB'A'))=CH$. \\
			Ta có $\left(\widehat{A'C,(ABC)}\right)=\widehat{A'CA}=60^\circ$. \\
			Xét tam giác $A'AC$ có 
			\[\tan\widehat{A'CA}=\dfrac{AA'}{AC} \Rightarrow AA'=AC\tan\widehat{A'CA}=3a\sqrt{3}.\]
			Xét tam giác $A'AB$ có $AB^2=A'B^2-AA'^2=4a^2 \Rightarrow AH^2=a^2$. \\
			Xét tam giác $AHC$ có $CH^2=AC^2-AH^2=8a^2 \Rightarrow CH=2a\sqrt{2}$. \\
			Vậy $\mathrm{d}(M,(ABB'A'))=\dfrac{2}{3}\cdot2a\sqrt{2}=\dfrac{4a\sqrt{2}}{3}$. 
		}
		{
			\begin{tikzpicture}[scale=0.7, font=\footnotesize, line join=round, line cap=round, >=stealth]
				\tkzDefPoints{0/0/B, 2/-2/A, 5/0/C};
				\coordinate (A') at ($ (A) + (0,5) $);
				\coordinate (B') at ($ (A') + (B) - (A) $);
				\coordinate (C') at ($ (A') + (C) - (A) $);
				\coordinate (M) at ($ (C)!0.333333333!(A')$);
				\coordinate (H) at ($ (B)!0.5!(A)$);
				\tkzDrawSegments(A,B A,C A',B' B',C' C',A' A',C A,A' B,B' C,C' A',B);
				\tkzDrawSegments[dashed](B,C C,H)
				\tkzDrawPoints[size=2](A,B,C,A',B',C',M,H);
				\tkzLabelPoints[above](A',B',C',M);
				\tkzLabelPoints[below](A,B,C,H);
				\tkzMarkRightAngle(C,H,A);
				\tkzMarkRightAngle(C,A,A');
				\tkzMarkRightAngle(B,A,A');
			
\end{tikzpicture}
		}
	}
\end{bt}

\begin{bt}%[1H8V5-3]
% nguon: Archive/SP_KhoLatex/ID6_old/19-OT-KhoangCach.tex #2851
%[Thi thử L2, Sở Vĩnh Phúc, 2021]%[Thành Đức Trung, dự án 12-EX-9] 
		Cho hình chóp $S.ABCD$ có đáy $ABCD$ là hình thoi cạnh bằng 
		$a$, $\widehat{BAD} = 120^{\circ}$. Mặt bên $SAB$ là tam giác đều và $(SAB)\perp(ABCD)$. 
		Tính khoảng cách từ $A$ đến $(SBC)$.  
	\loigiai
	{
		\immini
		{
			Ta có $\heva{&(SAB)\perp(ABCD)\\&(SAB)\cap(ABCD)=AB\\& SH \subset (SAB)\\&SH\perp AB} \Rightarrow SH \perp (ABCD)$.\\
			Qua $H$ kẻ $HK \perp BC$ tại $K$, khi đó ta có 
			\[\heva{& SH \perp BC\\& HK \perp BC} \Rightarrow BC \perp (SHK).\]
			Qua $H$ dựng $HI \perp  SK, (I \in SK)$, khi đó
			\[\heva{ & HI\perp BC \\ & HI \perp SK} \Rightarrow HI\perp (SBC) \Rightarrow \mathrm{d}\left(H, (SBC)\right)=HI.\]
		}
		{
			\begin{tikzpicture}[scale=1,>=stealth, font=\footnotesize, line join=round, line cap=round]
				\tkzDefPoints{0/0/A,-1.4/-1.6/B,2.5/-1.6/C}
				\coordinate (D) at ($(A)+(C)-(B)$);
				\coordinate (H) at ($(A)!0.5!(B)$);
				\coordinate (S) at ($(H)+(0,3.5)$);
				\coordinate (K) at ($(B)!0.25!(C)$);
				\coordinate (I) at ($(S)!0.5!(K)$);
				\tkzDrawSegments(S,C S,D S,B B,C C,D);
				\tkzDrawSegments[dashed](A,S A,B A,D S,H H,K S,K H,I);
				\tkzDrawPoints[fill=black,size=2](D,C,A,B,S,H,K,I);
				\tkzLabelPoints[above](S);
				\tkzLabelPoints[below](A,B,C,K);
				\tkzLabelPoints[left](H);
				\tkzLabelPoints[right](D,I);
				\tkzMarkRightAngle(H,K,B);
			
\end{tikzpicture}
		}
		\noindent Ta có $AH\cap(SBC)=B \Rightarrow \dfrac{\mathrm{d}\left(A, (SBC)\right)}{\mathrm{d}\left(H, (SBC)\right)}=\dfrac{BA}{BH}=2$. \\
		Suy ra $\mathrm{d}\left(A, (SBC)\right) = 2\mathrm{d}\left(H, (SBC)\right)=2HI$. \\
		Ta có $SH = \dfrac{a\sqrt{3}}{2}$, $HK = \sin\widehat{HBK} \cdot BH= \dfrac{a\sqrt{3}}{4}$.\\
		Xét tam giác $SHK$ vuông tại $H$ ta có
		\[\dfrac{1}{HI^2}=\dfrac{1}{SH^2}+\dfrac{1}{HK^2} \Leftrightarrow HI=\dfrac{SH\cdot HK}{\sqrt{SH^2 +HK^2}}=\dfrac{a\sqrt{15}}{10}.\]
		Vậy $ \mathrm{d}\left(A, (SBC)\right) = 2HI=\dfrac{a\sqrt{15}}{5}$.
	}
\end{bt}

\begin{bt}%[1H8H5-4]
% nguon: Archive/SP_KhoLatex/ID6_old/19-OT-KhoangCach.tex #2852
%[Thi thử, THPT Kim Sơn - Ninh Bình, 2021]%[Nguyễn Xuân Bảo, EX6]
	Cho hình lập phương $ABCD.A'B'C'D'$ có cạnh bằng $2a$. Tính khoảng cách giữa hai đường thẳng $AB'$ và $CD'$.
	\loigiai{
		\immini{ Vì $AB' \parallel C'D$ nên $AB' \parallel (CDD'C')$.\\
			Suy ra \[ \mathrm{d}(AB',CD')=\mathrm{d}(AB',(CDD'C'))=\mathrm{d}(A,(CDD'C'))=AD=2a,\] vì $AD \perp (CDD'C')$.
		}
		{
			\begin{tikzpicture}[scale=0.8, font=\footnotesize, line join=round, line cap=round, >=stealth]
				\def\bc{4} % cạnh BC
				\def\ba{2} % cạnh BA
				\def\h{4} % đường cao
				\def\gocB{35} % góc B của đáy
				\coordinate[label=below left:$B$] (B) at (0,0);
				\coordinate[label=above right:$A$] (A) at (\gocB:\ba);
				\coordinate[label=below:$C$] (C) at (\bc,0);
				\coordinate[label=right:$D$] (D) at ($(C)-(B)+(A)$);
				\coordinate[label=above left:$A'$] (A') at ($(A)+(90:\h)$);
				\coordinate[label=left:$B'$] (B') at ($(B)-(A)+(A')$);
				\coordinate[label=below right:$C'$] (C') at ($(C)-(A)+(A')$);
				\coordinate[label=right:$D'$] (D') at ($(D)-(A)+(A')$);
				\draw (B')--(B)--(C)--(D)--(D')--(A')--(B')--(C')--(D') (C)--(C') (C)--(D');
				\draw[dashed] (A')--(A)--(D) (A)--(B) (A)--(B');
				\foreach \diem in {A,B,C,D,A',B',C',D'}	\fill (\diem)circle(1.5pt);
			
\end{tikzpicture}
		}
	}
\end{bt}

\begin{bt}%[1H8H5-4]
% nguon: Archive/SP_KhoLatex/ID6_old/19-OT-KhoangCach.tex #2853
%[Khảo sát chất lượng lần 1, Sở Hải Phòng, 2021]%[Thành Đức Trung, dự án 12-EX-10] 
		Cho hình lập phương $ABCD.A'B'C'D'$ có cạnh bằng $1$. Tính khoảng cách giữa hai đường thẳng $AA'$ và $BD'$.
	\loigiai
	{
		\immini
		{
			Vì $AA' \parallel BB'$ nên $AA' \parallel (BB'D')$.\\
			Suy ra $\mathrm{d}(AA',BD')= \mathrm{d}(AA',(BB'D')) = \mathrm{d}(A',(BB'D'))$.\\
			Gọi $O$ là tâm của hình vuông $A'B'C'D'$.\\
			Vì $A'O \perp (BB'D)$ nên $\mathrm{d}(A',(BB'D'))=A'O = \dfrac{A'C'}{2}=\dfrac{\sqrt{2}}{2}$.\\
			Vậy khoảng cách giữa hai đường thẳng $AA'$ và $BD'$ bằng $\dfrac{\sqrt{2}}{2}$.
		}
		{
			\begin{tikzpicture}[scale=0.8,>=stealth, font=\footnotesize, line join=round, line cap=round]
				\tkzDefPoints{0/0/A',-1.1/-1.1/B',2.5/-1.1/C'};
				\coordinate (D') at ($(A')+(C')-(B')$);
				\coordinate (A) at ($(A')+(0,2.5)$);
				\coordinate (B) at ($(A)+(B')-(A')$);
				\coordinate (C) at ($(A)+(C')-(A')$);
				\coordinate (D) at ($(A)+(D')-(A')$);
				\draw (A)--(B) (A)--(D) (B)--(C)--(D) (C')--(C) (B)--(B') (D)--(D') (B')--(C')--(D'); 
				\draw[dashed] (A')--(B') (A')--(D') (A')--(A) (B)--(D') (B')--(C') (A')--(C') (B')--(D');   
				\tkzInterLL(A',C')(B',D')\tkzGetPoint{O};
				\foreach \p in {A,B,D,C,A',B',C',D',O}
				\fill (\p) circle (1pt);
				\tkzLabelPoints[above](A,D);
				\tkzLabelPoints[below](B',C',O);
				\tkzLabelPoints[left](B,A');
				\tkzLabelPoints[right](C,D');
			
\end{tikzpicture}
		}		
	}
\end{bt}

\begin{bt}%[1H8V5-4]
% nguon: Archive/SP_KhoLatex/ID6_old/19-OT-KhoangCach.tex #2854
%[Đề minh họa lần 2 - 2020]%[Mai Sương] 
	Cho hình chóp ${S.ABC}$ có đáy là tam giác vuông tại $A$, $AB=2a$, $AC=4a$, ${SA}$ vuông góc với mặt phẳng đáy và $SA=a$. Gọi $M$ là trung điểm của ${AC}$. Tính khoảng cách giữa hai đường thẳng $SM$ và $BC$.
	\loigiai{ 
		\immini
		{
			Gọi $N$ là trung điểm của ${AB}$, ta có $MN\parallel BC$ nên $BC\parallel (SMN)$.\\
			Do đó $$\mathrm{d}(BC,SM)=\mathrm{d}(BC,(SMN))=\mathrm{d}(B,(SMN))=\mathrm{d}(A,(SMN))=h.$$
			Tứ diện ${A.SMN}$ vuông tại $A$ nên ta có
			$$\dfrac{1}{h^2}=\dfrac{1}{AS^2}+\dfrac{1}{AM^2}+\dfrac{1}{AN^2}=\dfrac{1}{a^2}+\dfrac{1}{4a^2}+\dfrac{1}{a^2}=\dfrac{9}{4a^2}\Rightarrow h=\dfrac{2a}{3}.$$
			Vậy $\mathrm{d}(BC,SM)=\dfrac{2a}{3}$.
		}
		{
			\begin{tikzpicture}[scale=1,line join=round,line cap=round]
				\tkzDefPoints{0/0/A,-1.2/-1.7/B,4/0/C}
				\coordinate (S) at ($(A)+(0,3)$);
				\coordinate (M) at ($(A)!0.5!(C)$);
				\coordinate (N) at ($(A)!0.5!(B)$);
				\tkzDrawSegments(S,B S,C B,C)
				\tkzDrawSegments[dashed](A,C A,S A,B S,M S,N N,M)
				\foreach \p/\r in {A,B,C,S,M,N}	\fill (\p) circle (1pt); 
				\tkzLabelPoints[above](S)
				\tkzLabelPoints[below](B,M,N)
				\tkzLabelPoints[above right](A)
				\tkzLabelPoints[right](C)
			
\end{tikzpicture}
		}
	}
\end{bt}

\begin{bt}%[1H8V5-4]
% nguon: Archive/SP_KhoLatex/ID6_old/19-OT-KhoangCach.tex #2856
%[Thi thử, THPT Kinh Môn, Hải Dương, 2021]%[Toanvo, 12-EX-6]
	\immini
	{
		Cho hình lăng trụ đứng $ABC.A'B'C'$ có $ABC$ là tam giác vuông cân, $AB=AC=a$ và $AA'=a\sqrt{3}$. Tính khoảng các giữa hai đường thẳng chéo nhau $AB'$ và $BC'$. 
	}
	{
		\begin{tikzpicture}[scale=1, font=\footnotesize, line join=round, line cap=round, >=stealth]
			\path
			(0,0) coordinate (A')
			(5,0) coordinate (C')
			(3,-1.75) coordinate (B')
			(0,3) coordinate (A)
			($(B')+(A)-(A')$) coordinate (B)
			($(C')+(A)-(A')$) coordinate (C)
			;
			\pgfresetboundingbox
			\draw(A)--(B)--(C)--(A)--(A')--(B')--(C')--(C)(A)--(B')--(B)--(C');
			\draw[dashed](A')--(C');
			\foreach \p/\g in {C/0,C'/0,A/180,A'/180,B'/-90,B/90}{\node at ($(\p)+(\g:3mm)$) {$\p$}; \draw[fill=black] (\p) circle (1pt);}
		
\end{tikzpicture}
	}
	\loigiai
	{
		\immini
		{
			Từ giả thiết, ta có chóp $C'.A'ABB'$ có đáy là hình chữ nhật và $C'A'$ vuông góc với đáy.\\
			Gọi $O$ là tâm của đáy $AA'B'B$ và $I$ là trung điểm của $A'C'$.\\
			Trong $\triangle A'BC'$, ta có $O$ và $I$ lần lượt là trung điểm của $A'B$ và $A'C'$.\\
			Suy ra $OI$ là đường trung bình của $\triangle A'BC'$, nên $C'B\parallel IO$.\\
			Mà $IO\subset(AIB')$ nên $C'B\parallel (AIB')$.\\
			Suy ra $\mathrm{d}(AB',BC')=\mathrm{d}(BC',(IAB'))=\mathrm{d}(C',(IAB'))$.\\
			Do $I$ là trung điểm $A'C'$ nên $\mathrm{d}(C',(IAB'))=\mathrm{d}(A',(IAB'))=d$.
		}
		{
			\begin{tikzpicture}[scale=0.7, font=\footnotesize, line join=round, line cap=round, >=stealth]
				\path
				(0,0) coordinate (A')
				(5,0) coordinate (B')
				(-2,-2) coordinate (A)
				($(A)+(B')-(A')$) coordinate (B)
				(0,4) coordinate (C')
				($(A')!.5!(C')$) coordinate (I)
				($(A')!.5!(B)$) coordinate (O)
				;
				\pgfresetboundingbox
				\draw(A)--(B)--(B')--(C')--(A)(B)--(C');
				\draw[dashed](A)--(A')--(C')(A')--(B')(A)--(B')(B)--(A')(A)--(I)--(B')(I)--(O);
				\foreach \p/\g in {B/-45,B'/0,I/150,O/-100,A'/150,C'/90,A/150}{\node at ($(\p)+(\g:4mm)$) {$\p$}; \draw[fill=black] (\p) circle (1pt);}
			
\end{tikzpicture}
		}
		\noindent 
		Ta thấy $A'A$, $A'I$ và $A'B'$ đôi một vuông góc với nhau nên
		\[\dfrac{1}{d^2}=\dfrac{1}{A'A^2}+\dfrac{1}{A'I^2}+\dfrac{1}{A'B'^2}=\dfrac{1}{3a^2}+\dfrac{4}{a^2}+\dfrac{1}{a^2}=\dfrac{16}{3a^2}\Rightarrow d=\dfrac{a\sqrt{3}}{4}.\]
		Vậy khoảng cách cần tìm là $\mathrm{d}(AB',BC')=d=\dfrac{a\sqrt{3}}{4}$.
	}
\end{bt}

\begin{bt}%[1H8C6-2]
% nguon: Archive/SP_KhoLatex/ID6_old/1H3G2-5.tex #2904 
% KIEM_TRA: cần xem lại đáp số/dạng
%[Hồ Hà Đặng - dự án Latex hóa 11]
	Cho ba tia $Ox,Oy,Oz$ không đồng phẳng.  
	\begin{enumerate}
		\item  Đặt $\widehat{xOy}=\alpha,\widehat{yOz}=\beta,\widehat{zOx}=\gamma $. Chứng minh rằng $\cos \alpha +\cos \beta +\cos \gamma >-\dfrac{3}{2}$.  
		\item Gọi $Ox_1,Oy_1,Oz_1$ lần lượt là các tia phân giác của các góc $xOy,yOz,zOt$. Chứng minh rằng nếu $Ox_1$ và $Oy_1$ vuông góc với nhau thì $Oz_1$ vuông góc với cả $Ox_1$ và $Oy_1$.  
	\end{enumerate}
	\loigiai{
		Lấy $E_1,E_2,E_3$ lần lượt tuộc các tia $Ox,Oy,Oz$ sao cho $OE_1=OE_2=OE_3$.  
		Đặt $$\overrightarrow{OE_1}=\overrightarrow{e_1},\overrightarrow{OE_2}=\overrightarrow{e_2},\overrightarrow{OE_3}=\overrightarrow{e_3}.$$
		\begin{enumerate}
			\item  Do ba tia $Ox,Oy,Oz$ không đồng phẳng nên 
			\begin{eqnarray*}
				{\left(\overrightarrow{e_1}+\overrightarrow{e_2}+\overrightarrow{e_3}\right)}^2>0    &\Leftrightarrow& {\overrightarrow{e_1}}^2+{\overrightarrow{e_2}}^2+{\overrightarrow{e_3}}^2+2\left(\overrightarrow{e_1}\cdot\overrightarrow{e_2}+\overrightarrow{e_2}\cdot\overrightarrow{e_3}+\overrightarrow{e_3}\cdot\overrightarrow{e_1}\right)>0 \\
				\\	&\Leftrightarrow& 3OE_1^2+2OE_1^2\left(\cos \alpha +\cos \beta +\cos \gamma \right)>0.
			\end{eqnarray*}
			Vậy $\cos \alpha +\cos \beta +\cos \gamma >-\dfrac{3}{2}$. \\
			\item  Ta có $\overrightarrow{OE_1}+\overrightarrow{OE_2}\parallel Ox_1$; $\overrightarrow{OE_2}+\overrightarrow{OE_3}\parallel Oy_1$; $\overrightarrow{OE_3}+\overrightarrow{OE_1}\parallel Oz_1$ và \\
				$Ox_1\perp Oy_1 $. Do đó $$\left(\overrightarrow{OE_1}+\overrightarrow{OE_2}\right)\left(\overrightarrow{OE_2}+\overrightarrow{OE_3}\right)=0 
				\Leftrightarrow {\overrightarrow{OE_2}}^2+\overrightarrow{OE_1}\cdot\overrightarrow{OE_2}+\overrightarrow{OE_1}\cdot\overrightarrow{OE_3}+\overrightarrow{OE_2}\cdot\overrightarrow{OE_3}=0.$$ 
				\begin{center}
					\begin{tikzpicture}[scale=.7, line join=round, line cap=round,>=stealth]
					\tkzInit[xmin=-4,xmax=5.5,ymin=-4,ymax=5.5]
					\tkzClip
					\tkzDefPoints{0/0/O,4/0/Y,-2/-3/X,0/5/Z,2/0/E2,-1/-1.5/E1,0/2/E3}
					
					\draw[->] (O)--(Y) node[above]{$y$};
					\draw[->] (O)--(X) node[below]{$x$};
					\draw[->] (O)--(Z) node[above]{$z$};
					\tkzDefLine[bisector](X,O,Y) \tkzGetPoint{x1}
					\tkzDefLine[bisector](Y,O,Z) \tkzGetPoint{y1}
					\tkzDefLine[bisector](Z,O,X) \tkzGetPoint{z1}
					\coordinate (M) at ($(O)!.4!(x1)$);
					\coordinate (N) at ($(O)!.5!(y1)$);
					\coordinate (P) at ($(O)!.3!(z1)$);
					\draw[->] (O)--(M) node[below]{$x_1$};
					\draw[->] (O)--(N) node[above]{$y_1$};
					\draw[->] (O)--(P) node[above]{$z_1$};
					\tkzDrawPoints[fill=black](E1,E2,E3)
					\tkzLabelPoint[left](E1){$E_1$}
					\tkzLabelPoint[below](E2){$E_2$}
					\tkzLabelPoint[right](E3){$E_3$}
					\end{tikzpicture}
				\end{center}
				Ta có  
				\begin{eqnarray*}
					\left(\overrightarrow{OE_1}+\overrightarrow{OE_2}\right)\left(\overrightarrow{OE_3}+\overrightarrow{OE_1}\right)&=&{\overrightarrow{OE_1}}^2+\overrightarrow{OE_1}\cdot\overrightarrow{OE_2}+\overrightarrow{OE_2}\cdot\overrightarrow{OE_3}+\overrightarrow{OE_1}\cdot\overrightarrow{OE_3}
					\\&=&{\overrightarrow{OE_2}}^2+\overrightarrow{OE_1}\cdot\overrightarrow{OE_2}+\overrightarrow{OE_1}\cdot\overrightarrow{OE_3}+\overrightarrow{OE_2}\cdot\overrightarrow{OE_3}=0
				\end{eqnarray*} (do $OE_1=OE_2$ ). 
				Do đó $Ox_1\perp Oz_1$. \\Tương tự, ta cũng có $Oy_1\perp Oz_1$. 
			
		\end{enumerate}
		}\end{bt}

\begin{bt}%[1H8H1-3]
% nguon: Archive/SP_KhoLatex/ID6_old/DeThamKhao-02.tex #3214
%[Dự án ĐGNL-ĐGTD-2022 GĐ1]%[Lê Minh Thiện Anh]
	\immini[thm]{Cho hình lập phương $ABCD.A'B'C'D'$. Góc giữa hai đường thẳng $AC$ và $A'D$ bằng
	}
	{
		\begin{tikzpicture}[scale=.7, font=\footnotesize, line join=round, line cap=round, >=stealth]
		\coordinate[label=below left:$B$] (B) at (0,0);
		\coordinate[label=above left:$A$] (A) at (35:2);
		\coordinate[label=below:$C$] (C) at (4,0);
		\coordinate[label=right:$D$] (D) at ($(C)-(B)+(A)$);
		\coordinate[label=above left:$A'$] (A') at ($(A)+(90:4)$);
		\coordinate[label=left:$B'$] (B') at ($(B)-(A)+(A')$);
		\coordinate[label=below right:$C'$] (C') at ($(C)-(A)+(A')$);
		\coordinate[label=right:$D'$] (D') at ($(D)-(A)+(A')$);
		\draw (B')--(B)--(C)--(D)--(D')--(A')--(B')--(C')--(D') (C)--(C');
		\draw[dashed] (A')--(A)--(D)--cycle (C)--(A)--(B);
		\foreach \diem in {A,B,C,D,A',B',C',D'}	\fill (\diem)circle(1pt);
		\end{tikzpicture}
	}
	\textit{Kết quả:} \tikz\draw (0,0) rectangle (3,0.6);
	\loigiai{
		\immini{Do $ABCD.A'B'C'D'$ là hình lập phương nên $A'D$ song song với $B'C$.\\
			$\Delta ACB'$ đều $\Rightarrow \widehat{ACB'}=60^\circ$.\\
			Suy ra $\left(AC,A'D\right)=\left(AC,CB'\right)=\widehat{ACB'}=60^\circ$.}
		{
			\begin{tikzpicture}[scale=.7, font=\footnotesize, line join=round, line cap=round, >=stealth]
			\coordinate[label=below left:$B$] (B) at (0,0);
			\coordinate[label=below:$A$] (A) at (35:2);
			\coordinate[label=below:$C$] (C) at (4,0);
			\coordinate[label=right:$D$] (D) at ($(C)-(B)+(A)$);
			\coordinate[label=above left:$A'$] (A') at ($(A)+(90:4)$);
			\coordinate[label=left:$B'$] (B') at ($(B)-(A)+(A')$);
			\coordinate[label=below right:$C'$] (C') at ($(C)-(A)+(A')$);
			\coordinate[label=right:$D'$] (D') at ($(D)-(A)+(A')$);
			\draw (B')--(B)--(C)--(D)--(D')--(A')--(B')--(C')--(D') (B')--(C)--(C');
			\draw[dashed] (A')--(A)--(D)--cycle (C)--(A)--(B) (A)--(B');
			\foreach \diem in {A,B,C,D,A',B',C',D'}	\fill (\diem)circle(1pt);
			\end{tikzpicture}
		}
	}
\end{bt}

\begin{bt}%[1H8H6-1]
% nguon: Archive/SP_KhoLatex/ID6_old/DeThamKhao-03.tex #3283
%[Dự án ĐGNL-ĐGTD-2022 GĐ1]%[Duong Xuan Loi]
	\immini{
		Cho hình chóp tứ giác đều $S.ABCD$ có tất cả các cạnh bằng $a$. Gọi $M$ là trung điểm của $SD$ (tham khảo hình vẽ bên). Tính tang của góc giữa đường thẳng $BM$ và mặt phẳng $(ABCD)$. 
		
		\textit{Kết quả:} \tikz\draw (0,0) rectangle (3,0.6);
	}{
		\begin{tikzpicture}[scale=0.7, font=\footnotesize, line join=round, line cap=round,>=stealth]
			\path
			(0,0) coordinate (A)
			(-1.8,-1.6)coordinate (B)
			(1.7,-1.6)coordinate (C)
			($(A)+(C)-(B)$)coordinate (D)
			($(A)!0.5!(C)$)coordinate (O)
			($(O)+(0,3.3)$) coordinate (S)
			($(S)!0.5!(D)$) coordinate (M)			
			;
			\draw (S)--(B)--(C)--(D)--cycle (S)--(C);
			\draw[dashed] (S)--(A)--(D) (A)--(B)--(M);	
			\foreach \p/\q in {S/90,A/170,B/-90,C/-90,D/0,M/50}
			\fill[black] (\p) circle (1.0pt)node[shift={(\q:3mm)}]{$\p$};	
		\end{tikzpicture}
	}
	\loigiai{
		\immini{
			Vì hình chóp $S.ABCD$ đều nên $ABCD$ là hình vuông và $\break SO\perp (ABCD)$ với $O$ là tâm của $ABCD$.\\
			Gọi $H$ là trung điểm $OD$\\
			$\,\Rightarrow\, MH\parallel SO \,\Rightarrow\, MH \perp (ABCD)$.\\
			Suy ra, $BH$ là hình chiếu vuông góc của $BM$ trên $(ABCD)$, suy ra góc giữa $BM$ và $(ABCD)$ bằng góc giữa $BM$ và $BH$ và bằng góc $\widehat{MBH}$.\\
			Ta có
		}{
			\begin{tikzpicture}[scale=0.7, font=\footnotesize, line join=round, line cap=round,>=stealth]
				\path
				(0,0) coordinate (A)
				(-1.8,-1.6)coordinate (B)
				(1.7,-1.6)coordinate (C)
				($(A)+(C)-(B)$)coordinate (D)
				($(A)!0.5!(C)$)coordinate (O)
				($(O)+(0,3.3)$) coordinate (S)
				($(S)!0.5!(D)$) coordinate (M)
				($(O)!0.5!(D)$) coordinate (H)		
				;
				\draw (S)--(B)--(C)--(D)--cycle (S)--(C);
				\draw[dashed] (O)--(S)--(A)--(D) (C)--(A)--(B)--(D) (B)--(M)--(H);	
				\foreach \p/\q in {S/90,A/170,B/-90,C/-90,D/0,M/50,O/-90,H/-90}
				\fill[black] (\p) circle (1.0pt)node[shift={(\q:3mm)}]{$\p$};	
			\end{tikzpicture}
		}
		\[BH=\dfrac{3}{4}BD=\dfrac{3}{4}a\sqrt{2}=\dfrac{3a\sqrt{2}}{4};\]
		và
		\[MH=\dfrac{1}{2}SO=\dfrac{1}{2}\sqrt{SD^2-OD^2}=\dfrac{1}{2}\sqrt{a^2-\dfrac{a^2}{2}}=\dfrac{a\sqrt{2}}{4}.\]
		Tam giác $BMH$ vuông tại $H$, suy ra
		\[\tan{\widehat{MBH}}=\dfrac{MH}{BH}=\dfrac{1}{3}.\]
	}
\end{bt}

\begin{bt}%[1H8H2-2]
% nguon: Archive/SP_KhoLatex/ID6_old/DeThamKhao-07.tex #3298
%[Dự án ĐGNL-ĐGTD-2022 GĐ1]%[Dao-V- Thuy]
	Cho hình chóp $S.ABC$ có tam giác $ABC$ vuông cân tại $B$, $AB=BC=a$, $SA=a\sqrt{3}$, $SA \perp(ABC)$. Góc giữa hai mặt phẳng $(SBC)$ và $(ABC)$ có số đo bằng bao nhiêu độ?\\
	\textit{Kết quả:} \tikz\draw (0,0) rectangle (3,0.6);
	\loigiai{
		\immini{
			Ta có $(SBC) \cap (ABC)=BC$.\\
			Vì $A$ là hình chiếu của $S$ lên mặt đáy nên $SA\perp AB$. Mà $AB\perp BC$, suy ra $SB\perp BC$.\\
			Do đó $((SBC),(ABC))=(SB,AB) = \widehat{SBA}=\alpha$.\\
			Ta lại có $\tan \alpha = \dfrac{SA}{AB} = \sqrt{3} \Rightarrow \alpha = 60^{\circ}\Rightarrow((SBC),(ABC)) = 60^\circ$.
		}{
\begin{tikzpicture}[scale=1.0, font=\footnotesize, line join=round, line cap=round, >=stealth]
				\def\a{2}
				\path
				(0,0) coordinate (A)
				++(0:{\a*sqrt(2)}) coordinate (C)
				++(-150:\a) coordinate (B)
				($(A)+(90:{\a*sqrt(3)})$) coordinate (S)
				;
				\tkzDrawSegments(S,A S,B S,C A,B B,C)
				\tkzDrawSegments[dashed](A,C)
				\foreach \x/\y/\z/\t in {S/A/B/5pt, A/B/C/5pt} {
					\path 
					($(\y)!\t!(\z)$) coordinate (1)
					($(\y)!\t!(\x)$) coordinate (2)
					($(1)+(2)-(\y)$) coordinate (3)
					;
					\draw (1)--(3)--(2);
				}
				\foreach \x/\y in {B/-45,C/0,A/160,S/90}
				\draw[fill=black] (\x) circle (0.03cm) + (\y:0.3cm) node {$\x$};
\end{tikzpicture}}
	}
\end{bt}

\begin{bt}%[1H8H6-1]
% nguon: Archive/SP_KhoLatex/ID6_old/DeThamKhao-06.tex #3413
%[Dự án ĐGNL-ĐGTD-2022 GĐ1]%[Bùi Mạnh Tiến]
	Cho hình chóp $S.ABC$ có đáy $ABC$ là tam giác vuông tại $C$ với $AB=2a$. Tam giác $SAB$ đều và nằm trong mặt phẳng vuông góc với mặt đáy. Tính góc giữa đường thẳng $SC$ và $(ABC)$.\\
	\textit{Kết quả:} \tikz\draw (0,0) rectangle (3,0.6);
	\loigiai{
		\immini{
			Gọi $H$ là trung điểm của $AB$.\\
			Do $\triangle SAB$ đều và $(SAB)\perp (ABC)$ nên $SH\perp (ABC)$.\\
			Suy ra góc giữa $SC$ và $(ABC)$ là $\widehat{SCH}$.\\
			Ta có $\triangle SAB$ đều nên $SH=\dfrac{AB\sqrt{3}}{2}=a\sqrt{3}$.\\
			Có $\triangle CAB$ vuông tại $C$ nên $CH=\dfrac{1}{2}AB=a$.\\
			Xét $\triangle HSC$ vuông tại $H$ có $\tan\widehat{SCH}=\dfrac{SH}{CH}=\sqrt{3}\Rightarrow \widehat{SCH}=60^\circ$.\\
			Vậy góc giữa $SC$ và $(ABC)$ bằng $60^\circ$.
		}
		{
			\begin{tikzpicture}[scale=1,font=\footnotesize,line join = round, line cap = round, >= stealth]
				\def\x{4} \def\y{1.7} \def\z{3}
				\coordinate (A) at (0,0);
				\coordinate (B) at (0:\x);
				\coordinate (C) at (-60:\y);
				\coordinate (H) at ($(A)!.5!(B)$);
				\coordinate (S) at ($(H)+(90:\z)$);
				\draw 
				(S)--(A)--(C)--(B)--cycle
				(S)--(C)
				;
				\draw[dashed]
				(A)--(B) (S)--(H)--(C)
				;
				\foreach \p/\g in {S/90,A/180,B/0,C/-90,H/-45} \draw[fill] (\p) circle(.5pt)node [shift={(\g:.3)}] {$\p$};
				%SAC
				\draw ($(H)!5pt!(S)$)--($(H)!5pt!(S)+(H)!5pt!(A)-(H)$)--($(H)!5pt!(A)$);
				\draw ($(C)!5pt!(A)$)--($(C)!5pt!(A)+(C)!5pt!(B)-(C)$)--($(C)!5pt!(B)$);
				\tkzMarkSegments[mark=x](H,A H,B)
				\begin{scope}
					\clip (S)--(C)--(H)--cycle;
					\draw[] (C) circle (5mm);
				\end{scope}
			\end{tikzpicture}
		}
	}
\end{bt}

\begin{bt}%[1H8C6-2]
% nguon: Archive/SP_KhoLatex/ID6_old/DeThamKhao-06.tex #3415
%[Dự án ĐGNL-ĐGTD-2022 GĐ1]%[Bùi Mạnh Tiến]
	Cho hình chóp $S.ABCD$ có đáy là hình vuông cạnh $2a$, tam giác $SAB$ đều, góc giữa $(SCD)$ và $(ABCD)$ bằng $60^{\circ}$. Gọi $M$ là trung điểm của cạnh $AB$. Biết rằng hình chiếu vuông góc của đỉnh $S$ trên mặt phẳng $(ABCD)$ nằm trong hình vuông $ABCD$. Khoảng cách giữa hai đường thẳng $SM$ và $AC$ có dạng $\dfrac{a\sqrt{m}}{n}$ với $m$, $n\in \mathbb{N}$ và $m$ là số nguyên tố. Tổng $m+n^2$ có giá trị là\\
	\textit{Kết quả:} \tikz\draw (0,0) rectangle (3,0.6);
	\loigiai{
		\immini{Gọi $I$ là trung điểm cạnh $CD$, khi đó $\left\{\begin{aligned}
				&AB\perp SM \\
				&AB\perp MI
			\end{aligned}\right. \Rightarrow AB\perp (SMI)$. \\
			Do $CD \parallel AB$ nên $CD\perp (SMI) \Rightarrow ((SCD),(ABCD))=\widehat{SIM}$. \\
			Vẽ $SH\perp MI$ tại $H\in MI$ thì $SH\perp (ABCD)$. \\
			$\begin{aligned}\triangle SMI &\text{ có } SM^2=MI^2+SI^2-2MI\cdot SI\cos \widehat{SIM}\\
				&\Leftrightarrow 3a^2=4a^2+SI^2-2aSI\\
				& \Leftrightarrow SI^2-2aSI+a^2=0 \Leftrightarrow SI=a.
			\end{aligned}$}
		{\begin{tikzpicture}[line cap=round,line join=round,>=stealth,scale=1.1]
				\coordinate[label=left:{$A$}] (A) at (0,0);
				\coordinate[label=left:{$B$}] (B) at (-0.2,-1.5);
				\coordinate[label=right:{$C$}] (C) at ($(B)+(3.5,0)$);
				\coordinate[label=right:{$D$}] (D) at ($(A)+(C)-(B)$);
				\coordinate[label=left:{$M$}] (M) at ($(A)!1/2!(B)$);
				\coordinate[label=right:{$I$}] (I) at ($(C)!1/2!(D)$);
				\coordinate[label=above left:{$H$}] (H) at ($(I)!1/4!(M)$);
				\coordinate[label={$S$}] (S) at ($(H)+(0,3)$);
				\coordinate[label=below:{$N$}] (N) at ($(B)!1/2!(C)$);
				\draw (M)--(S)--(A)--(B)--(S)--(N)--(C)--(D)--(S)--(I) (B)--(N) (S)--(C);
				\draw[dashed] (N)--(M)--(I) (C)--(A)--(D) (S)--(H);
				\tkzMarkRightAngle[size=0.14](B,M,I)
				\tkzMarkRightAngle[size=0.14](S,H,I)
				\tkzDrawPoints[fill=black,size=1.2pt](S,A,B,C,D,M,N,I,H)
		\end{tikzpicture}}
		\noindent
		Theo định lý Pythagore đảo thì $\triangle SMI$ vuông tại $S \Rightarrow SH=\dfrac{SM\cdot SI}{MI}=\dfrac{a\sqrt{3}}{2}$. \\
		Gọi $N$ là trung điểm cạnh $BC$ ta có $AC \parallel MN$ \\
		$ \Rightarrow \mathrm{\,d}(AC,SM)=\mathrm{d}(AC,(SMN))=\mathrm{d}(C,(SMN))=\dfrac{3V_{SMNC}}{S_{\Delta SMN}}$. \\
		Ta có $V_{SMNC}=V_{S.MNB}=\dfrac{1}{3}SH\cdot\dfrac{1}{2}BM\cdot BN=\dfrac{1}{6}\cdot\dfrac{a\sqrt{3}}{2}\cdot a\cdot a=\dfrac{a^3\sqrt{3}}{12}$. \\
		Tam giác $SIC$ có $SC=\sqrt{SI^2+IC^2}=\sqrt{a^2+a^2}=a\sqrt{2}$. \\
		Tam giác $SBC$ có $SN^2=\dfrac{SB^2+SC^2}{2}-\dfrac{BC^2}{4}=2a^2 \Rightarrow SN=a\sqrt{2}$. \\
		Tam giác $SMN$ có nửa chu vi $p=\dfrac{SM+SN+MN}{2}=\dfrac{a\sqrt{3}+a\sqrt{2}+a\sqrt{2}}{2}$. \\
		Và diện tích $\triangle SMN$ là $S_{\triangle SMN}=\sqrt{p(p-SM)(p-SN)(p-BC)}=\dfrac{a^2\sqrt{15}}{4}$. \\
		Vậy $\mathrm{\,d}(AC,SM)=\dfrac{3V_{SMNC}}{S_{\Delta SMN}}=\dfrac{3\cdot \dfrac{a^3\sqrt{3}}{12}}{\dfrac{a^2\sqrt{15}}{4}}=\dfrac{a\sqrt{5}}{5}$.\\
		Suy ra $m=5$, $n=5\Rightarrow m+n^2=5+5^2=30$.
		}
\end{bt}

\begin{bt}%[1H8N1-3]
% nguon: Archive/SP_KhoLatex/ID6_old/18-OT-Goc.tex #3675
%[ĐGNL-ĐGTD-2022]%[Đình Nguyên]
	Cho hình chóp $S.ABCD$ có tất cả các cạnh bằng $a$. Gọi $I,J$ lần lượt là trung điểm của $SC$ và $BC$. Tính góc giữa $IJ$ và $CD$.\dapso{$60^\circ$}
	\loigiai{
		\immini{Vì $IJ\parallel SB$ (do $IJ$ là đường trung bình của tam giác $SBC$) và $CD\parallel AB$ nên 
			$$ (IJ,CD)=(SB,AB)=\widehat{SBA}=60^\circ,$$
			vì $\triangle SBA$ đều cạnh $a$.
		}
		{\begin{tikzpicture}[scale=.8, font=\footnotesize, line join=round, line cap=round, >=stealth]
				\tkzDefPoints{2.5/2.5/A, 0/0/D, 4/0/C}
				\tkzDefPointBy[translation = from D to C](A)
				\tkzGetPoint{B}
				\tkzDefMidPoint(C,A) \tkzGetPoint{O}
				\coordinate (S) at ($(O)+(0,4)$);
				\tkzDefMidPoint(S,C) \tkzGetPoint{I}
				\tkzDefMidPoint(B,C) \tkzGetPoint{J}
				\tkzDrawSegments(S,D S,B S,C B,C D,C I,J)
				\tkzDrawSegments[dashed](S,A D,A A,B)
				\tkzLabelPoints[right](B,C)
				\tkzLabelPoints[below](J)
				\tkzLabelPoints[above](S)
				\tkzLabelPoints[left](A,D,I)
				\tkzDrawPoints[fill=black](S,A,B,C,D,I,J)
			\end{tikzpicture}
		}
	}
\end{bt}

\begin{bt}%[1H8H1-3]
% nguon: Archive/SP_KhoLatex/ID6_old/18-OT-Goc.tex #3676
%[ĐGNL-ĐGTD-2022]%[Đình Nguyên]
	Cho hình chóp $S.ABCD$ có đáy là hình thoi cạnh $a$, $SA=a\sqrt 3$, $SA\perp BC$. Tính góc giữa hai đường thẳng $SD$ và $BC$.\dapso{$60^\circ$}
	\loigiai{
		\immini{Ta có 
			$\heva{&SA\perp BC\\&BC\parallel AD}\Rightarrow SA\perp AD$. Do đó
			$(SD,BC)=(SD,AD)=\widehat{SDA}$.\\
			Xét tam giác $SAD$ vuông tại $A$, ta có $\tan \widehat{SDA}=\dfrac{SA}{DA}=\dfrac{a\sqrt3}{a}=\sqrt{3}$. Suy ra $\widehat{SDA}=60^\circ$.\\
			Vậy góc giữa hai đường thẳng $SD$ và $BC$ bằng $60^\circ$.  
		}
		{\begin{tikzpicture}[scale=.8, font=\footnotesize, line join=round, line cap=round, >=stealth]
				\tkzDefPoints{1.5/1.5/A, 0/0/B, 3/0/C}
				\tkzDefPointBy[translation = from B to C](A)
				\tkzGetPoint{D}
				\tkzDefMidPoint(C,A) \tkzGetPoint{O}
				\coordinate (S) at ($(A)+(0,3)$);
				\tkzDrawSegments(S,D S,B S,C B,C D,C)
				\tkzDrawSegments[dashed](S,A D,A A,B)
				\tkzLabelPoints[right](D,C)
				%	\tkzLabelPoints[below](B)
				\tkzLabelPoints[above](S)
				\tkzLabelPoints[left](A,B)
				\tkzDrawPoints[fill=black](S,A,B,C,D)
			\end{tikzpicture}
		}		
	}
\end{bt}

\begin{bt}%[1H8V1-3]
% nguon: Archive/SP_KhoLatex/ID6_old/18-OT-Goc.tex #3677
%[ĐGNL-ĐGTD-2022]%[Đình Nguyên]
	Cho hình chóp $S.ABCD$ có đáy là hình vuông $ABCD$ cạnh $a$. Độ dài các cạnh bên của hình chóp bằng nhau và bằng hai lần độ dài cạnh hình vuông. Gọi $M$ và $N$ lần lượt là trung điểm của $AD$ và $SD$. Tính số đo của góc $\left(MN, SC\right)$.
	\loigiai{
		\immini{
			Do $M$ và $N$ lần lượt là trung điểm của $AD$ và $SD$ nên $MN\parallel SA$. Bởi vậy
			\[\left(MN, SC\right)=\left(SA, SC\right).\]
			Xét tam giác $SAC$ có $SA=SC=2a$, $AC=a\sqrt{2}$. Bởi vậy
			\[\cos \widehat{ASC}=\dfrac{SA^2+SC^2-AC^2}{2SA\cdot SC}=\dfrac{(2a)^2+(2a)^2-(a\sqrt{2})^2}{2\cdot 2a\cdot 2a}=\dfrac{3}{4}.\]
			Do đó 
			\[\left(MN, SC\right)=\left(SA, SC\right)=\widehat{ASC}=\arccos\left(\dfrac{3}{4}\right) .\]
		}
		{\begin{tikzpicture}[scale=1, font=\footnotesize, line join=round, line cap=round, >=stealth]
				\def\a{4}
				\def\h{4}	
				\path 
				(0,0)coordinate(B)
				(1.5,1)coordinate(A)
				(A)+(\a,0)coordinate(D)
				(B)+(\a,0)coordinate(C)
				($(A)!0.5!(C)$) coordinate(O)
				(O)+(0,\h)coordinate(S)
				($(A)!0.5!(D)$) coordinate(M)
				($(S)!0.5!(D)$) coordinate(N)
				;
				\draw[dashed]
				(B)--(A)--(D)--(B) (O)--(S)--(A)--(C)
				(M)--(N)
				;
				\draw
				(B)--(C)--(D) (B)--(S)--(C) (S)--(D)
				;
				\foreach \p/\g in{S/90,O/-90,A/160,B/-160,C/-20,D/20,M/140,N/50}
				\fill[black] (\p) circle (1.2pt) 
				($(\p)+(\g:3mm)$) node{$\p$};
			\end{tikzpicture}
		}	
	}
\end{bt}

\begin{bt}%[1H8V1-3]
% nguon: Archive/SP_KhoLatex/ID6_old/18-OT-Goc.tex #3678
%[ĐGNL-ĐGTD-2022]%[Đình Nguyên]
	Cho hình chóp $S.ABCD$ có đáy $ABCD$ là hình thang vuông tại $A$ và $D$, $AB$ là đáy lớn, biết $SA\perp AB$, $AD=DC=a$, $BC=a\sqrt2$, $SA=\dfrac{2a\sqrt3}{3}$. Tính góc giữa $SB$ và $DC$.\dapso{$30^\circ$}
	\loigiai{
		\immini{Kẻ $CI\perp AB$ với $I\in AB$. Khi đó $AICD$ là hình vuông cạnh $a$. Suy ra $\triangle BIC$ vuông tại $I$. Theo định lý Py-ta-go, ta có $$IB=\sqrt{BC^2-IC^2}=\sqrt{2a^2-a^2}=a.$$
			Suy ra $AB=AI+IB=a+a=2a$.\\
			Vì $AB\parallel DC$ nên $(SB,DC)=(SB,AB)=\widehat{SBA}$. Ta có
			$$\tan \widehat{SBA}=\dfrac{SA}{AB}=\dfrac{\dfrac{2a\sqrt{3}}{3}}{2a}=\dfrac{\sqrt{3}}{3}.$$
			Suy ra $\widehat{SBA}=30^\circ$. Vậy góc giữa $SB$ và $DC$ bằng $30^\circ$.
		}
		{
			\begin{tikzpicture}[scale=.8, font=\footnotesize, line join=round, line cap=round, >=stealth]
				\tkzDefPoints{1.5/2/A, 0/0/D, 6/0/K}
				\tkzDefPointBy[translation = from D to K](A)
				\tkzGetPoint{B}
				\coordinate (C) at ($(D)!0.5!(K)$);
				\coordinate (I) at ($(A)!0.5!(B)$);
				\coordinate (S) at ($(A)+(0,4)$);
				\tkzDrawSegments(S,B S,C S,D B,C C,D)
				\tkzDrawSegments[dashed](S,A A,B A,D A,C C,I)
				\tkzLabelPoints[right](C,B)
				\tkzLabelPoints[above](S, I)
				\tkzLabelPoints[left](A,D)
				\tkzMarkRightAngle(B,I,C)
				\tkzMarkRightAngle(S,A,B)
				\tkzDrawPoints[fill=black](S,A,B,C,D,I)
			\end{tikzpicture}
		}		
	}
\end{bt}

\begin{bt}%[1H8H6-1]
% nguon: Archive/SP_KhoLatex/ID6_old/18-OT-Goc.tex #3679
%[ĐGNL-ĐGTD-2022]%[Đình Nguyên]
	\immini{
		Cho hình chóp $S.ABCD$ có đáy $ABCD$ là hình vuông cạnh $a$, $SC$ vuông góc với đáy và $SC=a\sqrt{3}$. Tính $\tan$ góc giữa đường thẳng $SA$ và mặt phẳng $(SBC)$.
	}{
		\begin{tikzpicture}[line join=round, line cap=round, scale=0.7]
			%\draw[dashed,cyan] (-5,-5) grid (5,5);
			\def\d{5} \def\r{1.8} \def\h{4} \def\l{2.2}
			
			\coordinate[label={below left}:{\footnotesize $B$}] (B) at (-3,-3);
			\coordinate[label={below right}:{\footnotesize $A$}] (A) at ($(B)+(\d,0)$);
			\coordinate[label={below right}:{\footnotesize $C$}] (C) at ($(B)+(\l,\r)$);
			\coordinate[label={above right}:{\footnotesize $D$}] (D) at ($(C)+(\d,0)$);
			\coordinate[label={above}:{\footnotesize $S$}] (S) at ($(C)+(0,\h)$);
			\draw (S)--(B)--(A)--(D)--(S)--(A);
			\draw[dashed] (S)--(C)--(B) (C)--(D);
			\tkzMarkRightAngles(S,C,D S,C,B B,C,D)
			\tkzLabelSegment[right, yshift=-2mm](S,C){\footnotesize $a\sqrt{3}$}
			\tkzLabelSegment[below](A,B){\footnotesize $a$}
			\tkzDrawPoints[fill=black](S,A,B,C,D)
		\end{tikzpicture}
	}\dapso{$\dfrac12$}
	\loigiai{
		\immini{
			Ta có $\left. \begin{array}{l}
				AB\perp SC\\ 
				AB\perp BC
			\end{array}  \right\} \Rightarrow AB\perp (SBC)$.\\
			Suy ra hình chiếu của $SA$ lên $(SBC)$ là $SB$.\\
			$\Rightarrow \left( SA,(SBC) \right)=(SA,SB)=\widehat{ASB}$.\\
			Trong $\triangle SCB$ vuông tại $C$, ta có 
			$$SB=\sqrt{SC^2+CB^2}=\sqrt{4a^2}=2a.$$
			Trong $\triangle SBA$ vuông tại $B$, ta có
			$$\tan \widehat{BSA}=\dfrac{AB}{SB}=\dfrac{a}{2a}=\dfrac{1}{2}.$$ 
		}{
			\begin{tikzpicture}[line join=round, line cap=round, scale=0.8]
				%\draw[dashed,cyan] (-5,-5) grid (5,5);
				\def\d{5} \def\r{1.8} \def\h{4} \def\l{2.2}
				
				\coordinate[label={below left}:{\footnotesize $B$}] (B) at (-3,-3);
				\coordinate[label={below right}:{\footnotesize $A$}] (A) at ($(B)+(\d,0)$);
				\coordinate[label={below right}:{\footnotesize $C$}] (C) at ($(B)+(\l,\r)$);
				\coordinate[label={above right}:{\footnotesize $D$}] (D) at ($(C)+(\d,0)$);
				\coordinate[label={above}:{\footnotesize $S$}] (S) at ($(C)+(0,\h)$);
				\draw (S)--(B)--(A)--(D)--(S)--(A);
				\draw[dashed] (S)--(C)--(B) (C)--(D);
				\tkzMarkRightAngles(S,C,D S,C,B B,C,D A,B,C)
				\tkzLabelSegment[right, yshift=-2mm](S,C){\footnotesize $a\sqrt{3}$}
				\tkzLabelSegment[below](A,B){\footnotesize $a$}
				\tkzMarkAngles[size=0.6cm,arc=l](B,S,A)
				\tkzDrawPoints[fill=black](S,A,B,C,D)
			\end{tikzpicture}
		}
		\noindent Vậy $\tan$ góc giữa đường thẳng $SA$ và mặt phẳng $(SBC)$ là $\dfrac{1}{2}$.
	}
\end{bt}

\begin{bt}%[1H8H6-1]
% nguon: Archive/SP_KhoLatex/ID6_old/18-OT-Goc.tex #3680
%[ĐGNL-ĐGTD-2022]%[Đình Nguyên]
	Cho hình chóp $S.ABCD$ có đáy $ABCD$ là hình vuông cạnh $a$, $SA\perp (ABCD)$ và $SA=a$. Tính góc giữa đường thẳng $SB$ và $(ABCD)$.
	\loigiai{
		\immini{
			Ta có $SA\perp (ABCD)$, nên hình chiếu của $SB$ trên mặt phẳng $(ABCD)$ là $AB$.\\
			Vậy $(SB, (ABCD))=(SB,AB)=\widehat{SBA}$.\\
			Xét tam giác $SAB$ vuông tại $A$ có \\
			$\tan \widehat{SBA}=\dfrac{SA}{AB}=1\Rightarrow \widehat{SBA}=45^\circ$.}	
		{
			\begin{tikzpicture}[scale=0.9,>=stealth, font=\footnotesize, line join=round, line cap=round]
				\tkzDefPoints{0/0/A,-1.3/-1.6/B,2.5/-1.6/C}
				\coordinate (D) at ($(A)+(C)-(B)$);
				\coordinate (S) at ($(A)+(0,3)$);
				\coordinate (O) at ($(A)!0.5!(C)$);
				\tkzDrawPolygon(S,B,C,D)
				\tkzDrawSegments(S,C)
				\tkzDrawSegments[dashed](A,S A,B A,D A,C B,D)
				\tkzDrawPoints[fill=black](D,C,A,B,S)
				\tkzMarkRightAngles(S,A,B)
				\tkzLabelPoints[above](S)
				\tkzLabelPoints[below](A,B,C)
				\tkzLabelPoints[right](D)
		\end{tikzpicture}}
	}
\end{bt}

\begin{bt}%[1H8H6-1]
% nguon: Archive/SP_KhoLatex/ID6_old/18-OT-Goc.tex #3681
%[ĐGNL-ĐGTD-2022]%[Đình Nguyên]
	Cho hình chóp tam giác đều $S.ABC$ có cạnh đáy bằng $a$, đường cao $SH=\dfrac{a\sqrt{3}}{3}$. Tính góc giữa cạnh bên và mặt đáy của hình chóp.
	\loigiai{
		\immini{Vì $S.ABC$ là hình chóp tam giác đều, nên $H$ là tâm đường tròn ngoại tiếp $\triangle ABC$.\\
			Suy ra $CH$ là hình chiếu của $SC$ trên $(ABC)$, \\
			do đó $(SC,(ABC))=(SC,CH)=\widehat{SCH}$.\\
			Mặt khác $\tan \widehat{SCH}=\dfrac{SH}{CH}=\dfrac{a\sqrt{3}}{3}: \dfrac{a\sqrt{3}}{3}=1$ $\Rightarrow \widehat{SCH}=45^\circ$.\\
			Vậy góc cần tìm là $45^\circ$.  }{\begin{tikzpicture}[scale=1.0]
				\tkzInit[ymin=-3.0,ymax=3.5,xmin=-0.5,xmax=5.5]
				\tkzClip
				
				\tkzDefPoints{0/0/A,5/0/C,1.7/-2.5/B}
				\tkzCentroid(A,B,C) \tkzGetPoint{H}
				\coordinate[label=above:{$S$}] (S) at ($(H)+(0,3.5)$);
				
				\tkzDrawSegments(A,B B,C A,S S,C S,B)
				\tkzDrawSegments[dashed](A,C S,H C,H B,H A,H)
				
				\tkzMarkAngles[arc=l,size=0.7cm](S,C,H)
				\tkzLabelPoints[below right](B,H)
				\tkzLabelPoints[left](A)
				\tkzLabelPoints[right](C)
				
				\foreach \X in {S,A,B,C,H}
				{\draw[fill=black] (\X) circle (1.0pt);}
				
		\end{tikzpicture}}
	}
\end{bt}

\begin{bt}%[1H8V2-2]
% nguon: Archive/SP_KhoLatex/ID6_old/18-OT-Goc.tex #3682
%[ĐGNL-ĐGTD-2022]%[Đình Nguyên] 
	Cho hình chóp $S.ABCD$, đáy $ABCD$ là hình vuông cạnh bằng $a$, $ S A \perp(ABCD)$, $S A=a \sqrt{6}$.
	\begin{enumerate}
		\item Chứng minh $B C \perp(S A B)$, $B D \perp(S A C)$.
		\item Tính sin của góc tạo bởi đường thẳng $A C$ với mặt phẳng $(S B C)$.
	\end{enumerate}
	\loigiai{
		\immini
		{
			\begin{enumerate}
				\item Ta có $ \heva{& S A \perp(ABCD) \\ & ABCD\text{ là hình vuông}}\Rightarrow \heva{& SA\perp BC \\ & BC\perp AB}\Rightarrow B C \perp(S A B) $.\\
				Mặt khác $ \heva{& S A \perp(ABCD) \\ & ABCD\text{ là hình vuông}}\Rightarrow \heva{& SA\perp BD \\ & BD\perp AC}\Rightarrow B D \perp(S A C) $.
				\item Kẻ $ AH\perp SB $, $ H\in SB $.\\
				Do $B C \perp(S A B) $ nên $ BC\perp AH $.\\
				Dẫn đến $ AH\perp (SBC) \Rightarrow (AC,(SBC))=(AC,CH)=\widehat{ACH}$.
			\end{enumerate}
		}
		{
			\begin{tikzpicture}[line join = round, line cap = round,>=stealth,font=\footnotesize,scale=1]
				\tkzDefPoints{0/0/A}
				\coordinate (B) at ($(A)+(3,0)$);
				\tkzDefShiftPoint[A](-140:1.7){D}
				\coordinate (C) at ($(B)+(D)-(A)$);
				\coordinate (S) at ($(A)+(0,2.5)$);
				\coordinate (H) at ($(B)!0.45!(S)$);
				\tkzMarkRightAngles(S,A,B S,A,D A,H,S)
				\tkzDrawPolygon(S,B,C,D)
				\tkzDrawSegments(S,C C,H)
				\tkzDrawSegments[dashed](A,B A,S A,D A,H A,C)
				\tkzDrawPoints[fill=black](A,B,D,C,S,H)
				\tkzLabelPoints[above](S)
				\tkzLabelPoints[left](D)
				\tkzLabelPoints[above right](H)
				\tkzLabelPoints[right](B,C)
				\tkzLabelPoints[above left](A)
			\end{tikzpicture}
		}
		\noindent Do $ ABCD $ là hình vuông cạnh $ a $ nên $ AC=a\sqrt{2} $.\\
		Áp dụng hệ thức lượng trong $ \triangle SAB $, ta có $AH=\dfrac{SA\cdot AB}{\sqrt{SA^2+AB^2}}=\dfrac{a \sqrt{6}\cdot a}{\sqrt{6a^2+a^2}}=\dfrac{a\sqrt{42}}{7}$.
		Mặt khác $AH\perp (SBC)$ nên $ AH\perp HC$
		\[\Rightarrow \sin
		(AC,(SBC))=\sin\widehat{ACH}=\dfrac{AH}{AC}=\dfrac{\dfrac{a\sqrt{42}}{7}}{a\sqrt{2}}=\dfrac{\sqrt{21}}{7}.\]
	}
\end{bt}

\begin{bt}%[1H8V6-1]
% nguon: Archive/SP_KhoLatex/ID6_old/18-OT-Goc.tex #3683
%[ĐGNL-ĐGTD-2022]%[Đình Nguyên]
	Cho hình chóp $S.ABCD$ có đáy $ABCD$ là hình vuông tâm $O$ cạnh bằng $a$, $SO$ vuông góc với đáy. Gọi $M$, $N$ lần lượt là trung điểm của $SA$ và $BC$. Tính góc giữa đường thẳng $MN$ với mặt phẳng $(ABCD)$, biết $MN=\dfrac{a\sqrt{10}}{2}$.
	\loigiai{
		\immini{Gọi $H$ là trung điểm của $AO$ suy ra $MH\parallel SO$. Mà $SO\perp (ABCD)$ nên $MH\perp (ABCD)$, do đó góc giữa đường thẳng $MN$ và $(ABCD)$ bằng $\widehat{MNH}$. Theo định lí cô-sin trong $\triangle HNC$ ta có
			{\allowdisplaybreaks
				\begin{align*}
					HN&=\sqrt{NC^2+HC^2-2NC\cdot HC\cdot\cos 45^\circ}\\
					&=\sqrt{\left(\dfrac{a}{2}\right)^2+\left(\dfrac{3a\sqrt{2}}{4}\right)^2-2\cdot\dfrac{a}{2}\cdot\dfrac{3a\sqrt{2}}{4}\cdot\dfrac{\sqrt{2}}{2}}\\
					&=\dfrac{a\sqrt{10}}{4}.
			\end{align*}}
		}
		{\begin{tikzpicture}[>=stealth,line join=round,line cap=round,font=\footnotesize,scale=1]
				\coordinate[label=below left:$A$] (A) at (0,0);
				\coordinate[label=above right:$D$] (D) at (1.5,1);
				\coordinate[label=below right:$B$] (B) at (4,0);
				\coordinate[label=right:$C$] (C) at ($(B)-(A)+(D)$);
				\coordinate[label=below:$O$] (O) at ($(C)!.5!(A)$);
				\coordinate[label=above left:$S$] (S) at ($(O)+(0,4)$);
				\coordinate[label=left:$M$] (M) at ($(S)!.5!(A)$);
				\coordinate[label=below right:$N$] (N) at ($(B)!.5!(C)$);
				\coordinate[label=above left:$H$] (H) at ($(A)!.5!(O)$);
				\draw (S)--(A)--(B)--(C)--(S)--(A) (S)--(B);
				\draw[dashed] (A)--(D)--(C)--cycle (M)--(N)--(H)--cycle (O)--(S)--(D)--(B);
				\fill (A)circle(1.5pt) (B)circle(1.5pt) (C)circle(1.5pt) (D)circle(1.5pt) (S)circle(1.5pt) (H)circle(1.5pt) (M)circle(1.5pt) (N)circle(1.5pt) (O)circle(1.5pt);
		\end{tikzpicture}}
		\noindent{Xét} tam giác $MNH$ có $\cos\widehat{MNH}=\dfrac{NH}{NM}=\dfrac{1}{2}\Rightarrow \widehat{MNH}=60^\circ$. Vậy góc giữa đường thẳng $MN$ với mặt phẳng $(ABCD)$ bằng $60^\circ$.
	}
\end{bt}

\begin{bt}%[1H8H6-1]
% nguon: Archive/SP_KhoLatex/ID6_old/18-OT-Goc.tex #3684
%[ĐGNL-ĐGTD-2022]%[Đình Nguyên]
	\immini{
		Cho hình lập phương $ABCD.A'B'C'D'$. Góc giữa đường thẳng $BD'$ và mặt phẳng $(ADC')$ bằng $\alpha$. Tính $\tan \alpha$.
	}
	{
		\begin{tikzpicture}[>=stealth, font=\footnotesize, line join=round, line cap=round, scale=0.8]
			\tkzDefPoints{0/0/A', 3/0/D'}
			\coordinate (C') at ($(D')-(1,-0.6)$);
			\coordinate (B') at ($(A')-(1,-0.6)$);
			\tkzDefSquare(D',A')\tkzGetPoints{A}{D}
			\tkzDefSquare(C',B')\tkzGetPoints{B}{C}
			
			\tkzDrawSegments[dashed](B,C C,D C,C' D,C' A,C' B,D' D',C)
			\tkzDrawSegments[](A,B D,A A',B' A',D' A,A' B,B' D,D' C',B' C',D')
			\tkzLabelPoints[right](D')
			\tkzLabelPoints[above](B',C')
			\tkzLabelPoints[below left](A')
			\tkzLabelPoints[below](A,D)
			\tkzLabelPoints[above left](C)
			\tkzLabelPoints[left](B)
			\tkzDrawPoints[fill=black](A,B,C,D,A',B',C',D')
		\end{tikzpicture}
	}
	\loigiai{
		\immini{
			Gọi $O$, $M$ lần lượt là trung điểm của $BD'$, $CD'$. Ta có $D'C\perp DC'$ và $D'C\perp AD$ nên $D'C\perp \left(ADC'\right)$, suy ra 
			\[\alpha=\left( BD',(ADC') \right)=\left( OD', OM\right)=\widehat{ MOD'}.\]
			Do đó
			\[\tan \alpha =\dfrac{MD'}{MO}=\dfrac{CD'}{CB}=\sqrt{2}.\]
		}
		{
			\begin{tikzpicture}[>=stealth, font=\footnotesize, line join=round, line cap=round, scale=1.2]
				\tkzDefPoints{0/0/A', 3/0/D'}
				\coordinate (C') at ($(D')-(1,-0.6)$);
				\coordinate (B') at ($(A')-(1,-0.6)$);
				\tkzDefSquare(D',A')\tkzGetPoints{A}{D}
				\tkzDefSquare(C',B')\tkzGetPoints{B}{C}
				\tkzDefMidPoint(D,C')\tkzGetPoint{M}
				\tkzDefMidPoint(B,D')\tkzGetPoint{O}
				
				\tkzDrawSegments[dashed](B,C C,D C,C' D,C' A,C' B,D' D',C O,M)
				\tkzDrawSegments[](A,B D,A A',B' A',D' A,A' B,B' D,D' C',B' C',D')
				\tkzLabelPoints[right](D')
				\tkzLabelPoints[above](B',C')
				\tkzLabelPoints[below left](A')
				\tkzLabelPoints[below](A,D)
				\tkzLabelPoints[above left](C)
				\tkzLabelPoints[left](B)
				\tkzLabelPoints[above left](O)
				\tkzLabelPoints[right](M)
				\tkzDrawPoints[fill=black](A,B,C,D,A',B',C',D',O,M)
			\end{tikzpicture}
		}
	}
\end{bt}

\begin{bt}%[1H8H2-2]
% nguon: Archive/SP_KhoLatex/ID6_old/18-OT-Goc.tex #3685
%[ĐGNL-ĐGTD-2022]%[Đình Nguyên]
	Cho hình chóp $S.ABCD$ có đáy là hình vuông cạnh $a$, cạnh bên $SA \perp (ABCD)$ và $SA = a\sqrt{2}$.
	\begin{listEX}[1]
		\item Chứng minh $BC \perp (SAB)$.
		\item Chứng minh $(SAC) \perp (SBD)$.
		\item Tính góc giữa hai mặt phẳng $(SBC)$ và $(ABCD)$.
	\end{listEX}
	\loigiai{
		\begin{listEX}[1]
			\immini{
				\item Ta có $ABCD$ là hình vuông nên $BC \perp AB$. Vì $SA \perp (ABCD)$ nên $SA \perp BC$.\\
				Như vậy $\heva{&BC \perp AB \quad \textrm{ (cmt)}\\&BC \perp SA \quad \textrm{ (cmt)}\\&AB \cap SA = A\\&AB,SA \subset (SAB)} \Rightarrow BC \perp (SAB)$.
				\item Ta có $ABCD$ là hình vuông nên $BD \perp AC$.  Vì $SA \perp (ABCD)$ nên $SA \perp BD$.
			}{
				\begin{tikzpicture}[scale=1, font=\footnotesize, line join=round, line cap=round, >=stealth]
					\tkzDefPoints{0/0/A, 4/0/D, -1.5/-1.5/B, 0/3/S}
					\tkzDefPointBy[translation=from A to D](B)\tkzGetPoint{C}
					\tkzDrawSegments(B,C C,D S,B S,C S,D)
					\tkzDrawSegments[dashed](A,B S,A A,C B,D A,D)
					\tkzLabelPoints[above left](A)
					\tkzLabelPoints[below](B,C)
					\tkzLabelPoints[right](D)
					\tkzLabelPoints[above](S)
					\tkzDrawPoints(S,A,B,C,D)
					\tkzLabelSegment[right](A,B){$a$}
					\tkzLabelSegment[right](S,A){$a\sqrt{2}$}
					\tkzMarkAngles[size=.7cm](A,B,S)
				\end{tikzpicture}	
			}
			Như vậy $\heva{&BD \perp AC \quad \textrm{ (cmt)}\\&BD \perp SA \quad \textrm{ (cmt)}\\&AC \cap SA = A\\&AC,SA \subset (SAC)} \Rightarrow BD \perp (SAC)$ mà $BD \subset (SBD)$ nên $(SBD) \perp (SAC)$.
			\item Ta có $\heva{&(SBC) \cap (ABCD) = BC\\&BC \perp AB \quad \textrm{ (cmt)}\\&BC \perp SB \quad \textrm{ ($BC \perp (SAB)$)}} \Rightarrow ((SBC),(ABCD)) = \widehat{(SB,AB)} = \widehat{SBA}$.\\
			Xét $\triangle SAB$ vuông tại $A$, ta có $\tan\widehat{SBA} = \dfrac{SA}{AB} = \dfrac{a\sqrt{2}}{a} = \sqrt{2} \Rightarrow \widehat{SBA} \approx 54^\circ 44'8.2''$.\\
			Vậy góc giữa hai mặt phẳng $(SBC)$ và $(ABCD)$ gần bằng $54^\circ 44'8.2''$.
		\end{listEX}
	}
\end{bt}

\begin{bt}%[1H8H6-2]
% nguon: Archive/SP_KhoLatex/ID6_old/18-OT-Goc.tex #3686
%[ĐGNL-ĐGTD-2022]%[Đình Nguyên]
	Cho hình chóp $S.ABCD$ có đáy $ABCD$ là hình chữ nhật $AB=a,AD=a\sqrt{3}$. Cạnh bên $SA$ vuông góc với đáy và $SA=a$.
	\begin{enumerate}[a.]
		\item Tính góc giữa hai mặt phẳng $\left(SCD \right)$ và $\left(ABCD \right)$.
		\item Tính góc giữa hai mặt phẳng $\left(SBC \right)$ và $\left(SAD \right)$.
	\end{enumerate}
	
	\loigiai
	{\begin{center}
			\begin{tikzpicture}
				\draw [shift={(3.,8.5)},,color=violet,fill=violet,fill opacity=0.1] (0,0) -- (-110.55604521958348:0.5) arc (-110.55604521958348:-90.:0.5) -- cycle;
				\draw [shift={(6.5,6.)},,color=violet,fill=violet,fill opacity=0.1] (0,0) -- (144.46232220802563:0.5) arc (144.46232220802563:180.:0.5) -- cycle;
				
				\draw [dashed] (3.,6.)-- (1.5,4.5) (6.5,6.)-- (3.,6.) (3.,6.)-- (3.,8.5);
				\draw (1.5,4.5)-- (5.,4.5)-- (6.5,6.) (3.,8.5)-- (1.5,4.5) (3.,8.5)-- (5.,4.5) (3.,8.5)-- (6.5,6.) (1.5,8.5)-- (4.5,8.5);
				
				\draw  (3.,8.5) node[above] {$S$};
				\draw  (3.,6.) node[left] {$A$};
				\draw  (1.5,4.5) node[below left] {$B$};
				\draw  (5.,4.5) node[below left] {$C$};
				\draw  (6.5,6.) node[right] {$D$};
				\draw (4,8.5) node[above] {$d$};
				
			\end{tikzpicture}
		\end{center}
		\begin{enumerate}[a.]
			\item Ta có $(SCD)\cap (ABCD)=CD$
			
			$\left\{ \begin{aligned}
				& CD\perp SA \\ 
				& CD\perp AD 
			\end{aligned} \right.  \Rightarrow CD\perp (SAD)$.
			
			Ta có $(SAD)\cap (ABCD)=AD,(SAD)\cap (SCD)=SD$.
			
			$\Rightarrow \widehat{ \left(\left(SCD \right),\left(ABCD \right) \right)}=\widehat{\left(DA,SD \right)}=\widehat{SDA}=\varphi $.
			
			Ta có $\tan \varphi =\dfrac{SA}{AD}=\dfrac{a}{a\sqrt{3}}=\dfrac{1}{\sqrt{3}}\Rightarrow \varphi ={30^\circ}$.
			
			\item Ta có $\left\{\begin{aligned}
				& AD\subset \left(SAD \right) \\ 
				& BC\subset \left(SBC \right) \\ 
				& AD\parallel BC \\ 
			\end{aligned} \right.\Rightarrow \left(SAD \right)\cap \left(SBC \right)=d\parallel AD\parallel BC$.
			
			Vì $\left\{\begin{aligned}
				& SA\perp d \\ 
				& d\parallel AD \\ 
			\end{aligned} \right.\Rightarrow SA\perp d$, $\left\{\begin{aligned}
				& d\parallel AD \\ 
				& AD\perp AD \\ 
			\end{aligned} \right.\Rightarrow d\perp AB$ nên $(SAB)\perp d$.
			
			Ta có $(SAB)\cap (SBC)=SB,(SAB)\cap (SAD)=SA$ suy ra $\widehat{ASB}$ chính là góc giữa hai mặt phẳng $(SBC)$ và $(SAD)$.
			
			Tam giác $ASB$ vuông cân tại $A$ nên $\widehat{ASB}={45^\circ}$.
		\end{enumerate}
		
	}
	
\end{bt}

\begin{bt}%[1H8V6-2]
% nguon: Archive/SP_KhoLatex/ID6_old/18-OT-Goc.tex #3687
%[ĐGNL-ĐGTD-2022]%[Đình Nguyên]
	Cho tứ diện $S.ABC$ có đáy $ABC$ là tam giác đều cạnh $a$, $SA\perp (ABC)$ và $SA=\dfrac{3a}{2}$. Tính góc giữa hai mặt phẳng $(SBC)$ và $(ABC)$.
	\loigiai{
		\immini{
			Gọi góc giữa hai mặt phẳng $(SBC)$ và $(ABC)$ là $\alpha$.\\
			Gọi $M$ là trung điểm của $BC$. Do $\triangle ABC$ đều nên $AM\perp BC$.\hfill(1)\\
			Theo giả thiết $SA\perp(ABC)$, suy ra theo (1) ta có $SM\perp BC$.\hfill(2)\\
			Lại có $(SBC)\cap(ABC)=BC$.\hfill(3)\\
			Từ (1), (2) và (3) ta có $\alpha=\widehat{SMA}$.\\
			Ta có $AM=\sqrt{AC^2-CM^2}=\dfrac{a\sqrt{3}}{2}$.\\
			Xét tam giác $SAM$ vuông tại $A$, ta có $\tan\alpha=\dfrac{SA}{AM}=\dfrac{3}{\sqrt{3}}=\sqrt{3}$, suy ra $\alpha=60^\circ$.
		}{
			\begin{tikzpicture}[>=stealth,scale=.6,font=\footnotesize]
				\tkzDefPoints{0/0/A,0/4/S,4/0/B,1.5/-2/C}
				\tkzDefMidPoint(B,C)\tkzGetPoint{M}
				\tkzDrawSegments(S,A S,B S,C A,C B,C S,M)
				\tkzDrawSegments[dashed](A,B A,M)
				\tkzMarkAngle[mkpos=.2,size=.6](S,M,A)
				\tkzMarkSegments[mark=|](C,M M,B)
				\tkzMarkRightAngles(A,M,C S,M,B S,A,C S,A,B)
				\tkzDrawPoints(A,B,C,S,M)
				\tkzLabelPoints[left](A,S)
				\tkzLabelPoints[right](B,M)
				\tkzLabelPoints[below](C)
			\end{tikzpicture}
		}
	}
\end{bt}

\begin{bt}%[1H8V6-2]
% nguon: Archive/SP_KhoLatex/ID6_old/18-OT-Goc.tex #3688
%[ĐGNL-ĐGTD-2022]%[Đình Nguyên]
	Cho hình chóp $S.ABCD$ có đáy $ABCD$ là hình chữ nhật, $AB=a,AD=a\sqrt{3}$, $ SA\perp(ABCD)$.
	\begin{enumerate}[a)]
		\item Tính góc giữa hai mặt phẳng $(SCD)$ và $(ABCD)$ với $SA=a$.
		\item Tìm $x=SA$ để góc giữa hai mặt phẳng $(SCD)$ và $(ABCD)$ bằng $60^\circ$.
	\end{enumerate}
	\loigiai{
		\immini{
			\begin{enumerate}[a)]
				\item Vì $SA\perp (ABCD)\Rightarrow CD\perp SA$, theo giả thiết $CD\perp AD$ và $(SCD)\cap(ABCD)=CD$, suy ra $\Rightarrow \alpha=\left((SCD),(ABCD)\right)=\widehat{SDA}$.\\
				Xét tam giác $SAD$ vuông tại $A$ ta có\\
				$\tan\alpha=\dfrac{SA}{AD}=\dfrac{a}{a\sqrt{3}}=\dfrac{1}{\sqrt{3}}\Rightarrow \alpha=30^\circ$.
				\item Theo kết quả câu a ta có: $\tan60^\circ=\dfrac{x}{a\sqrt{3}}\\
				\Leftrightarrow x=a\sqrt{3}\tan60^\circ=3a$.
			\end{enumerate}
		}{
			\begin{tikzpicture}[>=stealth,scale=.7,font=\footnotesize]
				\tkzDefPoints{0/0/A,4/0/B,-1.5/-2/D,2.5/-2/C,0/4/S}
				\tkzDrawSegments(S,B S,C S,D B,C C,D)
				\tkzDrawSegments[dashed](S,A A,B A,D)
				\tkzMarkRightAngles(S,A,D S,A,B B,A,D)
				\tkzMarkAngles(A,D,S)
				\tkzDrawPoints(A,B,C,D,S)
				\tkzLabelPoints[left](S,D)
				\tkzLabelPoints[right](B,C)
				\tkzLabelPoints[below](A)
			\end{tikzpicture}
		}
	}
\end{bt}

\begin{bt}%[1H8V4-3]
% nguon: Archive/SP_KhoLatex/SP_Đề Theo Bài-Khối 11/1-B25.tex #3918

	Cho hình lập phương $ABCD.A'B'C'D'$ có cạnh bằng $a$. Tính số đo của góc giữa hai mặt phẳng $\left( BA'C \right)$ và $\left( DA'C \right)$. (kết quả được làm tròn đến một chữ số thập phân)
	\loigiai{
	\immini{
	Ta có: $AH\perp \left( BA'C \right)$, $AK\perp \left( DA'C \right)$ với $H$, $K$ lần lượt là trung điểm của $A'B$, $A'D$.\\
	Suy ra $\left( \widehat{\left( BA'C \right),\left( DA'C \right)} \right)=\left( \widehat{AH,AK} \right)$.\\
	Lại có $HK$ là đường trung bình của $\Delta A'BD$ nên $HK=\dfrac{1}{2}BD=\dfrac{a\sqrt 2} 2$\\
	Mặt khác: $AH=AK=\dfrac{a\sqrt 2} 2$\\
	Do đó $AH=AK=HK=a\sqrt 2$\\
	Suy ra $\Delta AHK$ đều, suy ra $\widehat{HAK}=60^\circ $.\\
	Vậy $\bigl( {\left( BA'C \right),\left( DA'C \right)} \bigr)=\widehat{HAK}=60^\circ $.}
	{
	\begin{tikzpicture}[line join=round,line cap=round,line width=.6pt,font=\footnotesize,scale=1]
		\coordinate (B) at (0,0);
		\coordinate (A) at (1,.8);
		\coordinate (C) at (4,0);
		\coordinate (D) at ($(C)-(B)+(A)$);
		\coordinate (A') at ($(A)+(90:4)$);
		\coordinate (B') at ($(B)-(A)+(A')$);
		\coordinate (C') at ($(C)-(A)+(A')$);
		\coordinate (D') at ($(D)-(A)+(A')$);
		\coordinate (H) at ($(A')!0.5!(B)$);
		\coordinate (K) at ($(A')!0.5!(D)$);
		\draw (B')--(B)--(C)--(D)--(D')--(A')--(B')--(C')--(D') (C)--(C');
		\draw[dashed] (A)--(D) (A')--(A)--(B) (A')--(B)--(D)--(A')--(C) (A)--(H)--(K)--(A);
		
		\foreach \x/\g in {A/150,B/-90,C/-90,D/30,A'/90,B'/180,C'/-30,D'/90,H/150,K/30}
		\draw [fill= black] (\x) circle (.05)+(\g:.3) node{$\x$};
		
	\end{tikzpicture}
	
	}
	}
\end{bt}

\begin{bt}%[1H8H6-1]
% nguon: Archive/SP-4Phan-Dot3/Data/CD_1_OTGK2_De06.tex #4176

	Cho hình chóp $S.ABCD$ có đáy hình vuông cạnh $a$, $SA \perp (ABCD)$ và $SA=a\sqrt{2}$. Tính số đo của góc tạo bởi đường thẳng $SC$ và mặt phẳng $(SAB)$.
	\loigiai{	\begin{center}
			\begin{tikzpicture}[scale=1, font=\footnotesize, line join=round,line cap=round, >=stealth]
				\def\a{5}
				\def\b{2}
				\def\h{3}
				\path 
				(0,0) coordinate (A)
				(A)+(\a,0) coordinate (B)
				(A)+(-2,-\b) coordinate (D)
				(D)+(\a,0) coordinate (C)
				($(A)!.5!(C)$) coordinate (O)
				(A)+(0,\h) coordinate (S)
				;
				\draw[dashed] (D)--(A)--(B) (A)--(C) (A)--(S)--(O) (B)--(D) ;
				\draw (D)--(C)--(B)--(S)--(C) (S)--(D) ;
				\foreach \i/\g in {A/-90,B/-90,C/-90,D/-90,O/-90,S/90} \fill (\i) circle(1pt) + (\g:0.3) node{$\i$};
			\end{tikzpicture}
		\end{center}
		Ta có $SA \perp (ABCD)$ nên $SA \perp BC$ và $ABCD$ là hình vuông nên $AB \perp BC$.\\
		$\Rightarrow BC \perp (SAB) \Rightarrow$ Góc tạo bởi đường thẳng $SC$ và mặt phẳng $(SAB)$ là $\widehat{CSB}$. \\
		Ta có $SC=\sqrt{SA^2+AC^2}=\sqrt{2a^2+a^2}=a\sqrt{3}$.\\
		Xét tam giác $SBC$ vuông tại $B \Rightarrow \sin \widehat{CSB}=\dfrac{BC}{SC}=\dfrac{a}{a\sqrt{3}}=\dfrac{\sqrt{3}}{3}$.\\
		$\Rightarrow \widehat{CSB}=35{,}3^\circ$.
	}
\end{bt}

\begin{bt}%[1H8V2-3]
% nguon: Archive/SP-4Phan-Dot3/Data/CD_1_OTGK2_De06.tex #4177

	Cho hình chóp tứ giác $S.ABCD$ có đáy là hình vuông tâm $O$ cạnh bằng $a$, $SA \perp (ABCD)$  và $SA=a$. Gọi $H$ là hình chiếu vuông góc của $O$ lên $SC$.  Tính độ dài của đoạn thẳng $BH$.
	\loigiai{	\begin{center}
			\begin{tikzpicture}[scale=1, font=\footnotesize, line join=round,line cap=round, >=stealth]
				\def\a{5}
				\def\b{2}
				\def\h{4}
				\path 
				(0,0) coordinate (A)
				(A)+(\a,0) coordinate (B)
				(A)+(-2,-\b) coordinate (D)
				(D)+(\a,0) coordinate (C)
				($(A)!.5!(C)$) coordinate (O)
				(A)+(0,\h) coordinate (S)
				($(S)!0.6!(C)$) coordinate (H)
				;
				\draw[dashed] (D)--(A)--(B) (A)--(C) (A)--(S) (B)--(D) (O)--(H);
				\draw (D)--(C)--(B)--(S)--(C) (S)--(D) (B)--(H)--(D);
				\foreach \i/\g in {A/-90,B/-90,C/-90,D/-90,O/-90,S/90,H/45} \fill (\i) circle(1pt) + (\g:0.3) node{$\i$};
				\draw pic[draw,angle radius=0.3cm]{right angle=O--H--S};
			\end{tikzpicture}
		\end{center}
		Ta có $SA \perp (ABCD)$ nên $SA \perp BC$ và $ABCD$ là hình vuông nên $AB \perp BC$.\\
		$\Rightarrow BC \perp (SAB) \Rightarrow BC \perp SB$.\\
		$\Rightarrow \triangle SBC$ vuông tại $B$. 
		Vì $ABCD$ là hình vuông nên $AC \perp BD$.\\
		Lại có $SA \perp (ABCD) \Rightarrow SA \perp BD$.\\
		$\Rightarrow BD \perp (SAC) \Rightarrow BD \perp SC$.\\
		Ta có $BD \perp SC$ và $OH \perp SC$.\\
		$\Rightarrow SC \perp (BHD) \Rightarrow BH \perp SC$.\\
		Dễ  dàng ta có $SB=a\sqrt{2}$ và $SC=a\sqrt{3}$.\\
		$\Rightarrow BH \cdot SC=BC \cdot SB \Leftrightarrow BH=\dfrac{BC\cdot SB}{SC}=\dfrac{a\cdot a\sqrt{2}}{a\sqrt{3}}=\dfrac{\sqrt{6}}{3}a$.}
\end{bt}

\begin{bt}%[1H8H6-2]
% nguon: Archive/SP-4Phan-Dot3/Data/CD_1_OTGK2_De05.tex #4185

	Cho hình chóp $S.ABCD$ có $SA\perp \left(ABCD\right)$, $ABCD$ là hình thoi cạnh $a, AC=a$, $SA=\dfrac{a}{2}$. Tìm số đo của góc nhị diện $[S,CD,A]$.
	\loigiai{   
		\begin{center}
			\begin{tikzpicture}[scale=0.6,>=stealth,line join=round,line cap=round,font=\footnotesize]
				\coordinate (A) at (0,0);
				\coordinate (B) at (-2,-3);
				\coordinate (D) at (7,0);
				\coordinate (C) at ($(B)+(D)-(A)$);
				\coordinate (S) at ($(A)+(0,5)$);
				\coordinate (H) at ($(C)!0.5!(D)$);
				\draw(S)--(B) (H)--(S)--(C) (S)--(D) (B)--(C)--(D);
				\draw[dashed,thin](S)--(A)--(C) (H)--(A)--(B)--(D) (A)--(D) ;
				\pic[draw,thin,angle radius=2mm] {right angle = S--A--D} pic[draw,thin,angle radius=2mm] {right angle = S--A--B} pic[draw,thin,angle radius=2mm] {right angle = C--O--D} pic[draw,thin,angle radius=2mm] {right angle = A--H--C};
				\foreach \i/\g in {S/90,A/-90,B/-90,C/-90,D/0,H/0}{\draw[fill=white](\i) circle (1.5pt) ($(\i)+(\g:3mm)$) node[scale=1]{$\i$};}
			\end{tikzpicture}
		\end{center}
Có $BD\perp AC$ (do $ABCD$ là hình thoi) và\\ $BD\perp SA$ (do $SA\perp (ABCD)$ nên $BD\perp\left(SAC\right)$.\\
Ta có $\left((SBD),(ABCD)\right)= \widehat{SOA}$. Mà $SA=AO= \dfrac{AC}{2}= \dfrac{a}{2}$.\\ 
Suy ra $\triangle SAO$ vuông cân tại $A \Rightarrow \widehat{SOA}=45^\circ$.\\
Hình chiếu của $SD$ lên $\left(ABCD\right)$ là $AD$ do đó góc giữa $SD$ và $\left(ABCD\right)$ là góc giữa $SD$ và $AD$ chính là góc $\widehat{SDA}$.\\
Ta có $\tan\widehat{SDA}=\dfrac{SA}{AD}=\dfrac{\dfrac{a}{2}}{a}=\dfrac{1}{2}\Rightarrow\widehat{SDA}\approx 26,5^{\circ}$.\\
Vậy góc giữa $SD$ và $\left(ABCD\right)$ nhỏ hơn $30^{^\circ}$.\\
Gọi $H$ là hình chiếu của $A$ trên $CD$. Khi đó,  $AH\perp CD$. Vì $SA\perp (ABCD)$ nên $SA\perp CD$. Suy ra $CD\perp (SAH)$. Khi đó, $SH\perp CD$. Như vậy, số đo của $[S,CD,A]$ bằng $\widehat{SHA}$.\\
Ta có $AH=\dfrac{a\sqrt{3}}{2},SA=\dfrac{a}{2}$ nên $\tan\widehat{SHA}=\dfrac{SA}{AH}=\dfrac{\dfrac{a}{2}}{\dfrac{a\sqrt{3}}{2}}=\dfrac{\sqrt{3}}{3}.$ \\
Vậy số đo của góc nhị diện $[S,CD,A]$ bằng $\widehat{SHA}=30^{^\circ}$.
}
\end{bt}

\begin{bt}%[1H8H6-7]
% nguon: Archive/SP-4Phan-Dot3/Data/CTST_1-OTHK2_De05.tex #4199

	Độ dốc của mái nhà, mặt sân, con đường thẳng là tang của góc tạo bởi mái nhà mặt sân, con đường thẳng đó với mặt phẳng nằm ngang. Độ dốc của đường thẳng dành cho người khuyết tật được quy định là không quá $\dfrac{1}{12}$. Hỏi theo đó, góc tạo bởi đường dành cho người khuyết tật và mặt phẳng nằm ngang không vượt quá bao nhiêu độ? (Làm tròn kết quả đến chữ số thập phân thứ hai).
	\loigiai{
		Gọi $\alpha$ là góc tạo bởi mái nhà mặt sân, con đường thẳng với mặt phẳng nằm ngang. \\
		Theo giả thiết, $\tan\alpha\leq\dfrac{1}{12}\Rightarrow \alpha\leq 4{,}76^\circ$.
	} 
\end{bt}

\begin{bt}%[1H8V2-2]
% nguon: Archive/SP-4Phan-Dot3/Data/CD_1_OTGK2_De04.tex #4209

Cho hình chóp $S.ABCD$ có đáy là hình vuông, cạnh bên $SA$ vuông góc với mặt phẳng đáy. Gọi $H$, $K$ lần lượt là hình chiếu vuông góc của điểm $A$ lên $SC$, $SD$. Chứng minh $SC \perp (AHK)$.
\loigiai{\begin{center}
			\begin{tikzpicture}[line join=round, line cap=round,>=stealth,font=\footnotesize,scale=1]
			\def\a{3.5}
			\def\b{2.5}
			\def\h{3.5}
			\path (0:0) coordinate (A)
			++(0:\a) coordinate (D)
			++(-135:\b) coordinate (C)
			($(A)+(C)-(D)$) coordinate (B)
			($(A)!1/2!(C)$) coordinate (O)
			($(A)+(90:\h)$) coordinate (S)
			($(S)!1/3!(C)$) coordinate (H)
			($(S)!0.4!(D)$) coordinate (K);
			\draw[dashed,thick] (D)--(A) (A)--(C) (B)--(D) (A)--(B) (S)--(A) (H)--(A)--(K);
			\draw[thick] (S)--(B)--(C)--(D)--(S)--(C) (H)--(K);
			\foreach \x/\g in {A/180,B/180,C/0,D/0,O/-90,S/180,H/70,K/0}
			\fill[black] (\x) circle (1pt) ($(\g:4mm)+(\x)$) node {$\x$};	
			\newcommand{\gv}[4][black]{\draw[thick] ($(#3)!8pt!(#2)$)--($(#3)!2!($($(#3)!8pt!(#2)$)!.5!($(#3)!8pt!(#4)$)$)$)--($(#3)!8pt!(#4)$);}
			\gv{A}{H}{C}
			\gv{A}{K}{D}
		\end{tikzpicture}
\end{center}
Vì $SA \perp (ABCD)$ nên $SA \perp CD$.\\
Ta có $\heva{&CD \perp SA\\&CD \perp AD} \Rightarrow CD \perp (SAD)$.\\
Mà $AK \subset (SAD)$ nên $CD \perp AK$.\\
Ta có
$\heva{&AK \perp CD\\&AK \perp SD}\Rightarrow AK \perp (SCD)  \Rightarrow AK \perp SC$.\\
Do đó $\heva{&SC \perp AK\\&SC \perp AH} \Rightarrow SC \perp (AHK)$.
} 

\end{bt}

\begin{bt}%[1H8V6-1]
% nguon: Archive/SP-4Phan-Dot3/Data/CD_1_OTGK2_De04.tex #4210

	Cho hình lăng trụ đứng $ABC.A'B'C'$ có đáy $ABC$ là tam giác đều cạnh $a$. Tính góc giữa $AC$ và mặt phẳng $ABB'A'$.
	\loigiai{	\begin{center}
			\begin{tikzpicture}[scale=1, font=\footnotesize, line join=round, line cap=round, >=stealth]
				\path
				(0,3) coordinate[label=left: A'] (A')
				(0,0) coordinate[label=left: A] (A)
				(3,0) coordinate[label=right: C] (C)
				(3,3) coordinate[label=right: C'] (C')
				(1,2) coordinate[label=left: B'] (B')
				(1,-1) coordinate[label=left: B] (B)
				($(A)!.5!(B)$) coordinate[label=left: H] (H);
				\draw (A')--(B')--(C')--cycle;
				\draw (B)--(B')--(C')--(C)--cycle;
				\draw (A')--(A)--(B)  ;		
				\draw[dashed] (A)--(C)--(H) ;
			\end{tikzpicture}
		\end{center}
		Gọi $H$ là trung điểm $AB$.\\
		Ta có $CH \perp AB$ nên $AH$ là hình chiếu vuông góc của $AC$ lên mặt phẳng $(ABB'A')$.\\		
		Do đó, $ (AC,(ABB'A'))=(AC,AH)=\widehat{CAB}$.\\		
		Mặt khác, tam giác $ABC$ đều nên $\widehat{CAB}=60^{\circ}$ hay $(AC,(ABB'A'))=60^{\circ}$.}
\end{bt}

\begin{bt}%[1H8H6-1]
% nguon: Archive/SP-4Phan-Dot3/Data/CTST_1-OTHK2_De01.tex #4227
%[Câu 2]
	Cho hình chóp $S.ABCD$ có đáy là hình chữ nhật có $BC=2a$. Biết $SA$ vuông góc với mặt phẳng $(ABCD)$ và $SA=2a \sqrt{3}$. Góc giữa đường thẳng $SD$ và mặt phẳng $(ABCD)$ bằng.
	\loigiai{
		\immini{Vì $SA \perp(ABCD)$ nên $AD$ là hình chiếu vuông góc của $SD$ trên mặt phẳng $(ABCD)$.\\
			Suy ra $(SD;(ABCD))=(SD; AD)=\widehat{SDA}$.\\
			Lại có  $\triangle SAD$ vuông tại $A$ nên suy ra\\
			$\tan \widehat{SDA}=\dfrac{SA}{AD}=\sqrt{3} \Rightarrow \widehat{SDA}=60^\circ$.\\
			Vậy $(SD,(A B C D))=60^\circ$.
		}{\begin{tikzpicture}[declare function={r=3;gocB=-140;},font=\footnotesize]
				\path (0,0) coordinate (A)
				(90:0.8*r) coordinate (S)
				(0:r) coordinate (D)
				(gocB:0.55*r) coordinate (B)
				($(B)+(D)-(A)$) coordinate (C)
				;
				\foreach \x/\y/\z in {S/A/B,S/A/D,B/A/D}{
					\path pic[draw=red,angle radius=5pt]{right angle= \x--\y--\z};
				}
				\path pic[draw,angle radius=9pt]{angle= S--D--A};
				\draw (S)--(B)--(C)--(D)--(S)--(C);
				\draw [dashed] (S)--(A)	(B)--(A)--(D);
				\foreach \t/\g in {A/140,B/210,C/-30,D/0,S/90}{
					\fill[ball color=red] (\t) circle (1pt) node[shift={(\g:7pt)}]{$ \t $};
				}
			\end{tikzpicture}
		}
	}
\end{bt}

\begin{bt}%[1H8V2-3]
% nguon: Archive/SP-4Phan-Dot3/Data/CD_1_OTGK2_De01.tex #4236

	Cho hình chóp $S.ABC$ có $SA=SB=SC=a$, $\widehat{ASB}=90^\circ$, $\widehat{BSC}=60^\circ$ và $\widehat{ASC}=120^\circ$. Gọi $I$ là trung điểm cạnh $AC$. Chứng minh $SI \perp (ABC)$.
	\loigiai{
		\immini{
			Hình chóp $S.ABC$ có $SA=SB=SC$ nên chân đường vuông góc kẻ từ $S$ tới đáy là tâm đường tròn ngoại tiếp tam giác $ABC$.\\
			Ta có $\triangle SAB$ vuông cân tại $S$, $SA=SB=a$\\
			 $ \Rightarrow AB=a\sqrt{2}$.\\
			Mặt khác, $\triangle SBC$ cân tại $S$ có $\widehat{BSC}=60^\circ $ \\
			$ \Rightarrow \triangle SBC$ đều.
			Suy ra $BC=SB=SC=a$.\\
			Lại có $\triangle SAC$ cân tại $S$ và $\widehat{ASC}=120^\circ$. Áp dụng định lý Côsin, ta có 
			$$AC=\sqrt{SA^2+SC^2-2\cdot SA\cdot SC \cdot \cos 120^\circ}=\sqrt{a^2+a^2-2\cdot a\cdot a \cdot \dfrac{-1}{2}}=a\sqrt{3}.$$
			Xét tam giác $\triangle ABC$, ta thấy $AB^2+BC^2=AC^2$.\\
			Suy ra, tam giác $ABC$ vuông tại $B$.\\
			Do đó, tâm đường tròn ngoại tiếp tam giác $\triangle ABC$ là trung điểm của cạnh $AC$.\\
			Vậy $SI \perp (ABC)$.
		}
		{
			\begin{tikzpicture}[scale=0.8, font=\footnotesize, line join=round, line cap=round, >=stealth]
				\coordinate (A) at (-2,0);
				\coordinate (B) at (2,-2);
				\coordinate (C) at (3,0);
				\coordinate (M) at ($(A)!0.5!(B)$);
				\coordinate (O) at ($(C)!2/3!(M)$);
				\coordinate (S) at ($(O)+(0,4)$);
				\coordinate (I) at ($(A)!(S)!(C)$);
				\draw(S)--(A) (S)--(B) (S)--(C) (B)--(C) (A)--(B) ;
				\draw[dashed,thin](A)--(C) (S)--(I);
				\foreach \i/\g in {S/90,A/180,B/-90,C/0,I/-90}{\draw[fill=black](\i) circle (1.5pt) ($(\i)+(\g:3mm)$) node[scale=1]{$\i$};}
			\end{tikzpicture}
		}
	}
\end{bt}

\begin{bt}%[1H8H6-7]
% nguon: Archive/SP-4Phan-Dot3/Data/CD_1_OTGK2_De01.tex #4237

	\immini[thm]{
		Một nhà kho được minh họa như hình bên. Biết kho có chiều cao bằng $8$ m, hai mái $EFIK$, $HGIK$ là hai hình chữ nhật bằng nhau, các bức tường tạo thành hình hộp chữ nhật $ABCD.EFGH$, $AB=10$ m, $AD=24$ m, $AE=7$ m. Khi đó góc tạo bởi hai mái nhà bằng bao nhiêu độ (kết quả làm tròn đến đơn vị độ)?
	}{
		\begin{tikzpicture}[scale=.75,>=stealth, font=\footnotesize, line join=round, line cap=round]
			\def \h{2}
			\path 
			(0:0) coordinate (A)
			++ (-120:2)coordinate (B)
			++(0:3) coordinate (C)
			($(A)+(C)-(B)$) coordinate (D)
			(A)++(90:\h) coordinate (E)
			(B)++(90:\h) coordinate (F)
			(C)++(90:\h) coordinate (G)
			(D)++(90:\h) coordinate (H)
			($(F)!.5!(G)$)++(90:1) coordinate (I)
			($(E)!.5!(H)$)++(90:1) coordinate (K)
			;
			\draw[fill=gray!50,opacity=200] (E)--(F)--(I)--(K)--cycle (H)--(G)--(I)--(K)--cycle;
			\draw[dashed] (B)--(A)--(D) (A)--(E)--(H);
%			\draw[->] (B)--++(-120:.7) node [right]{$x$};
%			\draw[->] (D)--++(0:1) node [below]{$y$};
%			\draw[->] (E)--++(90:1) node [left]{$z$};
			\draw (B)--(C)--(G)--(F)--cycle (C)--(D)--(H);
			\node[below right] at (A){$O$};
			\foreach \x/\g in {A/180,B/-135,C/-45,D/-45,E/150,F/180,G/-10,H/30,I/-90,K/6}
			\fill[black] 	(\x) circle (1pt)
			($(\g:3mm)+(\x)$) node {$\x$};
		\end{tikzpicture}
	}
	\loigiai{
	\immini{
		Vì $ABCD.EFGH$ là hình hộp chữ nhật nên ta có\\
		$AB=CD=EF=GH=10$,\\
		$AE=BF=CG=DH=7$,\\
		$AD=BC=EH=FG=24$.\\
		Chiều cao nhà kho $8$ m nên khoảng cách từ $I$ đến mặt phẳng $(EFGH)$ bằng $IM=1$.\\
		Khi đó góc tạo bởi hai mái nhà là góc phẳng nhị diện tạo bởi $(EFIK)$, $(HGIK)$ và $IK$.
		}{
			\begin{tikzpicture}[scale=.75,>=stealth, font=\footnotesize, line join=round, line cap=round]
				\def \h{2}
				\path 
				(0:0) coordinate (A)
				++ (-120:2)coordinate (B)
				++(0:3) coordinate (C)
				($(A)+(C)-(B)$) coordinate (D)
				(A)++(90:\h) coordinate (E)
				(B)++(90:\h) coordinate (F)
				(C)++(90:\h) coordinate (G)
				(D)++(90:\h) coordinate (H)
				($(F)!.5!(G)$)++(90:1) coordinate (I)
				($(E)!.5!(H)$)++(90:1) coordinate (K)
				($(F)!.5!(G)$) coordinate (M)
				;
				\draw[fill=gray!50,opacity=200] (E)--(F)--(I)--(K)--cycle (H)--(G)--(I)--(K)--cycle;
				\draw[dashed] (B)--(A)--(D) (A)--(E)--(H) (I)--(M);
				%			\draw[->] (B)--++(-120:.7) node [right]{$x$};
				%			\draw[->] (D)--++(0:1) node [below]{$y$};
				%			\draw[->] (E)--++(90:1) node [left]{$z$};
				\draw (B)--(C)--(G)--(F)--cycle (C)--(D)--(H);
				\node[below right] at (A){$O$};
				\foreach \x/\g in {A/180,B/-135,C/-45,D/-45,E/150,F/180,G/-10,H/10,I/30,K/6,M/40}
				\fill[black] 	(\x) circle (1pt)
				($(\g:4mm)+(\x)$) node {$\x$};
			\end{tikzpicture}
		}
		\noindent
	Theo đề bài $EFIK$, $HGIK$ là hai hình chữ nhật bằng nhau nên góc tạo bởi hai mái nhà bằng góc $\widehat{FIG}$.\\
	Xét tam giác $IGF$ cân tại $M$. Ta có 
	$\tan \widehat{MIG}=\dfrac{MG}{IM}=12\Rightarrow \widehat{MIG}\approx 85{,}24^\circ$.\\
	Suy ra $\widehat{FIG}=2\widehat{MIG}\approx 170{,}48^\circ$.\\
	Vậy góc tạo bởi hai mái nhà  (làm tròn đến đơn vị độ) bằng $170^\circ$.
	}
\end{bt}

\begin{bt}%[1H8C7-4]
% nguon: Archive/SP-4Phan-Dot3/Data/KNTT_1-OTHK2_De03.tex #4253

	Cho khối lăng trụ $ABC.A'B'C'$. Khoảng cách từ $C$ đến đường thẳng $BB'$ bằng $\sqrt 5 $,  khoảng cách từ $A$ đến các đường thẳng $BB'$ và $CC'$ lần lượt bằng $1$ và $2$,  hình chiếu vuông góc của $A$ lên mặt phẳng $\left(A'B'C'\right)$ là trung điểm $M$ của $B'C'$ và $A'M=\sqrt 5 $. Thể tích của khối lăng trụ đã cho bằng bao nhiêu? (kết quả được làm tròn đến hàng phần trăm).
	\loigiai{
		\begin{center}
			\begin{tikzpicture}[scale=1,line cap=round,line join=round,font=\footnotesize,>=stealth]
				\path
				(0,0) coordinate (A')
				($(A')+(-45:2)$) coordinate (B')
				($(A')+(0:5)$) coordinate (C')
				($(B')!0.5!(C')$) coordinate (M)
				($(M)+(90:3.5)$) coordinate (A)
				($(B')+(A)-(A')$) coordinate (B)
				($(C')+(A)-(A')$) coordinate (C)
				($(B)!0.5!(C)$) coordinate (N)
				($(B)!0.2!(B')$) coordinate (J)
				($(C)!0.5!(C')$) coordinate (K)
				($(J)+(C)-(K)$) coordinate (H)
				($(J)!0.5!(K)$) coordinate (F)
				(intersection of B--B' and A--N) coordinate (o)
				;
				\draw  (A)--(A')--(B')--(C')--(C)--(H)--(B)--(A) (B)--(B') (N)--(F) (A)--(J)--(K) (C)--(A)--(B)--(C) (A)--(o);
				
				\draw[dashed] (A')--(C') (A)--(M)--(A') (A)--(F) (A)--(K) (o)--(N);
				\foreach \p/\r in {A/90, B/-180, C/90, A'/-120,B'/-90,C'/-90,M/-90,J/140,F/-90,K/0,N/-60,H/90}
				\fill (\p) circle (1.25pt) node[shift={(\r:3mm)}]{$\p$};
			\end{tikzpicture}
		\end{center}
		\noindent Gọi $J$,  $K$ lần lượt là hình chiếu vuông góc của $A$ lên $BB'$ và $CC'$,  $H$ là hình chiếu vuông góc của $C$ lên $BB'$.\\
		Ta có $AJ\perp BB'$.\hfill (1)\\
		$AK\perp CC'\Rightarrow AK\perp BB'$.\hfill (2)\\
		Từ $(1)$ và $(2)$ suy ra $BB'\perp\left(AJK\right)$ $\Rightarrow BB'\perp JK$ $\Rightarrow JK\parallel CH$ $\Rightarrow JK=CH=\sqrt 5 $.\\
		Xét $\Delta AJK$ có $J{K^2}=A{J^2}+A{K^2}=5$ suy ra $\Delta AJK$ vuông tại $A$.\\
		Gọi $F$ là trung điểm $JK$ khi đó ta có $AF=JF=FK=\dfrac{\sqrt 5}{2}$.\\
		Gọi $N$ là trung điểm $BC$, có $NF\parallel BB'$ mà $BB'\perp\left(AJK\right)\Rightarrow NF\perp AF$.\\ Xét tam giác vuông $ANF$ ta có \\
		$\cos\widehat{NAF}=\dfrac{AF}{AN}$ $=\dfrac{\dfrac{\sqrt 5}{2}}{\sqrt 5}$ $=\dfrac{1}{2}$ $\Rightarrow\widehat{NAF}=60^\circ$  ($AN=A'M=\sqrt 5 $ vì $AN \parallel A'M$ và $AN=A'M$).\\
		Vậy ta có $S_{\Delta AJK}=\dfrac{1}{2}AJ\cdot AK$ $=\dfrac{1}{2}\cdot 1\cdot 2=1$\\ $\Rightarrow{S_{\Delta AJK}}=S_{\Delta ABC}\cdot \cos{60^\circ}$ $\Rightarrow{S_{\Delta ABC}}=\dfrac{S_{\Delta AJK}}{\cos{60^\circ}}=\dfrac{1}{\dfrac{1}{2}}=2$.\\
		Xét tam giác $AMA'$ vuông tại $M$ ta có $\widehat{MAA'}=\widehat{AMF}=30^\circ$ hay $AM=A'M\cdot \tan{30^\circ}$ $=\dfrac{\sqrt{15}}{3}$.\\
		Vậy thể tích khối lăng trụ là $V=AM\cdot S_{\Delta ABC}$ $=\dfrac{\sqrt{15}}{3}\cdot 2=\dfrac{2\sqrt{15}}{3}\approx 2{,}58$.}
\end{bt}

\begin{bt}%[1H8H2-5]
% nguon: Archive/SP-4Phan-Dot3/Data/CD_1_OTGK2_De03.tex #4261
 
	Cho hình chóp tam giác đều $S.ABC$ có cạnh đáy bằng $a\sqrt{3}$, cạnh bên bằng $2a$. Gọi $M$ là trung điểm của cạnh $BC$ và $O$ là trọng tâm tam giác $ABC$. Tính $\tan \widehat{\left(SM, (ABC)\right)}$.
	\loigiai{
		\begin{center}
			\begin{tikzpicture}[>=stealth,line join=round,line cap=round,scale=3.5]
				\draw (0,0)coordinate(A)--(-35:0.8)coordinate(B)--(2,0)coordinate(C);
				\draw[dashed]  ($(B)!.5!(C)$)coordinate(M1)--($(A)!2/3!(M1)$)coordinate(O)--(A)--(C) (O)--($(O)+(90:1)$)coordinate(S);
				\coordinate (M) at ($(B)!0.5!(C)$);
				\draw (A)--(S)--(C) (M)--(S)--(B);
				\path pic[angle radius=5,draw=blue,angle eccentricity=1.5] {right angle=S--O--A};
				\foreach \diem/\vitrin in {S/above,A/left,B/below left,C/right,O/below left,M/below right}	\fill (\diem)circle(0.5pt)node[\vitrin]{$\diem$};
			\end{tikzpicture}
		\end{center}

			Vì $S.ABC$ là chóp tam giác đều và $O$ là trọng tâm tam giác $ABC$ nên $SO\perp (ABC)$.\\
			Hình chiếu vuông góc của $SA$ trên $(ABC)$ là $SO$. Vậy góc giữa $SA$ và $(ABC)$ bằng góc giữa $SA$ và $AO$ và bằng góc $\widehat{SAO}$.\\
			Tam giác $SOA$ vuông tại $O$, có $AO=\dfrac{2}{3}\cdot \dfrac{a\sqrt{3}\cdot \sqrt{3}}{2}=a$.\\
			Vậy $SO^2=SA^2-AO^2=4a^2-a^2=3a^2 \Rightarrow SO=a\sqrt{3}$.\\
			Góc giữa đường thẳng $SM$ và $(ABC)$ là góc $\widehat{SMO}$ và
			$$\tan \widehat{SMO}=\dfrac{SO}{OM}=\dfrac{a\sqrt{3}}{\dfrac{a}{2}}=2\sqrt{3}.$$
}
\end{bt}

\begin{bt}%[1H8V2-3]
% nguon: Archive/SP-4Phan-Dot3/Data/CD_1_OTGK2_De03.tex #4263

	Cho hình chóp $S.ABCD$ có đáy $ABCD$ là hình vuông. Gọi $H$ là trung điểm của $AB$ và $SH\perp (ABCD)$, gọi $K$ là trung điểm của cạnh $AD$. Xác định góc của hai đường thẳng $BK$ và $SC$.
	\loigiai{
		\begin{center}
			\begin{tikzpicture}
				\def\a{3}
				\def\h{4}
				\path 	(0:0) coordinate (A)
				++(0:\a) coordinate (D)
				++(-130:\a/2) coordinate (C)
				($(A)+(C)-(D)$) coordinate (B)
				($(A)!0.5!(B)$) coordinate (H)
				($(H)+(90:\h)$) coordinate (S)
				($(A)!0.5!(D)$) coordinate (K)
				(intersection of A--C and B--D) coordinate (O);%giao điểm O
				\draw[dashed] 	(K)--(B)--(A)--(D)	(A)--(S)
				(S)--(H)--(C)	;
				\draw[thick] 			(B)--(C)--(D)
				(B)--(S)	(C)--(S)	(D)--(S);
				\foreach \x/\g in {A/135,B/-135,C/-45,D/45,S/90,H/135, K/-135}
				\fill[black] 	(\x) circle (1.5pt)
				($(\g:4mm)+(\x)$) node {$\x$};
				\draw pic[draw,angle radius=2mm]{right angle=S--H--A};
			\end{tikzpicture}
		\end{center}
		Xét hai tam giác $BCH$, $ABK$, ta có $\widehat{CBH}=\widehat{BAK}=90^{\circ}$, $BC=AB$, $BH=AK$.\\
		Suy ra $\triangle BCH=\triangle ABK \Rightarrow \widehat{ABK}=\widehat{BCH}$, mà $\widehat{ABK}+\widehat{CBK}=90^{\circ}$ nên $\break \widehat{CBK}+\widehat{BCH}=90^{\circ}$,
		suy ra $BK\perp CH$.\\
		Ta có $\heva{&BK\perp CH\\ &BK\perp SH} \Rightarrow BK\perp(SCH) \Rightarrow BK\perp SC$.
	}
\end{bt}

\begin{bt}%[1H8H6-1]
% nguon: Archive/SP-4Phan-Dot3/Data/CTST_1-OTHK2_De02.tex #4272

	Cho hình chóp $S.ABCD$ có đáy $ABCD$ là hình vuông tâm $O$ và tất cả các cạnh đều bằng $a$.
	\begin{enumerate}
		\item Chứng minh rằng $SO\perp (ABCD)$.
		\item Tính góc giữa đường thẳng $SA$ và mặt phẳng $(SBD)$.
	\end{enumerate}
	\loigiai{
		\begin{center}
			\begin{tikzpicture}[scale=1, font=\footnotesize, line join=round, line cap=round, >=stealth]
				\path 
				(0,0) coordinate (A)
				(-2,-1) coordinate (B)
				(4,0) coordinate (D)
				($(B)+(D)-(A)$) coordinate (C)
				($(A)!.5!(C)$) coordinate (O)
				(O)+(0,4) coordinate (S)
				;
				
				\draw (S)--(B)--(C)--(D)--cycle (S)--(C);
				\draw[dashed] (S)--(A) (B)--(A)--(D) (A)--(C) (S)--(O) (B)--(D);
				\foreach \p/\r in {A/120,B/180,C/-60,D/0,S/90,O/-90}
				\fill (\p) circle (1.5pt) node[shift={(\r:3mm)}]{$\p$};
			\end{tikzpicture}
		\end{center}
		\begin{enumerate}
			\item Ta có $\triangle SAC$ và $\triangle SBD$ là hai tam giác cân tại $S$.\\
			Suy ra $SO\perp AC$, $SO\perp BD$ mà $AC$ cắt $BD$ trong $(ABCD)$.\\
			Nên $SO\perp (ABCD)$.Ta có $\triangle SAC$ và $\triangle SBD$ là hai tam giác cân tại $S$.\\
			Suy ra $SO\perp AC$, $SO\perp BD$ mà $AC$ cắt $BD$ trong $(ABCD)$.\\
			Nên $SO\perp (ABCD)$.
			\item Ta có
			$\heva{&AO\perp BD\\&AO\perp SO}\Rightarrow AO\perp (SBD)$.\\
			Do đó, hình chiếu của $A$ lên $(SBD)$ là điểm $O$.\\
			Suy ra góc giữa đường thẳng $SA$ và mặt phẳng $(SBD)$ là góc $\widehat{ASO}$.\\
			Ta có $AO=\dfrac{AC}{2}=\dfrac{AB\sqrt{2}}{2}=\dfrac{a\sqrt{2}}{2}$.\\
			Mặt khác $SO=\sqrt{SA^2-AO^2}=\sqrt{a^2-\left(\dfrac{a\sqrt{2}}{2}\right)^2}=\dfrac{a\sqrt{2}}{2}$.\\
			Suy ra tam giác $SAO$ vuông cân tại $O$.\\
			Do đó, $\widehat{ASO}=45^\circ$.
		\end{enumerate}
		
	}
\end{bt}

\begin{bt}%[1H8V2-6]
% nguon: Archive/SP-4Phan-Dot3/Data/CD_1_OTGK2_De02.tex #4290

	Đại Kim tự tháp Giza có chiều cao $138{,}5$ m và độ dài các cạnh đáy là $230{,}33$ m. Xem Kim tự tháp này là một hình chóp tứ giác đều $S.ABCD$ có $O$ là tâm của đáy. Một đoàn thám hiểm đã đi đến được tâm $O$ của Kim tự tháp và nhận thấy ở mỗi mặt bên của kim tự tháp đều có một ô cửa trống, và từ vị trí của họ đến với ô cửa đó thì sẽ vuông góc với các mặt bên của kim tự tháp. Vậy khoảng cách từ vị trí đang đứng của đoàn thám hiểm đến ô cửa là bao nhiêu? (\textit{Kết quả làm tròn đến hàng thập phân})
	\loigiai{
		Theo đề bài, ta vẽ được hình sau
		\begin{center}
			\begin{tikzpicture}[line join=round,line cap=round,>=stealth,font=\footnotesize,scale=1]
				\def\a{4}
				\def\h{4}
				\path 	(0:0) coordinate (A)
				++(0:\a) coordinate (D)
				++(-140:\a/2) coordinate (C)
				($(A)+(C)-(D)$) coordinate (B)
				(intersection of A--C and B--D) coordinate (O)
				($(O)+(90:\h)$) coordinate (S)
				($(C)!.5!(D)$) coordinate (I)
				($(S)!.7!(I)$) coordinate (H)
				;
				\draw[dashed] (B)--(A)--(D)	(A)--(S) (A)--(C) (B)--(D) (S)--(O) (O)--(H) (O)--(I);
				\draw[thick] (B)--(C)--(D) (B)--(S)	(C)--(S) (D)--(S)--(I);
				\foreach \x/\g in {A/135,B/-135,C/-45,D/45,S/90,O/-90,I/-20,H/20}
				\fill[black] 	(\x) circle (1.5pt)
				($(\g:3mm)+(\x)$) node {$\x$};
				\draw pic[draw,angle radius=3mm]{right angle=D--O--S};
				\draw pic[draw,angle radius=3mm]{right angle=O--I--C};
				\draw pic[draw,angle radius=3mm]{right angle=O--H--I};
				\draw pic[draw,angle radius=3mm]{angle=S--I--O};
			\end{tikzpicture}
		\end{center}
		Ta có $SO$ vuông với đáy, nên $SO$ vuông với $CD$.\\
		Gọi $I$ là trung điểm $CD$, do $ABCD$ là hình vuông nên $OI$ vuông với $CD$.\\
		Suy ra $CD\perp(SIO)$.\\
		Gọi $H$ là hình chiếu của $O$ lên $SI$.\\
		Do $CD\perp(SIO)$ nên $OH\perp CD$ ($SC$ chứa trong $(SCD)$).\\
		Mà $OH\perp SI$ ($H$ là hình chiếu) ($SI$ chứa trong $(SCD)$).\\
		Nên $OH\perp(SCD)$.\\
		Vậy $H$ chính là ô cửa cần tìm.\\
		Do vậy khoảng cách từ đoàn thám hiểm đến ô cửa là $OH$.\\
		Bài toán trở thành tính độ dài $OH$.\\
		Theo định lý tam giác vuông, ta có
		\begin{eqnarray*}
			\dfrac{1}{OH^2}&=&\dfrac{1}{SO^2}+\dfrac{1}{OI^2}\\
			&=&\dfrac{1}{138{,}5^2}+\dfrac{1}{\left(\dfrac{230{,}33}{2} \right)^2} \\
			&\Rightarrow& OH\approx 88{,}6 \text{ (m).}
		\end{eqnarray*}
	}
\end{bt}

\begin{bt}%[1H8N7-3]
% nguon: Archive/SP-4Phan-Dot3/Data/KNTT_1-OTHK2_De02.tex #4296

    \immini{Cho hình chóp đều $S.ABCD$ có cạnh đáy bằng $a$, gọi $O$ là giao điểm của $AC$ và $BD$, $SO=3a$. Tính thể tích của khối chóp $S.ABCD$.
        }{
        \begin{tikzpicture} [font=\footnotesize, line join=round, line cap=round, >=stealth, scale=0.6]
            \def\a{4}
            \path (0:0) coordinate (A)
            ++(0:\a) coordinate (D)
            ++(-135:\a/2) coordinate (C)
            ($(A)+(C)-(D)$) coordinate (B)
            (intersection of A--C and B--D) coordinate (O)
            ($(O)+(90:\a)$) coordinate (S);
            \draw[dashed] 	(B)--(A)--(D)	(A)--(S)
            (A)--(C)	(B)--(D)	(S)--(O)	;
            \draw	(B)--(C)--(D)	(B)--(S)	(C)--(S)	(D)--(S);
            \foreach \x/\g in {A/180,B/-135,C/-45,D/0,S/90,O/-90}
            \fill (\x) circle (1pt)	($(\g:3mm)+(\x)$) node {$\x$};
            \pic[draw,angle radius=2mm]{right angle=D--O--S};%Theo chiều dương
        \end{tikzpicture}
    }
    \loigiai{
            Thể tích khối chóp $S.ABCD$ là $V=\dfrac{1}{3}S_{ABCD}\cdot SO=\dfrac{1}{3}\cdot a^2\cdot 3a=a^3$.
    }
\end{bt}

\begin{bt}%[1H8V7-2]
% nguon: Archive/SP-4Phan-Dot3/Data/CD_1_OTHK2_De04.tex #4303

	Cho lăng trụ tam giác $ABC.A'B'C'$ có đáy là tam giác đều cạnh bằng $5$, độ dài cạnh bên bằng $20$. Biết mặt phẳng $\left(BCC' B'\right)$ vuông góc với mặt phẳng đáy và góc $\widehat{B'BC}=30^{\circ}$. Tính thể tích khối chóp $A.CC'B'$.
	\loigiai{
		\begin{center}
			\begin{tikzpicture}[line join = round, line cap = round,>=stealth,font=\footnotesize,scale=.8,declare function={ab=3;ac=5;cao=4;gocA=-50;}]
				\path (0,0) coordinate (B)
				(gocA:ab) coordinate (A)
				(ac,0) coordinate (C)
				($(B)+(67:cao)$) coordinate (B')
				($(B')+(ac,0)$) coordinate (C')
				($(B')+(gocA:ab)$) coordinate (A')
				($(B)!0.5!(C)$) coordinate (M)
				($(B)!(B')!(C)$) coordinate (H)
				;
				\draw (A')--(B')--(C')--cycle (A')--(A) (B')--(B) (C')--(C)
				(B)--(A)--(C) (C')--(A)--(B')
				;
				\draw pic[draw,cap=butt,angle radius=4mm,"$30^\circ$",angle eccentricity=1.9]{angle=C--B--B'};
				\draw[dashed] (B)--(C)--(B') (A)--(M) (B')--(H);
				\foreach \x/\y/\z in {A/M/B,B'/H/B} \draw pic[draw,angle radius=2.5mm]{right angle=\x--\y--\z};
				\foreach \x/\gm in {B/180,A/-90,C/0,B'/90,A'/-50,C'/0,H/-90,M/60} \fill (\x) circle (1pt) ($(\x)+(\gm:3.5mm)$)node{$\x$};
			\end{tikzpicture}
		\end{center}
		Kẻ $B'H\perp BC$.\\
		Ta có $\heva{& (BCC'B')\perp(ABC) \\ & (BCC'B')\cap(ABC)=BC \\ & B'H\perp BC \\ & B'H\subset(BCC'B')}\Rightarrow B'H\perp(ABC)$.\\
		Gọi $M$ là trung điểm của $BC$ suy ra $AM\perp BC$.\\
		Ta có $\heva{& AM\perp BC \\ & AM\perp B'H}\Rightarrow AM\perp(BCC'B')\Rightarrow\mathrm{d}\left(A,(BCC'B')\right)=AM$.\\
		$AM=\sqrt{AC^2+MC^2}=\sqrt{5^2+\left(\dfrac{5}{2}\right)^2}=\dfrac{5\sqrt{5}}{2}$.\\
		Ta có $B'H=BB'\cdot\sin{30^\circ}=20\cdot\dfrac{1}{2}=10$.\\
		$S_{B'BC}=S_{CB'C'}=\dfrac{1}{2}\cdot B'H\cdot BC=\dfrac{1}{2}\cdot 10\cdot 5=25$.\\
		Thể tích khối chóp $A.CC'B'$:\\ $V_{A.CC'B'}=\dfrac{1}{3}\cdot AM\cdot S_{CC'B'}=\dfrac{1}{3}\cdot\dfrac{5\sqrt{5}}{2}\cdot 25=\dfrac{75\sqrt{5}}{6}$ (đvtt).
	}
\end{bt}

\begin{bt}%[1H8V5-3]
% nguon: Archive/SP-4Phan-Dot3/Data/CD_1_OTHK2_De04.tex #4304

	Cho hình chóp $S.ABCD$ có cạnh bên $SA$ vuông góc với mặt phẳng đáy, $ABCD$ là hình vuông cạnh bằng $4$. Biết góc giữa $SB$ và mặt đáy bằng $60^{\circ}$. Gọi $M$ là điểm nằm trên cạnh $CD$ sao cho $DM=2MC$. Tính khoảng cách từ điểm $M$ đến mặt phẳng $(SBD)$.
	\loigiai{
		\immini{
			Do $SA \perp(ABCD)$ nên $$\widehat{(SB,(ABCD))}=\widehat{SBA}=60^{\circ}.$$
			Tam giác $SAB$ vuông tại $A$ nên $$SA=AB \cdot \tan \widehat{SBA}=4\sqrt{3}.$$
			Gọi $I$ là giao điểm của $AM$ và $BD$. Ta có
			$$
			\dfrac{IM}{IA}=\dfrac{DM}{AB}=\dfrac{2}{3}=\dfrac{\mathrm{d}(M,(SBD))}{\mathrm{d}(A,(SBD))}.
			$$
			Suy ra $\mathrm{d}(M,(SBD))=\dfrac{2}{3} \mathrm{d}(A,(SBD))$.}
		{
			\begin{tikzpicture}[scale=1, font=\footnotesize, line join=round, line cap=round, >=stealth]
				\def\bc{4} \def\ba{2.25} \def\h{4} \def\gocB{30}
				\coordinate[label=below left:$D$] (B) at (0,0);
				\coordinate[label=above left:$A$] (A) at (\gocB:\ba);
				\coordinate[label=below:$C$] (C) at (\bc,0);
				\coordinate[label=right:$B$] (D) at ($(C)-(B)+(A)$);
				\coordinate[label=above:$S$] (S) at ($(A)+(90:\h)$);
				\draw (B)--(C)--(D)--(S)--cycle (S)--(C);
				\draw[dashed] (A)--(D) (S)--(A)--(B);
				\foreach \diem in {A,B,C,D,S}
				\fill (\diem)circle(1pt);
				\draw (A)--++(90:0.2)--++(0:0.2)--++(-90:0.2);
				\coordinate (M) at ($(C)!1/3!(B)$);
				\coordinate (O) at ($(C)!1/2!(A)$);
				\coordinate (H) at ($(O)!1/3.5!(S)$);
				\fill (M) circle (1pt) node[below]{$M$};
				\fill (O) circle (1pt) node[below]{$O$};
				\fill (H) circle (1pt) node[right]{$H$};
				\draw[dashed] (C)--(A)--(H) (S)--(O);
				% Giao điểm của AM và BD
				\draw[dashed,name path=AM] (A) -- (M);
				\draw[dashed,name path=BD] (B) -- (D);
				\path[name intersections={of=AM and BD, by=I}];
				\fill (I) circle (1pt) node[below]{$I$};
				\clip (A) -- (D) -- (S);
				\draw (D) circle (0.5cm);
				\draw ($(D)+(-0.25,0.25)$) node[left]{$60^\circ$};
			\end{tikzpicture}
		}\noindent
		Gọi $O$ là tâm đáy. Kẻ đường cao $AH$ của tam giác $SAO$.\\
		Khi đó $AH \perp(SBD)$.\\
		Suy ra $\mathrm{d}(A,(SBD))=AH=\dfrac{SA \cdot AO}{\sqrt{SA^2+AO^2}}=\dfrac{4\sqrt{3} \cdot \dfrac{4 \sqrt{2}}{2}}{\sqrt{(4\sqrt{3})^2+\left(\dfrac{4 \sqrt{2}}{2}\right)^2}}=\dfrac{4\sqrt{21}}{7}$.\\
		Do đó, $\mathrm{d}(M,(SBD))=\dfrac{2}{3} \mathrm{d}(A,(SBD))=\dfrac{2}{3} \cdot \dfrac{\sqrt{21}}{7}=\dfrac{8\sqrt{21}}{21}$.
	}
\end{bt}

\begin{bt}%[1H8V4-3]
% nguon: Archive/SP-4Phan-Dot3/Data/KNTT_1-OTGK2_De04.tex #4310

	Cho hình chóp tứ giác đều có cạnh đáy bằng $a$, chiều cao bằng $\dfrac{a}{2}$. Gọi $M$ là trung điểm $CD$ như hình vẽ. Tính số đo góc phẳng nhị diện $[S,CD,O]$.
	\begin{center}
		\begin{tikzpicture}[scale=0.7]
			\path 
			(0:0) coordinate (A)
			+(0:5) coordinate (D)
			+(-140:2.5) coordinate (B)
			($(B)+(D)-(A)$) coordinate (C)
			(intersection of A--C and B--D) coordinate (O)
			++(90:4) coordinate (S);
			\path (S)--(O) node[right, midway]{$\dfrac{a}{2}$};
			\path (B)--(C) node[below, midway]{$a$};
			\path ($(D)!.5!(C)$) coordinate (M);
			\draw[thin] pic[teal,draw, angle radius = 12pt]{angle = S--M--O};
			\foreach \x/\y/\z in {S/O/D,S/M/D,O/M/C}{
				\draw pic[red,draw, angle radius = 6pt]{right angle = \x--\y--\z};}
			\draw[dashed] 
			(A)--(B) (A)--(D) (A)--(S) (A)--(C) (B)--(D) (S)--(O)--(M);
			\draw 
			(D)--(C)--(B)
			(S)--(B) (S)--(C) (M)--(S)--(D);
			\foreach \x/\g in {A/135,D/0,C/-45,B/-135,O/-90,S/90,M/-45}
			\fill [blue] (\x) circle (1.5pt)
			+(\g:3mm) node {$\x$};
		\end{tikzpicture}
	\end{center}
	\loigiai{
		Ta có $\heva{&SM\perp CD\\ &OM\perp CD} \Rightarrow[S, CD, O]=\widehat{SMO}$.\\				
		Xét $\triangle SOM$ vuông tại $O$ có: $\tan \widehat{SMO}=\dfrac{SO}{OM} \Rightarrow \tan \widehat{SMO}=1\Rightarrow \widehat{SMO}=45^{\circ}$.
	}
\end{bt}

\begin{bt}%[1H8V7-9]
% nguon: Archive/SP-4Phan-Dot3/Data/KNTT_1-OTGK2_De04.tex #4311

Để chuẩn bị cho hoạt động cắm trại, bạn An tìm hiểu các mẫu lều cắm trại có kích thước như trong hình vẽ.
Bạn An muốn biết thể tích chênh lệch của hai lều nên thực hiện tính $V_1-V_2$, trong đó $V_1$, $V_2$ lần lượt là thể tích của mẫu lều cắm trại ở Hình a), b). Giá trị của $V_1-V_2$ bằng bao nhiêu decimét khối (làm tròn kết quả đến hàng đơn vị)?
\begin{center}
	\begin{tikzpicture}[line join=round,line cap=round,>=stealth,scale=1,font=\footnotesize]
		\tikzset{every node/.style={outer sep=-0.5mm}}
		\foreach \x/\y/\n in {0/0/A,3.5/1.8/C,0/1.8/B,3.5/0/D} \coordinate (\n) at (\x,\y);
		\coordinate (M) at ($(B)!0.5!(C)$);
		\coordinate (N) at ($(A)!0.5!(D)$);
		\coordinate (E) at ($(M)+(0,0.4)$);
		\coordinate (D') at ($(D)+(1.5,1.5)$);
		\coordinate (C') at ($(C)+(1.5,1.5)$);
		\coordinate (E') at ($(E)+(1.5,1.5)$);
		\coordinate (B') at ($(B)+(1.5,1.5)$);
		\fill[gray!50](A)--(B)--(B')--(E')--(C')--(D')--(D)--(A);
		\draw (A)--(B)node[midway,above,sloped]{$180$ cm}--(E)--(C)--(D)--(A)node[midway,below]{$350$ cm};
		\draw (B)--(B')--(E')--(C')--(D')--(D)node[midway,below,sloped]{$460$ cm} (E)--(E') (C)--(C');
		\draw[<->] (E)--(N)node[midway,fill=gray!50]{$220$ cm};
		\draw (N)node[yshift=-.75 cm]{Hình a)};
	\end{tikzpicture}\qquad
	\begin{tikzpicture}[line join=round,line cap=round,>=stealth,scale=1,font=\footnotesize]
		\tikzset{every node/.style={outer sep=-0.5mm}}
		\foreach \x/\y/\n in {0/0/A,3.15/2.1/C,0.55/2.1/B,3.7/0/D} \coordinate (\n) at (\x,\y);
		\coordinate (M) at ($(B)!0.5!(C)$);
		\coordinate (N) at ($(A)!0.5!(D)$);
		\coordinate (E) at ($(M)+(0,0.5)$);
		\coordinate (D') at ($(D)+(1.4,1.4)$);
		\coordinate (C') at ($(C)+(1.4,1.4)$);
		\coordinate (E') at ($(E)+(1.4,1.4)$);
		\coordinate (B') at ($(B)+(1.4,1.4)$);
		\fill[gray!50](A)--(B)--(B')--(E')--(C')--(D')--(D)--(A);
		\draw (A)--(B)--(E)--(C)--(D)--(A)node[midway,below]{$370$ cm};
		\draw (B)--(B')--(E')--(C')--(D')--(D)node[midway,below,sloped]{$430$ cm} (E)--(E') (C)--(C');
		\draw[<->] (E)--(N)node[midway,fill=gray!50]{$260$ cm};
		\draw (A)--($(A)+(-0.3,0)$) (B)--($(B)+(-0.85,0)$);
		\draw[<->] ($(A)+(-0.2,0)$)--($(B)+(-0.75,0)$)node[midway,above,sloped]{$210$ cm};
		\draw (B)--($(B)+(0,2.1)$) (C)--($(C)+(0,2.1)$);
		\draw[<->] ($(B)+(0,2)$)--($(C)+(0,2)$)node[midway,above,sloped]{$260$ cm};
		\draw (N)node[yshift=-.75 cm]{Hình b)};
	\end{tikzpicture}
\end{center}
\loigiai{
	\begin{center}
		\begin{tikzpicture}[line join=round,line cap=round,>=stealth,scale=1.25,font=\footnotesize]
			\tikzset{every node/.style={outer sep=-0.5mm}}
			\foreach \x/\y/\n in {0/0/A,3.5/1.8/C,0/1.8/B,3.5/0/D} \coordinate (\n) at (\x,\y);
			\coordinate (M) at ($(B)!0.5!(C)$);
			\coordinate (E) at ($(M)+(0,0.5)$);
			\coordinate (N) at ($(A)!0.5!(D)$);
			\draw (A)--(B)node[midway,sloped,above]{$180$}--(C)node[midway,below]{$350$}--(D)--(A)node[midway,below]{$350$} (B)--(E)--(C);
			\foreach \t/\g in {A/180,B/180,C/0,D/0,E/90}{
				\draw[fill=black] (\t) circle (.8pt) node[shift={(\g:8pt)}]{$ \t $};
			}
			\draw[<->] (E)--(M)node[midway,sloped,below]{$40$};
			\draw (N)node[yshift=-.75 cm]{Hình a)};
		\end{tikzpicture}\qquad
		\begin{tikzpicture}[line join=round,line cap=round,>=stealth,scale=1.25,font=\footnotesize]
			\tikzset{every node/.style={outer sep=-0.5mm}}
			\foreach \x/\y/\n in {0/0/A,3.15/2.1/C,0.55/2.1/B,3.7/0/D} \coordinate (\n) at (\x,\y);
			\coordinate (M) at ($(B)!0.5!(C)$);
			\coordinate (E) at ($(M)+(0,0.5)$);
			\coordinate (N) at ($(A)!0.5!(D)$);
			\draw (A)--(B)--(C)node[midway,below]{$260$}--(D)--(A)node[midway,below]{$370$} (B)--(E)--(C);
			\foreach \t/\g in {A/180,B/180,C/0,D/0,E/90}{
				\draw[fill=black] (\t) circle (.8pt) node[shift={(\g:8pt)}]{$ \t $};
			}
			\draw[<->] (B)--(0.55,0)node[midway,sloped,below]{$210$};
			\draw[<->] (E)--(M)node[midway,sloped,below]{$50$};
			\draw (N)node[yshift=-.75 cm]{Hình b)};
		\end{tikzpicture}
	\end{center}
	\begin{itemize}
		\item Ở hình a), ta có $S_{ABECD}=S_{ABCD}+S_{BEC}=180\cdot 350+\dfrac{1}{2}\cdot 40\cdot 350=70000$ (cm$^2$).\\
		Suy ra $V_1=S_{ABECD}\cdot h_1=70000\cdot 460=32200000$ (cm$^3$)$=32200$ (dm$^3$).
		\item Ở hình b), ta có $$S_{ABECD}=S_{ABCD}+S_{BEC}=\dfrac{1}{2}\cdot210\cdot(260+370)+\dfrac{1}{2}\cdot 50\cdot 260=72650 \ (\text{cm}^2).$$
		Suy ra $V_2=S_{ABECD}\cdot h_2=72650\cdot 430=31239500$ (cm$^3$)$=31239{,}5$ (dm$^3$).
	\end{itemize}
	Vậy $V_1-V_2=960{,}5\approx 961$ (dm$^3$).
}
\end{bt}

\begin{bt}%[1H8H1-3]
% nguon: Archive/SP-4Phan-Dot3/Data/CTST_1-OTGK2_De05.tex #4315

	Cho hình lăng trụ $ABC.MNP$ có $ABC$ là tam giác vuông tại $B$, $\widehat{ACB}=30^\circ$. Tính góc giữa hai đường thẳng $MN$ và  $AC$.
	\loigiai{\immini{
			Ta có $\left(MN, AC\right)=\left(AB, AC\right)=\widehat{BAC}=90^\circ-30^\circ=60^\circ$.}
		{\begin{tikzpicture}[>=stealth,line join=round,line cap=round]
				\draw (0,0)coordinate(A)--(-35:1.5)coordinate(B)--(3,0)coordinate(C)--($(C)+(90:3)$)coordinate(P)--($(B)-(C)+(P)$)coordinate(N)--($(A)-(B)+(N)$)coordinate(M)--(P) (A)--(M) (B)--(N);
				\draw[dashed]  (A)--(C);
				\foreach \diem/\vitrin in {A/left,B/below left,C/right,M/left,N/below left,P/right}	\fill (\diem)circle(1pt)node[\vitrin]{$\diem$};
	\end{tikzpicture}}}
\end{bt}

\begin{bt}%[1H8V2-2]
% nguon: Archive/SP-4Phan-Dot3/Data/CTST_1-OTGK2_De04.tex #4322

Cho tứ diện $ABCD$ có $AD\perp (ABC)$ với $\triangle ABC$ là tam giác đều cạnh bằng $2a$ và $AD=\alpha\sqrt{3}$. Gọi $M$ là trung điểm của $BC$. Điểm $I$ thuộc $AM$ thỏa $\dfrac{AI}{MI}=\dfrac{1}{2}$, $G$ là trọng tâm tam giác $BCD$. Chứng minh rằng $IG\perp (BCD)$. 
	\loigiai{
\begin{center}
	\begin{tikzpicture}
	\def\a{6} %Khai báo cạnh
	\def\h{4}
	\path 	(0:0) coordinate (A)
			++(0:\a) coordinate (B)
			++(-150:3*\a/5) coordinate (C)
			($(A)+(90:\h)$) coordinate (D)
			($(C)!0.5!(B)$) coordinate (M)
			($(A)!1/3!(M)$) coordinate (I)
			($(D)!2/3!(M)$) coordinate (G)
			($(D)!0.5!(M)$) coordinate (H);
	\draw	(A)--(C)--(B)
		(A)--(D)	(B)--(D)	(C)--(D)--(M);
	\draw[dashed] 	(A)--(B) (A)--(M) (I)--(G) (A)--(H);
	\foreach \x /\goc in {A/180,B/0,C/-135,D/90,M/-30,I/-120,G/60,H/60}
			\fill[black] (\x) circle (1.5pt)
			($(\x)+(\goc:3mm)$) node {$\x$};
	\draw pic[draw,angle radius=2mm]{right angle=B--A--D};%Theo chiều dương
\end{tikzpicture}
\end{center}
Ta có $\heva{&AD\perp (ABC)\\&BC\subset (ABC)}\Rightarrow AD\perp BC$.\\
Và $\heva{&\triangle ABC~\text{đều}\\&AM~\text{là đường trung tuyến}}\Rightarrow AM\perp BC$.\\
Khi đó $\heva{&BC\perp AD\\&BC\perp AM\\&AM\cap AD=A}\Rightarrow BC\perp (ADM)$.\\
Trong $(ADM)$, kẻ $AH\perp DM~(1)$.\\
Ta có $\heva{&BC\perp (ADM)\\&AH\subset (ADM)}\Rightarrow AH\perp BC~(2)$.\\
Từ $(1)$ và $(2)$ suy ra $AH\perp (BCD)~(3)$.\\
Vì $\triangle ABC$ đều nên $AM=a\sqrt{3}=AD$. Suy ra tam giác $ADM$ cân tại $M$.\\
Mặt khác, $AH\perp DM$ nên $H$ là trung điểm của $DM$.\\
Vì $G$ là trọng tâm tam giác $BCD$ nên $\dfrac{GM}{AM}=\dfrac{1}{3}\Rightarrow \dfrac{GM}{HM}=\dfrac{2}{3}$.\\
Trong tam giác $AHM$ có $\dfrac{GM}{AM}=\dfrac{IM}{AM}=\dfrac{2}{3}\Rightarrow IG\parallel AH~(4)$.\\
Từ $(3)$ và $(4)$ ta suy ra được $IG\perp (BCD)$.
	}
\end{bt}

\begin{bt}%[1H8V1-3]
% nguon: Archive/SP-4Phan-Dot3/Data/KNTT_1-OTGK2_De05.tex #4331
%%Tự luận 3
	\immini {
		Công ty sản xuất đồ chơi Electric X giao cho nhân viên thiết kế một mô hình khối hình hộp $ABCD.A'B'C'D'$ có $6$ mặt là hình vuông cạnh bằng $2\,024$ (cm) như hình vẽ. Giả sử $M$, $N$ lần lượt là trung điểm của cạnh $AA'$ và $A'B'$. Để ghi các thông số kĩ thuật thì công ty yêu cầu nhân viên tính số đo góc giữa hai đường thẳng $MN$ và $BD$ trước khi sản xuất hàng loạt. Hỏi số đo góc giữa hai đường thẳng $MN$ và $BD$ là bao nhiêu?
	}{
		\begin{tikzpicture}[scale=0.85, font=\footnotesize, line join=round, line cap=round, >=stealth]
			\path
			(0,0) coordinate (A) (3,0) coordinate (D)
			(-150:1.5) coordinate (B)++(3,0) coordinate (C);
			\foreach \p in {A,B,C,D}{\path (\p)++(0,2.5) coordinate (\p');}
			\path
			($(A)!1/2!(A')$) coordinate (M)
			($(A')!1/2!(B')$) coordinate (N)
			;
			\draw (B)--(C)--(C')--(B')--cycle (C)--(D)--(D')--(C') (A')--(B')--(C')--(D')--cycle;
			\draw[dashed] (A)edge(B)edge(D)edge(A') (N)--(M) (B)--(D);
			\foreach \x/\g in {A/45, B/-135, C/-45, D/0, A'/90, B'/180, C'/0, D'/90, M/0, N/180} 
			\fill[black] (\x) circle (1pt)+(\g:.3)node{$\x$};
		\end{tikzpicture}
	}		
	\loigiai{
		\begin{center}
			\begin{tikzpicture}[scale=1, font=\footnotesize, line join=round, line cap=round, >=stealth]
				\path
				(0,0) coordinate (A) (3,0) coordinate (D)
				(-150:1.5) coordinate (B)++(3,0) coordinate (C);
				\foreach \p in {A,B,C,D}{\path (\p)++(0,2.5) coordinate (\p');}
				\path
				($(A)!1/2!(A')$) coordinate (M)
				($(A')!1/2!(B')$) coordinate (N)
				($(A')!1/2!(D')$) coordinate (P);
				\draw (B)--(C)--(C')--(B')--cycle (C)--(D)--(D')--(C') (A')--(B')--(C')--(D')--cycle (B')--(D') (N)--(P);
				\draw[dashed] (A)edge(B)edge(D)edge(A') (N)--(M)--(P) (B)--(D);
				\foreach \x/\g in {A/45, B/-135, C/-45, D/0, A'/90, B'/180, C'/0, D'/90, M/0, N/180, P/90} \fill[black] (\x) circle (1pt)+(\g:.3)node{$\x$};
			\end{tikzpicture}
		\end{center}
		Ta có hình hộp $ABCD.A'B'C'D'$ có $6$ mặt là hình vuông, nên $ABCD.A'B'C'D'$ là hình lập phương cạnh $2\,024$ cm.\\
		Suy ra $B'D'=AB'=AD'=2\,024\sqrt2$.\\
		Ta có $BD\parallel B'D'$ nên $(MN,BD)=(MN,B'D')$.\\
		Gọi $P$ là trung điểm của $A'D'$, khi đó $NP$ là đường trùng bình của $\triangle A'B'D'$. Suy ra $NP\parallel B'D'$ và $NP=\dfrac{1}{2}B'D'=1\,012\sqrt2$ cm.\\
		Nên $(MN,B'D')=(MN,NP)=\widehat{MNP}$.\\
		Ta cũng có $MN$ và $MP$ lần lượt là đường trung bình của $\triangle A'AB'$ và $\triangle A'AD'$.\\
		Nên $MN=\dfrac{1}{2}AB'=1\,012\sqrt2$ cm và $MP=\dfrac{1}{2}AD'=1\,012\sqrt2$ cm.\\
		Suy ra $\triangle MNP$ đều, nên $\widehat{MNP}=60^\circ$.\\
		Vậy $(MN,BD)=60^\circ$.
	}
\end{bt}

\begin{bt}%[1H8V7-5]
% nguon: Archive/SP-4Phan-Dot3/Data/KNTT_1-OTGK2_De05.tex #4332
%%Tự luận 4
	Cho khối chóp $S.ABCD$ có đáy $ABCD$ là hình vuông cạnh bằng $2a$, $SA\perp \left(ABCD\right)$, $SA=\dfrac{2a}{\sqrt{3}}$.
	\begin{enumerate}
		\item Chứng minh tam giác $SCD$ vuông?
		\item Tính khoảng cách từ $A$ đến $(SCD)$.
		\item Tính thể tích của khối tứ diện $AHCD$?
	\end{enumerate}
	\loigiai{
		\begin{center}
			{\begin{tikzpicture}[scale=0.8, font=\footnotesize,line join=round, line cap=round, >=stealth]
					\def\a{4}
					\def\h{4}
					\path 	(0:0) coordinate (A)
					++(0:\a) coordinate (D)
					++(-130:\a/2) coordinate (C)
					($(A)+(C)-(D)$) coordinate (B)
					($(A)+(90:\h)$) coordinate (S)
					($(S)!(A)!(D)$) coordinate (H);
					\draw[dashed] 	(B)--(A)--(D)	(A)--(S) (C)--(A)--(H);
					\draw	(B)-- (C)--(D) (C)--(H)
					(B)--(S)	(C)--(S)	(D)--(S);
					\foreach \x/\g in {A/135,B/-135,C/-45,D/45,S/90,H/45}
					\fill[black] 	(\x) circle (1.5pt)
					($(\g:3mm)+(\x)$) node {$\x$};	
					\draw pic[draw,angle radius=2mm]{right angle=D--A--S}
					pic[draw,angle radius=2mm]{right angle=A--H--S};	
			\end{tikzpicture}}
		\end{center}
		\begin{enumerate}
			\item Ta có $\heva {&CD\perp AD\text { ($ABCD$ là hình vuông)}\\&CD\perp SA\  (SA\perp (ABCD)}\Rightarrow CD\perp (SCD)$.\\
			$\Rightarrow CD\perp SD$ hay $\triangle SCD$ vuông tại $D$.
			\item Gọi $AH$ là đường cao của tam giác $SAD$.\\
			Ta có $\heva{&AH\perp SD\\ &AH\perp CD}\Rightarrow AH\perp (SCD)$ nên $AH$ là khoảng cách từ $A$ đến $(SCD)$.\\
			Trong tam giác vuông $SAD$, ta có \[\dfrac{1}{AH^2}=\dfrac{1}{SA^2}+\dfrac{1}{AD^2}=\dfrac{3}{4a^2}+\dfrac{1}{4a^2}=\dfrac{1}{a^2}.\]
			Suy ra $AH=a$. 
			\item Ta có $HD=\sqrt{AD^2-AH^2}=\sqrt{4a^2-a^2}=a\sqrt{3}$.
			Từ đó ta có
			\[V_{C.AHD}=\dfrac{1}{3} S_{ADH} \cdot CD=\dfrac{1}{3} \cdot \dfrac{AH\cdot HD}{2} \cdot CD=\dfrac{1}{3} \dfrac{a\sqrt{3} \cdot a}{2} \cdot 2a=\dfrac{a^3 \sqrt{3}}{3}.\]
		\end{enumerate}		
	}
\end{bt}

\begin{bt}%[1H8H7-4]
% nguon: Archive/SP-4Phan-Dot3/Data/CD_1_OTHK2_De05.tex #4340

	Cho hình lăng trụ đứng $ABC.A'B'C'$ có đáy $ABC$ là tam giác vuông cân tại $A$, $BC=a\sqrt{2}$, $A'B$ tạo với đáy một góc bằng $60^{\circ}$. Tính thể tích của khối lăng trụ $ABC.A'B'C'$.
	\loigiai{\immini{Ta có $A'A\perp(ABC)$ nên $AB$ là hình chiếu của $A'B$ lên đáy $(ABC)$.\\
			Suy ra $(A'B,(ABC))=(A'B,AB)=\widehat{A'BA}=60^\circ$.\\
			Tam giác $ABC$ vuông cân tại $A$ và có $BC=a\sqrt{2}$ nên $AB=AC=a$.\\
			Tam giác $A'BA$ có $\tan60^\circ=\dfrac{A'A}{AB}\Rightarrow A'A=AB\tan60^\circ=a\sqrt{3}$.\\
			Thể tích của khối lăng trụ $ABC.A'B'C'$ là 
			$$V=A'A\cdot S_{ABC}=A'A\cdot\dfrac{1}{2}\cdot AB\cdot AC=\dfrac{a^3\sqrt{3}}{2}.$$
		}
		{\begin{tikzpicture}[>=stealth,line join=round,line cap=round,font=\footnotesize,scale=0.5]
				\path (0,0) coordinate (A)
				(2,-1.4) coordinate (B)
				(5,0) coordinate (C)
				($(A)+(0,5)$)coordinate (A')
				($(B)+(0,5)$)coordinate (B')
				($(C)+(0,5)$)coordinate (C')
				;
				\draw (B')--(B)--(A')--(A)--(B)--(C)--(C')--(A')--(B')--(C')--(A');
				\draw[dashed] (A)--(C);
				\foreach \p / \r in {A/135,B/-90,A'/135,B'/90,C/45,C'/45}
				\fill (\p) circle (1.2pt) node[shift={(\r:3mm)}]{$\p$};
	\end{tikzpicture}}}
\end{bt}

\begin{bt}%[1H8H5-4]
% nguon: Archive/SP-4Phan-Dot3/Data/CD_1_OTHK2_De06.tex #4356
%[Làm đề GHK, HK, 4 phần - 2025]%[Huỳnh Xuân Tín]
	Cho hình chóp tứ giác đều $S.ABCD$ có tất cả các cạnh đều bằng $a$ và có $O$ là giao điểm hai đường chéo của đáy. Tính khoảng cách giữa hai đường thẳng $AC$ và $SB$.
	\loigiai{
		\begin{center}
			\begin{tikzpicture}[=>stealth,line join=round,line cap=round, font=\footnotesize, scale=.8]
				\def\a{5}
				\def\goc{210}
				\def\b{3}
				\def\h{4.5}
				\path
				(0,0)coordinate (A)++(0:\a)coordinate (B)++(\goc:\b)coordinate (C)++(180:\a)coordinate(D)
				($(A)!.5!(C)$)coordinate (O)
				(O)++(90:\h)coordinate (S);
				\path 
				($(B)!(O)!(S)$)coordinate (H);
				\draw (S)--(B)--(C)--(D)--cycle (S)--(C); %(M)--(S)--(C);
				\draw[dashed] (O)--(S)--(A)--(C) (B)--(A)--(D)--(B) (O)--(H);
				\draw pic[draw,angle radius=2mm] {right angle = S--O--B};
				\draw pic[draw,angle radius=2mm] {right angle = O--H--B};
				\foreach \x/\g in{A/160,B/0,C/-90,S/90,O/-90,D/-90,H/60}
				\fill[black](\x)circle(1pt) ($(\x)+(\g:3mm)$)node{$\x$};
			\end{tikzpicture}
		\end{center}
		Kẻ $OH \perp SB$ tại $H$.\\
		Ta có $\heva{&AC\perp BD\\&AC \perp SO}\Rightarrow AC \perp (SBD)\Rightarrow AC \perp OH$. Khi đó $\mathrm{d}(AC,SB)=OH$.\\
		Có $\triangle SAC$ vuông tại $S$ suy ra $SO=\dfrac{1}{2}AC=\dfrac{a\sqrt{2}}{2}$.\\
		Khi đó $OH=\dfrac{SO\cdot OB}{SB}=\dfrac{\dfrac{a\sqrt{2}}{2}\cdot \dfrac{a\sqrt{2}}{2}}{a}=\dfrac{a}{2}$.\\
		Vậy $\mathrm{d}(AC,SB)=\dfrac{a}{2}$.				
	}
\end{bt}

\begin{bt}%[1H8C7-9]
% nguon: Archive/SP-4Phan-Dot3/Data/CD_1_OTHK2_De06.tex #4357
%[Làm đề GHK, HK, 4 phần - 2025]%[Huỳnh Xuân Tín]
	Người ta muốn xây một bể bơi hình hộp chữ nhật có thể tích là $125$ m$^3$. Đáy bể bơi là hình chữ nhật có chiều dài gấp ba lần chiều rộng. Tính chiều rộng của đáy bể bơi để khi thi công tiết kiệm được nguyên liệu nhất.
	\loigiai{
		\begin{itemize}
			\item Gọi chiều rộng đáy là $a$, chiều rộng là $b$, chiều dài là $3a$. Theo bài ra ta có $3a^2\cdot b =125\Rightarrow b=\dfrac{125}{3a^2}$.
			\item Diện tích toàn phần: $2(3a+a)\cdot b+a\cdot 3a=8a\cdot \dfrac{125}{3a^2}+3a^2=\dfrac{125\cdot 8}{3a}+3a^2$ (chỉ có 1 mặt đáy).
			\item Ta có: $\dfrac{125\cdot 8}{3a}+3a^2=\dfrac{125\cdot 4}{3a}+\dfrac{125\cdot 4}{3a}+3a^2\ge 3 \sqrt[3]{\dfrac{125\cdot 4}{3a}\cdot \dfrac{125\cdot 4}{3a}\cdot 3a^2}=10\sqrt[3]{12}$.\\ Dấu \lq\lq$=$\rq\rq \ xảy ra khi: $\dfrac{125\cdot 4}{3a}=3a^2\Rightarrow a=5\cdot \sqrt[3]{\dfrac{4}{9}}$.
		\end{itemize}	
	}
\end{bt}

\begin{bt}%[1H8V5-3]
% nguon: Archive/SP-4Phan-Dot3/Data/CD_1_OTHK2_De02.tex #4365

	Cho hình chóp $S.ABCD$ có đáy là hình vuông cạnh $1$, hai mặt phẳng $(SAB)$ và $(SBC)$ cùng vuông góc với mặt phẳng $(ABCD)$ và $SC=\sqrt 5$. Tính khoảng cách từ $D$ đến mặt phẳng $(SAC)$.
	\loigiai{
		\begin{center}
			\begin{tikzpicture}[line join=round, line cap=round,thick, scale=0.8]
				\coordinate (B) at (0,0);
				\coordinate (A) at (-2,-3);
				\coordinate (C) at (7,0);
				\coordinate (D) at ($(A)+(C)-(B)$);
				\coordinate (S) at ($(B)+(0,5)$);
				\coordinate (O) at ($(B)!0.5!(D)$);
				\coordinate (H) at ($(S)!(B)!(O)$);
				\draw(S)--(A) (S)--(D) (S)--(C) (A)--(D)--(C);
				\draw[dashed](S)--(B) (A)--(B) (B)--(C) (S)--(O) (B)--(H) (B)--(D) (A)--(C);
				\pic[draw,thin,angle radius=3mm] {right angle = S--B--D} pic[draw,thin,angle radius=3mm] {right angle = S--B--A} pic[draw,thin,angle radius=3mm] {right angle = O--H--B};
				\foreach \i/\g in {S/90,B/-90,A/-90,D/-90,C/0,O/30,H/30}{\draw[fill=black](\i) circle (1.0pt) ($(\i)+(\g:3mm)$) node[scale=1]{$\i$};}
			\end{tikzpicture}
		\end{center}
		Kẻ $BH\perp SO$ tại $H$.\\
		Ta có $\heva{&AC \perp SB\\&AC\perp BD} \Rightarrow AC\perp (SBD)\Rightarrow AC\perp BH$.\\
		Ta lại có $BH\perp SO\Rightarrow BH\perp (SAC)\Rightarrow \mathrm{d}(B,(SAC))=BH$.\\
		Ta có $SB=\sqrt{SC^2-BC^2}=\sqrt{(\sqrt 5)^2-1^2}=2$.\\
		Ta có $BH=\dfrac{1}{\sqrt{\dfrac{1}{SB^2}+\dfrac{1}{OB^2}}}$.\\ 
		Vậy $\mathrm{d}(B,(SAC))=\dfrac{2}{3}$.
		Ta có $DB$ cắt $(SAC)$ tại $O$.\\
		Suy ra, $\dfrac{\mathrm{d}(D,(SAC))}{\mathrm{d}(B,(SAC))}=\dfrac{DO}{BO}=1\Rightarrow \mathrm{d}(D,(SAC))=\mathrm{d}(B,(SAC))=\dfrac{2}{3}\approx 0{,}67$.	
	}
\end{bt}

\begin{bt}%[1H8V5-6]
% nguon: Archive/SP-4Phan-Dot3/Data/KNTT_1-OTGK2_De02.tex #4375

	Kim tự tháp ở Ai Cập có hình dạng là một hình chóp tứ giác đều có đáy là $262$ m và cạnh bên là $230$ m. Giả sử, từ một mặt bên của kim tự tháp ta cần đào một con đường ngắn nhất để đi đến tâm của đáy kim tự tháp, khi đó quãng đường ngắn nhất có độ dài khoảng bao nhiêu mét?
	\loigiai{
		\immini
		{Khoảng cách ngắn nhất từ một điểm tại mặt bên đến tâm $O$ của đáy chính là khoảng cách từ tâm $O$ đến mặt phẳng $(SCD)$.\\
		Dựng $OI\perp CD$ tại $I$, $OH\perp SI$ tại $H$.\\
		Ta có $\heva{&CD\perp OI\\&CD\perp SO} \Rightarrow CD\perp (SOI) \Rightarrow OI\perp OH$.\\
		Ta có $\heva{& OH\perp SI\\& OH\perp CD} \Rightarrow OH\perp (SCD)$.}
		{\begin{tikzpicture}[declare function={a=2.65;b=4.5;h=4;}, line join=round, line cap=round, >=stealth,scale=0.9]
			\path (0,0) coordinate (A)
			(-135:a) coordinate (B)
			(b,0) coordinate (D)
			($(D)-(A)+(B)$) coordinate (C)
			($(B)!0.5!(D)$) coordinate (O)
			($(O)+(0,h) $) coordinate (S)
			($(C)!.5!(D)$) coordinate (I)
			($(S)!0.65!(I)$) coordinate (H);		
			\path ($(D)-(A)+(B)$) coordinate (C);
			\draw[dashed] (S)--(A)--(B) (A)--(D) (A)--(C) (S)--(O) (B)--(D) (O)--(I) (O)--(H) ;
			\draw (B)--(C)--(D) (B)--(S) (D)--(S)--(C) (S)--(I);
			\draw pic[draw=black,angle radius=0.2cm] {right angle = O--I--C};
			\draw pic[draw=black,angle radius=0.2cm] {right angle = S--H--O};
			\foreach \t/\g in {A/150,B/-90,C/-90,D/0,S/90,O/-90,I/-20,H/10}{\draw[fill=black] (\t) circle (0.75pt) node[shift={(\g:7pt)}]{$\t$};}
		\end{tikzpicture}}
		\noindent
		Vậy $OH = \mathrm{d}\left(O,(SCD)\right)$.\\
		Ta có $OC=\dfrac{262\sqrt{2}}{2} \Rightarrow SO=\sqrt{SC^2-OC^2}=\sqrt{18578}$.\\
		Ta có $OH=\dfrac{OI\cdot SO}{\sqrt{OI^2+SO^2}} = \dfrac{131\cdot\sqrt{18578}}{\sqrt{131^2+18578}}$ (mét).
	}
\end{bt}

\begin{bt}%[1H8V7-9]
% nguon: Archive/SP-4Phan-Dot3/Data/KNTT_1-OTGK2_De02.tex #4376

	Từ một tấm bìa hình vuông, người ta cắt ở bốn góc của tấm bìa đó bốn hình vuông bằng nhau, mỗi hình vuông có cạnh bằng $6$ cm, rồi gập tấm bìa lại để được một chiếc hộp không nắp có dạng hình hộp chữ nhật.
	\begin{center}
		\begin{tikzpicture}[scale=0.7,x=1cm,y=1cm,line join = round, line cap = round]
			\coordinate[label=left:$ $] (D) at (0,0);
			\coordinate [label=left:$ $] (A) at (0,6);
			\coordinate [label=right:$ $] (B) at (6,6);
			\coordinate [label=right:$ $] (C) at (6,0);
			\coordinate [label=right:$ $] (E) at (1,0);
			\coordinate [label=right:$ $] (F) at (5,0);
			\coordinate [label=right:$ $] (G) at (1,6);
			\coordinate [label=right:$ $] (H) at (5,6);
			\coordinate [label=right:$ $] (I) at (0,5);
			\coordinate [label=right:$ $] (J) at (0,1);
			\coordinate [label=right:$ $] (M) at (6,1);
			\coordinate [label=right:$ $] (N) at (6,5);
			\draw(A)--(B)--(C)--(D)--(A);
			\draw[dashed] (G)--(E)(H)--(F)(I)--(N)(J)--(M);
		\end{tikzpicture}\qquad
		\begin{tikzpicture}[scale=0.7,x=1cm,y=1cm,line join = round, line cap = round]
			\coordinate[label=left:$ $] (D) at (0,0);
			\coordinate [label=left:$ $] (A) at (3,1);
			\coordinate [label=right:$ $] (B) at (6,0);
			\coordinate [label=right:$ $] (C) at (9,1);
			\coordinate [label=right:$ $] (E) at (0,2);
			\coordinate [label=right:$ $] (F) at (6,2);
			\coordinate [label=right:$ $] (G) at (3,3);
			\coordinate [label=right:$ $] (H) at (9,3);
			\draw(D)--(E)--(G)--(H)--(F)--(E)(D)--(B)--(C)(F)--(B)(C)--(H);
			\draw[dashed] (D)--(A)--(C)(A)--(G);
		\end{tikzpicture}
	\end{center}
	Tính cạnh của tấm bìa ban đầu, biết rằng thể tích của chiếc hộp bằng $864$ cm$^3$.
	\loigiai{
		Giả sử tấm bìa ban đầu có cạnh $x$ (cm), $(x>12)$. Khi đó đáy của chiếc hộp là hình vuông cạnh $x-12$ (cm) nên diện tích của đáy chiếc hộp là $(x-12)^2$ cm$^2$.\\
		Mà chiều cao của chiếc hộp là $6$ cm, suy ra thể tích của chiếc hộp bằng $(x-12)^2\cdot 6~\left(\mathrm{cm}^3\right)$.\\
		Theo đề bài, thể tích của chiếc hộp bằng $864$ cm$^3$ nên
		$$(x-12)^2\cdot 6=864 \Leftrightarrow(x-12)^2=144.$$
		Với $x>12$ ta có $x-12=12 \Leftrightarrow x=24$ (cm).\\
		Vậy độ dài cạnh của tấm bìa ban đầu là $24$ cm.
	}
\end{bt}

\begin{bt}%[1H8V2-2]
% nguon: Archive/SP-4Phan-Dot3/Data/CTST_1-OTGK2_De03.tex #4386
  % Câu 4-TL.
	Cho hình chóp $S.ABCD$ có đáy là hình vuông và $SA \perp (ABCD)$. Gọi $H, K$ lần lượt là hình chiếu vuông góc của $A$ lên $SB$ và $SD$. Chứng minh rằng $SC \perp (AHK)$.
	\loigiai{
		\begin{center}
			\begin{tikzpicture}[scale=1.2, thick, font=\small]
				% Định nghĩa tọa độ các đỉnh hình bình hành ABCD
				\coordinate (A) at (0,0);
				\coordinate (B) at (4,0);
				\coordinate (D) at (1,2);
				\coordinate (C) at ($(B)+(D)-(A)$); % C = B + D - A
				\coordinate (S) at (0,5); % Đỉnh chóp phía trên
				
				% Vẽ các cạnh đáy (hình bình hành)
				\draw[thick] (A) -- (B) -- (C);
				\draw[dashed] (C) -- (D) -- (A); % Cạnh CD và DA khuất
				
				% Vẽ các cạnh bên
				\draw[thick] (S) -- (A) (S) -- (B) (S) -- (C);
				\draw[dashed] (S) -- (D); % Cạnh SD khuất
				
				% Tính trung điểm H (SB) và K (SD)
				\coordinate (H) at ($(S)!0.5!(B)$); % Trung điểm SB
				\coordinate (K) at ($(S)!0.5!(D)$); % Trung điểm SD
				
				% Vẽ các đoạn AH, AK, HK và SC
				\draw[red, thick] (A) -- (H) ;
				\draw[red, dashed] (A) -- (K) ;
				\draw[red, dashed] (H) -- (K);
				% Đặt nhãn các điểm
				\foreach \p/\pos in {A/below left, B/below right, C/above right, D/above left, S/above}
				\node[\pos] at (\p) {$\p$};
				\node[above right=1pt] at (H) {H};
				\node[above right=1pt] at (K) {K};
				
				\foreach \x/\y/\z in {S/K/A,S/H/A}{
					\path (\y) pic[draw=blue,angle radius=5pt,pic text=.,fill=blue!25]{right angle = \x--\y--\z};
				}
			\end{tikzpicture}
		\end{center}
		Ta có $BC\perp AB$.
		$BC\perp SA$ (Vì $SA\perp(ABCD)\supset BC$)
		Suy ra $BC\perp(SAB)$. Mà $AH\subset(SAB)$. Do đó: $AH\perp BC$.
		Mặt khác $AH\perp SB\Rightarrow AH\perp(SBC)\Rightarrow SC\perp AH$.
		Chứng minh tương tự ta cũng có: $SC\perp AK$.
		Vậy $SC\perp(AHK)$.}
	\end{bt}

\begin{bt}%[1H8H6-2]
% nguon: Archive/SP-4Phan-Dot3/Data/CTST_1-OTGK2_De02.tex #4395

	Cho hình lăng trụ đứng $ABC.A'B'C'$ có đáy là tam giác vuông cân tại $B$, $AC=2a$ và $A'B=3a$. Tính góc phẳng nhị diện $\left[B', AC,B\right]$ (\textit{theo đơn vị độ và làm tròn đến hàng phần mười}).
	\loigiai{
		\immini{
			Gọi $I$ là trung điểm của $AC$.\\
			Vì tam giác $ABC$ vuông cân tại $B$ nên $BI \perp AC$. \quad(1)\\
			Mà $BB'\perp AC$ nên $AC \perp (BB'I)$.\\
			Suy ra $B'I \perp AC$.\quad(2)\\
			Từ (1) và (2) suy ra góc phẳng nhị diện $\left[B', AC,B\right]$ là góc $\widehat{B'IB}$.\\
			Vì tam giác vuông cân tại $B$ có $AC=2a$ nên $AB=a\sqrt{2}$ và $BI=\dfrac{AC}{2}=a$.\\ 
			Lại có $B'B=AA'=\sqrt{{A'B}^2-{AB}^2}=\sqrt{(3a)^2-\left(a\sqrt{2}\right)^2}=\sqrt{7} a$.\\
			Xét $\triangle BB'I$ vuông tại $B$, ta có 
			$$\tan \widehat{B'IB}=\dfrac{B'B}{BI}=\dfrac{\sqrt{7} a}{a}=\sqrt{7} \Rightarrow \widehat{B'IB}\approx 69{,}3^{\circ}.$$
		}
		{
			\begin{tikzpicture}[scale=1, font=\footnotesize,line join=round, line cap=round, >=stealth]
				\path 
				(0,0) coordinate (A)
				++(-40:2.5) coordinate (B)
				(3,0) coordinate (C)
				($(A)!1/2!(C)$) coordinate (I)
				;
				\foreach \i in{A,B,C}{
					\path ($(\i)+(0,3.5)$) coordinate (\i');
					\draw (\i)--(\i');};
				\draw (A)--(B)--(C) (A')--(B')--(C')--cycle (A')--(B);
				\draw[dashed] (A)--(C) (B)--(I)--(B');
				\pic[draw,angle eccentricity=1.8,angle radius=2mm]{right angle=A'--A--C}
				pic[draw,angle eccentricity=1.8,angle radius=2mm]{right angle=C--B--A}
				pic[draw,angle eccentricity=1.8,angle radius=2mm]{angle=B--I--B'}
				;
				\foreach \i/\g in {A/-120,B/-90,C/-90,A'/90,B'/80,C'/90,I/90}{\draw[fill=black](\i) circle (1pt) ($(\i)+(\g:4mm)$) node[scale=1]{$\i$};}
			\end{tikzpicture}
		}
	}
\end{bt}

\begin{bt}%[1H8V4-3]
% nguon: Archive/SP-4Phan-Dot3/Data/KNTT_1-OTGK2_De03.tex #4402

Chứng minh rằng 
$\left(1+\mathrm{e}^{2x}\right)=2x+\ln\left(1+\mathrm{e}^{-2x}\right)$.
	\loigiai{
Ta có $\ln \left(1+\mathrm{e}^{2x}\right)=\ln\left[\mathrm{e}^{2x}\left(1+\mathrm{e}^{-2x}\right)\right]=\ln\mathrm{e}^{2x}+\ln\left(1+\mathrm{e}^{-2x}\right)=2x+\ln\left(1+\mathrm{e}^{-2x}\right)$ (đpcm).
	}
\end{bt}

\begin{bt}%[1H8H7-7]
% nguon: Archive/SP-4Phan-Dot3/Data/KNTT_1-OTGK2_De03.tex #4404

	\immini{
		Cho hình hộp chữ nhật $ABCD.A'B'C'D'$ có thể tích $V$. Gọi $M$, $N$, $P$, $Q$ lần lượt là trung điểm các cạnh $AA'$, $BB'$, $CC'$, $DD'$.
	}{
		\begin{tikzpicture}[scale=0.6, font=\footnotesize,line join=round, line cap=round, >=stealth]
			\path 
			(0,0) coordinate (A)
			++(-150:2) coordinate (B)
			++(0:6) coordinate (C)
			($(A)+(C)-(B)$) coordinate (D)
			($(A)!1/2!(C)$) coordinate (O)
			($(A)!1/2!(A')$) coordinate (M)
			($(A')!1/2!(B)$) coordinate (N)
			($(A')!1/2!(C)$) coordinate (P)
			($(A')!1/2!(D)$) coordinate (Q)
			;
			\foreach \i in {A,B,C,D}{
				\coordinate (\i') at ($(\i)+(0,6)$);
			}
			\draw (A')--(B')--(C')--(D')--cycle;
			\draw (B)--(B') (C)--(C') (D)--(D')  (B)--(C)--(D);
			\draw[dashed,thin](B)--(A)--(A') (A)--(D) (A')--(B) (A')--(D) (A')--(C) (M)--(N)--(P)--(Q)--(M);
			\pic[draw,angle eccentricity=1.8,angle radius=2mm]{right angle=A'--A--D};
			\foreach \i/\g in {A'/90, B'/90, C'/-30, D'/90, A/150, B/-90, C/-90, D/-90, M/30, N/180, P/-90, Q/0}
			\fill[black] (\i) circle(1pt)+(\g:5mm)node[scale=1]{$\i$};
		\end{tikzpicture}
	}
Tính thể tích khối chóp cụt $MNPQ.ABCD$.
	\loigiai{
\begin{align*}
				V_{MNPQ.ABCD} &= \dfrac{1}{3} \left(S_{MNPQ} + S_{ABCD} + \sqrt{S_{MNPQ} \cdot S_{ABCD}}\right)\cdot A'A \\
				&=\dfrac{1}{3} \left(\dfrac{1}{4}S_{ABCD} + S_{ABCD} + \sqrt{\dfrac{1}{4}S_{ABCD}^2}\right)\cdot A'A \\
				&= \dfrac{1}{3} \cdot \dfrac{7}{4} S_{ABCD}\cdot A'A \\
				&=  \dfrac{7}{12} V.
			\end{align*}
}	
\end{bt}

\begin{bt}%[1H8V4-7]
% nguon: Archive/SP-4Phan-Dot3/Data/KNTT_1-OTGK2_De03.tex #4405

	Cho hình chóp tứ giác đều $S.ABCD$ có cạnh đáy bằng $2a$, cạnh bên bằng $a\sqrt{5}$. Tính cosin góc giữa $(SCD)$ với $(SBC)$.
\loigiai{
\immini{Ta có $(SCD) \cap (SBC) =SC$. Kẻ $DK \perp SC$ tại $K$. Ta có
\begin{itemize}
	\item [$\bullet$] $SC \perp DK \quad (1)$.
	\item [$\bullet$] $BD \perp (SAC) \Rightarrow SC \perp BD \quad (2)$
\end{itemize}
Từ (1) và (2) suy ra $SC \perp (BDK)$ nên $BK \perp SC$.\\
Do $DK \perp SC$ và $BK \perp SC$ nên góc giữa $(SCD)$ với $(SBC)$ bằng góc giữa $KD$ và $KB$.\\
Xét $\triangle KBD$. Ta có
\begin{itemize}
	\item [$\bullet$] $BD=2\sqrt{2}a$; $KB=KD=\dfrac{a\sqrt{13}}{2}$.
	\item [$\bullet$] $\cos (KB,KD)=\dfrac{\bigg|KB^2+KD^2-BD^2\bigg|}{2\cdot KB \cdot KD}=\dfrac{3}{13}.$
\end{itemize}}
{
\begin{tikzpicture}[line join=round, line cap=round,thick, scale=0.7]
	\coordinate (A) at (0,0);
	\coordinate (B) at (2,-2);
	\coordinate (D) at (5,0);
	\coordinate (C) at ($(B)+(D)-(A)$);
	\coordinate (O) at ($(A)!0.5!(C)$);
	\coordinate (S) at ($(O)+(0,5.5)$);
	\coordinate (K) at ($(S)!.5!(C)$);
	\draw(S)--(A) (S)--(B)--(K) (S)--(C) (A)--(B) (B)--(C);
	\draw[dashed,thin](A)--(C) (A)--(D) (C)--(D) (S)--(D) (S)--(O) (B)--(D)--(K);
	\pic[draw,thin,angle radius=2mm] {right angle = B--K--C} pic[draw,thin,angle radius=2mm] {right angle = D--K--C};
	\foreach \i/\g in {S/90,A/180,B/-90,C/-90,D/0,O/-90,K/0}{\draw[fill=white](\i) circle (1.5pt) ($(\i)+(\g:3mm)$) node[scale=.7]{$\i$};}
\end{tikzpicture}
}
Vậy góc giữa cosin $(SCD)$ với $(SBC)$ bằng $\dfrac{3}{13}$.
}
\end{bt}

\begin{bt}%[1H8V4-4]
% nguon: Archive/SP-4Phan-Dot3/Data/CD_1_OTHK2_De03.tex #4413

	Cho hình chóp $S.ABCD$ có đáy là hình vuông cạnh $a$, $SA \perp (ABCD)$ và \break $SA=a\sqrt{2}$. Điểm $M$ thuộc cạnh $AB$ và $AM=x$, với $0<x<a$. Mặt phẳng $(\alpha)$ qua $M$ và vuông góc với $SB$ tại $N$ cắt $SC$, $SD$ tại $P$ và $Q$. Tính diện tích của thiết diện $MNPQ$.
	\loigiai
		{\begin{itemize}
				\item Xét hình tính của $MNPQ$.
				\immini{
				Dựng $AH\perp SB$ tại $H$. \quad $(1)$\\
				Khi đó $\heva{& AH\perp SB \\ & AD\perp SB}$ mà $(\alpha)\perp SB\Rightarrow\heva{& (\alpha)\parallel AH \\ & (\alpha)\parallel AD.}$\\
				Ta có $\heva{& MN=(\alpha)\cap (SAB) \\ & (\alpha)\parallel AH\subset (SAB)}\Rightarrow MN\parallel AH$.\\
				Tương tự, 
					\begin{itemize}
						\item $\heva{& MQ=(\alpha)\cap (ABCD) \\ & (\alpha)\parallel AD\subset (ABCD)}\Rightarrow MQ\parallel AD$,
						\item  $\heva{& NP=(\alpha)\cap (SBC) \\ & (\alpha)\parallel BC\subset (SBC)}\Rightarrow NP\parallel BC$.
				\end{itemize}  }
			{\begin{tikzpicture}[scale=1, font=\footnotesize, line join=round, line cap=round,>=stealth]
						\path (0,0)coordinate (B) (3,0) coordinate (C) (4.5,1.3) coordinate (D) ($(B)+(D)-(C)$) coordinate (A) (0,2.5) coordinate (h) ($(A)+(h)$) coordinate (S) ($(B)!0.4!(A)$) coordinate (M) ($(C)!0.4!(D)$) coordinate (Q) ($(B)!0.25!(S)$) coordinate (N) ($(B)!5/2!(N)$) coordinate (H) ($(N)+(Q)-(M)$) coordinate(t) (intersection of N--t and S--C) coordinate (P);
						\draw (S)--(B)--(C)--(D)--cycle (S)--(C) (N)--(P)--(Q);
						\draw[dashed] (S)--(A)--(B) (H)--(A)--(D) (Q)--(M)--(N);
						\fill[opacity=0.2, color=gray] (M)--(N)--(P)--(Q)--cycle;
						\foreach \p/\g in {A/45,B/240,C/-60, D/0, S/90,M/-90,N/150,P/30,Q/-30,H/150} \fill[black] (\p) circle(1pt)+(\g:0.3) node{$\p$};
						\draw pic[draw, angle radius=0.2cm]{right angle= A--H--B} pic[draw, angle radius=0.2cm]{right angle= M--N--B};
				\end{tikzpicture}}\noindent
				Do đó, thiết diện cần tìm là hình thang $MNPQ$ (vì $MQ\parallel AD\parallel BC\parallel NP$).\\
				Ta có $AH\perp SB$, $AH\perp BC$ và $SB\cap BC=B$.\\
				Suy ra $AH\perp (SBC)\Rightarrow AH\perp NP$ (vì $NP\subset (SBC)$).\\
				Mặt khác, $MN\parallel AH$ nên $MN\perp NP$.\\
				Vậy $MNPQ$ là hình thang vuông tại $M$ và $N$.
				\item Tính diện tích của thiết diện.\\
				Tam giác $SAB$ vuông tại $A$ nên 
				\begin{itemize}
				\item $\dfrac{1}{AH^2}=\dfrac{1}{SA^2}+\dfrac{1}{AB^2}=\dfrac{1}{2a^2}+\dfrac{1}{a^2} \Rightarrow AH=\dfrac{a\sqrt{6}}{3}$ và $SB=a\sqrt{3}$,
				\item $AB^2=BH\cdot BS \Rightarrow BH=\dfrac{a^2}{a\sqrt{3}}=\dfrac{a}{\sqrt{3}}=\dfrac{a\sqrt{3}}{3}$.
				\end{itemize} 
				Ta có $MN\|AH \Rightarrow \dfrac{MN}{AH}=\dfrac{BM}{BA}=\dfrac{BN}{BH} \Leftrightarrow \dfrac{MN}{\dfrac{a\sqrt{6}}{3}}=\dfrac{a-x}{a}=\dfrac{BN}{\dfrac{a}{\sqrt{3}}}$.\\
				Suy ra $MN=\dfrac{(a-x)\sqrt{6}}{3}$ và $BN=\dfrac{a-x}{\sqrt{3}}$.\\
				Vì $NP\|BC \Rightarrow \dfrac{SN}{SB}=\dfrac{NP}{BC} \Leftrightarrow \dfrac{SB-BN}{SB}=\dfrac{NP}{BC}$\\
				$ \Rightarrow NP=\left(\dfrac{a\sqrt{3}-\dfrac{a-x}{\sqrt{3}}}{a\sqrt{3}}\right)\cdot a=\dfrac{2a+x}{3}$.\\
				Vậy $S_{MNPQ}=\dfrac{(MQ+NP)\cdot MN}{2}=\dfrac{\left(a+\dfrac{2a+x}{3}\right)\dfrac{a-x}{\sqrt{3}}}{2}=\dfrac{(5a+x)(a-x)}{6\sqrt{3}}$.
			\end{itemize}
		}	
	\end{bt}

\begin{bt}%[1H8V7-6]
% nguon: Archive/SP-4Phan-Dot3/Data/CD_1_OTHK2_De03.tex #4414

	Cho hình chóp $S.ABC$ có đáy $ABC$ là tam giác đều cạnh $a$. Hình chiếu vuông góc của $S$ trên đáy là điểm $H$ trên cạnh $AC$ sao cho $AH=\dfrac{2}{3}AC$; mặt phẳng $\left( SBC \right)$ tạo với đáy một góc $60^\circ$. Tính thể tích khối chóp $S.ABC$.
	\loigiai{
		\begin{center}
			\begin{tikzpicture}
				\def\a{4}
				\path 	(0:0) coordinate (A)
				(0:\a) coordinate (B)
				(-50:.5*\a) coordinate (C)
				($(B)!.5!(C)$)coordinate (M)
				($(C)!1/3!(M)$)coordinate (N)
				($(A)!2/3!(C)$)coordinate (H)
				($(H)+(90:\a)$) coordinate (S);
				\draw[dashed] 	(B)--(A) (A)--(M) (H)--(N);
				\draw[thick] 	(A)--(C)--(B) (S)--(A) (S)--(N)
				(S)--(B) (S)--(C) (S)--(H);
				\foreach \x/\g in {A/180,B/0,C/-90,S/90,M/-45,N/-45,H/-120}
				\fill[black] 	(\x) circle (1pt)
				($(\g:3mm)+(\x)$) node {$\x$};
				\draw pic[draw,angle radius=3mm]{right angle=A--M--C}
				pic[draw,angle radius=2mm]{right angle=H--N--C}
				pic[draw,angle radius=3mm]{right angle=S--H--A};
			\end{tikzpicture}
		\end{center}
		\begin{itemize}
			\item Gọi $M$ là trung điểm của $BC$,
			$N\in CM$ sao cho $\dfrac{CN}{CM}=\dfrac{CH}{CA}=\dfrac{1}{3} \Rightarrow HN \parallel AM$.
			\item Mà
			$\triangle ABC$ đều nên $AM\perp BC\Rightarrow HN\perp BC\Rightarrow BC\perp \left( SHN \right)$.\\ 
			Do đó $\left(\left( SBC \right),\left( ABC \right)\right)=\left(SN,HN\right)=\widehat{SNH}=60^\circ$.
			\item Do $\triangle ABC$ đều nên $AM=\dfrac{a\sqrt{3}}{2}\Rightarrow HN=\dfrac{1}{3}AM=\dfrac{a\sqrt{3}}{6}$.
			\item $\triangle SHN$ vuông tại $H$ có $SH=HN.\sin \widehat{SNH}=\dfrac{a\sqrt{3}}{6}\cdot \sin 60^\circ=\dfrac{a}{4}$.
			\item ${{V}_{S.ABC}}=\dfrac{1}{3}SH\cdot {{S}_{ABC}}=\dfrac{1}{3}\cdot \dfrac{a}{4}\cdot \dfrac{a^2\sqrt{3}}{4}=\dfrac{a^3\sqrt{3}}{48}$.
		\end{itemize}
	}
\end{bt}

\begin{bt}%[1H8H5-3]
% nguon: Archive/SP-4Phan-Dot3/Data/CD_1_OTHK2_De01.tex #4423

	Cho hình chóp $S.ABCD$ có đáy là hình vuông cạnh bằng $2$, cạnh bên $SA$ vuông góc với đáy và $SA=2\sqrt{3}$. Tính khoảng cách từ điểm $B$ đến mặt phẳng $(SCD)$.
	\loigiai{
		\immini{
			Vì $AB$ song song với $(SCD)$ nên khoảng cách từ $B$ đến $(SCD)$ bằng khoảng cách từ $A$ đến $(SCD)$.\\
			Vẽ $AH$ vuông góc với $SD$ tại $H$. Ta có $CD\perp AD$ và $CD\perp SA$, suy ra $CD\perp AH$.\\
			Mặt khác $AH\perp SD$ nên suy ra $AH\perp (SCD)$. Do đó khoảng cách từ $A$ đến $(SCD)$ bằng độ dài $AH$. Ta có $SD=\sqrt{SA^2+AD^2}=\sqrt{12+4}=4$.\\
			Từ hệ thức lượng $AH\cdot SD=SA\cdot AD$ ta tính được $$AH=\dfrac{SA\cdot AD}{SD}=\sqrt{3}.$$ Vậy khoảng cách cần tính bằng $\sqrt{3}$.}
		{\begin{tikzpicture}[scale=1, font=\footnotesize, line join=round, line cap=round,>=stealth]
			\path (0,0)coordinate (B)++(3,0) coordinate (C)++(1.3,1) coordinate (D) ($(B)+(D)-(C)$) coordinate (A) (0,3) coordinate (h) ($(A)+(h)$) coordinate (S) ($(S)!0.6!(D)$) coordinate (H);
			\draw (S)--(B)--(C)--(D)--cycle (S)--(C);
			\draw[dashed] (S)--(A)--(B) (H)--(A)--(D);
			\foreach \p/\g in {A/150,B/240,C/-60, D/0, S/90,H/45} \fill[black] (\p) circle(1pt)+(\g:0.3) node{$\p$};
			\foreach \x/\y/\z in {S/A/B,S/A/D,S/H/A,A/B/C,B/C/D,C/D/A,D/A/B}\draw pic[draw, angle radius=0.2cm]{right angle=\x--\y--\z};
		\end{tikzpicture}}
	} 
\end{bt}

\begin{bt}%[1H8V7-6]
% nguon: Archive/SP-4Phan-Dot3/Data/CD_1_OTHK2_De01.tex #4424

	Cho lăng trụ $ABC.A'B'C'$ có thể tích bằng $24$. Gọi $M$, $N$ và $P$ lần lượt là các điểm nằm trên các cạnh $A'B'$, $B'C'$ và $BC$ sao cho $M$ là trung điểm của $A'B'$, $B'N=\dfrac{3}{4} B'C'$ và $BP=\dfrac{1}{4} BC$. Đường thẳng $NP$ cắt đường thẳng $BB'$ tại $E$ và đường thẳng $EM$ cắt đường thẳng $AB$ tại $Q$. Tính thể tích của khối đa diện lồi $AQPC.A'MNC'$.
	\loigiai{
		Đặt $S$, $h$ lần lượt là diện tích đáy và chiều cao của lăng trụ đã cho ta có $S\cdot h=24$ và
		$V_{AQPC.A'MNC'}=V_{ABC.A'B'C'}-V_{BPQ.B'NM}$.\\
		Trong đó $BPQ.B'NM$ là chóp cụt có chiều cao $h$.\\
		Ta có $\dfrac{EB}{EB'}=\dfrac{EP}{EN}=\dfrac{EQ}{EM}=\dfrac{BP}{B'N}=\dfrac{BQ}{B'M} =\dfrac{PQ}{NM}=\dfrac{1}{3}$.\\
		Do đó hai tam giác $\triangle BPQ \backsim \triangle B'NM$ theo tỷ số $k=\dfrac{1}{3}$.
		\begin{center}
		\begin{tikzpicture}[scale=0.9, font=\footnotesize, line join=round, line cap=round, >=stealth]
			\def\a{4.5}
			\def\h{4.5}
			\path 	(0:0) coordinate (B')
					++(0:\a) coordinate (C')
					++(-135:3*\a/4) coordinate (A')
					($(A')+(90:\h)$) coordinate (A)
					($(B')+(90:\h)$) coordinate (B)
					($(C')+(90:\h)$) coordinate (C)
					($(A')!1/2!(B')$) coordinate (M)
					($(B')!3/4!(C')$) coordinate (N)
					($(B)!1/4!(C)$) coordinate (P)
					(intersection of N--P and B--B') coordinate (E)
					(intersection of E--M and B--A) coordinate (Q);
			\draw[dashed,thick] 	(B')--(C')		(M)--(N)--(P);
			\draw[thick]	(C)--(C') 	(B)--(B')	(A)--(A')	(M)--(E)--(P)--(Q)
							(B)--(E)
							(A)--(B)--(C)--cycle	(B')--(A')--(C');
			\foreach \x/\g in {A/180,B/180,C/0,A'/180,B'/180,C'/0, M/190, N/45, P/45, E/90, Q/-15}
					\fill[black] 	(\x) circle (1pt)
					($(\g:4mm)+(\x)$) node {$\x$};	
		\end{tikzpicture}
		\end{center}
		Suy ra\\
		\indent $S_{B'NM}=\dfrac{B'N}{B'C'}\cdot \dfrac{B'M}{B'A'} S=\dfrac{3}{4}\cdot \dfrac{1}{2} S=\dfrac{3}{8} S$;\\
		 \indent $S_{BPQ}=\left(\dfrac{1}{3}\right)^2 S_{B'NM}=\dfrac{1}{24} S$.\\
		Vì vậy
		$$
		V_{BPQ\cdot B'NM}=\dfrac{h}{3}\left(\dfrac{3}{8} S+\dfrac{1}{24} S+\sqrt{\dfrac{3}{8} S\cdot \dfrac{1}{24} S}\right)=\dfrac{13}{72} S \cdot h=\dfrac{13}{72} \cdot 24=\dfrac{13}{3}.
		$$
		$\Rightarrow V_{AQPC.A'MNC'}=24-\dfrac{13}{3}=\dfrac{59}{3}$.
	}
\end{bt}

\begin{bt}%[1H8H4-3]
% nguon: Archive/SP-4Phan-Dot3/Data/KNTT_1-OTGK2_De01.tex #4432
%[Câu 3]
Cho hình chóp $S.ABCD$ có đáy là hình vuông cạnh, cạnh bên $SA$ vuông góc với mặt phẳng đáy. Chứng minh rằng đường thẳng $BD$ vuông góc với mặt phẳng $(SAC)$.
\loigiai
{
\immini
{
Vì $\heva{ & SA\perp(ABCD) \\ & BD\subset(ABCD)}$ nên $SA\perp BD$. \\
Ta có $\heva{ & BD\perp AC\subset(SAC) \\ & BD\perp SA\subset(SAC) \\ & AC\cap SA=A} \Rightarrow BD\perp(SAC)$.
}
{
\begin{tikzpicture}[line join=round, line cap=round,thick]
\coordinate (A) at (0,0);
\coordinate (B) at (-2,-1);
\coordinate (D) at (3,0);
\coordinate (C) at ($(B)+(D)-(A)$);
\coordinate (S) at ($(A)+(0,2)$);
\draw(S)--(B) (S)--(C) (S)--(D) (B)--(C)--(D);
\draw[dashed,thin](S)--(A) (A)--(B) (A)--(D) (A)--(C) (B)--(D);
\pic[draw,thin,angle radius=3mm] {right angle = S--A--D} pic[draw,thin,angle radius=3mm] {right angle = S--A--B};
\foreach \i/\g in {S/90,A/-90,B/-90,C/-90,D/0}{\draw[fill=white](\i) circle (1pt) ($(\i)+(\g:3mm)$) node[scale=1]{$\i$};}
\end{tikzpicture}
}
}
\end{bt}
